\chapter{C1: the remaining source faces}
\label{ch:cert-4}
Formula tags are local to their certificate section. All source and testing fields remain separately labelled.

\section{\texorpdfstring{A simultaneous full-$D$/full-$Y$ identity and an extra-row conductor}{A simultaneous full-D/full-Y identity and an extra-row conductor}}
\label{component:pa73a17051258}

\paragraph{Scope and status.}
This is an independently reconstructed exact identity for the complete
original C1 parent, not a certificate admitting C14.3b. It retains the
same stationary source, $\Lambda=n=0$, $\beta\ge0$, all $j$ and $1-r$
openings, independent source/testing fields, and the remaining local
$v/u$-infinitesimal testing chart. No mixed inverse or comparison of
source and testing rates is used below.

Write $\mathcal P=F/4$, initially at $Z=1$. Set
\[
 a=1/\zeta,\quad t=\eta,\quad y=tx/a,\quad U=u^2,\quad V=v^2,
 \quad B=\beta+\tfrac12,\quad D_p=D+\tfrac12,
\]
\[
 \begin{gathered}
 E=tuW-v,\quad Q_t=Y-rv/t,\quad P=E/t-Q_t,\\
 \xi=X+hu,\quad R_D=v(U-x),\quad T_D=R_D-uv\xi/h.
 \end{gathered}
\]
Thus $Q_t=rE/t+\nu$ in the R15 notation. Put
\[
 q=\alpha+(m+B)(1-h)^2,\quad\mu=m-a^2q,
 \quad\lambda=\mu/(tax),\quad H=avW-u-a\xi,
\]
\[
 F_0=V-y,\quad F_b=F_0+vE,\quad F_t=F_0+(t/r)vQ_t,
 \quad T_0=A+\alpha+\tfrac12,\quad T_m=T_0+m,
 \quad e_4=r^2T_m/t^2.
\]
Here $F_t=\widetilde F$ and $F_b=\overline F$. Define pure-source
coefficients
\[
 b=Ar/t-m(1-h)-\alpha-\tfrac12,\qquad
 b'=b+e_4yt/(2r),
\]
\[
 \kappa=2[B(1-h)-D_pj/t-mr/t],\qquad
 \gamma=2b'/h+2D_pj/r-2tB(1-h)/r.
\]
They satisfy $b'=O(x)$, $t\kappa\sim-1$ and $\gamma\sim1$.

\paragraph{Complete parent, with no omitted hard row.}
Let
\[
 L=UV+(ax/t)V^2/U+tx^2/a-3xV,
\]
\[
 \mathscr C=\lambda L+\frac{\mu V}{t^2U}E^2
                +2\left(\frac{\mu v^3}{t^2U}-\lambda xv\right)E.
\]
The exact full-parent identity is
\begin{align*}
 \mathcal P={}&qH^2+\mathscr C
 +CP^2+AQ_t^2+D_p(P-hT_D)^2+2b'Q_tT_D\\
 &+\mathscr R+e_4UF_t^2-e_4yu(H+a\xi)F_t\\
 &+AUQ_t^2+B(1+V)\xi^2
                 -2(\mu/a)H\xi-\mu\xi^2\\
 &+\kappa uF_t\xi+\gamma uvQ_t\xi,                 \tag{JF1}
\end{align*}
where the entire joint radial row is
\[
 \mathscr R=\frac C{t^2}(F_b-rF_t)^2
       +\frac{r^2T_0}{t^2}F_t^2
       +\left(\frac{2mr}{t^2}+e_4x\right)F_bF_t.    \tag{JF2}
\]
In particular all $mVEQ_t/t$ terms are included in JF2; they are not
separate errors. The $mUVQ_t^2$ contribution is included in the complete
$e_4UF_t^2$ row, so it cannot also be retained independently.

A short independent way of certifying JF1 is useful. Without imposing
any source formulas, the difference between the original parent and
its displayed right side is exactly a linear combination of only
\[
                       vE,\ vQ_t,\ u\xi,\ UV,\ V,\ 1.
\]
The displayed expression has zero value and all five first derivatives
at the original contact $U=x,V=y,E=Q_t=\xi=0$.
Its six coefficients therefore vanish by the six original stationary
contact balances. This support calculation and the six contact identities
prove the identity for the complete original parent, including every state normal.

\paragraph{One strict radial pivot, with its exact companion.}
Set $\rho=m+rT_mx/2$. Then JF2 has the pure-source completion
\[
 \mathscr R=\frac C{t^2}
       [F_b-r(1-\rho/C)F_t]^2
       +\frac{r^2}{t^2}[T_0+2\rho-\rho^2/C]F_t^2.  \tag{JF3}
\]
Here $\rho=O(x)$ and $T_0+2\rho-\rho^2/C\sim1/2$;
all denominators are pure source and this remains uniform at both
source openings. Both $F_t/t$ and $F_b/t$ have bounded allocations
on JF3. This does not provide an independent strict $F_0/t$ pivot.
The non-small remaining coefficient is precisely $\gamma uvQ_t\xi$.
In particular a source split making $x/\sqrt t$ bounded does not
remove it, because $\gamma\sim1$.

\paragraph{The extra circuit row participates in the leading cancellation.}
Use $c^2=a/t$, $k^2=x$, $p=ckv/u$, and homogenize now in $Z$:
\[
 J_0=uv-x/c,\quad R_c=kv(tcW-Z),\quad
 R=\sqrt\lambda\,tk(cv-u)W,
 \quad \tau=tcW-Z.
\]
The generic exact circuit completion is
\[
 \mathscr C=\lambda J_0^2Z^2+R^2+
      \underbrace{\frac{2\lambda k}{p}R_c^2}_{\mathcal T^2}
      -\lambda xE^2.
\]
The full-$D$ conductor is
\[
 R_D=uJ_0Z+kR_c-(x/c)(H+a\xi),\qquad
 H+a\xi=avW-uZ.                                  \tag{JF4}
\]
The additional exact polarization is
\[
 E(H+a\xi)=uv\tau^2-t(u-cv)^2WZ.                  \tag{JF5}
\]
Consequently the extra circuit row and the principal cross left after
using JF4 satisfy
\[
 \mathcal T^2+\kappa h(x/c)E(H+a\xi)
 =\frac{x}{c}(2\lambda+\kappa h)uv\tau^2
      -\kappa ht\frac{x}{c}(u-cv)^2WZ.             \tag{JF6}
\]
This cancels the previously non-small $E$--angular Gram coupling.
JF6 is used as an exact identity, rather than as a separate positive
packet. In particular, its first coefficient has not been allocated
to a strict positive row. The sign of $WZ$ must be retained. The second summand is a
literal SOS on the pure-testing sign chart $WZ\ge0$ because
$-\kappa>0$, but not on the opposite sign chart.

\paragraph{An exact allocation of the negative small extra-row fraction.}
There is nevertheless a useful conductor:
\[
 x\frac{k}{p}R_c^2
          =U R_c^2-kuJ_0R_c\tau.                   \tag{JF7}
\]
If the testing state is normalized with $Z$ a unit and $W$ bounded,
then $\tau/Z=tc(W/Z)-1$ is bounded. Thus the right side is an
infinitesimal Gram on $\sqrt\lambda R_c$ and
$\sqrt\lambda J_0Z$, with coefficients $U$ and
$ku(\tau/Z)$. Since
$\sqrt\lambda R_c=R+\sqrt\lambda kE$, its additional coefficient
on $E/\sqrt t$ is $\sqrt{\lambda t}\,k=O(k)$.
Therefore a negative $O(x)$ fraction of $\mathcal T^2$ has a bounded
infinitesimal allocation on this state chart. No bound on $p$ is used.
The companion identity
\[
 F_b/t=(W/Z)J_0Z+(x/a)\tau                          \tag{JF8}
\]
may supply a strict $xZ$ radial pivot in a simultaneous state-chart
reassembly, since $\tau/Z\sim-1$ there. All divisions in JF7--JF8
are pure testing. Neither identity is an authorization to assume a
strict $E/\sqrt t$ pivot after the leading source block has become
semidefinite. That final simultaneous allocation is still required.

\paragraph{A direct conductor on the normalized $Z$-unit state chart.}
There is a simpler route on this state chart that does not perform JF4--JF6.
Put $e_E=E/\sqrt t$ and retain the original strict source completion
of the first line of JF1. Since $Q_t=E/t-P$ and $vZ=tuW-E$,
\[
 uvQ_t\xi=\sqrt t\,U(W/Z)e_E\xi
             -u(\xi/Z)e_E^2-uvP\xi.                \tag{JF9}
\]
When the testing state is normalized with $Z$ a unit, $W/Z$ and
$X/Z$ are bounded and $\xi/Z=X/Z+hu$ is bounded. Thus all three
coefficients in JF9 are infinitesimal in the bounded tensor ring.
The last uses the bounded pure-source allocation of $P$ on the
completed source roots. The first line of JF1 has source Schur pivot
$tS\sim1/2$, and $E/\sqrt t$ has a bounded allocation there,
including after the full $D_p$ shift. Consequently JF9 removes the
remaining $\gamma$ coupling on this state chart without producing
a semidefinite source frame.
The accompanying complete state-chart admission, if used, must still
list the retained positive rows and the strict angular block; JF9
alone does not cover states with normalized $Z$ infinitesimal.

\paragraph{Use of the full identity.}
The complete JF1 identity is the input to DG3/LD9 and the completed $J$-row conductor in Section~\ref{component:pe3bbfd3a8151}. Those arguments allocate the remaining normal and retain the full angular and radial packets.


\section{\texorpdfstring{C1.4: unbounded $L=x\zeta$, bounded-unit $\sigma=1/\eta$}{C1.4: unbounded L=xzeta, bounded-unit sigma=1/eta}}
\label{component:pd4b92bf7de71}

\label{pd4b92bf7de71:sec:c14-large-l-unit}
This section proves this full original-parent subdomain, retaining every
state normal and every independent source/testing rate. It makes no claim
for infinitesimal $\sigma$. The inherited common remote theorem leaves
$u$ bounded and $w=1/v$ infinitesimal; both $u$ infinitesimal and $u$ a
bounded unit are covered below. All calculations use the exact NC source,
with $d=1$, $n=\lambda=0$, and its original nonnegative $D,\beta$ faces.

\paragraph{Source chart and bounded products.}
Put $x=k^2$, $z=q^2=L^{-1}$, $r=bz$, $Q=1-b$, $\theta=1-h$,
$K=1+x$, $R=1+xbz$, $M=R-xz\theta$, and
\[
 J=1+xh-\sigma(1-bz),\qquad E=1-xh/\sigma.
\]
Here $x,z$ are positive infinitesimals, $0<b<1$, and $\sigma$ is a positive
bounded unit. The NC source gives exactly
\begin{align*}
 P&=\sigma+K+\sigma z(1+bx+2xh)+\sigma bx^2hz^2+x^2h\theta z,\\
 V_0&=z(1+xh)E+(1-bz)(1+xhz),\\
 C&=cz,&c&=\frac{K(1+\sigma bz)E}{2(1-bz)P},\\
 m&=\mu x^2z,&\mu&=\frac{K(1+\sigma bz)h(1+\sigma z)}{2\sigma MP},\\
 A&=a/z,&a&=\frac{V_0}{2bP},\\
 \alpha&=\ell/z,&\ell&=\frac{(1+\sigma bz)[(x/\sigma)Rh(1+\sigma z)+ME]}{2QP},\\
 D&=\frac{K(1+\sigma z)(1+\sigma bz)}{2JP}-\frac12,&
 \beta&=\frac{K(1+\sigma bz)(1+xhz)}{2\theta P}-\frac12.
\end{align*}
In particular $c,\mu,ab,\ell Q,DJ,\beta\theta$ are bounded,
$c,\mu,ab,\ell Q,DJ$ are positive units, and $ab-\ell Q\in z\mathcal O$.
At $x=z=0$ their residues are
\[
\begin{gathered}
c_0=\frac1{2(1+\sigma)},\quad
\mu_0=\frac h{2\sigma(1+\sigma)},\quad
(ab)_0=(\ell Q)_0=\frac1{2(1+\sigma)},\\
(DJ)_0=\frac{\sigma^2}{2(1+\sigma)},\quad
(\beta\theta)_0=\frac{(1+\sigma)h-\sigma}{2(1+\sigma)}.
\end{gathered}
\]
Source nonnegativity gives $0<\sigma\leq1$ and
$h\geq\sigma/(1+\sigma)$ at residue. Thus $h$ is a positive unit.
The openings $b,Q,J,\theta$ remain independent; $\beta=0$ is retained.
The face $D=0$ cannot occur here since $DJ$ is a positive unit.

Write $D_u=1+u^2$, $D_w=1+w^2$ and
\begin{gather*}
 M_w=W-wX+uY-uwZ,\quad R_w=uW-Y,\quad
 J_w=uwW+uX-wY-Z,\quad B_w=X+uZ,\\
 K_w=J_w^2+D_u(Y-wZ)^2+D_wB_w^2-D_uD_wZ^2.
\end{gather*}
The complete reciprocal parent is
\[
 F_w=4mM_w^2+4czD_wR_w^2+4(a/z)D_uD_wY^2
 +4DJ_w^2+4(\ell/z)D_u(Y-wZ)^2+4\beta D_wB_w^2+2K_w.
\]
Choose bounded positive-unit caps $a_*\leq a$, $\ell_*\leq\ell$,
$D_*\leq D$, and a bounded nonnegative cap $\beta_*\leq\beta$.
Keep as nonnegative squares the four centered excesses
\begin{gather*}
4(a-a_*)D_uD_wN_y^2,\quad4(\ell-\ell_*)D_uN_y^2,\\
4(D-D_*)(J_w+JZ)^2,\quad4(\beta-\beta_*)D_wN_x^2,\\
N_x=X+uhZ,\qquad N_y=(Y-bwZ)/q.
\end{gather*}
These roots vanish at the exact source contact
\[
u=k,\quad w=q\sqrt\sigma,\quad
(W,X,Y,Z)=(\sqrt\sigma/(kq),-kh,bq\sqrt\sigma,1).
\]
Set $d_J=(D-D_*)J$ and $b_\theta=\beta\theta$.
Their coefficients and the centered excesses preserve the original parent.

\paragraph{Bounded full-normal remainder.}
Define $G=wZ/q$, $R_h=G-uqW$ and
\[
c_N=8D_u[abD_w-\ell Q],\qquad c_G=4D_u[ab^2D_w+\ell Q^2].
\]
The $A,\alpha$ remainder is
$4D_u(a_*D_w+\ell_*)N_y^2+c_NN_yG+c_GG^2$.
Here $c_N\in(z,w^2)\mathcal O$ and $c_G$ is a positive unit.
Keep $c_GR_h^2+c_NN_yR_h$ in the bounded Gram and use
$G^2-R_h^2=2uwWZ-u^2zW^2$. The rest is the bounded polynomial
\begin{align*}
G_0={}&4D_u(a_*D_w+\ell_*)N_y^2+c_NuqWN_y+c_G(2uwWZ-u^2zW^2)\\
&+4mM_w^2+4czD_wR_w^2+4D_*J_w^2-8d_JZJ_w-4d_JJZ^2\\
&+4\beta_*D_wB_w^2+8(b_\theta-\beta_*\theta)uD_wZB_w
 -4(b_\theta-\beta_*\theta)u^2\theta D_wZ^2+2K_w.
\end{align*}
All its coefficients are bounded. Put $N_z=Z-uwW/\sigma$ and substitute
$Z=N_z+uwW/\sigma$, $X=N_x-uhZ$, $Y=qN_y+bwZ$.
For $t=u^2$, $\xi=w^2$, define
\[
 T_0=1+th+b\xi,\quad M_0=1+t(h-Q)\xi/\sigma,\quad
 R_f=1-b\xi/\sigma,\quad J_f=1-T_0/\sigma.
\]
In the formulas below $c_N,c_G$ mean their even-polynomial expressions.
The exact expansion is
\[
G_0=\phi W^2+wL_xWN_x+uqL_yWN_y+uwL_zWN_z+H(N_x,N_y,N_z),
\]
where
\begin{align*}
\phi={}&4\mu x^2zM_0^2+4cz(1+\xi)tR_f^2+c_Gt(2\xi/\sigma-z)\\
 &+4D_*t\xi J_f^2-8d_Jt\xi J_f/\sigma-4d_JJt\xi/\sigma^2
 +4b_\theta\theta(1+\xi)t^2\xi/\sigma^2\\
 &+2[t\xi J_f^2+(1+t)tQ^2\xi^2/\sigma^2
 +(1+\xi)t^2\theta^2\xi/\sigma^2-(1+t)(1+\xi)t\xi/\sigma^2],\\
L_x={}&-8\mu x^2zM_0+t[(8D_*+4)J_f-8d_J/\sigma
 +(8b_\theta+4\theta)(1+\xi)/\sigma],\\
L_y={}&8\mu x^2zM_0-8cz(1+\xi)R_f+c_N-(8D_*+4)\xi J_f
 +8d_J\xi/\sigma-4(1+t)Q\xi/\sigma,\\
L_z={}&8\mu x^2z(h-Q)M_0-8cz(1+\xi)bR_f+2c_G
 -(8D_*+4)T_0J_f-8d_JJ_f+8d_J(T_0-J)/\sigma\\
 &+8b_\theta\theta(1+\xi)t/\sigma+4(1+t)Q^2\xi/\sigma
 +4(1+\xi)t\theta^2/\sigma-4(1+t)(1+\xi)/\sigma.
\end{align*}
For a compact explicit specification of the entire quadratic block, put
\[
\begin{gathered}
m_N=-wN_x+uqN_y+uw(h-Q)N_z,\quad
r_N=-qN_y-bwN_z,\\
j_N=uN_x-wqN_y-T_0N_z,\quad
b_N=N_x+u\theta N_z.
\end{gathered}
\]
Then every quadratic normal coefficient is retained by the finite formula
\begin{align*}
H={}&4D_u(a_*D_w+\ell_*)N_y^2+4mm_N^2+4czD_wr_N^2
 +(4D_*+2)j_N^2-8d_JN_zj_N-4d_JJN_z^2\\
&+(4\beta_*+2)D_wb_N^2+8(b_\theta-\beta_*\theta)uD_wN_zb_N
 -4(b_\theta-\beta_*\theta)u^2\theta D_wN_z^2\\
&+2D_u(qN_y-QwN_z)^2-2D_uD_wN_z^2.
\end{align*}
The exact source substitutions give, at $(t,\xi)=(x,\sigma z)$,
\[
\phi=\phi_t=\phi_\xi=L_x=L_y=L_z=0.
\]
Moreover $\phi(0,\xi)=4\mu x^2z$,
$L_x(0,\xi)=-8\mu x^2z$, and $\phi(t,\sigma z)\in z\mathcal O[t]$.
The checker verifies all six source jets and these supports separately.

\paragraph{Exact six-root lift in the critical strip $u$ infinitesimal.}
Monic division, first by $\xi-\sigma z$ and then by $t-x$, gives bounded
polynomial coefficients satisfying
\begin{align*}
\phi&=z(t-x)^2 A+t(t-x)(\xi-\sigma z)B+t(\xi-\sigma z)^2C,\\
L_x&=(t-x)(tA_x+zB_x)+t(\xi-\sigma z)C_x,\\
L_y&=(t-x)A_y+(\xi-\sigma z)B_y,\\
L_z&=(t-x)A_z+(\xi-\sigma z)B_z.
\end{align*}
For example, the constant coefficient of
$L_x(t,\sigma z)/(t-x)$ equals $8\mu xz$, giving its $(t,z)$ support.
No small or mixed inverse enters these monic quotients. Also $A(0)=4\mu$.
Take the six finite separated roots
\[
 N_x,\quad N_y,\quad N_z,\quad R_h,\quad
 A_\theta=q(t-x)W,\quad V=uw(t-x)W.
\]
The complete scalar lift follows from
\begin{align*}
 z(t-x)^2W^2&=A_\theta^2,\\
 t(t-x)(\xi-\sigma z)W^2&=u\sigma A_\theta R_h-\sigma VN_z,\\
 t(\xi-\sigma z)^2W^2&=\sigma^2(qR_h-wN_z)^2.
\end{align*}
All first normals lift by the following finite list, which includes their
whole coefficient support:
\begin{align*}
wt(t-x)W&=uV,&wz(t-x)W&=qwA_\theta,\\
wt(\xi-\sigma z)W&=u\sigma qwR_h-u\sigma w^2N_z,\\
uq(t-x)W&=uA_\theta,&uq(\xi-\sigma z)W&=\sigma zR_h-\sigma qwN_z,\\
uw(t-x)W&=V,&uw(\xi-\sigma z)W&=\sigma qwR_h-\sigma w^2N_z.
\end{align*}
Together with the displayed $H$ these give an exact bounded six-root Gram,
not merely a restriction to the contact line. At residue its first five
diagonal pivots are respectively
\[
2+4\beta_*,\quad4(a_*+\ell_*),\quad
4[D_*+d_J(1+\sigma)],\quad \frac2{1+\sigma},\quad4\mu_0.
\]
They are positive units through all allowed openings. The only potentially
nonzero remaining cross entry is $N_zV$. The exact relation
\[
 V^2=\sigma tA_\theta^2+u\sigma(t-x)A_\theta R_h
       -\sigma(t-x)VN_z
\]
introduces an arbitrarily large finite positive $V^2$ pivot with only
infinitesimal changes to all other entries. Its real Schur complement
therefore makes the residue strictly positive definite. The exact bounded
Gram is positive definite by the standard finite-dimensional local
criterion. All excess squares are nonnegative with separated coefficients.

\paragraph{Away from the angular contact: $u$ a bounded unit.}
The same source chart has a positive-unit angular gap at every such real
residue. A convenient raw finite frame is
$E_t=qW,F_t=Y/q,G_t=wZ/q,R_t=wW,X,Z$; products obey $E_tG_t=R_tZ$.
The residue parent divided by $2/(1+\sigma)$ is exactly
\[
u^2E_t^2+D_u\{F_t^2/b+(F_t-G_t)^2/Q\}
 +\frac{(uR_t+uX-Z)^2}{J_0}
 +\frac{(X+uZ)^2}{\theta}-(1+\sigma)D_uZ^2,
\quad J_0=1-\sigma.
\]
Eliminate the positive $F_t$ and $X$ directions and allocate
the common product between the remaining two quadratic blocks. Up to a
common positive factor the two blocks have diagonal products $AB,CD_0$
and total off-diagonal magnitude $E_0$, with $U=u^2>0$ and
\[
\begin{gathered}
\delta=1-\sigma+\theta U,\qquad
A=U,\quad B=1+U,\quad C=U/\delta,\\
D_0=(1+U)[(1+U)/\delta-(1+\sigma)],\quad
E_0=\sqrt U(1+U)/\delta.
\end{gathered}
\]
Both diagonal products are strictly positive, since
$1+U-(1+\sigma)\delta=\sigma^2+U[(1+\sigma)h-\sigma]>0$.
All square roots in this sentence are real-residue quantities used to
choose real allocation constants, not mixed tensor radicals.
The required strict inequality is $E_0<\sqrt{AB}+\sqrt{CD_0}$.
Set $N=U(2\sigma-\theta U)$. Directly,
\[
E_0^2-AB-CD_0=(1+U)N/\delta,\quad
N^2-4U^2[1+U-(1+\sigma)\delta]
 =U^3(\theta^2U-4h).
\]
If $N\leq0$ strictness is immediate. If $N>0$, then
$\theta^2U<2\sigma\theta\leq2h<4h$, again giving strictness.
The single-opening faces follow by their constrained limits. At the only
double opening $J=\theta=0$, so $\sigma=h=1$, the exact limit is
$E_0-\sqrt{CD_0}\longrightarrow\sqrt U$, whereas $AB-U=U^2>0$.
Here is a full-frame square completion which also justifies those cap
limits. For clarity first take $u>0$ and select a finite real
$\rho$ in the nonempty interval
\[
\max\{-\sqrt{AB},E_0-\sqrt{CD_0}\}<\rho
 <\min\{\sqrt{AB},E_0+\sqrt{CD_0}\}.
\]
Its nonemptiness follows from the strict allocation inequality above.
Put
\[
r_\rho=D_u/u-\rho\delta/U,\qquad
s_\rho=D_u[2\rho/u-(1+\sigma)]-\rho^2\delta/U>0.
\]
The preceding complete residue form equals
\begin{align*}
&UE_t^2+D_uG_t^2-2\rho E_tG_t
 +\frac{D_u}{bQ}(F_t-bG_t)^2\\
&\quad+\frac{\delta}{J_0\theta}
 \left[X+\frac{U\theta R_t+u(J_0-\theta)Z}{\delta}\right]^2
 +\frac U\delta(R_t-r_\rho Z)^2+s_\rho Z^2.
\end{align*}
The equality uses only $E_tG_t=R_tZ$ and is checked exactly.
For $u<0$ take the correspondingly signed $\rho$; all sign choices are
real constants in this allocation. All triangular minimizer coefficients
are bounded through every opening because $U$ is a unit and
$\delta=J_0+U\theta$. Their positive hard weights stay bounded below.
At the double face choose $|u|<|\rho|<\sqrt{U(1+U)}$, with matching
sign; then $s_\rho$ has a positive-unit limit. This proves uniform
coercivity of the constrained residue, not just determinant positivity.

For completeness the exact centered cap rows in the raw frame are
\[
 F_t-bG_t,\quad
 uR_t+uX-qwF_t+(J-1)Z,\quad X+uhZ.
\]
Only one of $b,Q$ can tend to zero. The other two rows have hard parts
on $(R_t,X)$ with determinant $u$, a unit; each single row has a unit
hard entry. Their simultaneous kernel is solved exactly by
\[
F_t=bG_t,\qquad X=-uhZ,\qquad
R_t=\{(u^2h-J+1)Z+qwbG_t\}/u.
\]
More generally, prescribed changes $(e_F,e_D,e_\beta)$ in these three
rows are produced by
$\Delta F_t=e_F$, $\Delta X=e_\beta$,
$\Delta R_t=(e_D-u e_\beta+qw e_F)/u$.
This is a bounded separated right inverse: its only inversion is the
pure-testing unit $u$. For the finite-cap existence argument fix $u$
at its ordinary real residue, so the same inverse is an ordinary real
finite-dimensional right inverse. No order on the tensor product is
used. It lifts the limiting constrained kernel to the exact centered
kernel while keeping every opening independent. The
bounded residual form restricted to that kernel is the strictly positive
constrained form just computed. An elementary block Schur complement
then provides finite sufficiently large caps on every divergent row;
bounded rows retain their full coefficients. This is the finite-cap
lemma used in the UH proof. It yields a strictly positive bounded
full-frame Gram. Exact coefficients differ from their residues
infinitesimally, completing this strip.
All divisions by $\delta$ and radicals in this paragraph belong only
to the ordinary real residue calculation and its face limits. The exact
tensor Gram uses the polynomial raw frame, real allocation constants,
and finite source caps; it introduces no mixed inverse or radical.

Consequently this section closes the $L$ unbounded, $\sigma$ bounded-unit
subdomain of C1.4 under the original parent, including all state normals,
all independent opening rates, and $\beta=0$. It does not close the
$\sigma$ infinitesimal subdomain or claim a whole-row count change.


\section{\texorpdfstring{Small $L$, bounded $\eta$: exact source openings and the remote complement}{Small L, bounded eta: exact source openings and the remote complement}}
\label{component:p6d2f659c0fde}

\label{p6d2f659c0fde:sec:r8-c1-small-L-bounded-eta}

Work in the original C1 source, with $d=1$, $x\in\Inf_{>0}$,
$\zeta$ noninfinitesimal, $L=x\zeta\in\Inf_{>0}$, and $\eta$
bounded. Native $r=1-a\eta$, $h=1-\theta$ remain strictly positive.
This section admits every original projective test except possibly
$u,v$ both infinitesimal. It does not admit the remaining two-soft
local chart.

\paragraph{Exact source restrictions, including the total-scale intersection.}
Put $\rho=r+\eta$ and $j=\eta J=\eta(1+xh)-1+r>0$.
Since $a<1+xh$, $\rho$ is a positive bounded unit. The exact beta
numerator is
\[
 2\beta\theta P_d=K\rho(xh+L)-\theta P_d.       \tag{SB1}
\]
It implies $\theta=O(x+L)=O(L)$, since $x/L=1/\zeta$ is bounded.
Consequently
\[
 h,M,E,Q,P_d,P_4\sim1,\qquad B_5\sim1.           \tag{SB2}
\]
For the last assertion, use the exact expression
$B_5=K(1+L\eta)\rho-jP_d$: its residue is
$r+\eta-(\eta-1+r)=1$. Thus original $D=0$ cannot occur on
this source face; literal $\beta=0$ remains included.
The original formulas, with openings canceled before taking residues,
give uniformly
\[
\begin{gathered}
 C(1-r)\sim\rho/2,\quad Ar/\eta\sim1/2,
 \quad Dj\sim1/2,\\
 mL/x^2\sim\rho/2,\quad
 \alpha/[L(L+x)]\sim\rho/2,\quad
 (2\beta+1)\theta/(x+L)\sim\rho.
\end{gathered}                                         \tag{SB3}
\]
In particular $C,D$ are strong, $m,\alpha$ weak, and $A$ is strong
precisely when $\eta$ is a unit. The four exact large-weight
alternatives are
\[
\begin{array}{c|c}
\text{weight unbounded}&\text{equivalent opening}\\\hline
C&1-r\in\Inf\\
D&j\in\Inf\\
A&r\in\Inf\\
\beta&\theta/(x+L)\in\Inf.
\end{array}                                                   \tag{SB4}
\]
For the $A$ equivalence, $r\in\Inf$ forces $\eta$ to be a unit
because $j>0$; if $r$ is a unit, $\eta/r$ is bounded.
All coefficients are bounded exactly when the four quantities in
the right column of (SB4) are positive bounded units. That entire
interior is already admitted by the positive-total-scale theorem.
After intersection with $\delta_{\rm tot}\in\Inf_{>0}$, at least
one of (SB4) therefore occurs. These are exhaustive necessary source
alternatives, not assertions of their independent realizability.

\paragraph{The unit-$\eta$ remote tests: two complete three-state blocks.}
Suppose first $\eta$ is a unit, with real residue $e>0$, and write
$r_0=\operatorname{res}r\in[0,1]$. For an interior opening,
the source residues are
\[
 C_0=\frac{r_0+e}{2(1-r_0)},\quad
 A_0=\frac e{2r_0},\quad D_0=\frac1{2(e+r_0-1)}.       \tag{SB5}
\]
Retain the original $m,\alpha,\beta$ packets throughout. Bounded
selected weights retain their actual coefficients; unbounded ones
are capped at sufficiently large finite real numbers and their
excesses retained as exact squares.

For $u$ a bounded nonzero unit and $v$ infinitesimal, use the full
raw frame
\[
 p=uW,\quad Y,\quad R=vZ,\quad B_t=X+uZ.
\]
The exact core is (CE2), and the $D$ event is
$J=p-Y+uvB_t-D_uR$. Every Gram coefficient is bounded after
capping, and no state normal is omitted. Set $U=\operatorname{res}u^2>0$.
The residue, divided by four, is $B_t^2/2$ plus
\[
\begin{split}
 f_U={}&C_0(p-Y)^2+A_0(1+U)Y^2
       +D_0\{p-Y-(1+U)R\}^2\\
 &+\tfrac12\{(p-Y)^2+(1+U)UR^2-2(1+U)pR\}.
\end{split}                                            \tag{SB6}
\]
Its first two principal minors and determinant are, respectively,
\[
\begin{gathered}
\Delta_1=\frac{e(e+r_0)}{2(1-r_0)j_0},\qquad
\Delta_2=\frac{e^2(1+U)(e+r_0)}{4r_0(1-r_0)j_0},\\
\Delta_3=\frac{Ue(1+U)^2(e+r_0)^2}{8r_0(1-r_0)j_0}>0,
\qquad j_0=e+r_0-1.
\end{gathered}                                      \tag{SB7}
\]

For $u$ infinitesimal and $v$ a bounded nonzero unit, use the full
frame $(p=uW,X,Y,R=vZ)$. The same exact parent has residue,
divided by four, $(1+V)X^2/2$ plus
\[
\begin{split}
 f_V={}&C_0(1+V)(p-Y)^2+A_0(1+V)Y^2
       +D_0(p-Y-R)^2\\
       &+\tfrac12\{(p-Y)^2+VY^2-2pR\},
 \qquad V=\operatorname{res}v^2>0.
\end{split}                                            \tag{SB8}
\]
The leading two minors are positive, and
\[
 \det f_V=\frac{V(1+V)(e+r_0)^2}{8r_0(1-r_0)j_0}>0. \tag{SB9}
\]

These arguments include every opening in (SB4). To see explicitly
why the finite cap remains valid, impose the kernels of the divergent
events and inspect the remaining two-state Gram. The determinants
are
\[
\renewcommand{\arraystretch}{1.7}
\begin{array}{c|c|c}
\text{opening}&\det(f_U|_{\ker})&\det(f_V|_{\ker})\\\hline
r_0=0,\ e>1&\dfrac{Ue^2(1+U)}{4(e-1)}&\dfrac{Ve}{4(e-1)}\\
r_0=1&\dfrac{U(1+U)^2(e+1)}4&\dfrac{V(e+1)}{4e}\\
j_0=0,\ 0<r_0<1&
 \dfrac{U(1+U)^2}{4r_0}&\dfrac{V(1+V)}{4r_0(1-r_0)}.
\end{array}                                                   \tag{SB10}
\]
The respective kernels are $Y=0$, $p=Y$, and
$p-Y=(1+U)R$ for $f_U$ or $p-Y=R$ for $f_V$.
All first pivots are positive. At the only double opening,
$r_0=j_0=0,e=1$, the two kernels leave the positive scalar
$U(1+U)R^2/2$, respectively $VR^2/2$. There are no other
multiple openings since $e>0$. The elementary finite-penalty
Schur lemma therefore supplies finite real caps simultaneously
for every divergent selected event.

For completeness the remaining projective ends also use original
homogeneous degree. If $u$ is unbounded, (CE4)--(CE5) for bounded
$v$, and (CE7)--(CE8) for unbounded $v$, are strict with any positive
finite $C,A$ caps. If $v$ is unbounded and $u$ bounded, put
$w=1/v$ and use $(p=uW,X,Y,Z)$. The exact $D$ event is
$w(p-Y)+uX-Z$. At $w=0$ the core is
$(1+u^2)(X^2+Y^2)$, while the three caps give
$C(p-Y)^2+A(1+u^2)Y^2+D(uX-Z)^2$; these are strict on
all four roots, including $u=0$. Each term suppressed in this
residue has a bounded infinitesimal coefficient in this same
full frame. Literal $w=0$ is therefore included. The both-unit
tests are covered by the inherited noncommuting-residue theorem.
Together these statements exhaust the complement of $u,v\in\Inf$.

\paragraph{The infinitesimal-$\eta$ remote tests.}
Here $r\sim1$, $C,D$ are unbounded, and $y=L\eta$ is infinitesimal.
Define the pure-source harmonic coefficient
$\gamma=yCD/(C+D)$. Uniformly through the $a,J$ openings,
\[
 \gamma\sim L/2,\qquad A/y\sim1/(2L),\qquad
 \gamma A/y\sim1/4,\qquad y/\gamma\sim2\eta\in\Inf. \tag{SB11}
\]
The strong-II remote proof (RR3)--(RR6), which uses only $C,D$
unbounded, covers every noninfinitesimal $v$ and every $u$.
For $u$ noninfinitesimal and $v$ infinitesimal, use (RR7)--(RR9)
with zero alpha cap and $k=-1$. In its notation let
$a=u/\sqrt{1+u^2}$, $b=1/\sqrt{1+u^2}$,
$\ell=vY$, $x_t=bX$, $R_0=x_t+a(Z-\ell)$, and $q=\sqrt yY$.
Then
\[
 M_{-1}=\tfrac12R_0^2+\tfrac12(b\ell)^2,\qquad
 v(Z-\ell)=avR_0-bv(b\ell)-b\sqrt y\,V_h,       \tag{SB12}
\]
where $V_h=v(aX-bZ)/\sqrt y$. Normalize the harmonic pair as
$P_h=\sqrt\gamma V_h$, $Q_h=\sqrt{A/y}\,q$. Its residue
determinant is $1-b_0^2>0$ by (SB11), where $b_0=\operatorname{res}b<1$;
the off-diagonal entry is $b_0$. The sole negative term in (RR9),
$-v^2(Z-\ell)^2/2$, is an infinitesimal Gram perturbation of
$(R_0,b\ell,P_h)$ by (SB12), since $\sqrt{y/\gamma}$ is
infinitesimal. Thus the full four-root Gram is strict. Retain the
original positive $\alpha(Z-\ell)^2$, the C1 beta event, and every
other unused packet. The homogeneous endpoint $u=\infty$ has
$a=\pm1,b=0$ and is included; literal $v=0$ also causes no change.

\paragraph{The actual remaining obstruction.}
Only the simultaneous infinitesimal-test chart remains, intersected
with (SB4) and all previous admissions. It has two soft state
directions, not the single soft ray of the admitted $L$-unit proof.
Indeed the original parent at $u=v=0$ is exactly
\[
 F= (4m+4\beta+2)X^2+(4C+4A+4D+2)Y^2+4\alpha Z^2. \tag{SB13}
\]
Since $m,\alpha$ are weak, the two $W,Z$ directions remain soft
after the hard $X,Y$ normals are normalized. On the contact graph
$X=-huZ$, $Y=ruW$, the exact two-soft scalar is (WS2)--(WS3)
with $n=0$, the WS letters $h,r$ replaced by native $r,h$, and
the additional positive term
$\beta\theta^2U(1+V)Z^2$ in $F_g/4$. Its two negative
monomials are still $-2uv(b_0+b_1U)WZ$. No strict finite Gram
on the existing hard roots supplies the missing $W,Z$ pivots.
A weighted contact circuit and its complete first-normal allocation
are still needed. Neither the exact source alternatives nor the
remote theorem is presented as that local completion.


\section{\texorpdfstring{A C1 extension: $L$ a unit and $\eta$ infinitesimal}{A C1 extension: L a unit and eta infinitesimal}}
\label{component:p080aa62de16d}

\label{p080aa62de16d:sec:r7-c1-L-unit-eta-small}

\begin{theorem}
The entire original C1 source subdomain
\[
 x=k^2\in\Inf_{>0},\qquad L=x\zeta\text{ a positive bounded unit},
 \qquad \eta\in\Inf_{>0}
\]
has a full-parent tensor-SOS certificate for every independent
projective testing pair.  All original openings and the retained
boundaries $D=0$ or $\beta=0$ remain allowed whenever they occur
in this source subdomain.
\end{theorem}

The proof is an exact adapter to the strong-II local and remote
certificates.  Its additional event is retained with a centered
square; it is not deleted from the parent.

\paragraph{The native source in the shared contact geometry.}
Keep the native notation of (NC1)--(NC4), but in this section denote
its $r,h,a$ by $r_n,h_n,a_n$ to avoid collisions.  Normalize $d=1$.
The exact native contact with $Z_c=1$ is
\[
 (u_c,v_c;W_c,X_c,Y_c,Z_c)
 =\left(\sqrt x,\sqrt{x\eta\zeta};
 \sqrt{\zeta/\eta},-\sqrt x\,h_n,
 \sqrt x\,r_n\sqrt{\zeta/\eta},1\right).
\]
It therefore has the same contact geometry used in (LII2), with
\[
 z=x,\quad y=L\eta,\quad B=L,\quad H=Lr_n=B-a_ny,
 \quad r=h_n=1-\theta,\quad t=a_n,\quad
 j=1+zr-t=J_n.                                           \tag{LC1}
\]
In particular $z,y$ are positive infinitesimals.  Here $r_n\sim1$,
$0<Lr_n<1$, and hence $b_0:=\operatorname{res}B$ satisfies
$0<b_0\le1$.  We use $b_0$ only for this source residue.

Set, as in the local strong-II proof,
\[
 c=Cyt,\quad d_p=Dyj,\quad
 \mu=m/z^2,\quad a_A=A/y,\quad
 \sigma=\alpha(1-H),\quad \tau=\beta(1-r)=\beta\theta.
\]
These are bounded source quantities, and $c,d_p,\mu,a_A,\sigma$
are positive units.  Writing $r_0=\operatorname{res}r$, exact
substitution in (NC2)--(NC4), with opening factors canceled first,
gives
\[
\begin{gathered}
 c,d_p\sim\frac{b_0}{2(1+b_0)},\qquad
 a_A\sim\frac1{2b_0(1+b_0)},\\
 \mu\sim\frac{r_0}{2b_0(1+b_0)},\qquad
 \sigma\sim\frac{b_0^2}{2(1+b_0)},\\
 \operatorname{res}\tau=
 \frac{b_0}{2(1+b_0)}-\frac{1-r_0}{2}.                    
\end{gathered}
\tag{LC2}
\]
Since $\tau\ge0$, the last identity implies
$r_0\ge1/(1+b_0)>0$, proving the asserted unit lower bound for
$\mu$.  Also $t+j=1+zr\sim1$, and $c,d_p$ being units proves
$C,D$ unbounded uniformly through either opening $t=0$ or $j=0$
in residue.  Thus $D=0$ is in fact absent from this specific
subdomain by the exact source implications; it has not been excluded
by an additional hypothesis.  Literal $\beta=0$ is included.

The source-harmonic quantity $yCD/(C+D)$ and the cap behavior of
$\alpha$ have precisely the residues (RR1)--(RR2), with
$h=b_0$.  The remote proof (RR3)--(RR13) therefore applies unchanged,
retaining the extra positive packet
$4\beta D_v(X+uZ)^2$ throughout.  This covers every test outside
the local region where $u,v$ are both infinitesimal, including the
original projective endpoints.

\paragraph{The additional event and its exact local remainder.}
For the local region use exactly the graph, roots, harmonic pair,
and caps (LII2)--(LII4), with the source quantities (LC1)--(LC2).
Put
\[
 R=v(U-z)W,\qquad
 E_\beta=\xi+u\zeta+R/B
          =(X+uZ)-\frac{z(1-r)v}{B}W.                    \tag{LC3}
\]
This root vanishes at the actual contact.  Retain the entire
positive packet $4\beta D_vE_\beta^2$, in addition to the first,
second and fourth packets of (LII5); there is no $n$ packet since
$n=0$.  The exact remainder from the $\beta$ event is
\[
 \begin{split}
 \Delta_\beta={}&
 \frac{8\tau z}{B}D_vvW(\xi+u\zeta)\\
 &+\frac{4\tau z}{B^2}D_vV\{2U-z(1+r)\}W^2.
 \end{split}                                             \tag{LC4}
\]
All its coefficients are bounded even if $\beta$ is unbounded:
only $\tau=\beta(1-r)$ occurs.  The exact full remainder is (LII6)
with its $n$ line deleted, plus (LC4).  Thus no state normal,
boundary coefficient, or first contact jet has been omitted.

The $R_hW$ gauge (LII7) is unchanged, because (LC4) contains no
$R_h$.  The coefficient representation (LII8), its exact contact
and first-normal vanishings (LII9)--(LII10), and its bounded monic
divisions (LII12)--(LII13) all remain valid.  Explicitly, (LC4)
adds $8\tau z(1+V)/B$ to $P_\xi$ and $P_\zeta$, adds zero to
$P_N$, and adds its printed scalar line to $\phi$.  Hence the
essential support $P_\xi\in(U,z)$ and $\deg_U\phi\le2$ remain
literal polynomial statements.

The only changed data in (LII11) are
\[
\begin{split}
 d_s&=\frac{4z^2}{B^2}
   \left\{\mu(B+zr)^2+\frac{r^2}{2}-\tau(1+r)\right\},\\
 e_s&=\frac{4z^2}{B^2}
   \left\{\frac{r^2}{2}-\tau(1+r)\right\}.
\end{split}                                             \tag{LC5}
\]
Both quotients by $z^2$ are bounded; $e_s$ need not be positive.
Its contribution to the final Gram has an infinitesimal coefficient
because it is multiplied by $V$.  The decisive exact residue
cancellation, using (LC2), is
\[
 \begin{split}
 \operatorname{res}\left\{
 \mu B^2+\frac{r^2}{2}-\tau(1+r)\right\}
 &=\frac{b_0r_0}{2(1+b_0)}+\frac{r_0^2}{2}
 -\left\{\frac{b_0}{2(1+b_0)}-\frac{1-r_0}{2}\right\}(1+r_0)\\
 &=\frac1{2(1+b_0)}.
 \end{split}                                             \tag{LC6}
\]
Consequently
\[
 \operatorname{res}(d_s/z^2)=\frac2{b_0^2(1+b_0)},\qquad
 \operatorname{res}A_\xi=\frac4{b_0(1+b_0)}.              \tag{LC7}
\]
The second equality is unchanged from (LII18): the new $P_\xi$
term is a bounded multiple of $z$ independent of $U$, so it has
zero $U$-quotient residue.  The four hard-normal residues (LII17)
are unchanged as well.

Apply the exact allocations (LII15) to all the coefficients,
including (LC4).  They give the same bounded six-root Gram
$(N,\xi,\zeta,R_h,R,T)$ as in (LII16).  Its finite $(\xi,R)$
block is exactly (LII19), with $h=b_0$, and has strictly positive
Schur complement $2/[b_0(1+b_0)^2]$.  The same exact conductor
(LII20) supplies an arbitrarily large real $T^2$ pivot while
changing every other residue entry by zero.  The complete bounded
Gram is strict.  Restoring the four retained packets, including
$4\beta D_vE_\beta^2$, proves the local tensor SOS for all
independent rates and literal affine axes.  The remote admission
above completes the theorem.

\paragraph{Verification and exact frontier effect.}
The accompanying checker reconstructs the entire original C1 parent
from this remainder and all four retained packets, verifies the new
centered-event identity, the unchanged contact-ideal support, every
new source residue, the finite block, and the conductor.  Stationarity
itself is the unchanged native source identity, not an added
restriction.  The result removes the whole $L$-unit, $\eta$-infinitesimal
subdomain from C1.4.  It does not assert the other $L$-unit source
regimes or the entire C1.4 row.


\section{\texorpdfstring{A sharper source bound on the remaining $L$-unit, $\eta$-large face}{A sharper source bound on the remaining L-unit, eta-large face}}
\label{component:p57d751e20564}

\label{p57d751e20564:sec:r7-c1-L-unit-eta-large-frontier}

This section proves a source restriction and a leading quadratic
identity.  It does not claim the remaining full-parent admission.
Use the unchanged native notation with $x\in\Inf_{>0}$,
$L=x\zeta$ a positive unit, and $\eta$ unbounded.  Write
$\theta=1-h$ and $K=1+x$.

In the exact polynomial $B_5=a_2\eta^2+a_1\eta+a_0$, substitution
$\zeta=L/x$ gives
\[
 a_2=-xh\{KL+x\theta(1+xh)\},\qquad a_1=xb_1,
                                                               \tag{LH1}
\]
where
\[
\begin{split}
 b_1={}&-Lhrx-Lh-Lr+L-h^2rx^2+h^2rx-3h^2x\\
       &-2hrx+hx-3h+1,
\end{split}
\]
and
\[
 a_0=Kr+(1-r)\{1+L+2xh+rxL+rx^2h\}.
\]
Both $b_1$ and $a_0$ are bounded uniformly in every source opening,
while $KL+x\theta(1+xh)$ has a positive-unit lower bound.
The original $B_5\ge0$ inequality therefore implies
\[
 x\eta h\le O(x)+O(\eta^{-1})\in\Inf,
 \qquad\boxed{E=1-x\eta h\sim1.}                       \tag{LH2}
\]
This is stronger on this face than the global bound $E\ge1/18$.
It does not require $x\eta$ bounded.  In particular it is not
legitimate to infer $x\eta^2h=O(1)$ from (LH2); the exact available
bound is $x\eta^2h=O(x\eta)+O(1)$.

The corresponding opening-cleared leading source weights are
\[
 C_0=\frac1{2L(1-r)},\qquad
 A_0=\frac{1+L(1-r)}{2Lr},\qquad
 \alpha_0=\frac L{2(1-Lr)}.                              \tag{LH3}
\]
The ratios of the actual weights to these positive expressions
tend to one.  In reciprocal Bob coordinates put $H_t=wZ$ and
$P=uW$.  The leading three-state quadratic has the exact identity
\[
\begin{split}
 &C_0(P-Y)^2+A_0Y^2+\alpha_0(Y-H_t)^2
       +\tfrac12Y^2-PH_t\\
 &\quad=C_0\{P-Y-L(1-r)H_t\}^2
       +(A_0+\alpha_0+\tfrac12)\{Y-LrH_t\}^2.
\end{split}                                             \tag{LH4}
\]
All radical or inverse coefficients in this identity are pure source.
Its two roots vanish at the exact contact to leading order.
The original $m,D,\beta$ packets, the small corrections to (LH3),
and every normal coupling still require an exact simultaneous
allocation.  Neither (LH2) nor (LH4) is a tensor-SOS certificate for
that remaining source face.

\paragraph{Why the completed C2 IH finite block cannot be copied unchanged.}
The common contact geometry is still exact, but its source pivot
hypotheses change. To make this failure concrete, take a positive
pure-source infinitesimal $s$ and the native rational source
\[
 L=\tfrac12,\qquad x=s^2,\qquad\eta=s^{-1},\qquad
 r_n=s,\qquad h_n=4s,\qquad\theta=1-4s.              \tag{LH5}
\]
All native strict gaps are positive for sufficiently small $s$.
Opening-cleared substitution into the exact formulas gives
\[
 C\sim1,\quad sA\sim\tfrac32,\quad
 D/s^2\sim\tfrac32,\quad m/x^2\sim4,\quad
 \alpha\sim\tfrac14,\quad\beta/s\sim\tfrac12.       \tag{LH6}
\]
Thus this diagnostic also has an unbounded original coefficient;
it is not removed merely by the all-coefficients-bounded argument.

Write $B=L$, $H=Br_n$, $r_g=h_n$, $z=x$, $y=B\eta$.
In the C1 version of the IC7 scalar bookkeeping, the beta event
adds $4\beta(1-r_g)^2$ to its mixed coefficient $K$. Explicitly set
\[
\begin{split}
 K_1={}&4AH^2+8DHr_g+4\alpha(1-H)^2
       +4m(H+r_g-1)^2\\
       &+2(H+r_g-1)^2+4\beta(1-r_g)^2,\\
 k_1={}&K_1-4B^2m/z^2+8AH^2/y.
\end{split}                                                   \tag{LH7}
\]
The same exact substitution yields $k_1\to-1$. Consequently
$k_1/(Bs)$ is unbounded, whereas the completed strong-C2 IH
finite block requires this coefficient to be bounded in its
$sT$ normalization. The C2 premise $r_g/s\le c_s$ also fails:
here $r_g/s=4$ and $c_s=(1+B-H)/B\to3$.
This exhibits the precise failed transfer premise. It is not a
counterexample to positivity or a proof that a generic remaining
source is nonempty; this individual rational curve has its separate
inherited admission. A new beta-aware soft normalization or circuit
is required before that uniform C2 argument can be applied.


\section{\texorpdfstring{C1 with both $L$ and $\eta$ units}{C1 with both L and eta units}}
\label{component:p686209b87229}

\label{p686209b87229:sec:r7-c1-L-eta-unit}

\begin{theorem}
For the original C1 source with $x\in\Inf_{>0}$ and both
$L=x\zeta$ and $\eta$ positive bounded units, the entire original
parent is tensor-SOS for every independent projective testing pair.
All original positive openings and literal $\beta=0$ are included.
\end{theorem}

The exact source formulas force $D$ strong on this subdomain, so
$D=0$ is absent by source implication.  No opening is removed by
an extra hypothesis, and the proof does not invoke a table or atlas.

\paragraph{The source and the only testing strip left by four events.}
Use the exact shared geometry (LC1), and write $\rho=r_n$ for the
native shape parameter and $r=1-\theta$ for the geometric one.
Thus
\[
 B=L,\quad y=B\eta,\quad H=B\rho=B-yt,\quad
 t=(1-\rho)/\eta,\quad j=1+zr-t,
 \quad z=x\in\Inf_{>0}.
\]
Here $B,y,\eta$ are units.  For the boundary coefficient formulas
below set
\[
 P=1+B+y,\qquad P_4=1+B(1-\rho),\qquad S=H+y=P-P_4.
                                                               \tag{LU1}
\]
All three quantities are positive units, uniformly through
$\rho\to0$, $1-\rho\to0$, $j\to0$, and $1-H\to0$.
The native formulas imply that $D,C,A,\alpha$ are strong.
One forward cyclic relabeling sends these four consecutive events
to $C,A,\alpha,m$ and sends testing coordinates to
$(u',v')=(v,1/u)$.  The full-parent four-consecutive theorem
(CE11)--(CE16) therefore covers every test except
\[
                    u\in\Inf,\qquad v\text{ a bounded nonzero unit}.
                                                               \tag{LU2}
\]
The cycle is the inverse of (CE9)--(CE10), so its literal endpoint
scope is the original homogeneous parent's scope.

It remains to prove (LU2).  All roots and coefficients used below
are valid for both testing signs, including literal $u=0$.

\paragraph{A bounded full-parent remainder with every opening retained.}
Use the graph, harmonic roots and gauge (LII2)--(LII8), with
$k=\sqrt y$ now a source unit.  The bounded source data are
\[
 c=Cyt,\quad d_p=Dyj,\quad
 \mu=m/z^2,\quad \sigma=\alpha(1-H),\quad
 \tau=\beta(1-r),\quad A_H=AH.
\]
They are bounded; $c,d_p,\mu,\sigma,A_H$ are units.  In residue,
with the boundary symbols of (LU1), their exact values are
\[
 c=\frac S{2P},\quad d_p=\frac{BP_4}{2P},\quad
 \mu=\frac{Sr(1+y)}{2B^2P},\quad
 \sigma=\frac{BS}{2P},\quad A_H=\frac{yP_4}{2P},\quad
 \tau=\frac S{2P}-\frac{1-r}{2}.                         \tag{LU3}
\]
These are residue identities, not substitutions of source and
testing fields.  They are obtained after multiplying away every
possibly small opening in the exact native formulas.
Since $\tau\ge0$, $r\ge P_4/P$ in residue and is a positive unit.
Also $t+j\sim1$, so $\ell=y(C+D)$ has a positive lower bound,
$q=D/(C+D)$ is bounded in $[0,1]$, and
$g=yCD/(C+D)$ is a positive unit.

Choose positive bounded caps $\ell_*\le\ell$,
$A_*\le A$, and $\alpha_*\le\alpha$.  Choose a bounded cap
$0\le\beta_*\le\beta$ as specified below.  Retain the five
contact-centered positive packets
\[
\begin{gathered}
 4(\ell-\ell_*)N^2,\qquad 4CV\eta_t^2,\qquad
 4(A-A_*)D_uD_v\eta_t^2,\qquad
 4(\alpha-\alpha_*)D_u(v\eta_t-\zeta)^2,\\
 4(\beta-\beta_*)D_v(\xi+u\zeta+R/B)^2,
 \qquad \eta_t=k(N+qR_h),\quad R=v(U-z)W.
\end{gathered}
\tag{LU4}
\]
The exact remainder is (LII6), with its $n$ line deleted, its $A$
packet replaced by
\[
 4A_*D_uD_v\eta_t^2
 +\frac{8A_H}B D_uD_vuW\eta_t
 +\frac{4A_HH}{B^2}D_uD_vUW^2,
                                                               \tag{LU5}
\]
and with (LC4) plus $4\beta_*D_v(\xi+u\zeta+R/B)^2$ added.
This specifies every coefficient of the full remainder.  Only
bounded source products occur, even at an unbounded original weight.

The same dependent-root gauge and exact contact equations give
(LII9)--(LII13).  The new constant scalar coefficients are
\[
 \frac{d_s+e_sV}{z^2}=
 \frac4{B^2}\left[
 \mu(B+zr)^2+
 \left\{\frac{r^2}{2}-\tau(1+r)+\beta_*\right\}(1+V)
 \right].                                                   \tag{LU6}
\]
The support $P_\xi\in(U,z)$ and scalar degree at most two in $U$
remain exact.  Apply all of (LII15).  As $v$ is a unit, eliminate
$T$ by the exact pure-factor identity
\[
                        T=(uk/v)R.                         \tag{LU7}
\]
Its coefficient is a bounded infinitesimal tensor.  Thus the entire
remainder has a bounded Gram on
$(N,\zeta,R_h,\xi,R)$.  No source/testing ratio or mixed inverse
has been introduced.  We now verify its complete residue.

\paragraph{A three-square identity supplies all three state/angle normals.}
The following boundary computation is an exact rational identity
before taking any opening limit.  Define
\[
\begin{aligned}
 C_0&=\frac{\rho+\eta}{2(1-\rho)P},&
 A_0&=\frac{\eta P_4}{2\rho P},\\
 D_0&=\frac{P_4}{2(\eta-1+\rho)P},&
 \alpha_0&=\frac{B^2(\rho+\eta)}{2(1-B\rho)P}.
\end{aligned}
\]
These may have unbounded openings; only the final capped identity
below is used as a bounded Gram.  For an indeterminate $p$, put
\[
 e=Y-\rho p,\qquad f=Z-vp/B,\qquad T_v=(V-y)p.
\]
The entire leading three-state form is exactly
\[
\begin{split}
 &4C_0D_v(p-Y)^2+4A_0D_vY^2
 +(4D_0+2)(p-Y-vZ)^2+(4\alpha_0+2)(vY-Z)^2-2D_vZ^2\\
 &\quad=\frac{4D_0}{B^2}(T_v+Be+Bvf)^2
       +4\alpha_0(f-ve)^2+4(C_0+A_0+\tfrac12)D_ve^2.
\end{split}                                                    \tag{LU8}
\]
This identity holds for every $v$, not only near the actual Bob
angle.  In the contact allocation at $z=u=0$,
$e=k(N+qR_h)$ and $T_v=-B(kR_h+v\zeta)$.
Subtract the residue contributions of the retained packets (LU4).
The harmonic identity then gives the complete three-normal block
\[
 \boxed{4\ell_*N^2+4gR_h^2
 +4(A_*+\tfrac12)D_vy(N+qR_h)^2
 +4\alpha_*\{\zeta-vk(N+qR_h)\}^2.}                     \tag{LU9}
\]
It has strictly positive residue on $(N,\zeta,R_h)$ because
$\ell_*,g,\alpha_*$ are positive units.  No normal has been
suppressed.  Every rational cancellation used to obtain (LU9)
has bounded source products on both sides, so the identity extends
continuously through every opening listed after (LU1).

\paragraph{The remaining two-root block and a finite cap.}
There is no finite residue coupling between these three roots and
$(\xi,R)$.  Such terms carry $u$, $z$, or the coefficient in (LU7).
For completeness, the exact first-normal support gives this block
without dropping a term.  To first order in $(U,z)$ write
$P_\xi=UA(V)+zZ(V)$; both coefficient functions are affine in $V$,
and
\[
 Z(V)=\frac{D_v}{B}(-4r+8\tau-8\beta_*).
\]
The exact contact equation gives $A(y)+Z(y)=0$ in residue.
After the first allocation of (LII15), the coefficient of $\xi R$
is therefore $A(V)-(V-y)(A'+Z')=-Z(V)$.
Using (LU3) and (LU6), the complete block is
\[
 M_{\xi R}=
 \begin{pmatrix}
 2D_v(1+2\beta_*)&\dfrac{2D_v}{B}(P_4/P+2\beta_*)\\[1mm]
 \dfrac{2D_v}{B}(P_4/P+2\beta_*)&
 \dfrac2{B^2}\{(P_4/P+2\beta_*)D_v-(S/P)r(V-y)\}
 \end{pmatrix}.                                           \tag{LU10}
\]
Here all letters denote real residues when positivity is assessed.

If $\beta$ is bounded, take $\beta_*=\beta$.  The native relation
$(1-r)(2\beta+1)=S/P$ follows from (LU3), and (LU10) has determinant
\[
 \frac{4S D_v}{B^2P}(1+2\beta)r(1+y)>0.                 \tag{LU11}
\]
Its first diagonal entry is positive, including $\beta=0$.
If $\beta$ is unbounded, take $\beta_*$ to be a sufficiently large
positive real cap.  Increasing it adds the positive square
$4\beta_*D_v(\xi+R/B)^2$.  On its kernel $\xi=-R/B$,
the uncapped block equals
\[
          \frac{2S}{B^2P}\{1+ry+(1-r)V\}R^2>0.           \tag{LU12}
\]
The elementary real two-by-two Schur argument therefore supplies
a finite cap making (LU10) strictly positive.  Unboundedness of the
original $\beta$ permits that cap.

Together (LU9)--(LU12) prove strict positive residue of the full
five-root Gram.  The bounded Gram lemma gives its tensor SOS;
restoring (LU4) proves (LU2).  The cyclic four-event complement
completes every testing chart and original projective endpoint.

\paragraph{Verification scope.}
The checker verifies the exact three-square state identity, all three
capped normals, both explicit two-root positivity formulas, the scalar
coefficient and support assertions, and prints the full extracted
first-normal and scalar polynomials.  The exact source contact and
monic contact divisions are the unchanged full-parent argument of
(LII7)--(LII15), with the bounded modifications (LU4)--(LU6).
This theorem closes the $L$-unit, $\eta$-unit subdomain of C1.4;
the $\eta$-unbounded subdomain is not included.


\section{\texorpdfstring{An interior small-$L$ C1 source: a full three-circuit closure}{An interior small-L C1 source: a full three-circuit closure}}
\label{component:pe3ef590c01d6}
% Edited scoped mathematical packet, 22 October 2026.

\label{pe3ef590c01d6:sec:r9-c1-small-interior}

Use the normalized native C1 parent, $d=1$, with $x=k^2>0$
infinitesimal, and set
\[
 \rho=\eta+r,\quad y=x\eta\zeta,\quad L=x\zeta,\quad
 h=1-\theta,\quad s=\sqrt{\eta\zeta},\quad
 \lambda=\rho/(2\eta).
\]
The reference interior has
\[
 \eta,\zeta,r,1-r,j_0=\eta+r-1
       \quad\text{positive bounded units}.                    \tag{IC1}
\]
The formulas below are the exact scalar and full-normal dictionary.
The opening-uniform RO proof supplies their normalized bounds when
$r$ or $1-r$ tends to zero.  The displayed reference $X$ normal
uses $\sqrt\beta$ when $\beta>0$; RO replaces it throughout by
$\sqrt{\beta+1/2}$, which includes literal $\beta=0$.
The full scalar graph coefficients $A_w,A_z,b_0,b_1$ are printed
explicitly in the native dictionary SN1--SN2.
For IC4--IC21, first assume $Z\ne0$ and normalize $Z=1$ inside
the testing field.  Homogenizing the resulting state-linear roots
restores every $Z$; the literal $Z=0$ state has the direct core
certificate.  The reciprocal identities below are already homogeneous.
\paragraph{A bounded-ratio frame retaining all three state normals.}
First suppose the pure-testing ratio $t=v/u$ is bounded, possibly
infinitesimal. Define
\[
 \begin{gathered}
 E=\eta uW-v,\quad
 b={r(1-r)\over\eta(\eta+2r-1)},\quad
 N=Y-ruW-bE,\\
 \xi=X+hu,\qquad N_x=\sqrt\beta\,\xi,\quad h=1-\theta.
                                                               \end{gathered} \tag{IC4}
\]
The state change is exact and uses only the pure-testing inverse $u$:
$W=(v+E)/(\eta u)$, $Y=r(v+E)/\eta+bE+N$, $X=-hu+\xi$.
Every normal vanishes at the actual source contact. 

\paragraph{The exact scalar, with the negative defect absorbed.}
Set $U=u^2,V=v^2$, $s=\sqrt{\eta\zeta}$, $k=\sqrt x$, and introduce
the positive separated roots
\[
 A_t=uv,\quad B_t={k v^2\over s u},\quad
 C_t=xs,\quad G=kv,\qquad F_v=V-y.
                                                               \tag{IC6}
\]
They satisfy $A_tB_tC_t=G^3$ and coincide at actual contact
$(u,t)=(k,s)$. The pure-testing pole in $B_t$ is inside the root;
none of the coefficient arrays below has a negative power of $t$.

Use $A_w,A_z,b_0,b_1$ from (SN1)--(SN2), including the beta term.
The scalar at $E=N=\xi=0$ is exactly
\[
 \mathcal S={VA_w\over\eta^2U}+A_z
                  -{2V\over\eta}(b_0+b_1U)
       =c_FF_v^2+f_AA_t^2+f_BB_t^2+f_CC_t^2-f_GG^2.    \tag{IC7}
\]
Here the complete coefficients are
\[
\begin{aligned}
 c_F={}&{C(1-r)^2+Ar^2+(\alpha+\tfrac12)r^2+2mr\over\eta^2},\\
 a_{UV}={}&{Ar^2\over\eta^2}+{h^2\over2}+2Dh+\beta\theta^2
                                     -{2b_1\over\eta},\\
 a_{UUV}={}&h^2(D+\tfrac12),\qquad
 a_{UVV}={r^2(m+A+\alpha+\tfrac12)\over\eta^2},\\
 q_V={}&{C(1-r)^2+Ar^2+(D+\tfrac12)(1-r)^2\over\eta^2}
                                     +D-{2b_0\over\eta},\\
 q_U={}&\alpha+(m+\beta+\tfrac12)\theta^2,\\
 f_A={}&a_{UV}+a_{UUV}U+a_{UVV}V,\qquad
 f_B={m\zeta\over\eta x},\\
 f_C={}&{\alpha-c_Fy^2+q_UU\over\eta\zeta x^2},\qquad
 f_G=-{q_V+2c_Fy\over x}.                              \end{aligned} \tag{IC8}
\]

\paragraph{Finite scalar Taylor arrays.}
Let $\epsilon_i=f_i-\lambda$ and $\epsilon_G=f_G-3\lambda$,
$h_c=\epsilon_A+\epsilon_B+\epsilon_C-\epsilon_G$.
Equations (IC8) show that $h_c=h_{00}+aU+cV$ with bounded
pure-source $a,c$. Exact contact value gives
\[
 h_c=a(u^2-k^2)+c(u^2t^2-k^2s^2).                     \tag{IC11}
\]
For $d_u=u-k,d_t=t-s$, its finite second Taylor identity is
\[
\begin{split}
 h_c={}&2k(a+cs^2)d_u+2k^2sc\,d_t
 +(a+ct^2)d_u^2+2kc(t+s)d_ud_t+k^2c\,d_t^2.            \end{split} \tag{IC12}
\]
The coefficient of $d_t^2$ is infinitesimal. For any affine function
$f=f_{00}+a_fU+c_fV$ among (IC8), the exact first difference is
\[
 f-f(k,s)=(u-k)(u+k)(a_f+c_ft^2)
                    +(t-s)k^2(t+s)c_f.                \tag{IC13}
\]
These are the required bounded $u$ arrays and infinitesimal $t$ arrays;
no appeal to unspecified Taylor membership is needed.
The two exact first-jet balances are
\[
 2(\epsilon_{A0}-\epsilon_{C0})+k h_{c,u0}=0,
 \quad 2(\epsilon_{B0}-\epsilon_{C0})+s h_{c,t0}=0.      \tag{IC14}
\]
They follow from the original stationary contact, and are also checked
directly using (NC4). The extra square $c_FF_v^2$ has zero value and
first jets there. Apply the accompanying three-circuit contact-lift
lemma with the exact conductors
\[
 G(u-k)=k(A_t-G),\qquad G(t-s)=s(B_t-G).                \tag{IC15}
\]

\paragraph{Explicit finite allocations of every first normal.}
Use $f_w,f_z,k_w,k_z$ in (SN2), and define the polynomials
\[
\begin{aligned}
 P_Y(U,V)&=2(k_w/\eta+k_z),\\
 P_E(U,V)&={2\over\eta^2}(A_w/U-mV/U)
              -{2\over\eta}(b_0+b_1U)+bP_Y(U,V).
\end{aligned}                                                   \tag{IC16}
\]
The quotient $A_w/U-mV/U$ is polynomial, as its displayed numerator
shows. For either polynomial $P$ use the exact divisions
\[
 P=P_{00}+UP_U+VP_V,\quad
 P_U={P-P(0,V)\over U},\quad
 P_V={P(0,V)-P(0,0)\over V}.                           \tag{IC17}
\]
The normalized bounds on these exact arrays are supplied by RO. The complete first normals
of the original $F/4$ in the frame (IC4) are exactly
\[
\begin{aligned}
 \mathcal L_N&=vP_Y,\\
 \mathcal L_E&=2mv t^2/\eta^2+vP_E,\\
 \mathcal L_x&={1\over\sqrt\beta}\left[
 2u\theta(m+\beta+\tfrac12)
 -{2m\over\eta}ut^2(1+rU)+2uVZ_X\right],\\
 Z_X&=(\beta+\tfrac12)\theta-(D+\tfrac12)(1+hU)
                         +(D+\tfrac12)(1-r)/\eta .     \end{aligned} \tag{IC18}
\]
These identities follow by differentiating the entire original parent
after the full state substitution (IC4).

Their actual root allocations are
\[
\begin{aligned}
 vP&={P_{00}\over k}G+u(P_U+t^2P_V)A_t,\\
 {2m\over\eta^2}vt^2&={2ms t\over\eta^2k}B_t,\\
 \mathcal L_x&={2u\theta(m+\beta+\tfrac12)
                    \over\sqrt\beta\,xs}C_t
 -{2ms(1+rU)\over\eta k\sqrt\beta}B_t
 +{2vZ_X\over\sqrt\beta}A_t.                          \end{aligned} \tag{IC19}
\]

The exact original contact makes each complete normal zero at
$(u,t)=(k,s)$. For a normal $\ell=\sum b_iA_i+b_GG$, put
$B=\sum b_i+b_G$, and use the explicit monic divisions
\[
 B=(u-k){B(u,t)-B(k,t)\over u-k}
       +(t-s){B(k,t)-B(k,s)\over t-s}.                 \tag{IC21}
\]

\paragraph{Reciprocal testing-ratio identities.}
Return to homogeneous state coordinates and set
\[
 p=uW,\ q=vZ,\ E=\eta p-q,\quad
 Y=rp+bE+N,\quad B=X+uZ=\xi+\theta uZ.
\]
Use the complete raw root list
\[
 E,\ N,\ N_x=\sqrt\beta\xi,\quad
 P=\sqrt m\,vW,\quad Q=\sqrt\alpha Z,\quad
 R_t=u^2W,\quad S_t=uvW.                              \tag{IC22}
\]
\[
 \sqrt x\,p=\sqrt{x/m}\,(u/v)P,\quad
 Y=rp+bE+N,\quad q=\eta p-E,                          \tag{IC23}
\]
For completeness, with $B=0$ the exact higher-angle part of the
$C,A,D$ packets and half core, after the quadratic base is removed, is
\[
\begin{split}
 &U\{AY^2-2Dq(p-Y-q)+q^2/2-pq\}\\
 &\quad+V\{C(p-Y)^2+(A+\tfrac12)Y^2\}
 +(A+\tfrac12)UVY^2+(D+\tfrac12)U^2q^2.               \end{split} \tag{IC24}
\]
The terms involving $B$ are equally explicit. Relative to $B=0$, the
$D$ packet adds $2DuvB(p-Y-D_uq)+DUVB^2$, and the half core adds
$uvB(p-Y-D_uq)+(D_v+UV)B^2/2$.
Here
\[
 B={N_x\over\sqrt\beta}+{\theta u\over\sqrt\alpha}Q,
 \quad uvp=uS_t,
\]
The opening-specific proof bounds their coefficients.  Write
$M=vW-X+uvY-uZ$. The three remaining
original packets keep their entire roots:
\[
\begin{aligned}
 \sqrt m\,M&=P+\sqrt m\,uvY-\sqrt m\,B,\\
 \sqrt\alpha\sqrt{D_u}(vY-Z)
  &=\sqrt{D_u}\{\sqrt\alpha(rS_t+v(bE+N))-Q\},\\
 \sqrt\beta\sqrt{D_v}B
  &=\sqrt{D_v}\{N_x+\sqrt{\beta\theta^2/\alpha}\,uQ\}.
\end{aligned}                                                   \tag{IC26}
\]


\section{\texorpdfstring{Small $L$: a complete native normal identity and the exact soft defect}{Small L: a complete native normal identity and the exact soft defect}}
\label{component:pc15b7caff0a3}
% Edited scoped mathematical packet, 22 October 2026.

\label{pc15b7caff0a3:sec:r8-small-L-native-normal}

Normalize $d=1$, use the native C1 source $\lambda=n=0$, and keep
all other original event weights, including literal $\beta=0$.
Put $U=u^2,V=v^2,D_u=1+U,D_v=1+V$, $\theta=1-h$, and
\[
             X=-huZ+\xi,\qquad Y=ruW+\nu.
\]
Here $\nu$ denotes this state normal only.  Direct expansion of the
original event packets gives
\[
 \frac{F_g}{4}=A_wW^2+A_zZ^2-2uv(b_0+b_1U)WZ,
\]
with the complete graph coefficients
\[
\begin{aligned}
 A_w={}&[(C+D+\tfrac12)(1-r)^2+Ar^2]U+Ar^2U^2\\
 &+\{m+[2mr+C(1-r)^2+(A+\alpha+\tfrac12)r^2]U
                  +r^2(m+A+\alpha+\tfrac12)U^2\}V,\\
 A_z={}&\alpha+[\alpha+(m+\beta+\tfrac12)\theta^2]U\\
 &+\{D+[2Dh+\tfrac12h^2+\beta\theta^2]U
                     +(D+\tfrac12)h^2U^2\}V,\\
 b_0={}&m\theta+D(1-r)+\alpha r+\tfrac12,\\
 b_1={}&mr\theta+(D+\tfrac12)h(1-r)+(\alpha+\tfrac12)r.
\end{aligned}
\]
These identities are homogeneous in $W,Z$ and need no state-ratio
bound.  The first-normal and hard-block formulas that follow are
obtained by the same substitution, or by differentiating the original
parent twice in $X,Y$.
The entire original parent is
exactly
\[
\begin{split}
 F/4={}&F_g/4+2\xi(vWf_w+uZf_z)+2\nu(uWk_w+vZk_z)\\
 &+H_{xx}\xi^2+2H_{xy}\xi\nu+H_{yy}\nu^2.
\end{split}                                               \tag{SN1}
\]
where
\[
\begin{aligned}
 f_w={}&-m(1+rU)+(D+\tfrac12)(1-r)U,\\
 f_z={}&m\theta-(D+\tfrac12)V(1+hU)
                    +(\beta+\tfrac12)D_v\theta,\\
 k_w={}&mV(1+rU)-CD_v(1-r)+AD_uD_vr
             -(D+\tfrac12)(1-r)+(\alpha+\tfrac12)D_urV,\\
 k_z={}&-m\theta U+(D+\tfrac12)(1+hU)
                              -(\alpha+\tfrac12)D_u,\\
 H_{xx}={}&m+\beta D_v+(D+\tfrac12)UV+D_v/2,\\
 H_{yy}={}&mUV+CD_v+AD_uD_v+D+\tfrac12
                              +(\alpha+\tfrac12)D_uV,\\
 H_{xy}={}&-(m+D+\tfrac12)uv.
\end{aligned}                                                   \tag{SN2}
\]
There is no discarded normal or excess packet in this identity.
Its two exact first-normal contact equations are
\[
 \zeta f_w(x,y)+f_z(x,y)=0,\qquad
 k_w(x,y)+\eta k_z(x,y)=0,\qquad y=x\eta\zeta.             \tag{SN3}
\]
They follow by substituting (NC4), not by setting a testing angle
equal to a source angle in the desired tensor assertion.

On the small-$L$, bounded-$\eta$ source, with $x,L$ infinitesimal
and $\zeta$ noninfinitesimal, the beta-opening identity gives:
\[
 2\beta\theta P_d=K(r+\eta)(xh+L)-\theta P_d,
 \qquad \beta\theta=O(L),\quad \beta\theta^2=O(L^2).       \tag{SN4}
\]
In particular, after the pure-source normalization
$\widehat\xi=\sqrt{1+\beta}\,\xi$, its cross coefficient is
$\beta\theta/\sqrt{1+\beta}=O(L)$, including every unbounded
beta opening.  The complete normal allocations are the retained
RO/JO/HO and DG/LD/CJ packets.


\section{\texorpdfstring{Beta-positive extension and the exact remaining local source branch}{Beta-positive extension and the exact remaining local source branch}}
\label{component:p909aa5765721}

\label{p909aa5765721:sec:r13-beta-positive-extension}

Retain the same original C14.1 source as JR1--JR8, with $x,z,\sigma$
positive infinitesimals, but allow $\beta\ge0$. Write
\[
 b_\theta=\beta\theta,\qquad H=h/\sigma,\qquad
 M_\beta=b_\theta/\sigma^2.
\]
This section proves two precise extensions. First, if $M_\beta$ is
bounded, the \emph{entire source domain is admitted for every test}.
This includes the beta-zero theorem. Second, if $H$ is bounded but
$M_\beta$ is unbounded, every bounded nonzero-unit $u$ is admitted
at the residual infinitesimal-$w$ tests. On this second source domain
only $u$ infinitesimal remains. Neither statement assumes $d_0$ a
unit: the second includes $D=0$. The $H$-unbounded domain and C14.2
are not admitted here.

\paragraph{The generic joint radial identity.}
The exact six-root identity JR5 continues to hold if its coefficients
are given the following beta corrections:
\begin{gather*}
 k_1\ \mapsto\ k_1+4b_\theta\theta,\qquad
 R_x\ \mapsto\ R_x+8b_\theta,\\
 Z_0\ \mapsto\ Z_0+\frac{4b_\theta\theta t}{\sigma^2},\qquad
 Z_x\ \mapsto\ Z_x+8ub_\theta/\sigma,\\
 S\ \mapsto\ S+4b_\theta\theta/\sigma^2.
\tag{JB1}
\end{gather*}
Here the uncorrected formulas on the right are the polynomials in
JR2--JR3 evaluated at the actual $h,D,m,p,l$, not at a beta-zero
source. The scalar $B$ remains its full monic quotient. The additional
retained centered excess is $4\beta D_wN_x^2$. The complete generic
coefficient is, explicitly,
\[
 S=\frac{2\{h^2(2D+1)(t+x)+4Dh-2bp\sigma^2z
                       +2b_\theta\theta+h^2\}}{\sigma^2}.
\tag{JB2}
\]
All four cancellations JR4 are exact with these changes. They were
checked symbolically with $b_\theta$ free. Thus no radial pole remains
unallocated in this extension.

\paragraph{The source estimates needed by the extension.}
The source equation BE3 gives exactly
\[
 H=H_0/B+\frac{2P}{B}\frac{b_\theta}{\sigma}.
\tag{JB3}
\]
Consequently $M_\beta$ bounded implies $H\sim1$ and $\theta\sim1$.
The sphere budget then implies $D/\sigma^2\sim1/2$, since
$xb_\theta/\sigma^2=xM_\beta$ is infinitesimal. Also $\mu\sim1/2$.

More generally, suppose only $H$ bounded. Then $\res H\ge1$;
no exact inequality $H\ge1$ is imposed. One has $\theta\sim1$,
and $b_\theta/\sigma$ is bounded. Write
\[
 d=\res(D/\sigma^2),\qquad
 c=\res(xb_\theta\theta/\sigma^2).
\]
BE4--BE6 give $d,c\ge0$ and $d+c=1/2$. Formula JB3 gives
$\res(b_\theta/\sigma)=(\res H-1)/2$. The bounded source formulas
also give $\mu\sim H/2$, $p,l\sim1/2$, $C=O(z)$, and
$p-l\in z\mathcal O_{\rm source}$. These statements retain both
source-sphere endpoints and both $b,Q$ openings.

\paragraph{An exact square joining the beta pole to its angular reserve.}
Suppose first $u$ is a bounded nonzero unit, so that $1/u$ is an
allowed bounded pure-testing inverse. Put
\[
 a_\beta=b_\theta\theta,\qquad
 C_\beta=4a_\beta/\sigma^2,\qquad T_c=t\Gamma+V,
 \qquad P_\beta=uZ-(x/\sigma)wW.
\]
The identity $T_c=\sigma uP_\beta$ has no $u$ denominator.
Adding $(C_\beta/t)$ times JR6's zero, and completing the resulting
$T_c$ square, gives the exact entire-parent identity
\[
 G=4a_\beta P_\beta^2+\widehat G,
\tag{JB4}
\]
where $G$ is the complete generic JR5 and
\begin{align*}
 \widehat G=G
 &-C_\beta t\Gamma^2
 -(C_\beta+C_\beta x/t)\Gamma V
 -C_\beta\sigma A_\theta^2\\
 &-C_\beta\sigma(t-x)A_\theta R/u.
\tag{JB5}
\end{align*}
Every coefficient of $\widehat G$ is bounded when $H$ is bounded
and $u$ a unit. The only apparent angular reserve subtraction is
$4a_\beta/\sigma$. Its exact source comparison is
\[
 \mu-a_\beta/\sigma
 =\kappa H_0/B+
    \frac{b_\theta}{\sigma}\left(\frac{2\kappa P}{B}-\theta\right),
\tag{JB6}
\]
where $\kappa$ is the BE8 source factor. The bracket is infinitesimal,
so this difference has residue $1/2$. Thus the beta pole is paid for
by the actual shared $m$ packet. No leading coefficient is substituted
in the exact identity JB4.

\paragraph{Bounded $M_\beta$: all residual tests.}
When $M_\beta$ is bounded and $u$ a unit, the residue of $\widehat G$
is exactly the beta-zero residue in JR7 before the finite $V$ pivot:
$H\sim1$, $d=1/2$, $c=0$. Add the same pure-testing pivot
$\lambda(t)=2(t^2+t+1)/(1+t)^2$ and apply JR8. Restore the positive
square in JB4 and every centered excess. This covers every unit $u$.

When $u$ is infinitesimal, use the generic six-root form directly,
without $1/u$. Its residue has five positive diagonal pivots equal to
two, with the only cross term
\[
 (2+4M_0)\Gamma V,\qquad M_0=\res M_\beta\ge0.
\]
Add JR6's zero with the bounded pure-source coefficient
$\lambda=(1+2M_\beta)^2$. All its non-$V^2$ corrections have zero
residue. The $(\Gamma,V)$ Schur complement is
$(1+2M_0)^2/2>0$. This includes literal $u=0$ and either sign.
The exact remote join is the CE9--CE10 and CE20--CE21 join printed
after JR8; it applies unchanged here because $D,\beta$ were unused.
This proves the full all-tests admission under the sole extra
pure-source condition $\beta\theta/\sigma^2$ bounded.

\paragraph{Unbounded $M_\beta$, bounded $H$: all unit-Alice tests.}
For any real nonzero residue $u_0$, set $t_0=u_0^2$. In the notation
of JB3, and writing $H$ now for its finite real residue, the only
nontrivial $\widehat G$ residue block is
\begin{align*}
 Q={}&2(1+t_0)N_x^2+[4d+2H^2t_0(1+t_0)]\Gamma^2\\
 &-4u_0(1+Ht_0)N_x\Gamma-4Hu_0N_xV\\
 &+[2H^2(1+t_0)-4c/t_0]\Gamma V.
\tag{JB7}
\end{align*}
Its angular $(A_\theta,R)$ block is the strict matrix from JR7 and
its $N_y$ pivot is $2(1+t_0)$. All other residue cross entries vanish.
The beta square in JB4 equals $(C_\beta/t)T_c^2$, with
$C_\beta$ unbounded. On its exact residue kernel $V=-t_0\Gamma$,
JB7 reduces to
\[
 2(1+t_0)N_x^2-4u_0N_x\Gamma+2\Gamma^2,
 \qquad \det=4.
\tag{JB8}
\]
Here the cancellation uses exactly $d+c=1/2$. It remains strict at
$d=0$ and $c=0$. The row $T_c=V+t\Gamma$ has a bounded separated
right inverse, obtained by changing $V$ alone. The elementary finite
penalty lemma therefore chooses a positive real $K$ such that
$\widehat G+(K/t)T_c^2$ has strictly positive residue. This is possible
because $t_0>0$. The exact retained excess is
$((C_\beta-K)/t)T_c^2$: $C_\beta-K$ is a positive pure-source
coefficient, and its square root separates from the pure-testing
$1/\sqrt t$. No ordering of $C_\beta/t$ minus a real number is used.
This proves the entire
parent for all bounded nonzero-unit $u$, with infinitesimal or zero
$w$, all signs and complex states.

The new exact remaining obstruction on the bounded-$H$ source branch
is therefore $M_\beta$ unbounded and $u$ infinitesimal. JB5 has
$x/u$ and $x/u^2$ terms there, so its bounded-Gram argument cannot be
imported. The square $P_\beta$ itself is still legal and has no
$u$ inverse; a joint allocation with the remaining angular terms is
required to remove those coefficient denominators without comparing
independent source and testing rates. No such allocation is claimed.

\paragraph{Denominator and boundedness audit.}
Every source inverse in the roots or exact coefficients is a pure-source
inverse. On the large-$L$, bounded-$H$ domains, $P,J,B$ and the original
geometry coefficient $M=1+xbz-xz\theta$ are units; $1-r$ and $\theta$
are units as well. Possible $b,Q$ inverses occur only in the original
source weights and the retained positive centered excesses, while their
opening-cleared products $p,l$ are bounded. The normalized quantities
$H$, $D/\sigma^2$, $b_\theta/\sigma$, and
$xb_\theta\theta/\sigma^2$ are bounded by JB3 and the BE sphere.
The corrections in JB5 therefore leave bounded coefficients: the
remaining $\Gamma^2$ terms involve $D/\sigma^2$, $Dh/\sigma^2$ and
$H^2$; the $\Gamma V$ coefficient involves those same bounded source
products and $(xb_\theta\theta/\sigma^2)/t$; the angular-radial
coefficient involves $b_\theta\theta/\sigma$ times $(t-x)/u$.
The factors $1/u,1/t$ are used only when $u$ is a testing unit.
All $u=0$ and infinitesimal-$u$ assertions for bounded $M_\beta$
instead use the unaltered generic six-root identity and its bounded
pure-source pivot $(1+2M_\beta)^2$. The everywhere bounded testing
denominator $1+t=1+u^2$ in JR8 is also pure testing. No mixed inverse,
mixed radical, or source/testing rate comparison occurs in either
admitted domain or its remote join.


\section{\texorpdfstring{A complete joint radial allocation on the literal beta-zero face}{A complete joint radial allocation on the literal beta-zero face}}
\label{component:pb1104217f67b}

\label{pb1104217f67b:sec:r13-joint-radial-beta-zero}

\paragraph{Quantified domain and conclusion.}
Retain the original physical C1.4 source and all its residual-complement
conditions. Suppose $x=k^2$, $z=L^{-1}$ and $\sigma=\eta^{-1}$ are
positive infinitesimals, $0<b<1$, and the original coefficient $\beta=0$.
There are no further source hypotheses. In particular $h/\sigma\sim1$
and $D/\sigma^2\sim1/2$ follow from BE3--BE7, rather than being imposed.
The conclusion below covers every bounded independent real test $u$,
every infinitesimal or zero independent reciprocal test $w$, all their
signs, and real or complex states. The previously proved remote
complement supplies the other tests on this large-$L$ source chart.
Thus this proves the literal $\beta=0$ part of C14.1. The adjoining
Section~\ref{p909aa5765721:sec:r13-beta-positive-extension} extends the actual admission
to all sources with $\beta\theta/\sigma^2$ bounded; the present result
is its beta-zero corollary and supplies the residue calculation used
there. C14.2 is not included.

The calculation joins the entire radial scalar quotient with every
radial first normal \emph{before} deciding boundedness. Their exact
cancellations remove the BR3 pole, including its separate
$8(ab)b/\sigma$ summand through $Q=1-b$ infinitesimal.

\paragraph{Notation and the original contact identities.}
Use the complete reciprocal parent of BE10 and the large-$L$ source
notation. Thus $m=\mu x^2z$, $C=cz$, $A=a/z$, and
$\alpha=\ell/z$. In expressions below $p=ab$ and $l=\ell Q$ are
the two opening-cleared products; $p$ is not the source parameter $b$.
Set $t=u^2$, $\xi=w^2$, $\xi_0=\sigma z$, $\delta=\xi-\xi_0$,
$g=h-Q$, $D_u=1+t$, $D_w=1+\xi$, and $q=\sqrt z$.
The source identities established in the handoff give
\[
 \phi=\phi_t=\phi_\xi=L_x=L_y=L_z=0
 \quad\hbox{at }(t,\xi)=(x,\xi_0),
\]
as well as $\phi(0,\xi)=4m$ and $L_x(0,\xi)=-8m$.
Here $\phi,L_x,L_y,L_z$ are the \emph{complete} EN2 polynomials,
with $D_*=D$, $\beta_*=0$ and $\beta=0$.

Choose the fixed caps $a_* = \ell_* =1/4$. They are smaller than
the actual $a,\ell$, since $p,l\sim1/2$ and $0<b,Q<1$.
Retain the two exact nonnegative centered excesses
\[
 E_{\rm cap}=4(a-a_*)D_uD_wN_y^2+4(\ell-\ell_*)D_uN_y^2.
\]
Define six separated roots
\[
\begin{aligned}
N_x&=X+uhZ,&N_y&=(Y-bwZ)/q,\\
\Gamma&=\sigma Z-uwW,&R&=wZ/q-uqW,\\
A_\theta&=q(t-x)W,&V&=uw(t-x)W.
\end{aligned}\tag{JR1}
\]
No test/source quotient is used in their definition. Their inverses
and radicals are pure source.

\paragraph{Exact quotient specification.}
The contact identities give the exact divisions
\begin{align*}
 \phi&=z(t-x)^2A+t(t-x)\delta B+t\delta^2C_r,\\
 L_x&=(t-x)(tA_x+zB_x)+t\delta C_x,\\
 L_y&=(t-x)A_y+\delta B_y,\\
 L_z&=(t-x)A_z+\delta B_z.
\end{align*}
All quotients are monic polynomial divisions in one testing variable.
Their complete support can equivalently be specified by
\begin{gather*}
 A=4\mu+\frac{2h^2(2D+1)}\sigma t,\qquad
 A_x=-\frac{4h(2D+1)}\sigma,\qquad B_x=8\mu x,\\
 A_y=8(p-l)+4z\{(1+2m)g+2Dh+2p\sigma\},\\
 S:=A_z/\sigma-B=S_0+
       \frac{2h^2(2D+1)}{\sigma^2}(t-x).
\tag{JR2}
\end{gather*}
The same monic divisions give the much shorter exact expression
\[
 S_0=\frac{2\{2h^2(2D+1)x+4Dh-2bp\sigma^2z+h^2\}}{\sigma^2}.
\tag{JR2a}
\]
Thus $S_0\sim2$ follows immediately from $h/\sigma\sim1$ and
$D/\sigma^2\sim1/2$. This removes any dependency on the longer R10
reduction modulo the original beta numerator. The source contact jets
are still needed for the quotient remainders to vanish.
For a denominator-free quotient specification, $C_r$ is obtained by
dividing $\phi$ twice by $\xi-\xi_0$ and then dividing the quotient
by $t$; $B$ is the quotient of
$t^{-1}\phi_\xi(t,\xi_0)$ by $t-x$; and $A_z$ is the quotient of
$L_z(t,\xi_0)$ by $t-x$. These divisions introduce no mixed inverse.

The remaining exact coefficients are most compactly written as
\begin{align*}
 K_r:=\sigma^2C_r={}&k_0+k_1t+k_2\xi,\\
 k_0={}&2b\{(2C+1)b+4D+4p\sigma
             +2(2C+2D+1)\sigma(bz-1)\},\\
 k_1={}&2(1+2m)g^2+8Dhb+8pb\sigma,\\
 k_2={}&2b^2(2C+2D+1),\\
 R_x:=\sigma C_x={}&-4\{2Db+(1+2m)g\},\\
 R_y:=\sigma B_y={}&4\{t[2Dh+2p\sigma+(1+2m)g]
      +b(2C+2D+1)(\xi+\sigma z)\\
 &\hspace{23mm} +(2C+1)(b-\sigma)+2D(1-\sigma)+2p\sigma\},\\
 E_r={}&2(C+D)(bz-1)+2p+bz-1,\qquad
 R_\gamma=-4(E_r+2pt),\\
 Z_0={}&\frac{4D+2t(4Dh+h^2)+2t^2h^2(2D+1)}{\sigma^2},\\
 Z_x={}&-\frac{4u}{\sigma}\{2D+h+th(2D+1)\}.
\tag{JR3}
\end{align*}
In particular all of these coefficients are bounded when $u$ is bounded
and $w$ infinitesimal. The $\sigma^{-1}$ and $\sigma^{-2}$ in $Z_0,Z_x$
are controlled by the \emph{proved pure-source} bounds on $h,D$.

\paragraph{The joint cancellation, including all first normals.}
The exact full normal block EN3--EN4 obeys
\begin{align*}
 H_z/\sigma^2+\xi C_r-(\xi/\sigma)B_z&=Z_0,\\
 B_z-2\sigma C_r&=bR_\gamma,\\
 \frac4\sigma\{2DT_{\rm test}+bD_w(1+2C)+tg(1+2m)\}-B_y&=R_\gamma,\\
 \frac{4u}{\sigma}\{-2DT_{\rm test}-hD_u-(1+2m)\xi g\}
                         -u\xi C_x&=Z_x.
\tag{JR4}
\end{align*}
Here $T_{\rm test}=1+th+b\xi$. These are polynomial identities after clearing
pure-source powers of $\sigma$. No coefficient is discarded: $K_r$
contains the full BR3 term, and the last three lines keep its radial
couplings. The cancellation explains why a scalar-only extraction of
BR1 followed by boundedness checking cannot see this certificate.

Put
\begin{align*}
 H_x&=2(1+t+\xi)+4Dt+4m\xi,\\
 H_y&=4D_u(a_*D_w+\ell_*)+4mtz+4CzD_w+(4D+2)\xi z+2D_uz,\\
 c_G&=4D_u\{pbD_w+lQ\},\qquad
 c_N=8D_u\{pD_w-l\}.
\end{align*}
The \emph{entire original parent}, without restriction to a state
subspace, is exactly $F_w=E_{\rm cap}+G$, where
\begin{align*}
 G={}&H_xN_x^2+H_yN_y^2+Z_0\Gamma^2
        +(c_G+zK_r)R^2+AA_\theta^2\\
 &-4uqw(2m+2D+1)N_xN_y+Z_xN_x\Gamma
       +qwR_\gamma N_y\Gamma+qwbR_\gamma R\Gamma\\
 &+uqwR_xN_xR+(c_N+zR_y)N_yR
       +qwB_xN_xA_\theta+uA_yN_yA_\theta\\
 &+uA_xN_xV+u\sigma B A_\theta R+S\Gamma V.
\tag{JR5}
\end{align*}
For example, the scalar radial conductor is
$t\delta^2W^2=\sigma^2(qR-w\Gamma/\sigma)^2$;
the radial $Z$ first normal is
$uw\delta W(\Gamma/\sigma)=qwR\Gamma-(\xi/\sigma)\Gamma^2$.
Their joint use gives the first two identities of JR4. The other
conductors are those of the verified large-$L$ proof with
$N_z=\Gamma/\sigma$. These identities prove JR5 exactly and include
every first normal and every hard quadratic coupling.

\paragraph{One bounded positive Gram for every bounded Alice test.}
All entries of JR5 are bounded separated tensors. Indeed
\[
 h/\sigma\sim1,\quad D/\sigma^2\sim1/2,\quad
 \mu\sim1/2,\quad C=O(z),\quad p,l\sim1/2,
 \quad p-l\in z\mathcal O_{\rm source},
\]
and $\sigma B=A_z-\sigma S\sim4$ uniformly for bounded $u$.
The displayed formula for $A_z$ obtained from $L_z$ has only the
bounded terms $h^2/\sigma$, $Dh/\sigma$, $p,l$, and bounded multiples
of $z$; hence this last assertion does not order any mixed rate.

For completeness, the possibly needed opening-uniform divisibility
$p-l\in z\mathcal O_{\rm source}$ follows directly from the exact
source formulas, not from a residue subtraction. Write $H=h/\sigma$.
The two numerators over their common unit denominator $2P$ are
\begin{align*}
 U_A={}&z(1+x\sigma H)(1-xH)+(1-bz)(1+x\sigma Hz),\\
 U_\ell={}&(1+\sigma bz)\big[xH(1+xbz)(1+\sigma z)
       +\{1+xbz-xz(1-\sigma H)\}(1-xH)\big].
\end{align*}
They are bounded polynomials in the source letters and both equal
one at $z=0$. Consequently $(U_A-U_\ell)/z$ is a bounded polynomial.
Neither numerator has a $b$ or $Q$ denominator. This proves the claimed
divisibility through both openings and is the source argument in EN6's
adjoining radial paragraph.

The exact syzygy
\[
 V^2=\sigma t A_\theta^2+u\sigma(t-x)A_\theta R-(t-x)V\Gamma
\tag{JR6}
\]
allows the addition of zero times the pure-testing function
\[
 \lambda(t)=\frac{2(t^2+t+1)}{(1+t)^2}>0.
\]
Specifically add
$\lambda[V^2-\sigma tA_\theta^2-u\sigma(t-x)A_\theta R+(t-x)V\Gamma]$
to JR5. There is no source/test denominator in $\lambda$.

At any real residue $u_0$, write $t_0=u_0^2$. The $N_y$ pivot is
$2(1+t_0)$. The $(A_\theta,R)$ block is
\[
 \begin{pmatrix}2&2u_0\\2u_0&2(1+t_0)\end{pmatrix},
 \qquad \det=4.
\]
The remaining $(N_x,\Gamma,V)$ block is
\[
 \begin{pmatrix}
 2(1+t_0)&-2u_0(1+t_0)&-2u_0\\
 -2u_0(1+t_0)&2(1+t_0+t_0^2)&1+t_0+\lambda(t_0)t_0/2\\
 -2u_0&1+t_0+\lambda(t_0)t_0/2&\lambda(t_0)
 \end{pmatrix}.
\tag{JR7}
\]
Its first two leading pivots are positive, with leading $2\times2$
determinant $4(1+t_0)$. The last Schur complement is exactly
\[
 \lambda(1+t_0^2)+\frac{2t_0^2}{1+t_0}
       -\frac12(1+t_0+\lambda t_0/2)^2
   =\frac{4t_0+3}{2(1+t_0)^4}>0.
\tag{JR8}
\]
Thus the residue Gram is strictly positive definite for every finite
real $u_0$, including zero and both signs. Crucially it does not depend
on the residue of $b$: the result remains strict through $b\to0$
and $Q\to0$. The finite bounded positive-Gram criterion already used
by the project then gives the exact legal tensor-SOS.

\paragraph{Boundaries and scope.}
JR1--JR6 are identities, including $u=0$ and $w=0$; no root is inverted
at a zero. Their sign coverage is literal. Apply the real identity to
the real and imaginary parts of the state for complex states.
The exact remote join is the common-parent theorem in
Section~\ref{component:p89fd3edb3edd},
Section~\texttt{sec:r6-c1-common-remote}. First, $A,\alpha$ are
unbounded on this source, so its adjacent-event theorem transported by
(CE9)--(CE10) covers every $v$ that is not projectively unbounded.
For projectively unbounded $v$ and projectively unbounded $u$, put
$p_t=1/u$ and apply the complete raw frame (CE20), whose strict
residue is (CE21). Its exact source premises (CE17), without the
unneeded $m$ lower bound in this $p_t$ chart, hold with $q=\sqrt z$:
$C\sim z/2$, $A\ge1/(4z)$, and $\alpha\ge1/(4z)$.
The $C$ lower bound may likewise be chosen as $z/4$. The unselected
$D,\beta$ events are retained. The identities include literal
$p_t=0$, $w=0$ and all signs. These two remote cases and JR1--JR8
exhaust all projective tests on the printed source face. No source
domain with $L$ a unit, or $\beta>0$, is admitted by this proof.

\paragraph{Corroboration.}
The accompanying focused checker verifies the complete joint polynomial
allocation, the four cancellations, every conductor, the short exact
$S$ coefficient, and the rational Schur-complement identity. The
original six contact identities are explicitly identified imported
results, not silently replaced by numerical residue checks.


\section{\texorpdfstring{Small-$L$, unbounded eta: uniform source weights}{Small-L, unbounded eta: uniform source weights}}
\label{component:p66922bb7f235}

\label{p66922bb7f235:sec:r13-small-L-eta-large-bounds}

This is a source theorem, not an additional all-tests admission.
Use the original physical C1 source with $x,L=x\zeta>0$
infinitesimal, $\zeta$ noninfinitesimal and $\eta$ unbounded.
Put $\sigma=1/\eta$, $S=\sigma+L$, and retain the original
$0<r,h<1$, $\theta=1-h$, $M>0$, and $D,\beta\ge0$.
The native gap gives $E=1-xh/\sigma\ge1/18$. Set
\[
 e={xh\over\sigma}=1-E,\quad
 P=\sigma P_d,\quad J=1+xh-\sigma(1-r),\quad
 H_0=1+L(1-r),\quad B=P+x(J-\sigma),
\]
\[
 \mathcal H=R H_0+(1-r)xh(2+xr)
 -e\{2+x(1+h+r\theta)\}+Ke^2,
 \quad R=1+xr,\quad K=1+x.                           \tag{ES1}
\]
No $E\sim1$ assumption is made here. The original source formulas
give the two exact identities
\[
 2\theta P\beta=hB-\sigma H_0,\qquad
 2P(JD+x\theta\beta)=\sigma^2\mathcal H.             \tag{ES2}
\]
They are the beta-shared identities with the reciprocal-$L$ factors
cleared. In particular their use does not require bounded $1/L$.

The positive denominator is exactly
\[
 P=\sigma(1+L+2xh+rxL+rx^2h)+KL+x^2h\theta.
                                                               \tag{ES3}
\]
Since $x/L=1/\zeta$ is bounded and $x^2h=x\sigma e$,
all terms beyond $\sigma+L$ in (ES3), divided by $S$, are
infinitesimal. Consequently
\[
 {P\over S}\sim1,\qquad J\sim1,\qquad
 {B\over S}-(1+x/S)\in\Inf,\qquad
 H_0\sim1,\qquad \mathcal H-E^2\in\Inf.             \tag{ES4}
\]
Here $x/S$ is bounded. Both $B/S$ and $\mathcal H$ are
positive bounded units, uniformly through all original openings;
for the latter use the actual bound $E\ge1/18$.

Define nonnegative pure-source sphere coordinates
\[
 d_0^2={2PJD\over\sigma^2\mathcal H},\qquad
 b_0^2={2Px\theta\beta\over\sigma^2\mathcal H}.
\]
Then
\[
 d_0^2+b_0^2=1,\qquad
 D={\sigma^2\mathcal H\over2PJ}d_0^2,\qquad
 x\theta\beta={\sigma^2\mathcal H\over2P}b_0^2.
                                                               \tag{ES5}
\]
Thus $D=O(\sigma^2/S)=O(\sigma)$ and
$x\theta\beta=O(\sigma^2/S)$, including either literal coefficient
boundary. The first identity of (ES2) becomes
\[
 eB=xH_0+\sigma\mathcal H b_0^2.                     \tag{ES6}
\]
Put $\kappa=x/S$ and $\nu=\sigma/S$. Multiplying
$e(B/S)=\kappa H_0+\nu\mathcal H b_0^2$ by $e$ and using
$e\nu=\kappa h$ gives
\[
 e^2(B/S)=\kappa\{eH_0+h\mathcal H b_0^2\}.
                                                               \tag{ES7}
\]
All factors in braces are bounded and nonnegative. Hence
$e^2=O(x/S)$, and (ES6) also gives $x/S=O(e)$.
In particular $\zeta$ unbounded implies $e$ infinitesimal,
since $x/S\le1/\zeta$. The unit-$\zeta$ cases may have
$e$ a positive bounded unit; (SN7) and (SN12) are not imported.

The other original coefficients have the opening-cleared bounds
\[
 {CS(1-r)\over E}\sim\tfrac12,\qquad
 {ASr\over E}\sim\tfrac12,\qquad
 m={K(1+\sigma r)eS\over2M\zeta P},                  \tag{ES8}
\]
\[
 \alpha={(1+\sigma r)\{LR eS+L^2ME\}\over2QP}
       =O(L),\qquad Q=1-Lr\sim1.                   \tag{ES9}
\]
For $A$ use $P_4=(1+xh)E+(1-r)(xh+L)\sim E$.
The bound on $\alpha$ follows because $0<M<R$ is bounded and
$L/S\le1$. Thus $C,A$ are unbounded while $D,\alpha$ are
infinitesimal on the entire stated source domain. No bound on $m$
is asserted through $M=0$, and no bound on beta is asserted through
$\theta=0$.

Finally put $\mu=K(1+\sigma r)S/(2M\zeta P)>0$. The
original $m$ packet has the exact shared-sphere split
\[
 m=m_0+m_1b_0^2,\qquad
 m_0={\mu xH_0\over B},\qquad
 m_1={\mu\sigma\mathcal H\over B}.                  \tag{ES10}
\]
The remaining source opening $M$ is displayed explicitly. These
identities provide valid weights for C14.4 and the boundedness step
missing from its earlier checklist. A full-parent allocation for every
residual infinitesimal-$u$ test is still required. The narrow checker
verifies (ES2), (ES3), (ES6)--(ES10) algebraically; the inequalities
above are the ordered-field argument. No mixed inverse, mixed order,
or source/testing rate comparison is used.


\section{\texorpdfstring{The actual-$j$ C1 opening: centering the whole divergent event}{The actual-j C1 opening: centering the whole divergent event}}
\label{component:p48da253a150f}

\label{p48da253a150f:sec:r13-c1-j-opening}

\begin{theorem}
Use the original physical native C1 source with $d=1$, $x=k^2>0$
infinitesimal, and
\[
 \eta,\ \zeta,\ r,\ 1-r\quad\hbox{positive bounded units},\qquad
 j=\eta(1+xh)-1+r\quad\hbox{positive infinitesimal}.       \tag{JO1}
\]
Then the complete original parent is tensor SOS for every independent
projective test and every state. Literal $\beta=0$ is retained.
There is no assumption that $j_0=\eta+r-1$ is positive: $j_0=0$
and a slightly negative $j_0$ are included. Simultaneous infinitesimal
$r$ and $j$ are outside this theorem.
\end{theorem}

All scalar and Gram identities in this section use the normalized
parent $F/4$; multiplication of every square root by $2$ restores $F$.

The divergence $D\asymp1/j$ must be centered by its full event,
including its testing-angle dependence. A normal centered only by its
quadratic limit would lose control when $j\ll x$. The following exact
frame removes that problem before any residue is taken.

\paragraph{Source bounds and contact-exact coordinates.}
Write $\theta=1-h=x\pi$, where $\pi>0$ is bounded by (SB1),
$B_x=\beta+1/2\ge1/2$, and $y=x\eta\zeta$. Put
\[
 R_0=1+j-\eta=r+\eta xh,\qquad
 g(U)=R_0-\eta hU,\qquad U=u^2,\ V=v^2.
                                                               \tag{JO2}
\]
Thus $g(x)=r$ exactly. The native source formulas, with
$r=1+j-\eta(1+xh)$, imply
\[
 \begin{gathered}
 C,A=O(1),\quad Dj=\tfrac12+O(x),\quad
 m/x\sim1/(2\zeta),\quad
 \alpha/x^2\sim\zeta(1+\zeta)/2,\\
 B_x\theta/x=K(\eta+r)(h+\zeta)/(2P_d),\qquad
 \theta/\sqrt\alpha,\quad B_x\theta^2/\alpha=O(1).
 \end{gathered}                                                \tag{JO3}
\]
In particular $D>0$, so its square root and inverse are pure source
and legitimate. Set
\[
 \begin{gathered}
 p=uW,\quad q=vZ,\quad E_t=\eta p-q,\qquad
 N_d=\sqrt D\,(J_t+jp),\\
 J_t=uW+uvX-Y-vZ,\qquad
 N_x=\sqrt{B_x}(X+huZ),\quad
 n=N_d/\sqrt D,\quad \xi=N_x/\sqrt{B_x}.
 \end{gathered}                                                \tag{JO4}
\]
The exact inverse state formulas are polynomial and homogeneous:
\[
 Y_0=(1+j)p-q-n,\qquad
 Y=Y_0-hUq+uv\xi,\qquad X=-huZ+\xi,\qquad
 J_t=n-jp.                                             \tag{JO5}
\]
All three normals vanish at the actual source contact. At their zero
graph, $q=\eta p$ and $Y=g(U)p$, so the original native contact
graph is recovered at $U=x$ without changing any source contact jet.

For later use introduce positive reference weights
\[
 \widetilde C={1+j\over2(\eta-j)},\quad
 \widetilde A={\eta\over2R_0},\quad
 \widetilde D={1\over2j}.
\]
Since $R_0,1-R_0,\eta$ are units, exact source numerator division gives
\[
 C-\widetilde C=O(x),\qquad A-\widetilde A=O(x),\qquad
 Dj-\tfrac12=O(x).                                    \tag{JO6}
\]
The last assertion is deliberately about $Dj$, not $D$. The canceled
denominators of (JO3),(JO6), at $x=j=0$, contain only the unit
factors $\eta,\zeta,1-\eta$ and constants. No comparison of $j/x$
has been made.

\paragraph{The complete bounded hard block.}
The reference quadratic base has the exact identity
\[
 \begin{aligned}
 &\widetilde C(p-Y_0)^2+\widetilde A Y_0^2+
 \widetilde D(p-Y_0-q)^2+\tfrac12(p-Y_0)^2-pq\\
 &=\widetilde\gamma E_t^2+
 \widetilde H\{(1-\widetilde b)E_t-n\}^2,               \\
 \widetilde H&={\eta(j+R_0)\over2R_0(1-R_0)j},\quad
 \widetilde b={R_0(1-R_0)\over\eta(j+R_0)},\quad
 \widetilde\gamma={1+j\over2\eta(j+R_0)} .
 \end{aligned} \tag{JO7}
\]
No $p$ term survives. Furthermore
\[
 1-\widetilde b={j(1+j)\over\eta(j+R_0)},\qquad
 \widetilde H/D\sim1.                                 \tag{JO8}
\]
On $(E_t,N_d)$ this is a bounded hard Gram with diagonal residue
$(1/(2\eta(1-\eta)),1)$ and vanishing off-diagonal residue.
The actual-base error has coefficients $O(x)$ in $(p,E_t,N_d)$.
For its apparently singular part use
\[
 (D-\widetilde D)(n-jp)^2,\quad
 {D-\widetilde D\over D}=O(x),\quad
 {(D-\widetilde D)j\over\sqrt D}=O(x\sqrt j),\quad
 (D-\widetilde D)j^2=O(xj).                           \tag{JO9}
\]
These are consequences of (JO3),(JO6) and hold for every relative
source rate. All other weights in the base are bounded.

For bounded $t=v/u$ at $u,v>0$ both infinitesimal, put $Z=1$.
Then $p=(v+E_t)/\eta$ in (JO5). The higher-angle terms have
bounded infinitesimal hard coefficients, and the $m$ hard coefficient
$m t^2$ is infinitesimal. Combining beta with its half-core part gives
the third hard pivot on $N_x$. Thus the full hard residue is
\[
                 \operatorname{diag}(1/(2\eta(1-\eta)),1,1).
                                                               \tag{JO10}
\]
The exact higher-angle identity printed below also proves these hard
bounds directly; no state normal is suppressed here.

\paragraph{An exact scalar three-circuit with bounded polynomial coefficients.}
Use $s=\sqrt{\eta\zeta}$ and the same separated roots
\[
 A_t=uv,\quad B_t={kv^2\over su},\quad C_t=xs,\quad
 G=kv,\quad F_v=V-y.
\]
On the zero graph of (JO4), the exact scalar is
\[
 \mathcal S=c_FF_v^2+f_A(U,V)A_t^2+f_BB_t^2+
                       f_C(U)C_t^2-f_GG^2,            \tag{JO11}
\]
where, with $g=g(U)$,
\[
 \begin{aligned}
 c_F&={2mR_0+C(1-R_0)^2+(A+\alpha+\tfrac12)R_0^2\over\eta^2},\\
 q_U&=\alpha+(m+B_x)\theta^2,\\
 q_V&={C(1-R_0)^2+AR_0^2+(D+\tfrac12)j^2\over\eta^2}
       -{2m\theta\over\eta}-{(2\alpha+1)R_0\over\eta}-\tfrac12,\\
 f_B&={m\zeta\over\eta x},\quad
 f_C={\alpha-c_Fy^2+q_UU\over\eta\zeta x^2},\quad
 f_G=-{q_V+2c_Fy\over x},                              \\
 f_A&={mV(g^2-2\eta h)\over\eta^2}-{2m\theta g\over\eta}
       +C(1+V)\{2h(1-R_0)/\eta+h^2U\}\\
 &\quad+\{A(1+V)+(\alpha+\tfrac12)V\}
                   \{g^2/\eta^2-2R_0h/\eta+h^2U\}\\
 &\quad-(2\alpha+1)(g/\eta-h)+B_x\theta^2-\tfrac12 .
 \end{aligned} \tag{JO12}
\]
In particular $f_A$ is an explicit bounded polynomial of degree at
most two in $U$ and one in $V$; the divergent $D$ is absent.
The only $D$ factor elsewhere is $Dj^2=O(j)$. These formulas
follow by substituting $W=v/(\eta u)$, $X=-hu$, $Y=g(U)v/\eta$
in the entire parent; the new checker differentiates that entire
substituted parent as well.

The source bounds and (JO12) give
\[
 c_F\sim1/(2\eta^2),\qquad
 f_A,f_B,f_C\sim\lambda=1/(2\eta).
\]
Exact scalar contact gives
$f_G=f_A(x,y)+f_B+f_C(x)$, proving that $f_G$ is bounded
and $f_G\sim3\lambda$ without estimating a singular quotient
$q_V/x$. Every nonconstant coefficient in $f_A,f_C$ is bounded;
for $q_U/x^2$ use (JO3).

For clarity the required finite Taylor step works for these higher
polynomial coefficients too. Let $H(u,t)=f_A(u^2,u^2t^2)+f_B+
f_C(u^2)-f_G$. It is a finite polynomial in $u,t$ with bounded
source arrays and $H(k,s)=0$. Write $\Delta_u=u-k$,
$\Delta_t=t-s$. A finite identity, obtained by monic polynomial
division, is
\[
 H=H_u(k,s)\Delta_u+H_t(k,s)\Delta_t+
       H_{uu}\Delta_u^2+H_{ut}\Delta_u\Delta_t+
       H_{tt}\Delta_t^2,                              \tag{JO13}
\]
where one may take
\[
 \begin{aligned}
 H_{uu}&={H(u,t)-H(k,t)-(u-k)H_u(k,t)\over(u-k)^2},\\
 H_{ut}&={H_u(k,t)-H_u(k,s)\over t-s},\\
 H_{tt}&={H(k,t)-H(k,s)-(t-s)H_t(k,s)\over(t-s)^2}.
 \end{aligned}
\]
These displayed divisions mean their finite polynomial quotients,
not mixed inverses. The first two arrays are bounded and $H_{tt}$
is infinitesimal: every $t$-dependent monomial of $H(k,t)$
contains $k^2$. The coefficient first differences similarly have a
bounded $u$ quotient and infinitesimal $t$ quotient. Exact first
angular stationarity supplies the two unchanged contact-jet balances
of the three-circuit lemma. The chain rule is legitimate here because
$g(x)=r$ and the original first state derivatives vanish at contact.
Consequently (JO11) has the strict finite three-circuit Gram and the
strict $F_v$ pivot, including the complete soft defect.

\paragraph{Every complete first normal.}
The following bounded polynomials are explicit:
\[
 \begin{aligned}
 P_Y/2={}&m\{V(1+gU)/\eta-\theta U\}
       -C D_v(1-g)/\eta+A D_uD_vg/\eta\\
       &+(D+\tfrac12)j/\eta+
                  (\alpha+\tfrac12)D_u(gV/\eta-1),\\
 P_E={}&{2mV\over\eta^2}\{g+1+j+g(1+j)U\}
       -{2m\theta\over\eta}\{1+(1+j)U\}\\
       &-{2CD_vj(1-g)\over\eta^2}
        +{2AD_uD_vg(1+j)\over\eta^2}
        +{2(D+\tfrac12)j^2\over\eta^2}\\
       &+{2(\alpha+\tfrac12)D_u(1+j)\over\eta}
                                      (gV/\eta-1),\\
 Z_X={}&B_x\theta+P_Y/2-(D+\tfrac12)j/\eta .
 \end{aligned}                                                \tag{JO14}
\]
Their only surviving $D$ factors are $Dj$ and $Dj^2$; $D$
cancels entirely from $Z_X$. Direct full-parent differentiation gives
\[
 \begin{aligned}
 \mathcal L_d&=-vP_Y/\sqrt D,\\
 \mathcal L_E&=2mvt^2/\eta^2+vP_E,\\
 \mathcal L_x&={1\over\sqrt{B_x}}
 \{2u\theta(m+B_x)-2mut^2(1+gU)/\eta+2uVZ_X\}.
 \end{aligned}                                                \tag{JO15}
\]
No normal is missing. For $P=P_Y,P_E$ use the exact finite
division $P=P_{00}+UP_U+VP_V$. Contact implies respectively
\[
 \begin{aligned}
 P_{Y,00}/x&=-P_{Y,U}(x,y)-\eta\zeta P_{Y,V}(y),\\
 P_{E,00}/x&=-2(m/x)\eta\zeta/\eta^2
                  -P_{E,U}(x,y)-\eta\zeta P_{E,V}(y).
 \end{aligned}                                                \tag{JO16}
\]
Thus these coquotients are bounded by finite formulas with no $j$
inverse. Allocate $vP$ by
$vP=(P_{00}/k)G+u(P_U+t^2P_V)A_t$; allocate the additional
$2mvt^2/\eta^2$ on $B_t$ with coefficient
$2mst/(\eta^2k)$. The $X$ line allocates exactly on $(C_t,B_t,A_t)$
with coefficients
\[
 {2u\theta(m+B_x)\over\sqrt{B_x}xs},\quad
 -{2ms(1+gU)\over\eta k\sqrt{B_x}},\quad
 {2vZ_X\over\sqrt{B_x}}.                               \tag{JO17}
\]
Every displayed coefficient is infinitesimal; for the first use
$\theta(m+B_x)/x=O(1)$, and for the others use $m/x=O(1)$.
The $N_d$ row has the additional bounded factor $1/\sqrt D$.
Each sum of coefficient functions vanishes at contact. Its finite
$u$ difference is bounded and its $t$ difference at $u=k$ is
infinitesimal, because every $t$ dependence carries $u$ or $k$.
The two unchanged conductors therefore place every whole normal on the
strict circuit roots with infinitesimal coefficients. Joining (JO10)
proves the full parent on the bounded-ratio chart.

\paragraph{The reciprocal testing ratio: a complete polynomial remainder.}
For $v/u$ unbounded, return to homogeneous state coordinates and use
\[
 E_t,\ N_d,\ N_x,\ P_m=\sqrt m\,vW,\ Q_\alpha=\sqrt\alpha Z,
                  \ R_t=u^2W,\ S_t=uvW.              \tag{JO18}
\]
The entire base error from (JO9) has zero Gram residue on this list:
its $p^2$ coefficient is $O(x)$ and
$\sqrt x\,p=\sqrt{x/m}(u/v)P_m$. Every coefficient is a product
of a bounded source factor and a bounded testing factor; $u/v$ is
pure testing.

Here is the full higher-angle identity used to control the other terms.
Write $T_s=-hUq+uv\xi$, so $Y=Y_0+T_s$, with $Y_0$ in
(JO5). After removing the actual quadratic base and the three entire
positive packets $mR_m^2$, $\alpha D_u(vY-Z)^2$ and
$B_xD_v(X+uZ)^2$, the exact remainder is
\[
 \begin{aligned}
 \mathcal R={}&C\{V(p-Y_0)^2-2D_v(p-Y_0)T_s+D_vT_s^2\}\\
 &+A\{(D_uD_v-1)Y_0^2+2D_uD_vY_0T_s+D_uD_vT_s^2\}\\
 &+\tfrac12D_uVY^2-UqY+(h-\tfrac12)Uq^2-uvq\xi .
 \end{aligned}                                                \tag{JO19}
\]
This is a quadratic in the full state, with bounded coefficients
$C,A,h$ and no $D$. It is checked as a whole polynomial. Insert
$q=\eta p-E_t$, $n=N_d/\sqrt D$, $\xi=N_x/\sqrt{B_x}$.
Every term contains $U,V$ or $uv$. Cross terms with $p$ have
infinitesimal coefficients on (JO18), using $Up=uR_t$,
$Vp=vS_t$ and $uvp=uS_t$. Hard quadratic terms have zero
residue. Terms $U^aV^bp^2$ are allocated on $R_t^2$ when
$a>0$, and on $S_t^2$ when $a=0,b>0$. Their two only nonzero
residues, verified from (JO19), are
\[
                         \kappa_R=\eta/2,\qquad
                         \kappa_S=1/2.               \tag{JO20}
\]
Every remaining such coefficient contains another $U$ or $V$.
These identities explicitly exhaust the scalar and every normal
coupling of (JO19).

The three retained positive packets have leading roots
$P_m,-Q_\alpha,N_x$. Their remaining root coefficients are
infinitesimal because $m,\alpha$ are infinitesimal and
$\theta/\sqrt\alpha$, $B_x\theta^2/\alpha$ are bounded:
\[
 \sqrt m\,R_m=P_m+\sqrt m\,uvY-\sqrt m(\xi+\theta uZ),
 \quad
 \sqrt{B_xD_v}(X+uZ)=
 \sqrt{D_v}\{N_x+\sqrt{B_x\theta^2/\alpha}\,uQ_\alpha\}.
\]
The alpha packet is $\sqrt{D_u}(\sqrt\alpha\,vY-Q_\alpha)$.
Thus the complete parent has a bounded Gram on (JO18) with strictly
positive diagonal residue
\[
             (1/(2\eta(1-\eta)),1,1,1,1,\eta/2,1/2).
                                                               \tag{JO21}
\]
This proves the reciprocal-ratio chart and completes the local proof.

\paragraph{All-tests join and receipt.}
The inherited small-$L$, bounded-$\eta$ remote theorem applies to
(JO1), including the exact original source openings and literal beta
zero. It covers the complement of both-infinitesimal testing angles.
The original exact sign gauges, literal axes, homogeneous state
extension and real/imaginary decomposition apply as in IC1 and RO.
All coefficients and radicals used above are pure-field quantities;
no mixed inverse or comparison of source and testing rates occurs.

The two new checkers record 25 and 4 exact assertions. They verify
the original full-parent substitution, every complete first normal,
the scalar identity, all newly used source denominator factors,
the reference hard cancellation, the whole reciprocal higher-angle
identity and both strict soft pivots. The finite polynomial bounds
and full-rates circuit allocation are the proof above, rather than
an inference from those checks. This admits precisely (JO1) and its
exact symmetry images. It does not remove the simultaneous
$r,j$ opening or any unbounded-$\zeta$ or infinitesimal-$\eta$ source.


\section{\texorpdfstring{Small-$L$ C1: a weighted closure through both $r$ openings}{Small-L C1: a weighted closure through both r openings}}
\label{component:p3df6e6bdd8f4}
% Edited scoped mathematical packet, 22 October 2026.

\label{p3df6e6bdd8f4:sec:r13-c1-r-openings}

\begin{theorem}
Use the original physical native C1 source, normalized by $d=1$, and
retain all original intersections. Suppose that $x=k^2$ is positive
infinitesimal and
\[
 \eta,\quad\zeta,\quad j_0=\eta+r-1
       \quad\hbox{are positive bounded units}.                 \tag{RO1}
\]
Then the complete original parent is tensor SOS at every independent
projective test and every state. The original $0<r<1$ may approach
either endpoint at any positive rate, and the literal $\beta=0$ face
is included. No lower positive real bound on $r$, $1-r$, or $\beta$
is imposed. This extends the admitted IC1 domain; it does
not cover an infinitesimal $j_0$, infinitesimal $\eta$, or unbounded
$\zeta$.
\end{theorem}

All scalar and Gram identities in this section use the normalized
parent $F/4$; multiplication of every square root by $2$ restores $F$.

The proof reuses the unchanged exact scalar and full-normal identities
(IC7)--(IC8), (IC16)--(IC19), and the three-circuit lemma. Its new
content is a uniform weighting of the two divergent event openings,
and a different $X$ normal that retains bounded and zero beta.

\paragraph{Opening-cleared source arrays.}
Put $\rho=\eta+r$, $y=x\eta\zeta$, $h=1-\theta$,
$c=C(1-r)$ and $a=Ar$. The beta numerator (SB1) gives
$\theta=x\pi$ with $\pi>0$ bounded. We do not require $\pi$
infinitesimal. The original source formulas give, uniformly for
all $0<r<1$,
\[
\begin{gathered}
 c=\rho/2+O(x),\quad a=\eta/2+O(x),\quad
 D=1/(2j_0)+O(x),\quad h=1+O(x),\\
 m/x\sim\rho/(2\zeta),\qquad
 \alpha/x^2\sim\rho\zeta(1+\zeta)/2,\qquad
 B_x\theta/x={K\rho(h+\zeta)\over2P_d},\quad
 B_x=\beta+\tfrac12\ge\tfrac12.                      
\end{gathered} \tag{RO2}
\]
These are finite canceled numerator identities, followed by residue
evaluation. Every canceled denominator at $x=0$ is a product of
positive units from $\eta,\zeta,j_0,j_0+r$ and $\rho$, with no
factor $r$, $1-r$, or $\pi$. In particular $D>0$ on this domain;
the original $D=0$ face has empty intersection here, as already proved
in (SB2). Both $\theta/\sqrt\alpha$ and
$B_x\theta^2/\alpha$ are bounded, not necessarily infinitesimal.

The exact source formulas for the two cleared strong weights are
\[
 c={K\rho E\over2P_d},\qquad a={\eta P_4\over2P_d}.
\]
Their order-$x$ differences in (RO2), and the following arrays, have
opening-free finite canceled formulas:
\[
 {P_{Y,00}\over x},\quad {P_{E,00}\over x},\quad
 {\theta(m+B_x)\over x},\quad
 {\alpha+(m+B_x)\theta^2\over x^2},\quad {q_V\over x}.
                                                               \tag{RO3}
\]
Here $P_Y,P_E,q_V$ are precisely (IC16) and (IC8). To check why no
unbounded event remains in these arrays, replace every occurrence of
$C(1-r)^2$ by $c(1-r)$, every $Ar^2$ by $ar$, every
$C(1-r)$ by $c$, and every $Ar$ by $a$ before division. No other
occurrence of $C$ or $A$ survives in the scalar or first-normal
polynomials. The accompanying new checker prints the denominators of
(RO2)--(RO3), retains arbitrary bounded $\pi$, and checks both
$r=0$ and $r=1$ after cancellation. This is a finite symbolic
criterion for boundedness on the stated field domain, not a numerical
endpoint approximation.

\paragraph{The normalized hard frame.}
Define the pure-source coefficients
\[
 b={r(1-r)\over\eta(j_0+r)},\quad
 H_0={\eta(j_0+r)\over2r(1-r)j_0},\quad
 \gamma={\rho\over2\eta(j_0+r)}.
\]
Both $1/\sqrt{H_0}$ and $1/\sqrt{B_x}$ are bounded; $H_0$
itself may be unbounded. Put
\[
 E_t=\eta uW-vZ,\quad
 N_y=\sqrt{H_0}(Y-ruW-bE_t),\quad
 N_x=\sqrt{B_x}(X+huZ).                              \tag{RO4}
\]
All three normals vanish at the actual source contact. For the
comparison weights
\[
 C_0={\rho\over2(1-r)},\quad A_0={\eta\over2r},\quad
 D_0={1\over2j_0},
\]
the exact quadratic base is
\[
 C_0(p-Y)^2+A_0Y^2+D_0(p-Y-q)^2+\tfrac12(p-Y)^2-pq
       =H_0(Y-rp-b(\eta p-q))^2+\gamma(\eta p-q)^2.  \tag{RO5}
\]
Thus the divergent normal has already been normalized before taking
residue. The cancellation factors needed for its entire hard block are
\[
 {C_0\over H_0}={r\rho j_0\over\eta(j_0+r)},\quad
 {A_0\over H_0}={(1-r)j_0\over j_0+r},\quad
 {D_0\over H_0}={r(1-r)\over\eta(j_0+r)}.             \tag{RO6}
\]
For example, after $p=(v+E_t)/\eta$ and $Z=1$, the $E_t$
coefficients in $p-Y$ and $Y$ are respectively
\[
 {1-r\over\eta}-b={(1-r)j_0\over\eta(j_0+r)},\qquad
 {r\over\eta}+b={r\rho\over\eta(j_0+r)}.            \tag{RO7}
\]
The strong $C,A$ factors therefore cancel even on those complete
first-state directions. For instance an error
$\delta_c(p-Y)^2/(1-r)$, with $\delta_c=O(x)$, has a
$N_y^2$ coefficient $\delta_c/((1-r)H_0)=O(x)$,
an $E_t^2$ coefficient $O(x)(1-r)$, and an $E_tN_y$
coefficient $O(x)/\sqrt{H_0}$. The $A$ error has the same
properties using $rH_0$ and the second formula in (RO7).
The original higher-angle factors have the same cleared coefficients,
multiplied by $U=u^2$ or $V=v^2$.

Suppose first that $u,v>0$ are infinitesimal and the pure-testing
ratio $t=v/u$ is bounded. After the state substitution (RO4), the
complete hard quadratic on $(E_t,N_y,N_x)$ consequently has bounded
entries and residue
\[
                         \operatorname{diag}(\gamma,1,1).       \tag{RO8}
\]
Here $\gamma$ is a positive bounded unit even when either $r$
opening is infinitesimal. The additional $m t^2$ contribution is
infinitesimal. The beta and half-core coefficients jointly give the
third pivot because $m/B_x$ is infinitesimal. These estimates concern
the full hard block, not only its restriction to the contact graph.

\paragraph{The complete scalar and all first normals.}
Normalize $Z=1$, homogenizing at the end as in IC1. Use the unchanged
positive separated roots
\[
 A_t=uv,\quad B_t={kv^2\over s u},\quad C_t=xs,\quad
 G=kv,\quad F_v=V-y,\qquad s=\sqrt{\eta\zeta}.
\]
The scalar at all three normals zero is exactly (IC7). Its coefficient
residues are still
\[
 c_F\sim{\rho\over2\eta^2},\quad
 f_A,f_B,f_C\sim\lambda={\rho\over2\eta},\quad
 f_G\sim3\lambda.                                   \tag{RO9}
\]
They do not depend on the residue of $\pi$. Every coefficient of
$U,V$ in (IC8) is bounded after the clearing described above;
for the additional term in $f_C$ use
$B_x\theta^2/x^2=(B_x\theta/x)\pi$. The exact contact
value and the two first-jet identities (IC11)--(IC14) are unchanged.
Their finite arrays therefore obey precisely the hypotheses of the
three-circuit lemma, with $s$ a unit. It gives the same strict finite
circuit Gram and the strict $F_v$ pivot. In particular the negative
soft defect has been allocated by the entire scalar (IC7), not
dropped by comparing only its quadratic base.

The complete first normals are (IC18), with the $N$ line divided
by $\sqrt{H_0}$ and the $X$ line divided by $\sqrt{B_x}$
instead of $\sqrt\beta$. Their exact root allocations are (IC19)
with those same replacements. Every coefficient remains infinitesimal:
the $N_y$ factors are bounded source factors times the old cleared
polynomials; the $X$ coefficient on $C_t$ is
\[
 {2u\theta(m+B_x)\over\sqrt{B_x}\,xs},
\]
which is infinitesimal by (RO3), and its coefficients on $B_t,A_t$
are respectively $O(k)$ and $O(v)$. In the latter coefficient
$Z_X=B_x\theta-(D+\tfrac12)(1+hU)
 +(D+\tfrac12)(1-r)/\eta$ is bounded. The $E_t$ line uses
$m/x,P_{E,00}/x$ and bounded polynomial arrays, as in IC1.
There is no negative power of $t$ in any coefficient.

Each coefficient is polynomial in $u,t$ with bounded pure-source
arrays; every term depending on $t$ carries $u,k$ or $m/k=O(k)$.
The exact monic differences (IC21) hence have bounded $u$ quotient
and infinitesimal $t$ quotient after evaluation at $u=k$. The
unchanged exact source contact cancels the bare $G$ term in each
complete normal. The conductors
$G(u-k)=k(A_t-G)$ and $G(t-s)=s(B_t-G)$ therefore allocate
all normals with infinitesimal coefficients on the strict circuit
roots. Joining this allocation to (RO8) proves the entire parent
on the bounded-ratio chart.

\paragraph{The reciprocal ratio chart.}
For $v/u$ unbounded use the homogeneous root list
\[
 E_t,\quad N_y,\quad N_x,\quad P_m=\sqrt m\,vW,
 \quad Q_\alpha=\sqrt\alpha Z,\quad R_t=u^2W,
 \quad S_t=uvW.                                      \tag{RO10}
\]
The error of the actual quadratic base relative to (RO5) is an
$O(x)$ quadratic in $(p,E_t,N_y)$ after the opening cancellations
(RO6)--(RO7). Every $\sqrt x\,p$ factor is allocated exactly by
\[
 \sqrt x\,p=\sqrt{x/m}\,(u/v)P_m,
\]
with a bounded source unit $x/m$ and an infinitesimal pure-testing
ratio $u/v$. Thus the entire error has zero Gram residue on (RO10).
All coefficients in the exact higher-angle identity (IC24) remain
bounded after $Y=rp+bE_t+N_y/\sqrt{H_0}$. This follows term
by term from (RO6), $b=O(r(1-r))$, $Ar=a$, and $C(1-r)=c$.
Its only nonzero scalar residues are the two positive diagonal entries
\[
                 \kappa_R=\rho\eta/2,\qquad
                 \kappa_S=\rho/2.                  \tag{RO11}
\]
No opening has been substituted before cancellation.

It remains essential to combine beta with its half-core part. Set
$B_t^{\rm state}=X+uZ$ and $J_0=p-Y-D_uvZ$. Relative to
$B_t^{\rm state}=0$, the exact $D$, beta and half-core increment is
\[
 2(D+\tfrac12)uvB_t^{\rm state}J_0
 +(D+\tfrac12)UV(B_t^{\rm state})^2
 +B_xD_v(B_t^{\rm state})^2.                          \tag{RO12}
\]
This whole-polynomial identity is checked in the new checker. Use
\[
 B_t^{\rm state}={N_x\over\sqrt{B_x}}
               +{\theta u\over\sqrt\alpha}Q_\alpha.
\]
The first two terms of (RO12) have zero Gram residue on (RO10);
$D$ is bounded, $1/\sqrt{B_x}$ and $\theta/\sqrt\alpha$
are bounded, and every coefficient carries $u$ or $v$.
The three remaining positive packets have their entire roots
\[
\begin{aligned}
 \sqrt m\,R_m&=P_m+\sqrt m\,uvY-\sqrt m\,B_t^{\rm state},\\
 \sqrt\alpha\sqrt{D_u}(vY-Z)
   &=\sqrt{D_u}\{\sqrt\alpha[rS_t+v(bE_t+N_y/\sqrt{H_0})]
                         -Q_\alpha\},\\
 \sqrt{B_x}\sqrt{D_v}B_t^{\rm state}
   &=\sqrt{D_v}\{N_x+\sqrt{B_x\theta^2/\alpha}\,uQ_\alpha\}.
\end{aligned}                                                    \tag{RO13}
\]
Each is its leading root plus an infinitesimal-coefficient combination
of (RO10). The square root in the last line is pure source and bounded;
its factor $u$ supplies the needed infinitesimal even when beta is
zero or bounded. The full parent therefore has a bounded Gram with
strict diagonal residue
\[
                 (\gamma,1,1,1,1,\kappa_R,\kappa_S).
                                                               \tag{RO14}
\]
This proves the remaining ratio chart.

\paragraph{All-tests join and exact scope.}
The earlier small-$L$, bounded-$\eta$ remote theorem includes (RO1)
with every original opening, so it covers the complement of
$u,v$ both infinitesimal. Literal axes have the elementary full-core
identities recorded in IC1. Its exact sign gauges apply unchanged.
Homogenizing the $Z=1$ identity gives polynomial state identities,
including $Z=0$; applying them to real and imaginary parts covers
complex states. All inverse and radical factors above are pure
source or pure testing. This completes the claimed full-parent
all-tests admission.

The newly removed source region is the portion of (RO1) outside IC1,
namely either $r$ or $1-r$ infinitesimal, with all total-scale and
earlier-complement intersections retained. Remaining C14.3 source
configurations still include infinitesimal $\eta$, unbounded
$\zeta$, and nonunit or nonpositive $j_0$ while the actual
$j=\eta(1+xh)-1+r$ remains positive. No conclusion on those
regions or on C14.4 follows from this theorem.

\paragraph{A local algebraic extension below beta zero.}
The preceding two local ratio-chart arguments have a stronger precise
scope that is useful when another event is restored later. Assume the
exact native stationary coefficient formulas (NC1)--(NC4),
$x>0$ infinitesimal, (RO1), $0<r,h<1$, and
\[
              B_x=\beta+\tfrac12\ge b_*>0
              \quad\hbox{for some real }b_* .                 \tag{RO15}
\]
Original source physicality and $\beta\ge0$ are not required for
this local assertion. Then the full C1-form parent is tensor SOS at
all independent $u,v$ both infinitesimal, including their literal
zero axes. In particular $-1/2<\beta<0$ is allowed whenever
$\beta+1/2$ is a positive bounded unit.

Indeed $P_d\sim1$ before any $h$ estimate, since all its source
arguments are bounded. The exact formula
$B_x\theta/x=K\rho(h+\zeta)/(2P_d)$ has bounded right side,
and (RO15) therefore implies $\theta/x$ bounded. All of (RO2)--(RO14)
follows, with $1/\sqrt{B_x}\le1/\sqrt{b_*}$. In particular
$m,\alpha>0$, $C,A>0$, and $D\sim1/(2j_0)>0$ on this
domain. No remaining use of $\beta\ge0$ occurs in the local
scalar, first-normal or hard-block proof: its complete contribution
is combined with the half-core as (RO12). At $u=0$ the signed
core contributes $D_v(X^2+Y^2)/2$, so its $X^2$ part plus
the beta packet is $B_xD_vX^2$. At $v=0$ the same combination
is $B_x(X+uZ)^2$ and the remaining half-core is
$(uW-Y)^2/2$. Thus both axes are also exact SOS under (RO15).
Only this local assertion is extended. The earlier all-tests remote
admission is not imported for negative beta.


\section{\texorpdfstring{The simultaneous $r,j$ opening: a harmonic full-event frame}{The simultaneous r,j opening: a harmonic full-event frame}}
\label{component:pbf6de1d7328a}

\label{pbf6de1d7328a:sec:r13-c1-simultaneous-rj}

\begin{theorem}
Use the original physical native C1 source, with $d=1$ and $x=k^2>0$
infinitesimal. Suppose $\eta,\zeta$ are positive bounded units and
$r$ and the actual $j=\eta(1+xh)-1+r$ are both positive
infinitesimals. Then the complete original parent is tensor SOS for
every independent projective test and every state. Literal $\beta=0$
is retained. No relative rate among $x,r,j$ or between source and
testing parameters is imposed.
\end{theorem}

All scalar and Gram identities in this section use the normalized
parent $F/4$; multiplication of every square root by $2$ restores $F$.

The essential change is to combine the whole positive $A,D$ events
before choosing the scalar graph. The weighted $uvX$ term is retained
inside a hard root. Solving for $p=uW$, rather than for $q=vZ$,
then gives a polynomial inverse with a pure-source unit denominator.

\paragraph{Source bounds and the harmonic frame.}
Write $U=u^2$, $V=v^2$, $D_u=1+U$, $D_v=1+V$,
$\theta=1-h=x\pi$, $B_x=\beta+1/2$, and
\[
 T=1+xh,\qquad e=\eta T=1-r+j,\qquad y=x\eta\zeta .
 \tag{HO1}
\]
The bounded-$\eta$ source bound gives $\pi=O(1)$.
Here $e,\eta,h,1-r\sim1$. Exact native numerator cancellation gives
\[
 \begin{gathered}
 C-C_* =O(x),\qquad Ar-e/2=O(x),\qquad Dj-1/2=O(x),\\
 C_*={1+j\over2(1-r)},\qquad A_*={e\over2r},\qquad D_*={1\over2j},\\
 m/x\sim1/(2\zeta),\qquad
 \alpha/x^2\sim\zeta(1+\zeta)/2,\\
 B_x\theta/x=O(1),\qquad
 q_U/x^2=O(1),\quad q_U=\alpha+(m+B_x)\theta^2 .
 \end{gathered}
 \tag{HO2}
\]
The checker cancels the exact rational formulas before setting
$x=r=j=0$: every denominator in the seven displayed bounded
coquotients factors there into constants and powers of the retained
unit $\zeta$. This proves boundedness at arbitrary simultaneous
source rates, rather than just nonvanishing at a selected boundary.
In particular $A,D>0$ and both are unbounded.
Also $\theta/\sqrt\alpha$ and $B_x\theta^2/\alpha$ are bounded,
including $\beta=0$.

Define pure-source positive quantities
\[
 \mathsf H=A+D,\qquad \rho={D\over\mathsf H},\qquad
 \mathsf K={AD\over\mathsf H},\qquad
 \delta={Ar+Dj\over\mathsf H}.
 \tag{HO3}
\]
Both $\mathsf H$ and $\mathsf K$ are unbounded,
$0<\rho<1$, and
$\mathsf K^{-1}=A^{-1}+D^{-1}=O(r+j)$.
Put $p=uW$, $q=vZ$, $J_0=p+uvX-q$, and use the full-event roots
\[
 \begin{aligned}
 E&=\sqrt{\mathsf K}\{J_0-(r-j)p\},\\
 N&=\sqrt{\mathsf H}\{Y-\rho J_0-\delta p\},\\
 N_x&=\sqrt{B_x}(X+huZ),\qquad
 \ell=E/\sqrt{\mathsf K},\quad n=N/\sqrt{\mathsf H},
 \quad \xi=N_x/\sqrt{B_x}.
 \end{aligned}
 \tag{HO4}
\]
All three vanish at the actual source contact. Their exact inverse is
\[
 \begin{aligned}
 p&={(1+hU)q-uv\xi+\ell\over e},&
 Y&=rp+\rho\ell+n,&X&=-huZ+\xi,\\
 J_t&=-jp+(1-\rho)\ell-n,&J_t&=p+uvX-Y-q .
 \end{aligned}
 \tag{HO5}
\]
The denominator $e$ is pure source. In particular no inverse of
$1+hU$ occurs, even if the testing and source rates are unrelated.
The useful exact weighted identities are
\[
 {A\rho^2\over\mathsf K}=\rho,\qquad
 {A\over\mathsf H}=1-\rho,\qquad
 {D(1-\rho)^2\over\mathsf K}=1-\rho,\qquad
 {D\over\mathsf H}=\rho .
 \tag{HO6}
\]

\paragraph{The complete hard block and its reference cancellation.}
The $A,D$ base has the exact identity
\[
 \begin{split}
 AY^2+DJ_t^2={}&N^2+E^2
 +2p\left\{{Ar+Dj\over\sqrt{\mathsf H}}N
                 +\sqrt{\mathsf K}(r-j)E\right\}\\
 &+(Ar^2+Dj^2)p^2 .
 \end{split}
 \tag{HO7}
\]
Both displayed cross coefficients and the last coefficient are
infinitesimal: the squared $E$ cross coefficient is $O(r+j)$
and the squared $N$ cross coefficient is $O(\mathsf H^{-1})$.
For bounded testing ratio $t=v/u$, (HO5) has bounded coefficients
on the hard roots, and the only potentially large $m$ coefficient is
$mt^2/\mathsf K$, which is infinitesimal.
The higher $A$ packet is controlled by (HO6); for example
\[
 {Ar^2\over\mathsf K}=r^2+{Ar\over Dj}rj=O((r+j)^2),
 \qquad {Ar\rho\over\mathsf K}=r .
 \tag{HO8}
\]
Consequently its hard coefficients carry an additional $U$ or $V$.
The bounded $C$ and half-core packets have infinitesimal coefficients
on $E,N$, and combining the beta packet with its half-core part gives
the third hard pivot. The full hard residue on $(E,N,N_x)$ is
$\operatorname{diag}(1,1,1)$.

A stronger exact base statement is needed for unbounded $t$.
Set $\widetilde q=q-uvX=ep-\ell$, so that
$J_t=p-Y-\widetilde q$. The reference quadratic form
\[
 \mathcal B_*=
 C_*(p-Y)^2+A_*Y^2+D_*(p-Y-\widetilde q)^2
                    +\tfrac12(p-Y)^2-p\widetilde q
 \tag{HO9}
\]
has no soft $p$ coordinate after substituting
$Y=rp+\rho\ell+n$ and $\widetilde q=ep-\ell$:
\[
 \mathcal B_*=
 (A_*+C_*+\tfrac12)(\rho\ell+n)^2
                 +D_*\{(1-\rho)\ell-n\}^2 .
 \tag{HO10}
\]
This is an exact contact cancellation, checked before taking residues.
Since $A_*/A=1+O(x)$ and $D_*/D=1+O(x)$, (HO6) shows that
its bounded hard Gram on $(E,N)$ has residue $\operatorname{diag}(1,1)$.
Replacing the reference weights by the actual $C,A,D$ creates a
quadratic error with coefficients $O(x)$ in $(p,E,N)$.
For the singular weights this follows from
$(A-A_*)Y^2=O(x)\,AY^2$ and
$(D-D_*)J_t^2=O(x)\,DJ_t^2$, using (HO6),(HO7);
neither an $x/r$ nor an $x/j$ estimate is used.

\paragraph{The exact bounded-ratio scalar circuit.}
At $E=N=N_x=0$ and $Z=1$, let
\[
 w(U)={1+hU\over e},\qquad
 p=w(U)v,\quad Y=rw(U)v,\quad X=-hu .
 \tag{HO11}
\]
Since $w(x)=1/\eta$, this graph agrees exactly with the native
source contact. The source contact value and all first state and
angular jets are therefore inherited by the chain rule.

For compact scalar coefficients set
$a=C(1-r)^2$, $b=Ar^2$, $d=(D+1/2)j^2$,
$c=(\alpha+1/2)r^2$, and
\[
 L_A(U)=1+2h+(2h+h^2)U+h^2U^2 .
 \]
The complete scalar is exactly
\[
 \mathcal S=c_F(V-y)^2+f_A(U,V)A_t^2+f_BB_t^2+
                         f_C(U)C_t^2-f_GG^2 ,
 \tag{HO12}
\]
on the separated roots
$A_t=uv$, $B_t=kv^2/(su)$, $C_t=xs$, $G=kv$,
where $s=\sqrt{\eta\zeta}$. The coefficients are
\[
 \begin{aligned}
 c_F&={2m(h+r)+a+b+c\over e^2},\\
 q_V&={a+b+d\over e^2}-{2m\theta\over e}
                       -{(2\alpha+1)r\over e}-\tfrac12,\\
 f_B&={m\eta\zeta\over e^2x},\qquad
 f_C={\alpha-c_Fy^2+q_UU\over\eta\zeta x^2},\qquad
 f_G=-{q_V+2c_Fy\over x},\\
 f_A&={mV\over e^2}
 \{(h+r)^2+2hr+2hr(h+r)U+h^2r^2U^2\}
       -{2m\theta\over e}(h+r+hrU)\\
 &\quad+{a\over e^2}(2h+h^2U)(1+V)
       +{b\over e^2}L_A(U)(1+V)
       +{d\over e^2}(2h+h^2U)+{c\over e^2}L_A(U)V\\
 &\quad-{(2\alpha+1)r\over e}(1+h+hU)
                              +B_x\theta^2-\tfrac12 .
 \end{aligned}
 \tag{HO13}
\]
The only opening weights here are $Ar^2$ and $Dj^2$, both
infinitesimal. Thus $f_A$ is a bounded polynomial of degree at most
two in $U$ and one in $V$, and all nonconstant coefficients of
$f_C$ are bounded by (HO2). Direct source reduction gives
\[
 c_F\sim\tfrac12,\qquad f_A,f_B,f_C\sim\tfrac12 .
 \tag{HO14}
\]
Exact scalar contact gives
$f_G=f_A(x,y)+f_B+f_C(x)\sim3/2$, without separately estimating
$q_V/x$.

Apply the finite three-circuit lemma with $\lambda=1/2$.
The polynomial Taylor arrays are the monic quotients in (JO13),
now applied to
$f_A(u^2,u^2t^2)+f_B+f_C(u^2)-f_G$.
Their $u$ and mixed arrays are bounded; the $t,t$ array and the
coefficient $t$ first differences are infinitesimal, since every
$t$ dependence at $u=k$ contains $k^2$.
The two exact angular contact jets give the two circuit balances.
This supplies the strict three-circuit Gram and the strict
$V-y$ pivot for the entire scalar, including its negative soft term.

\paragraph{All complete first normals.}
The following formulas come from differentiating the entire original
parent after (HO5), with $Z=1$. Put
$a_E=r/e+\rho$, $b_E=(1-r)/e-\rho$,
$c_E=1-\rho-j/e$. The bounded polynomials are
\[
 \begin{aligned}
 P_N/2={}&m\{wV(1+rU)-\theta U\}
 -CD_v(1-r)w+AD_uD_vrw\\
 &+(D+\tfrac12)jw
             +(\alpha+\tfrac12)D_u(rwV-1),\\
 P_E/2={}&{mV\over e^2}(h+r+hrU)-{m\theta\over e}
 +ma_E\{wV(1+rU)-\theta U\}\\
 &+CD_v(1-r)wb_E+AD_uD_vrwa_E
 -(D+\tfrac12)jwc_E
            +(\alpha+\tfrac12)D_u(rwV-1)a_E,\\
 Z_X={}&{m\theta(1+rU)\over e}+B_x\theta\\
 &-{CD_v(1-r)^2w+AD_uD_vr^2w+(D+\tfrac12)j^2w
                 +(\alpha+\tfrac12)D_ur(rwV-1)\over e}.
 \end{aligned}
 \tag{HO15}
\]
Their only opening-weight factors are $Ar,Dj,Ar^2,Dj^2$.
Every coefficient is bounded uniformly in the source rates and
$0<\rho<1$. The three complete first normals are
\[
 \begin{aligned}
 \mathcal L_N&={vP_N\over\sqrt{\mathsf H}},\\
 \mathcal L_E&={2mvt^2/e^2+vP_E\over\sqrt{\mathsf K}},\\
 \mathcal L_x&={1\over\sqrt{B_x}}\left\{
 2u\theta(m+B_x)-2mwut^2(1+rU)
           \left(1+{V(1+rU)\over e}\right)+2uVZ_X\right\}.
 \end{aligned}
 \tag{HO16}
\]
For either $P=P_N,P_E$, use its exact finite monic decomposition
$P=P_{00}+UP_U+VP_V$.
State contact gives
\[
 \begin{aligned}
 P_{N,00}/x&=-P_{N,U}(x,y)-\eta\zeta P_{N,V}(y),\\
 P_{E,00}/x&=-2(m/x)\eta\zeta/e^2
                         -P_{E,U}(x,y)-\eta\zeta P_{E,V}(y).
 \end{aligned}
 \tag{HO17}
\]
Thus these coquotients are bounded without any new opening inverse.
Allocate $vP=(P_{00}/k)G+u(P_U+t^2P_V)A_t$.
The extra $E$ term allocates on $B_t$ with coefficient
$2mst/(e^2k\sqrt{\mathsf K})$.
The three $X$ terms allocate respectively on $C_t,B_t,A_t$,
with coefficients
\[
 {2u\theta(m+B_x)\over xs\sqrt{B_x}},\qquad
 {-2msw(1+rU)\{1+V(1+rU)/e\}\over k\sqrt{B_x}},\qquad
 {2vZ_X\over\sqrt{B_x}} .
 \tag{HO18}
\]
Every displayed allocation coefficient is infinitesimal.
The coefficient sums vanish at contact by the exact state jets;
their finite $u$ first differences are bounded and their $t$
first differences at $u=k$ are infinitesimal.
This follows directly from the bounded polynomial arrays in
(HO15): every $t$ dependence carries $u$ or $k$.
The two unchanged circuit conductors therefore place every whole
normal on the strict circuit roots with infinitesimal coefficients.
Together with the complete hard block, this proves all bounded $t$,
including infinitesimal $t$, without a missing state normal.

\paragraph{The reciprocal ratio: retain both soft directions.}
If $v/u$ is unbounded, work homogeneously and use seven roots
\[
 E,\quad N,\quad N_x,\quad
 P_m=\sqrt m\,vW,\quad Q_\alpha=\sqrt\alpha Z,\quad
 R_t=uvZ,\quad S_t=uvW .
 \tag{HO19}
\]
The base error following (HO10) has zero Gram residue on this list:
use the exact separated identity
$\sqrt x\,p=\sqrt{x/m}(u/v)P_m$.
This absorbs every $O(x)$ coefficient in $(p,E,N)$, since $x/m$
is a source unit and $u/v$ is a pure-testing infinitesimal.

After removing the actual quadratic base (HO9), with actual $C,A,D$,
and the three entire positive packets
$mM_t^2$, $\alpha D_u(vY-Z)^2$,
$B_xD_v(X+uZ)^2$, the full remaining quadratic is exactly
\[
 \begin{split}
 \mathcal R={}&CV(p-Y)^2+A(U+V+UV)Y^2
                      +\tfrac12D_uVY^2-\theta UqY\\
 &+(h-\tfrac12)Uq^2+\tfrac12h^2U^2q^2
       -uv\xi\{Y+(1+hU)q\}+\tfrac12UV\xi^2 .
 \end{split}
 \tag{HO20}
\]
This identity retains the whole $D$ event in the quadratic base,
so no coefficient containing a bare divergent $D$ remains.
It is checked as one complete homogeneous polynomial.

Here is a finite allocation of every term. The exact identity
\[
 up={(1+hU)R_t-u^2v\xi+uE/\sqrt{\mathsf K}\over e}
 \tag{HO21}
\]
has bounded coefficients and uses no mixed inverse.
The $C$ packet has leading root
$\sqrt C(1-r)S_t$; its other coefficients are
$-\sqrt C\,\rho v/\sqrt{\mathsf K}$ and
$-\sqrt C\,v/\sqrt{\mathsf H}$ on $E,N$, both infinitesimal.
All three $A$ higher packets have zero residue. For example,
\[
 \sqrt A\,uY=\sqrt A\,r(up)+\sqrt\rho\,uE+
                               \sqrt{1-\rho}\,uN,\qquad
 \sqrt A\,vY=\sqrt A\,rS_t+\sqrt\rho\,vE+
                               \sqrt{1-\rho}\,vN .
 \tag{HO22}
\]
Here $\sqrt A\,r=O(\sqrt r)$ and (HO8),(HO21)
control every coefficient; the $uvY$ packet is $v$ times the first.
The half-$VY^2$ term is smaller by the same identities.
The $-\theta UqY$ term is $-\theta R_t(uY)$ and has zero residue.
The next two terms are $(h-1/2)R_t^2$ and
$h^2U R_t^2/2$.
Finally, the cross term is
$-\xi v(uY)-\xi v(1+hU)R_t$ and all its coefficients,
as well as those of $UV\xi^2/2$, are infinitesimal.
Thus this allocation exhausts the complete quadratic (HO20);
its only nonzero soft residues are $1/2$ on each of $R_t,S_t$.

The three retained positive packets have leading roots
$P_m,-Q_\alpha,N_x$:
\[
 \begin{aligned}
 \sqrt m\,M_t&=P_m+\sqrt m\,uvY-\sqrt m\,\xi
                  -\sqrt{m/\alpha}\,\theta uQ_\alpha,\\
 \sqrt{\alpha D_u}(vY-Z)
             &=\sqrt{D_u}(\sqrt\alpha\,vY-Q_\alpha),\\
 \sqrt{B_xD_v}(X+uZ)
             &=\sqrt{D_v}\{N_x+
                      \sqrt{B_x\theta^2/\alpha}\,uQ_\alpha\}.
 \end{aligned}
 \tag{HO23}
\]
Every other coefficient here is infinitesimal by (HO2),(HO21),
because $\theta\sqrt{m/\alpha}=O(\sqrt x)$.
The complete bounded Gram on (HO19) therefore has strictly positive
diagonal residue
\[
             \operatorname{diag}(1,1,1,1,1,1/2,1/2).
 \tag{HO24}
\]
This proves the reciprocal chart for the full parent and every state.

\paragraph{All-tests join and receipt.}
The inherited small-$L$, bounded-$\eta$ remote theorem covers the
complement of both-infinitesimal testing angles on this original
physical source domain. Its exact source openings and literal beta
zero remain available. The original exact sign gauges, literal axes,
homogeneous state extension and real/imaginary decomposition join
as in IC1, RO and JO. All roots, radicals and inverses above are
pure-field or separated; every bound is a product of a source bound
and a testing bound.

The separate checker records 28 exact assertions, including the
full-parent scalar substitution, all three complete first normals,
the two full-event maps, the harmonic completion, the exact reference
kernel cancellation, the whole reciprocal remainder, its two strict
soft pivots and the retained-unit denominator factorizations.
Those checks support the finite bounds and circuit allocation proved
above; they are not being used as a substitute for them.
This theorem removes precisely the simultaneous positive
infinitesimal $r,j$ opening with unit $\eta,\zeta$.

\begin{corollary}[The entire unit-$\eta,\zeta$ small-$L$ region]
Within the original physical native small-$L$ C1.4 source, suppose
$x>0$ is infinitesimal and $\eta,\zeta$ are positive bounded units.
Then the full parent is tensor SOS at every independent projective
test and every state, with all original source intersections retained.
\end{corollary}
\begin{proof}
The positive actual gap $j=\eta(1+xh)-1+r$ is bounded.
If $j$ is a unit, then $j_0=j-\eta xh$ is a positive unit,
so RO applies, including either $r$ opening. If $j$ is infinitesimal,
then $1-r\sim\eta$ is a unit. A unit $r$ is exactly the JO domain;
an infinitesimal $r$ is exactly the HO domain. These alternatives
exhaust the stated source, require no comparison of $j$ with $x$,
and each supplies its complete local, remote and endpoint join.
\end{proof}


\section{\texorpdfstring{Two exact conductors for the remaining small-$\eta$ chart}{Two exact conductors for the remaining small-eta chart}}
\label{component:peea9e9c546c5}
% Edited scoped mathematical packet, 22 October 2026.

Use the small-$\eta$, unit-$\zeta$ source: $x,\eta>0$
infinitesimal, $\zeta$ a positive bounded unit, and
$\beta\ge0$.  Put
\[
 a=1/\zeta,\quad c^2=a/\eta,\quad k^2=x,\quad
 B_x=\beta+\tfrac12,\quad q=\alpha+(m+B_x)(1-h)^2,
 \quad\mu=m-a^2q,\quad\lambda=\mu/(\eta ax).
\]
For homogeneous testing states use the exact roots
\[
\begin{gathered}
 \xi=X+huZ,\quad E=\eta uW-vZ,\quad
 D_0=u-cv,\quad J=(uv-x/c)Z,\\
 R_D=v(u^2-x)Z,\quad R_c=kv(\eta cW-Z),\quad
 \mathcal Q=\sqrt q(avW-uZ-a\xi).
\end{gathered}
\]
Direct multiplication gives
\[
 R_D-uJ=(x/c)D_0Z,\qquad
 \mathcal Q/\sqrt q=-D_0Z+(c/k)R_c-a\xi.
\]
Hence the complete principal $\xi$ row cancels exactly to
\[
 -2\eta\lambda ckR_c\xi+2\eta\lambda c(R_D-uJ)\xi
 =-2\eta\lambda x\mathcal Q\xi/\sqrt q
                    -2\eta\lambda xa\xi^2.             \tag{SC5}
\]
Together with the singular rank's $+\mu\xi^2$, the last term
leaves $-\mu\xi^2$, because $\mu=\eta\lambda ax$.
The latter coefficient is infinitesimal relative to $B_x\xi^2$.
The complete angular block on $(\mathcal Q,\sqrt{B_x}\xi)$
has strictly positive residue: the original source expressions give
\[
 B_xq/x^2\sim
 \frac{B_x\zeta(1+\zeta)}2+\frac{(1+\zeta)^2}4
       >\frac14\sim(\eta\lambda)^2.
\]
For unbounded $B_x$ the normalized off-diagonal tends to zero.
For bounded $B_x\ge1/2$, the displayed difference is bounded
below by $\zeta(3+2\zeta)/4>0$.  The source substitutions are
the regular DG/LD dictionary; the direct derivative and margin
calculation is also printed in the refinement proof.
All inverses and radicals in these coefficient expressions are
pure source.  The full radial and remaining normal allocations are
the retained JF/DG/LD/CJ packets; SC5 supplies their angular block.


\section{\texorpdfstring{C14.3b: a complete source subchart from a damped two-circuit completion}{C14.3b: a complete source subchart from a damped two-circuit completion}}
\label{component:p024c821bfb59}


\begin{theorem}
Normalize $d=1$ and retain the unchanged physical C14.3b source, with $x,\eta>0$
infinitesimal, $\zeta$ a positive bounded unit, every original $j$
and $1-r$ opening, $\beta\ge0$ including literal $\beta=0$, and all inherited intersections.
On the additional \emph{pure-source} subchart
\[
                         x/\eta\text{ bounded},              \tag{DG1}
\]
the entire original parent is finite tensor SOS for every independent
projective angle and state test. No relation between a source rate
and a testing rate is assumed. The proof below treats the whole local
corner; the inherited remote theorem gives its complement. The source
subchart $x/\eta$ unbounded is not admitted by this theorem.
\end{theorem}

\paragraph{The sign reduction and the complete frame.}
Use the sign gauges to take $u,v\ge0$. For real states with $WZ\le0$
the independent half-core sign lemma proved in this revision gives a
direct tensor-SOS decomposition of the original signed half core;
all remaining original event packets are positive source multiples
of squares. Thus only $WZ>0$ needs the following construction. No
normalization or inverse of a state coordinate is used here.

Use all homogeneous JF1 definitions, and abbreviate $Q=Q_t$,
$T=T_D$, $B=B_x$, $a=1/\zeta$, $x=k^2$, $a=\eta c^2$. In particular
\[
 E=\eta uW-vZ=\eta(P+Q),\quad H=avW-uZ,
 \quad \tau=\eta cW-Z,\quad D_0=u-cv,
\]
\[
 J=(uv-x/c)Z,\quad
 R=-\sqrt\lambda\,\eta kD_0W,\quad
 f_t=F_t/\eta,\quad \mathcal Q=\sqrt q(H-a\xi),
 \quad N_x=\sqrt B\xi.
\]
The full radial packet $\mathscr R$ is JF2, with the strict two-root
completion JF3. Its roots will be denoted $R_b,f_t$; both normalized
pivots have positive unit residue. The angular packet is exactly
\[
 \mathscr A=\mathcal Q^2
       -\frac{2\mu}{a\sqrt{qB}}\mathcal QN_x
       +(1-\mu/B)N_x^2.                              \tag{DG2}
\]
Its normalized Schur pivot is strict by the unchanged R15/R16 source
margin, including $\beta=0$ and unbounded $B$. Those previously
proved angular and radial blocks are used once, as their complete
packets.

\paragraph{Damp the remaining normal and complete both circuit companions.}
Set the positive source unit
\[
 G_0=2\eta\lambda-2\eta\sim1,\qquad
 \epsilon=\gamma-G_0/h,
\]
\[
 \widehat R=R+\frac{G_0k}{2\sqrt\lambda}Q,
 \qquad \widehat J=\sqrt\lambda J+
                         \frac{G_0u}{2\sqrt\lambda}Q,
 \qquad c_U=A-\frac{G_0^2}{4\lambda}.
\]
Use the exact full-$D$ conductor and polarization displayed in Chapter~\ref{ch:mechanisms} to conduct only
$(G_0/h)uvQ\xi$, leaving $\epsilon uvQ\xi$ unaltered. Complete
\emph{both} $R$ and $J$, incorporating the entire native $AUQ^2$
row into the latter completion. Since
$G_0/(2\lambda)=\eta-\eta/\lambda$, the complete original parent
has the following exact identity:
\[
\begin{split}
 F/4={}&\mathscr A+\mathscr R+S_0
    -\lambda x(\eta P+\eta Q/\lambda)^2
    +\widehat R^2+\widehat J^2\\
 &+c_UUQ^2+e_4UF_t^2+BV\xi^2
    +\frac{2x}{c}uv\tau^2
    +\frac{G_0x}{c}D_0^2WZ\\
 &-e_4yuHF_t+\kappa uF_t\xi
                 +\frac{G_0x}{c}PH+\epsilon uvQ\xi,       
\end{split}\tag{DG3}
\]
where
\[
 S_0=CP^2+AQ^2+D_p(P-hT)^2+2b_0QT,
 \qquad b_0=b'-G_0/2.                                 \tag{DG4}
\]
This is a simultaneous completion of the entire remaining normal and
the coupled circuit. It retains every full-$D$, full-$Y$ and higher
normal of JF1. There is no second bare radial pivot and no extracted
uncentered $m$ event. The coefficient of the extra $uv\tau^2$ row
is exactly $2x/c>0$, because $2\lambda-G_0/\eta=2$.

\paragraph{Two strictly positive source deficits replace the failed soft pivot.}
Put $K=CD_p/(C+D_p)$, $s=x+\eta$, and
\[
                         \Delta_0=h^2AK-b_0^2.
\]
The original source formulas imply exact finite divisions
\[
 \Delta_0=x d_x+\eta d_\eta,\qquad
 \frac{c_U}{\eta}=x e_x+\eta e_\eta,                \tag{DG5}
\]
with bounded source coefficients having the following residues:
\[
 \begin{array}{c|c}
 d_x&\pi/2+1/4+1/(4a)\\
 d_\eta&\sigma/2+3/4\\
 e_x&\pi+1+1/(2a)\\
 e_\eta&\sigma+3/2
 \end{array}
 \qquad h=1-x\pi,\quad r=1-\eta\sigma.             \tag{DG6}
\]
All four residues are strictly positive. Here $\pi,\sigma\ge0$
are bounded by the unchanged physical source estimates. Thus
\[
 \Delta_0/s\text{ is a positive bounded unit},\qquad
 c_U/(\eta s)\text{ is a positive bounded unit}.     \tag{DG7}
\]
These assertions compare source quantities only and do not use DG1.

For clarity, the divisions in DG5 are opening-safe. Write
$C_n=(1-r)C$, $D_j=jD_p$ and $\Theta_B=B(1-h)$, all given by their
unchanged unit-denominator source formulas. The identity
$\eta\Theta_B-D_j=-C_n$ gives
$b'-\gamma h/2=-hC_n/r$. The gap-cleared harmonic coefficient is
\[
 \eta K=\frac{C_nD_j}
 {C_n(1+xh-\sigma)+D_j\sigma}.
\]
Its denominator is a positive unit, even at either opening. The
quantities $\eta\lambda$, $b'$, $A/\eta$ and $\gamma h$ also have
unit-denominator expressions after their displayed monomial factors
are canceled. Both functions on the left of DG5 vanish at
$x=\eta=0$. Define their $\eta$ coefficients by their exact $x=0$
values divided by $\eta$, and then divide the remainder by $x$.
This gives DG5 with only source-unit denominators; differentiation at
the origin gives DG6. No assumption on $x/\eta$ enters these exact
finite divisions.

\paragraph{The source roots control the complete soft direction.}
There is the exact pure-source completion
\[
 S_0=N_s^2+N_T^2+N_Q^2,
\]
\[
 N_s=\sqrt{C+D_p}\left(P-\frac{hD_p}{C+D_p}T\right),
 \quad N_T=h\sqrt K\left(T+\frac{b_0}{h^2K}Q\right),
 \quad N_Q=\sqrt{\frac{\Delta_0}{h^2K}}Q.            \tag{DG8}
\]
As $\eta K\sim1/2$, DG7 gives
$N_Q/\{\sqrt{\eta s}\,Q\}$ a positive source unit. The inverse
identities show
\[
 T=\frac{N_T}{h\sqrt K}-\frac{b_0}{h^2K}Q,
 \qquad
 P=\frac{N_s}{\sqrt{C+D_p}}
                      +\frac{hD_p}{C+D_p}T.         \tag{DG9}
\]
Consequently $P$ has bounded coefficients on $(N_s,N_T,N_Q)$:
the first two coefficients are $O(\sqrt\eta)$ and the last is
$O(\sqrt{\eta/s})$, which is bounded because $\eta\le s$.
This is the complete $P$ control after the full $D$ shift; it is not
an assumption that the old $E/\sqrt\eta$ pivot remains strict.

The negative square in DG3 is an infinitesimal Gram on these three
roots. Its root is
\[
 \sqrt{\lambda x}(\eta P+\eta Q/\lambda)
   =\eta\sqrt{\lambda x}\,P+\eta\sqrt{x/\lambda}\,Q.
\]
Since $\eta\lambda\sim1/2$, the first coefficient is
$O(\sqrt{\eta x})$ on the bounded allocation of $P$; the second,
after normalizing $Q$ by $N_Q$, is
$O(\eta\sqrt{x/s})$. Both are infinitesimal, without DG1.

\paragraph{Every remaining complete normal is small on this frame.}
The angular-radial term in DG3 is
\[
 -e_4y\eta u\left(\mathcal Q/\sqrt q
                             +aN_x/\sqrt B\right)f_t.
\]
Its coefficients are $O(u)$ and $O(xu)$, using
$e_4\eta^2\sim1/2$ and $q/x^2$ a positive bounded unit.
The term $\kappa uF_t\xi$ has coefficient
$(\eta\kappa)u/\sqrt B$, again infinitesimal.
The new $PH$ term has coefficients
\[
 \frac{G_0x}{c\sqrt q}\quad\hbox{and}\quad
 \frac{G_0ax}{c\sqrt B}
\]
on $P\mathcal Q$ and $PN_x$; these are respectively
$O(\sqrt\eta)$ and $O(x\sqrt\eta)$, so DG9 allocates the whole
term infinitesimally.

Finally $\epsilon=O(s)$ by the same gap-free source expressions.
The complete, still unconducted normal has coefficient
\[
 \frac{\epsilon uv}{\sqrt{\eta s}\sqrt B}
         =O\!\left(\sqrt{s/\eta}\,uv\right)          \tag{DG10}
\]
on $N_QN_x$. DG1 is used \emph{exactly here}: it makes $s/\eta$
bounded, and hence the coefficient infinitesimal for independent
infinitesimal $u,v$. No testing comparison is imposed.

Thus the remainder on
\[
                  (N_s,N_T,N_Q,R_b,f_t,\mathcal Q,N_x)
\]
has a bounded Gram with strictly positive residue. The bounded-Gram
lemma makes it finite tensor SOS. Add the whole positive rows
$\widehat R^2$, $\widehat J^2$, $c_UUQ^2$, $e_4UF_t^2$,
$BV\xi^2$, $(2x/c)uv\tau^2$ and
$(G_0x/c)D_0^2WZ$. The last is a literal tensor square on the current
pure-testing sign chart; all radicals are source-only or testing-only.
This proves the complete local parent, not merely its scalar
restriction.

\paragraph{Equality, endpoints and the global join.}
At original contact $P,Q,T,J,R,\tau,F_t,\xi$ vanish, so all newly
completed rows vanish too. DG3 preserves all contact normals exactly.
The row formulas have no angle ratio or state-coordinate inverse;
$\sqrt{uv}$ is defined as zero on a literal axis. Thus they include
both literal local axes. The inherited gauges cover every angle sign.
For a complex state, split the real-coefficient Hermitian parent into
its real and imaginary evaluations and apply either the $WZ\le0$
core lemma or this $WZ>0$ proof separately to each real state. There
is no assumption that both real and imaginary state components have
a unit $Z$ or $W$. Combine the local proof with the unchanged remote
all-state theorem to obtain the all-tests assertion on DG1.

The hypothesis DG1 is used only in the final allocation DG10. The exact DG3 identity and preceding source bounds also apply without it; the completed $J$-row conductor in Section~\ref{component:pe3bbfd3a8151} supplies the remaining allocation for arbitrary $x/\eta$.


\section{\texorpdfstring{The damped frame also enlarges the pure-testing state chart}{The damped frame also enlarges the pure-testing state chart}}
\label{component:p8cfada778490}
% Edited scoped mathematical packet, 22 October 2026.

In any local chart with $u$ infinitesimal, suppose $Z\ne0$ and
the pure-testing ratio $|X/Z|$ is unbounded.  Then
\[
 \frac{1+v^2}{2}\{(X+uZ)^2-(1+u^2)Z^2\}>0.             \tag{DS3}
\]
Indeed, after division by $Z^2>0$, the square $(X/Z+u)^2$
is unbounded and $1+u^2$ is bounded.  The factor $1+v^2$
is positive.  This expression is the complete $X$ square and the
negative $Z$ term of the signed quarter-core.  Its other two
terms are $(uW+uvX-Y-vZ)^2/2$ and
$(1+u^2)(vY-Z)^2/2$, both nonnegative.  Thus the signed core
itself is a nonnegative pure-testing scalar and has a square root
in the testing field.  Every original event remains a square with
its nonnegative source coefficient.  This proves a direct separated
SOS certificate for the complete parent, with arbitrary $W,Y$.

The state chart $Z=0$ and the sign chart $WZ\le0$ have the
direct core certificates already stated in the geometric reduction.
Consequently the C14.3 local certificate only needs the remaining
case $Z\ne0$, $WZ>0$, $X/Z$ bounded.  The completed-J
conductor CJ treats that whole case without bounding $W/Z$ or
$Y/Z$.  Apply the appropriate real chart to each real and imaginary
part to include complex states.


\section{\texorpdfstring{The entire unit-$\eta$ large-$\zeta$ source, through the actual-$j$ opening}{The entire unit-eta large-zeta source, through the actual-j opening}}
\label{component:p60eca5442075}

\label{p60eca5442075:sec:r17-lz-unit-eta}

\begin{theorem}
Retain the original physical C1 source and every original contact,
strict-gap, nonnegative-weight, total-scale and earlier-exclusion
intersection. Suppose $x=k^2,L=x\zeta$ are positive infinitesimals,
$\zeta$ is unbounded and $\eta$ is a positive bounded unit.
Then the complete original parent is finite tensor SOS for every
independent projective test and every real or complex state.
No positivity or noninfinitesimal lower bound is imposed on
$j_0=\eta+r-1$: only the actual original gap
$j=\eta(1+xh)-1+r>0$ is used. Both $r$ openings, simultaneous
$r,j$ openings and literal $\beta=0$ are included.
\end{theorem}

This theorem treats the whole $\eta$-unit source portion of C14.3a.
Every identity below uses the normalized quarter-parent $\mathcal P=F/4$.
The complete parent identity JF1--JF3 from R16 is algebraic in the
original stationary source; none of its identities assumes unit
$\zeta$. We retain that identity and change its source normalization.

\paragraph{Opening-cleared source bounds.}
Put
\[
 a=1/\zeta,\quad x=aL,\quad y=\eta L,\quad
 \theta=1-h=L\pi,\quad B=\beta+\tfrac12,\quad D_p=D+\tfrac12,
 \quad \rho_0=r+\eta.
\]
The original beta numerator gives bounded $\pi>0$;
$h,M,E,Q,P_d,P_4\sim1$ as in SB1--SB3. In particular
$\rho_0$ is a positive bounded unit. All formulas here preserve
actual $j>0$, rather than replacing it by $j_0$.
Define, exactly as in JF1,
\[
\begin{gathered}
 q=\alpha+(m+B)\theta^2,\quad \mu=m-a^2q,\quad
 \lambda=\mu/(\eta a^2L),\\
 T_0=A+\alpha+\tfrac12,\quad T_m=T_0+m,\quad
 e_4=r^2T_m/\eta^2,\\
 b'=Ar/\eta-m\theta-\alpha-\tfrac12+rT_mL/2,\\
 \kappa=2[B\theta-D_pj/\eta-mr/\eta],\qquad
 \gamma=2b'/h+2D_pj/r-2\eta B\theta/r.
\end{gathered}                                                    \tag{EU1}
\]
Finite canceled source division gives
\[
\begin{gathered}
 C(1-r)\sim\rho_0/2,\quad Ar\sim\eta/2,\quad
 D_pj\sim\rho_0/2,\quad m/(a^2L)\sim\rho_0/2,\\
 \alpha/L^2\sim\rho_0/2,\quad B\theta/L\sim\rho_0/2,
 \quad q/L^2-\rho_0(1+\pi)/2\in\operatorname{Inf},\\
 \mu/(a^2L)\sim\rho_0/2,\quad rT_m\sim\rho_0/2,
 \quad b'/L\sim-\rho_0/4,\quad \eta\kappa\sim-\rho_0.
\end{gathered}                                                    \tag{EU2}
\]
The canceled denominators at $a=L=0$ are only $1,2,4$;
there is no $r$, $1-r$, $j$, $\theta$ or $j_0$ factor.
Thus these are uniform through every opening in the theorem.
In particular $C,A,D_p$ are bounded below by positive real constants,
although any of them may be unbounded. Also $q/L^2$ and $\lambda$
are positive bounded units, $\mu>0$, and $b'=O(L)$.
The positive lower bound for $D_p$ uses bounded positive actual $j$
and the third estimate, not positivity of $j_0$.

\paragraph{Use the complete homogeneous parent.}
For nonnegative testing angles $u,v$ put $U=u^2,V=v^2$ and define
\[
\begin{gathered}
 E_t=\eta uW-vZ,\quad Q_t=Y-rvZ/\eta,\quad
 P=E_t/\eta-Q_t,\quad \xi=X+huZ,\\
 R_D=v(U-x)Z,\quad T_D=R_D-uv\xi/h,\quad
 H=avW-uZ-a\xi,\\
 F_b=(V-y)Z+vE_t,\qquad
 F_t=(V-y)Z+(\eta/r)vQ_t.
\end{gathered}                                                    \tag{EU3}
\]
All inverses in these roots are pure source. The exact full parent is
\[
\begin{split}
\mathcal P={}&qH^2+\mathscr C
 +CP^2+AQ_t^2+D_p(P-hT_D)^2+2b'Q_tT_D+\mathscr R\\
 &+e_4UF_t^2-e_4yu(H+a\xi)F_t
 +AUQ_t^2+B(1+V)\xi^2-2(\mu/a)H\xi-\mu\xi^2\\
 &+\kappa uF_t\xi+\gamma uvQ_t\xi .
\end{split}                                                       \tag{EU4}
\]
This is the complete JF1 identity, including every hard row and every
first normal. Set
\[
 \rho=m+rT_mx/2=O(x),\qquad
 K_r=T_0+2\rho-\rho^2/C.
\]
Then $K_r/T_m\sim1$ and $K_r$ is bounded below by a positive
real constant. The complete radial row is exactly
\[
\mathscr R=R_b^2+R_t^2,\qquad
 R_b={\sqrt C\over\eta}[F_b-r(1-\rho/C)F_t],\qquad
 R_t={r\sqrt{K_r}\over\eta}F_t.                         \tag{EU5}
\]
It follows that
\[
 {F_b\over\eta}={R_b\over\sqrt C}
               +{1-\rho/C\over\sqrt{K_r}}R_t          \tag{EU6}
\]
is a bounded pure-source allocation even when $r$ is infinitesimal.
No independent bare $(V-y)Z$ pivot has been inserted.

\paragraph{A strict source frame without either gap inverse bound.}
Use the three original weighted triangle roots
\[
 N_C=\sqrt C\,P,\quad N_A=\sqrt A\,Q_t,\quad
 N_D=\sqrt{D_p}(P-hT_D).                               \tag{EU7}
\]
Their positive squares are exactly the three original source-triangle
terms in EU4. They give the bounded pure-source allocations
\[
 P={N_C\over\sqrt C},\quad Q_t={N_A\over\sqrt A},\quad
 T_D={N_C/\sqrt C-N_D/\sqrt{D_p}\over h},\quad
 E_t=\eta(N_C/\sqrt C+N_A/\sqrt A).                    \tag{EU8}
\]
Consequently the entire $2b'Q_tT_D$ term is an infinitesimal Gram
on this strict diagonal frame. This does not cap or remove a divergent
original event and does not assert boundedness of its coefficient.

\paragraph{The complete two-normal cancellation.}
The source definitions in EU1 give the exact identity
\[
             {\eta\kappa\over r}+\gamma=2b'/h-2m.
\]
Since $F_t=F_b-vE_t+(\eta/r)vQ_t$, both complete normal rows in EU4
therefore satisfy
\[
\boxed{\quad
 \kappa uF_t\xi+\gamma uvQ_t\xi
 =\kappa uF_b\xi
     +uv\{(2b'/h-2m)Q_t-\kappa E_t\}\xi .\quad}       \tag{EU9}
\]
This is the crucial simultaneous allocation. It cancels the potentially
unbounded $1/r$ coefficients before any estimate is made. It neither
estimates the two original rows separately nor uses their contact
vanishing as a substitute for the identity.

Set $N_x=\sqrt B\xi$ and $\mathcal Q=\sqrt q H$.
By EU2, EU6 and EU8, every coefficient of EU9 on
$(N_C,N_A,R_b,R_t,N_x)$ is infinitesimal: the first uses
$\eta\kappa=O(1)$ and an extra $u$; the second uses
$b',m\in\operatorname{Inf}$ and an extra $uv$; the third uses
$\eta\kappa=O(1)$ and an extra $uv$.
The factors $1/\sqrt B$, $1/\sqrt A$, $1/\sqrt C$ and
$1/\sqrt{K_r}$ are bounded. No testing-state normalization occurs.

\paragraph{Every remaining scalar and hard term.}
The complete angular block on $(\mathcal Q,N_x)$ is
\[
 \mathcal Q^2-{2\mu\over a\sqrt{qB}}\mathcal QN_x
                         +(1-\mu/B)N_x^2.
\]
Its off-diagonal coefficient is $O(a)$ and $\mu/B$ is infinitesimal:
\[
 {\mu\over a\sqrt{qB}}
   =a\,{\mu/(a^2L)\over\sqrt{q/L^2}\sqrt B}.
\]
Thus its residue is the identity, including literal $\beta=0$.
The radial-angular cross has the exact normalized expression
\[
 -e_4yu(H+a\xi)F_t
 =-{rT_mL\over\sqrt{K_r}}u
       \left({\mathcal Q\over\sqrt q}
                         +{aN_x\over\sqrt B}\right)R_t.          \tag{EU10}
\]
Both coefficients are infinitesimal because $rT_m=O(1)$,
$L/\sqrt q=O(1)$ and $u$ is infinitesimal.

The coupled circuit $\mathscr C$ is exactly the sum of the following
three whole positive squares, followed by $-\lambda xE_t^2$:
\[
 \sqrt\lambda(uv-x/c)Z,\qquad
 \sqrt\lambda\,\eta k(cv-u)W,\qquad
 \sqrt{2\lambda/c}\,k\sqrt{uv}(\eta cW-Z),
 \qquad c^2=a/\eta,\quad k^2=x.                       \tag{EU11}
\]
The last negative term is an infinitesimal Gram on EU8 since
$\lambda$ is bounded and $x$ infinitesimal. All radicals in EU11
are separated pure-source and pure-testing radicals.

Retain EU11's three squares, $e_4UF_t^2$, $AUQ_t^2$ and
$BV\xi^2$ as whole positive terms. Every other term of EU4 has now
been assigned explicitly to the seven-root frame
\[
                 (N_C,N_A,N_D,R_b,R_t,\mathcal Q,N_x).             \tag{EU12}
\]
The Gram on this frame has identity residue and bounded coefficients.
The bounded-Gram lemma proves finite tensor SOS for the complete
parent whenever both testing angles are infinitesimal. Every estimate
is a product of a source bound and a testing bound; there is no
source/testing order comparison and no inverse of a mixed tensor.

\paragraph{All tests, endpoints and states.}
The roots and the whole-parent identity are homogeneous in all four
state coordinates. In particular no $Z$ inverse or state restriction
was used. The positive circuit rows EU11 are defined on literal
$u=0$ or $v=0$ and their squared identity is polynomial there.
The inherited exact sign gauges give all local sign charts.
The unchanged small-$L$, bounded-$\eta$ remote theorem covers the
complement of both infinitesimal testing angles on this exact source
domain, retaining all its projective endpoints and openings.
Applying the real homogeneous identity to real and imaginary state
parts proves the complex-state assertion. All of the source unit-$\eta$
portion of C14.3a is therefore all-tests admitted.

\paragraph{The complementary source.}
The remaining C14.3a sources have $\eta$ infinitesimal, hence $r\sim1$, with all original $1-r,j$ openings retained. Their complete parent is treated by the LD identities and the completed $J$-row conductor in Sections~\ref{component:pc720c1e06aca} and~\ref{component:pe3bbfd3a8151}.


\section{\texorpdfstring{C14.1b: the entire parent on bounded testing state quotients}{C14.1b: the entire parent on bounded testing state quotients}}
\label{component:pff0a29ebf284}

\label{pff0a29ebf284:sec:r17-c1-bounded-state}

\begin{theorem}
Retain every original C14.1b source hypothesis and intersection:
$x,z,\sigma>0$ are infinitesimal, $H=h/\sigma$ is unbounded,
$0<h<1$, and all original stationary-contact, strict-gap,
coefficient-sign and residual-complement conditions remain in force.
Normalize $d=1$ by its positive pure-source value, and retain the
original positive-infinitesimal total-scale intersection.
For every independently chosen bounded $u$, infinitesimal or zero $w$,
and real testing state with $Z\ne0$ and $f=W/Z$ bounded, the complete
original reciprocal parent is a finite tensor SOS. Both testing signs
and zero angles are included. The literal $Z=0$ chart is covered by
the original-core sign exit.
\end{theorem}

This is a full-parent state-subchart admission, not a complete C14.1b
source admission. After the sign exit, the remaining real states have
$WZ>0$ and unbounded $W/Z$. A complex state is covered here when
each of its real and imaginary parts satisfies this state chart or
the sign exit. No assertion for other complex states is implicit.

\paragraph{Exact parent entry.}
Use SP1--SP6 of the companion sphere-pole reduction, including all its
generic JR/JB coefficients and retained centered excesses. In particular
$a_*=\ell_*=1/4$ are the fixed JR1 caps; their exact nonnegative
centered excesses remain in the parent. Also
$P_\beta=uZ-(x/\sigma)wW$, $X_e=X+(xh/\sigma)wW$,
$A_\theta=q(u^2-x)W$, $q^2=z$, and $R=wZ/q-uqW$.
Set
\[
 s=x/\sigma,\quad L_h=\lambda/H,\quad
 K_A=4\kappa H-\lambda,\quad \rho_\lambda=\lambda x/\sigma.
\]
The SP3 identity $\lambda=2H+\lambda_0$, with $\lambda_0$ bounded,
implies $\lambda>0$ unbounded and $L_h\sim2$.
Also $s$ is infinitesimal by SP11, and
$\rho_\lambda=2xh/\sigma^2+\lambda_0s$ is bounded and positive.
The exact source cancellation SP10 makes $K_A$ bounded. More precisely,
\[
 \operatorname{res}K_A=2-4\operatorname{res}h,\qquad
 \operatorname{res}\chi=4(1-\operatorname{res}h).
 \tag{BS1}
\]
These follow from $\alpha_\kappa=2H(2\kappa-1)$,
$K_A=\alpha_\kappa-\lambda_0$:
$\operatorname{res}\alpha_\kappa=-2\operatorname{res}h$ and
$\operatorname{res}\lambda_0=-2(1-\operatorname{res}h)$.
The polynomial source divisibility in SP10 proves these assertions
uniformly in $h$, rather than multiplying an unspecified infinitesimal
error by $H$.

Use SP7 with the \emph{whole} coefficient $\lambda$. The last line of
SP6 becomes exactly
\[
 K_A A_\theta^2+\chi uA_\theta R+
 \underbrace{\lambda\sigma P_\beta^2+
       \rho_\lambda(u^2-x)(w^2-\sigma z)W^2}_{\mathscr P}.
 \tag{BS2}
\]
This step retains the signed contact product; it has not discarded
one of its negative terms.

\paragraph{An exact bounded conductor for the contact product.}
All following state quotients belong solely to the testing field.
Put
\[
 U_h=\sqrt h\,uZ,\qquad
 Q_\lambda=\sqrt{\rho_\lambda}\,uwW,\qquad
 M_\lambda=\sqrt\lambda\,xqW.
\]
The roots are homogeneous in the testing state and explicitly separated.
At $Z\ne0$, the identity $W=fZ$ gives
\[
 \sqrt h\,P_\beta=(1-szf^2)U_h-\sqrt h\,sqfR.
 \tag{BS3}
\]
Expanding all four terms of the contact product gives the exact identity
\[
\begin{split}
 \mathscr P={}&L_h\{(1-szf^2)U_h-\sqrt h\,sqfR\}^2
                  +Q_\lambda^2+M_\lambda^2\\
 &-L_hszf^2U_h^2
       -L_hs^2\{\sqrt h\,qfR+zf^2U_h\}^2 .
\end{split}                                                   \tag{BS4}
\]
Both negative lines are retained exactly. Their coefficients, expanded
as a Gram on $(U_h,R)$, are infinitesimal sums of separated tensors:
$s,z,q$ are source infinitesimals, $\sqrt h$ is bounded, and $f$
is bounded pure testing. The first square has residue $2U_h^2$.
No inverse of $1-szf^2$, or of any other mixed quantity, is used.
The two remaining displayed squares are retained with their exact
positive source weights, through either sphere endpoint.

\paragraph{Every coupling to the collapsing Gamma direction.}
The only $\Gamma$ cross terms in SP6 are
$-4\nu N_x(u\Gamma)$ and
$R_\gamma(N_y+bR)(qw\Gamma)$. They both have exact conductors.
Define the bounded infinitesimal separated expression
\[
 B_f=(\sigma-xzf^2)/\sqrt h.
\]
Indeed $\sigma/\sqrt h=\sqrt{\sigma}\sqrt{\sigma/h}$ is infinitesimal
because $H$ is unbounded, and
$x/\sqrt h=(x/\sigma)(\sigma/\sqrt h)$ is infinitesimal.
Direct substitution of the original roots gives
\[
 u\Gamma=B_fU_h-uqfA_\theta-u^2qfR,                 \tag{BS5}
\]
\[
 qw\Gamma=qzfB_fU_h-uz^2f^2A_\theta
                  +z(\sigma-ufw-u^2zf^2)R.         \tag{BS6}
\]
Every coefficient on the right is infinitesimal and separated for
bounded $u$ and $f$, including arbitrary relative source and testing
rates. Thus both full first-normal couplings, and the $R\Gamma$
coupling, become infinitesimal bounded Gram entries on
$(N_x,N_y,U_h,A_\theta,R)$. No free positive pivot on $\Gamma$
has been imposed.

The entire remaining diagonal in $\Gamma$ is
\[
 \left(\frac{4D}{\sigma^2}+4bpzu^2\right)\Gamma^2,
 \tag{BS7}
\]
a sum of two exact squares with nonnegative pure-source weights.
It is retained separately. The other already-positive extracted square
$2(2D+1)(uX_e)^2$, both centered cap excesses, and
$4\beta D_wN_x^2$ are also retained exactly.

\paragraph{All remaining full-parent entries.}
After BS2--BS7, the only other cross terms are exactly SP5 without
its two $\Gamma$ terms:
\[
\begin{split}
 &-4uqw(2m+2D+1)N_xN_y
 +uqwR_xN_xR+(c_N+zR_y)N_yR\\
 &\hspace{12mm}+qwB_xN_xA_\theta+uA_yN_yA_\theta .
\end{split}                                                   \tag{BS8}
\]
Their coefficients are infinitesimal. The opening-uniform source
bounds are $p,l\sim1/2$, $p-l\in z\mathcal O_c$,
$m,D,b_\theta,C$ bounded, $B_x=8\kappa e$ infinitesimal, and
$A_y=O(z)$; $u$ and $f$ are bounded and $w$ infinitesimal.
The three corresponding hard/radial diagonals have residues
\[
 [2(1+w^2)+4mw^2]_0=2,\quad
 (H_y)_0=2(1+u_0^2),\quad
 (c_G+zK_r)_0=2(1+u_0^2).                          \tag{BS9}
\]
Thus no unallocated full-parent entry is hidden in the scalar argument.

\paragraph{Two exhaustive pure-source cases.}
If $h$ is infinitesimal, BS1 gives $(K_A)_0=2$ and $\chi_0=4$.
Retain $Q_\lambda^2+M_\lambda^2$ outside the remaining Gram.
In the root order $(N_x,N_y,U_h,A_\theta,R)$ its exact bounded
Gram has residue
\[
 \operatorname{diag}(2,2(1+u_0^2),2)\ \oplus\
 \begin{pmatrix}2&2u_0\\2u_0&2(1+u_0^2)\end{pmatrix}.
 \tag{BS10}
\]
The last determinant is exactly $4$, so this residue is strictly
positive for every bounded testing residue $u_0$, including zero.

If $h$ is a positive unit, the potentially signed $A_\theta$ diagonal
is instead removed using the exact root identity
\[
 A_\theta=\frac{qfu}{\sqrt h}\,U_h
                       -\frac1{\sqrt\lambda}M_\lambda.        \tag{BS11}
\]
Both coefficients are infinitesimal. Keep $M_\lambda^2$ in the Gram
and retain $Q_\lambda^2$ separately. Since $K_A,\chi,u$ are bounded,
all the $A_\theta$ terms now contribute only infinitesimal entries.
In the root order $(N_x,N_y,U_h,R,M_\lambda)$ the complete residue is
\[
 \operatorname{diag}(2,2(1+u_0^2),2,2(1+u_0^2),1)\succ0.
 \tag{BS12}
\]
This includes the $\theta=1-h$ opening. The two cases require no
comparison of a testing parameter with $h$.

The finite bounded positive-Gram criterion supplies legal tensor SOS
in either case. Its strict real residue can depend on the real number
$u_0$. Every error is a finite sum of separated bounded factors with
an infinitesimal source or testing factor; the usual finite two-factor
coefficient splitting applies. There is no mixed Cholesky operation.

\paragraph{Endpoints, signs, and scope.}
Every identity used above is exact at $u=0$ or $w=0$. The sources
$D=0$ and either sphere endpoint were never inverted. At the additional
$D=0$ null point $u=w=W=X=Y=0$, all roots in the strict Gram vanish
and the retained BS7 contribution is exactly $4DZ^2$.
Thus the certificate respects that null point and does not restore
the invalid old soft-coordinate cap.

The argument itself allows signed bounded $f$. For $Z=0$ the
original-core sign lemma applies since $uwWZ=0$; all original event
weights remain nonnegative pure source. That exit also covers
$uwWZ\le0$ regardless of the size of $W/Z$. After the exact testing
sign gauges, the only real-state complement on the remaining tests
is therefore $WZ>0$ and unbounded $W/Z$.
For completeness the reciprocal core at $Z=0$ is the manifest sum
$(uwW+uX-wY)^2+D_uY^2+D_wX^2$.
At the literal reciprocal endpoint $w=0$ it is exactly
$K_w=D_u(X^2+Y^2)$. Thus this endpoint needs no finite-$v$
Schur formula or continuity argument.

The focused checker verifies BS2--BS8 as an identity of the entire
SP6 parent, both Gamma conductors, the contact-product expansion,
BS11, both strict residue calculations, and the exact extra zero.
It does not assert that the remaining unbounded-state quotient is
covered, or count this result as a whole-source admission.


\section{\texorpdfstring{A full-parent pole extraction and a larger comparison-model state chart}{A full-parent pole extraction and a larger comparison-model state chart}}
\label{component:pfeca652aef94}

\label{pfeca652aef94:sec:r17-c1-sphere}

\paragraph{Status.}
The first result below is a new exact reduction of the complete generic
JR5/JB1 parent. On C14.1b it removes every unbounded $h/\sigma$
coefficient from that parent and leaves one explicitly stated bounded
contact-product conductor problem. It is not a positivity proof of the
resulting bounded form. The second result extends the R14 JU8
comparison-model certificate to a strictly larger pure-testing state
chart and to $Z=0$. Neither result admits a complete C1 source branch.

\paragraph{Notation and imported exact identity.}
Use the generic JR/JB coefficients, the exact retained centered excesses,
and the source BE2--BE8 definitions. Put
\[
\begin{gathered}
 t=u^2,\quad \xi=w^2,\quad \Delta=t-x,\quad
 \delta=\xi-\sigma z,\quad q^2=z,\\
 H=h/\sigma,\quad e=xh/\sigma,\quad b_\theta=\beta\theta,
 \quad\theta=1-h.
\end{gathered}
\]
Here $p=ab$, $l=\ell(1-b)$ are the opening-cleared source products,
and $\mu=m/(x^2z)=\kappa H$ is an \emph{exact} consequence of BE7--BE8.
The complete parent is
$F_w=E_{\rm cap}+4\beta D_wN_x^2+G$, with $G$ exactly JR5/JB1.
No state normal or retained excess is removed below.

Define the separated, denominator-free-in-the-tests roots
\[
 P_\beta=uZ-\frac{x}{\sigma}wW,\qquad
 X_e=X+ewW,\qquad T=t\Gamma+V.
\]
They obey the exact identities
\[
 T=\sigma uP_\beta,\qquad N_x-hP_\beta=X_e,
 \tag{SP1}
\]
and, using JR6,
\[
 t^2\Gamma^2+(t+x)\Gamma V+\sigma tA_\theta^2
 =T^2-u\sigma\Delta A_\theta R.
 \tag{SP2}
\]
In particular the \emph{whole} quadratic $h^2$ packet and its
$N_x$ first normal combine with the corresponding $N_x^2$ term into
$2(2D+1)(uX_e)^2$. This identity includes $u=0$, and does not invert it.

Put
\begin{align*}
 \nu&=\frac{2D-2b_\theta+h}{\sigma},\\
 \lambda&=\frac{8Dh+4b_\theta\theta+2h^2}{\sigma}
                    -4bp\sigma z=2H+\lambda_0,\\
 \lambda_0&=\frac{4D(1+h)}{\sigma}-2\theta\nu-4bp\sigma z,\\
 \chi&=8(pb+l(1-b))-4h+4eh(2D+1)-8Dh\\
 &\hspace{8mm}+4z\{4Dbh+4pb\sigma+(1+2m)(h-1+b)^2
                                      +2b_\theta\theta\}.
 \tag{SP3}
\end{align*}
The complete coefficient $B$ of JR/JB satisfies exactly
\[
 \sigma B-\frac{2h^2(2D+1)\Delta}{\sigma}=\lambda+\chi.
 \tag{SP4}
\]

For compactness retain the following explicit, full collection of
unchanged cross terms:
\begin{align*}
 \mathscr C={}&-4uqw(2m+2D+1)N_xN_y
 +qwR_\gamma N_y\Gamma+qwbR_\gamma R\Gamma\\
 &+uqwR_xN_xR+(c_N+zR_y)N_yR
 +qwB_xN_xA_\theta+uA_yN_yA_\theta.
 \tag{SP5}
\end{align*}
Every generic coefficient in SP5 is its complete JR/JB value. In
particular this line keeps both hard normals and every radial first
normal, including the generic beta correction in $R_x$.

The entire parent core, with no restriction on states, is exactly
\begin{align*}
G={}&2(2D+1)(uX_e)^2
 +[2(1+\xi)+4m\xi]N_x^2+H_yN_y^2\\
 &+\left(\frac{4D}{\sigma^2}+4bpzt\right)\Gamma^2
 +(c_G+zK_r)R^2-4u\nu N_x\Gamma+\mathscr C\\
 &+4\kappa H A_\theta^2
 +\lambda u(\Gamma P_\beta+A_\theta R)
 +\chi uA_\theta R.
 \tag{SP6}
\end{align*}
This first form already isolates all possible $H$ poles in one
shared scalar packet. It is an identity, not an entrywise comparison.

\paragraph{A single signed scalar contact product.}
There is a second exact conductor,
\[
 A_\theta^2+u\Gamma P_\beta+uA_\theta R
 =\sigma P_\beta^2+
       \frac{x}{\sigma}\Delta\delta W^2.
 \tag{SP7}
\]
Consequently SP6 can be written
\begin{align*}
 G={}&2(2D+1)(uX_e)^2+2hP_\beta^2+\mathscr B
               +\rho\Delta\delta W^2,\\
 \rho={}&\frac{2xh}{\sigma^2},\\
 \mathscr B={}&[2(1+\xi)+4m\xi]N_x^2+H_yN_y^2
 +\left(\frac{4D}{\sigma^2}+4bpzt\right)\Gamma^2
 +(c_G+zK_r)R^2\\
 &-4u\nu N_x\Gamma+\mathscr C+\alpha_\kappa A_\theta^2
 +\lambda_0u(\Gamma P_\beta+A_\theta R)+\chi uA_\theta R,\\
 \alpha_\kappa={}&2H(2\kappa-1).
 \tag{SP8}
\end{align*}
This is the promised exact full-parent reduction. The two squares
displayed before $\mathscr B$ have nonnegative pure-source weights.
The term $\rho\Delta\delta W^2$ is signed; it has not been called a
square or bounded below by a mixed-rate comparison.

\paragraph{Opening-uniform source bounds on C14.1b.}
The cancellation in the hard first normal is essential. BE3 gives
\[
 \nu=\frac{2D}{\sigma}
       +\frac{H_0-ez(J-\sigma)}{P}.
 \tag{SP9}
\]
Thus $\nu$ is bounded even when $H$ is unbounded. Equations SP3 show
that $\lambda_0$ and $\chi$ are bounded, since $D=O(\sigma^2)$,
$e$ is infinitesimal, $0<h<1$, and $b_\theta=O(h+\sigma)$.
No lower bound on $h$ or $\theta$ is used.

The apparently unbounded coefficient $\alpha_\kappa$ is bounded as
well, by an exact divisibility. Let
\[
 N_\kappa=(1+x)(1+\sigma bz)(1+\sigma z)-MP,
 \quad F_x=\frac{N_\kappa|_{\sigma=0}}x,
 \quad F_\sigma=\frac{N_\kappa-N_\kappa|_{\sigma=0}}\sigma.
\]
Both displayed quotients are polynomials in the bounded source
letters $x,z,b,h,\sigma$. The source $M$ here is the geometry factor
$1+xbz-xz\theta$, not $M_\beta$. Exactly,
\[
 \alpha_\kappa=\frac{2(hF_\sigma+eF_x)}{MP}.
 \tag{SP10}
\]
Since $M,P$ are positive units this proves the asserted bound without
using a rate comparison between a source and a test.

On C14.1b, $H$ unbounded and BE7 force $\sigma/x$ unbounded:
if it were bounded then both terms of BE7 would be bounded.
Thus $x/\sigma$ is infinitesimal. The source sphere gives
\[
 \rho=\frac{2xH_0}{\sigma B}
                  +\frac{2\mathcal H}{B}b_0^2,
 \tag{SP11}
\]
so $\rho$ is bounded through both sphere endpoints. Its residue is
$2\operatorname{res}(b_0^2)$ in the large-$L$ chart.
The coefficients in SP8, including the coefficients in SP5, are now
all bounded for bounded $u$ and infinitesimal $w$. Here
$C=O(z)$, $p,l$ are bounded and $p-l\in z\mathcal O_c$ uniformly
through $b$ and $1-b$ openings; the latter follows from the same
opening-cleared source numerators used in JR, written with $e$ rather
than imposing bounded $H$. Also $B_x=8\mu x=8\kappa e$ is bounded.

One may introduce the legal root
$B_\delta=\delta W/q$. Then the last term of SP8 is exactly
$\rho A_\theta B_\delta$. The remaining conductor problem is now
precise: allocate this bounded signed cross term together with
$\mathscr B$ and the two retained squares, using
\[
 \frac{xq}{\sigma}B_\delta
       =uqR+qA_\theta-wP_\beta,
 \qquad u\delta W=\sigma qR-w\Gamma.
 \tag{SP12}
\]
Neither relation permits cancellation of $u$, $w$, a state coordinate,
or the infinitesimal $x/\sigma$ while retaining a bounded coefficient
array. In particular no $B_\delta^2$ pivot has been proved. This is
an exact bounded full-parent reduction on C14.1b, not its admission.
Nor is $\mathscr B$ asserted to be positive by itself: for example,
at the algebraic specialization $z=0$ one has exactly
$\alpha_\kappa=-2h/(1+x+\sigma)$, whose residue is negative when
$h$ is a unit. Thus the two extracted squares, $\mathscr B$, and the
signed contact product still require a joint allocation.

The algebraic identities SP1--SP12 also apply in the common BE source
domain with $L$ a positive unit. The preceding claim that every entry
of SP5 is bounded uses large $L$; on C14.2 the separate $r\to1$
opening can make $C$ unbounded and is not silently removed here.
On C14.1a SP8 remains exact but SP11 does not in general bound
$\rho$; the following result addresses its comparison model instead.

\paragraph{A larger state chart for the complete JU8 model.}
Retain the C14.1a assumptions and notation of R14.1--R14.2:
$u$ infinitesimal, $\tau$ bounded, $\varepsilon$ infinitesimal,
$x=c\varepsilon$, $\sigma=\tau\varepsilon$, and $0<c\le1/2$.
Suppose $Z\ne0$ and the \emph{pure-testing} quotient
$uW/Z$ is bounded. Unlike R14, $W/Z$ itself may be unbounded.
Put $E=qW/Z$, $U=|u|$, $V=|E|$, $a=UV$, and use the exact
circuit roots $d_i=A_i-G$ of R14.2. Then $uE=q(uW/Z)$ is
infinitesimal, although no size is assigned to $E$.

When $c$ is infinitesimal, R14.3 still applies because its coefficient
is $uE/d_0$, which is infinitesimal on this larger chart.
When $c$ is a positive unit, the new identity is
\[
 uEK=\operatorname{sgn}(uE)\left\{
       -\frac{a}{\sqrt{2c}}d_1
       +\frac{a}{\sqrt{2c}}d_2
       +\frac{U\sqrt\tau}{\sqrt2\,c}d_3\right\},
 \qquad K=\sigma-u^2E^2.
 \tag{SP13}
\]
Every coefficient is infinitesimal. Crucially, the coefficient
$V\sqrt\varepsilon$ in R14.4 has disappeared; no bound on $V$ is
needed. The identity follows by substituting
$A_1=\sqrt{2c}U^2V^2$, $A_2=\sqrt{2\sigma\tau}U$,
$A_3=\sqrt{2\tau}xV$, $G=\sqrt{2\tau x}UV$,
and $\sigma c=\tau x$.

Substitute SP13 (or R14.3) in the exact R14.1 identity. Its radial
diagonal is $1+u^2E^2$ of residue one; its other diagonal entries
are the identity on the same fourteen circuit roots. The remaining
coupling is $\tau uA_\theta R$, with infinitesimal coefficient.
The finite bounded positive-Gram criterion therefore proves the
\emph{whole JU8 model}, including its first radial normal, on this
larger chart. All radicals are taken on their explicitly separated
pure-source and pure-testing factors, as in R14.2.

At the omitted literal chart $Z=0$, the whole model has the direct SOS
\[
 \frac{\mathcal L}{2\kappa_s}
   =(uqW)^2+(uwW)^2+(1+2\tau)\{q(u^2-x)W\}^2.
 \tag{SP14}
\]
The sole still-uncovered model state chart is therefore
$uW\ne0$ and $0<|Z/(uW)|$ infinitesimal, a comparison entirely in
the testing field. The full JU3 correction is still unallocated
even on the newly covered model chart. No full-parent admission is
inferred from SP13--SP14.

\paragraph{Verification and boundaries.}
The accompanying checker verifies the complete SP6 and SP8 polynomial
identities relative to the imported full generic JR/JB parent, every
new conductor and source cancellation, the polynomial divisions in
SP10, and SP13--SP14. It does not replay the old contact-jet proof or
turn a symbolic equality into a positivity theorem. All identities
include zero tests and all signs, retain the $D=0$ extra zero, and
apply separately to real and imaginary state parts. In particular,
at $u=w=W=X=Y=0$ SP8 gives exactly $4DZ^2$.


\section{\texorpdfstring{The damped completion at large zeta with bounded L over eta}{The damped completion at large zeta with bounded L over eta}}
\label{component:pc720c1e06aca}

\label{pc720c1e06aca:sec:r17-large-zeta-damped}

\begin{theorem}
Normalize $d=1$ and retain the unchanged physical C14.3a source,
its actual stationary contact, every strict source gap,
$\beta\ge0$ including literal $\beta=0$, the original total-scale
intersection and every earlier exclusion. Suppose $x,L=x\zeta,\eta$
are positive infinitesimals, $\zeta$ is unbounded, and
\[
                         L/\eta\text{ is bounded}.       \tag{LD1}
\]
Then the entire original parent is finite tensor SOS at every independent
projective angle and state test, real or complex. This is a whole-state,
all-tests source-subchart admission. No condition is imposed on the
testing angle ratio or on the source ratio $a/\eta$, where $a=1/\zeta$.
The source subchart with $\eta$ infinitesimal and $L/\eta$ unbounded
is not admitted here.
\end{theorem}

The proof uses the exact complete damped identity DG3, but its source
scale is $L+\eta$ rather than $x+\eta$. This change must be proved:
one cannot invoke the unit-$\zeta$ estimates by substituting an
infinitesimal $a$ into their displayed $1/a$ coefficients. All formulas
below use the actual physical source and the original contact.

\paragraph{Source variables and opening-cleared coefficients.}
Put
\[
 a=1/\zeta,\quad x=aL,\quad y=\eta L,\quad
 h=1-L\pi,\quad r=1-\eta\sigma,\quad B=\beta+\tfrac12.
\]
The original source estimates give bounded $\pi,\sigma\ge0$, and
actual $j>0$ says
$0<\sigma<1+xh$, $j=\eta(1+xh-\sigma)>0$.
Neither opening has a positive-residue lower bound.
The exact native source formulas displayed below give
\[
 \frac{m}{a^2L}\sim\frac12,\quad
 \frac{\alpha}{L^2}\sim\frac12,\quad
 \frac{q}{L^2}-\frac{1+\pi}{2}\in\operatorname{Inf},\quad
 \frac{\mu}{a^2L}\sim\frac12,
\]
\[
 \eta\lambda\sim\frac12,\quad
 \eta K\sim\frac12,\quad
 \eta\kappa\sim-1,\quad \eta^2e_4\sim\frac12,
 \quad b'=O(L),                                     \tag{LD2}
\]
where $q=\alpha+(m+B)(1-h)^2$, $\mu=m-a^2q$,
$\lambda=\mu/(\eta a^2L)$ and $K=CD_p/(C+D_p)$.

Here is a completely regular specification of the new source
expressions. Write
\[
 K_x=1+x,\ R_s=1+xr,\ T_s=1+xh,\ M_s=R_s-aL\pi,
 \ E_s=1-x\eta h,\ Q_s=1-Lr,
\]
\[
 P_d=1+2xh+rx^2h+rxL+L
                  +\eta\{K_xL+x^2hL\pi\},\qquad
 P_4=T_sE_s+(1-r)(xh+L).
\]
Define the opening-cleared quantities
\[
 C_n=\frac{K_x(r+\eta)E_s}{2P_d},\qquad
 D_j=\frac{K_x(1+y)(r+\eta)}{2P_d},\qquad
 A_\eta=\frac{P_4}{2rP_d},
\]
\[
 m_n=\frac{K_x(r+\eta)h(1+y)}{2M_sP_d},\quad
 \alpha_n=\frac{(r+\eta)\{aR_sh(1+y)+M_sE_s\}}
                         {2Q_sP_d},\quad
 B_n=\frac{K_x(r+\eta)(1+ah)}{2P_d}.                 \tag{LD3}
\]
Thus $C_n=(1-r)C$, $D_j=jD_p$, $A_\eta=A/\eta$,
$m_n=m/(a^2L)$, $\alpha_n=\alpha/L^2$ and
$B_n=B(1-h)/L$. In particular
\[
 q_n:=q/L^2=\alpha_n+m\pi^2+B_n\pi,\qquad
 \Lambda:=\eta\lambda=m_n-Lq_n,
\]
\[
 \eta K=\frac{C_nD_j}
                  {C_n(T_s-\sigma)+D_j\sigma}.       \tag{LD4}
\]
There is no $a^{-1}$ in (LD3)--(LD4). Their denominators at
$a=L=\eta=0$ have residues $1$, $2$ or $1/2$, independently
of $\pi,\sigma$. The remaining definitions are
\[
 T_m=A+\alpha+\tfrac12+m,\quad
 b'=A_\eta r-mL\pi-\alpha-\tfrac12+rT_mL/2,
\]
\[
 G_f:=\gamma h=2b'+2D_jh/r-2\eta LB_nh/r,
 \quad G_0=2\Lambda-2\eta\sim1,\quad b_0=b'-G_0/2.
                                                               \tag{LD5}
\]
These are the actual native-source expressions used in the completion.

\paragraph{Positive source deficits at the new scale.}
Set $s=L+\eta$ and retain DG's definitions
\[
 \Delta_0=h^2AK-b_0^2,
 \qquad c_U=A-\frac{G_0^2}{4\lambda},
 \qquad\epsilon=\gamma-G_0/h.
\]
The regular source formulas above give exact finite divisions
\[
 \Delta_0=Ld_L+\eta d_\eta,\qquad
 \frac{c_U}{\eta}=Le_L+\eta e_\eta,\qquad
 G_f-G_0=Lf_L+\eta f_\eta,                           \tag{LD6}
\]
with bounded coefficients and residues
\[
\begin{array}{c|cc}
 &L\text{ coefficient}&\eta\text{ coefficient}\\\hline
 \Delta_0&\pi/2+1/4&\sigma/2+3/4\\
 c_U/\eta&\pi+1/2&\sigma+3/2\\
 G_f-G_0&\pi+1/2&\sigma+2
\end{array}                                                     \tag{LD7}
\]
(as usual the bounded $\pi,\sigma$ on the right can be replaced
by their real residues). To justify the finite divisions explicitly,
each expression on the left of (LD6) vanishes at $L=\eta=0$
\emph{for arbitrary $a$}. Subtract its exact value at $L=0$, divide
that difference by $L$, and divide the $L=0$ value by $\eta$.
Clearing the source-unit denominators proves each divisibility in a
polynomial ring. Evaluation of these quotients at $a=L=\eta=0$
gives (LD7). There is consequently no hidden pure-$a$ remainder.

All four coefficients for the first two rows have strictly positive
residue. Hence
\[
 \Delta_0/s\text{ and }c_U/(\eta s)
       \text{ are positive bounded source units},\qquad
 \epsilon=O(s).                                      \tag{LD8}
\]
This source calculation itself does not use (LD1).

\paragraph{The same complete damped identity, without a state chart.}
Use DG's unshifted $H=avW-uZ$; the strict angular root is
$\mathcal Q=\sqrt q(H-a\xi)$, so this $H$ differs from the
shifted letter $H$ in EU3. Retain every other homogeneous JF/DG
definition. In particular
\[
 E=\eta(P+Q),\quad N_x=\sqrt B\xi,\quad
 \tau=\eta cW-Z,\quad D_0=u-cv,\quad c^2=a/\eta.
\]
After the pure-testing sign gauges, the half-core lemma supplies all
$WZ\le0$ states. On $WZ>0$, the exact identity DG3 is valid for
this source without any change:
\[
\begin{split}
 F/4={}&\mathscr A+\mathscr R+S_0
       -\lambda x(\eta P+\eta Q/\lambda)^2
       +\widehat R^2+\widehat J^2\\
 &+c_UUQ^2+e_4UF_t^2+BV\xi^2
       +\frac{2x}{c}uv\tau^2+\frac{G_0x}{c}D_0^2WZ\\
 &-e_4yuHF_t+\kappa uF_t\xi
       +\frac{G_0x}{c}PH+\epsilon uvQ\xi.             
\end{split}\tag{LD9}
\]
Here $S_0=CP^2+AQ^2+D_p(P-hT)^2+2b_0QT$,
$\widehat R=R+G_0\sqrt xQ/(2\sqrt\lambda)$ and
$\widehat J=\sqrt\lambda J+G_0uQ/(2\sqrt\lambda)$.
The identity conducts only $G_0uvQ\xi/h$ and completes both
circuit companions exactly, retaining every first normal and hard
coupling. Its derivation uses only $x=k^2$, $c^2=a/\eta$ and the
homogeneous JF conductor, not a bounded reciprocal of $a$ or $c$.

\paragraph{Strict source, radial and angular frames.}
Use the exact source roots
\[
 N_s=\sqrt{C+D_p}\left(P-\frac{hD_p}{C+D_p}T\right),
 \quad N_T=h\sqrt K\left(T+\frac{b_0}{h^2K}Q\right),
 \quad N_Q=\sqrt{\frac{\Delta_0}{h^2K}}Q.            \tag{LD10}
\]
Their squares sum to $S_0$, and
$N_Q/[\sqrt{\eta s}\,Q]$ is a positive source unit.
The inverse source identities give bounded coefficients for $P$
on this frame: the coefficients of $N_s,N_T$ are $O(\sqrt\eta)$
and that of $N_Q$ is $O(\sqrt{\eta/s})$. This uses only
$\eta\le s$, so no old $E/\sqrt\eta$ pivot is imported.

The complete radial packet has its inherited strict roots $R_b,f_t$,
with $f_t=F_t/\eta$. Indeed $\rho=m+rT_mx/2=O(x)$ and
$K_r\sim1/2$. The complete angular block is
\[
 \mathscr A=\mathcal Q^2-
       \frac{2\mu}{a\sqrt{qB}}\mathcal QN_x
                  +(1-\mu/B)N_x^2.
\]
It has identity residue, since
\[
 \frac{\mu}{a\sqrt{qB}}
       =\frac{a[\mu/(a^2L)]}{\sqrt{q/L^2}\sqrt B}=O(a),
 \qquad\mu/B\in\operatorname{Inf}.                 \tag{LD11}
\]
This includes unbounded $B$ and literal $\beta=0$.

\paragraph{The entire remainder has infinitesimal coefficients.}
The negative-square root in (LD9) has, on the source frame,
the same coefficient bounds as DG:
\[
 \eta\sqrt{\lambda x}\,P+
           \eta\sqrt{x/\lambda}\,Q,
 \qquad O(\sqrt{\eta x}),\quad
           O(\eta\sqrt{x/s}).                       \tag{LD12}
\]
They are infinitesimal; here $x=aL$ and no comparison of $a$ with
$\eta$ is needed. For the radial-angular term use
$H=\mathcal Q/\sqrt q+aN_x/\sqrt B$ and $F_t=\eta f_t$.
Since $e_4y\eta=(e_4\eta^2)L=O(L)$ and $\sqrt q\asymp L$,
its two coefficients are $O(u)$ and $O(xu)$.
The $\kappa$ term has coefficient
$(\eta\kappa)u/\sqrt B$, which is also infinitesimal.

The new $PH$ term improves at this source end:
\[
 \frac{G_0x}{c\sqrt q}=O(\sqrt{a\eta}),\qquad
 \frac{G_0ax}{c\sqrt B}=O(x\sqrt{a\eta}).            \tag{LD13}
\]
Thus the bounded $P$ allocation makes both of its angular coefficients
infinitesimal. Finally, the complete unprocessed normal is
\[
 \frac{\epsilon uv}{\sqrt{\eta s}\sqrt B}N_QN_x
   \quad\hbox{up to a bounded positive source-unit factor}.
\]
Its coefficient is
\[
                  O\!\left(\sqrt{1+L/\eta}\,uv\right),          \tag{LD14}
\]
which is infinitesimal under the bounded-source-ratio
hypothesis (LD1). This is the only use of that hypothesis.

Consequently the complete remainder on
$(N_s,N_T,N_Q,R_b,f_t,\mathcal Q,N_x)$ has a bounded strict-residue
Gram. The bounded-Gram lemma gives its finite tensor SOS. Every other
term in (LD9) is a whole positive square: in particular $c_U>0$,
$G_0>0$, and $uv,WZ$ are nonnegative pure-testing factors on this
sign chart. No mixed inverse or radical is taken.

\paragraph{All tests and exact remaining source end.}
No state normalization, state-coordinate inverse or testing ratio
occurs in this proof. The division-free circuit rows and
$\sqrt{uv}=0$ include both literal local axes. The half-core exit
and the exact angle sign gauges cover every real state; splitting
the real-coefficient Hermitian form covers complex states. The
unchanged small-$L$, bounded-$\eta$ remote theorem covers all other
projective testing pairs, including literal infinities. This proves
the theorem with every original source opening retained.

The restriction LD1 is used only in LD14. The complete LD9 identity, source bounds and all preceding normal allocations are inputs to the completed $J$-row conductor in Section~\ref{component:pe3bbfd3a8151}, which treats the full infinitesimal-$\eta$ source at arbitrary $L/\eta$. The unit-$\eta$ source is the EU theorem of Section~\ref{component:p60eca5442075}.


\section{\texorpdfstring{The completed $J$ row conducts the shared remaining normal}{The completed J row conducts the shared remaining normal}}
\label{component:pe3bbfd3a8151}

\label{pe3bbfd3a8151:sec:r18-completed-j}

\paragraph{Scope and unchanged inputs.}
Retain the complete actual-parent identity DG3, equivalently LD9,
with $d=1$, every original source opening and residual intersection,
and $\beta\ge0$ including literal $\beta=0$.  On C14.3b use
$s=x+\eta$ and the unit $a=1/\zeta$; on C14.3a use $s=L+\eta$
with $x=aL$ and $a$ infinitesimal.  In both cases $s$ and $\eta$
are positive infinitesimals.  The exact DG5--DG9 and LD3--LD10 source identities give
\[
 \Lambda=\eta\lambda\sim\tfrac12,\qquad
 \epsilon=O(s),\qquad
 N_Q=\sqrt{\frac{\Delta_0}{h^2K}}Q,
 \qquad
 Q=\frac{\omega}{\sqrt{\eta s}}N_Q,                 \tag{CJ1}
\]
where $\omega$ is a bounded positive source unit.  The complete strict
frame is
\[
 \mathcal F=(N_s,N_T,N_Q,R_b,f_t,\mathcal Q,N_x),
 \qquad N_x=\sqrt B\,\xi,\quad B=\beta+\tfrac12.
                                                               \tag{CJ2}
\]
All DG3/LD9 residuals other than $\epsilon uvQ\xi$ already have
infinitesimal Gram coefficients on this frame without any bound on
$s/\eta$.  The two residuals in the exact split CJ6a below have that same
property whenever $Z\ne0$ and $X/Z$ is bounded.  Thus the only
expression to be allocated on that state chart is
\[
 N=\epsilon U(W/Z)E\xi,
 \qquad E=\eta(P+Q)=\eta uW-vZ.                    \tag{CJ3}
\]
There is no restriction on $Y/Z$ or $W/Z$ below.

\paragraph{An exact identity from the already retained completed row.}
Let $\phi=\xi/Z=X/Z+hu$, and retain exactly DG's completed square
\[
 \widehat J=\sqrt\lambda\,(uv-x/c)Z
       +\frac{G_0u}{2\sqrt\lambda}Q,
 \qquad G_0=2\eta\lambda-2\eta,\qquad c^2=a/\eta.
\]
The defining identity $vZ=\eta uW-\eta(P+Q)$ gives
\[
 \boxed{\quad
 \frac{\widehat J}{\sqrt\lambda}
 =\eta U W-\frac{x}{c}Z
                -\eta u\left(P+\frac Q\lambda\right).
 \quad}                                                       \tag{CJ4}
\]
Multiply CJ4 by $\epsilon\phi E/\eta$.  The complete remaining
normal has the exact conductor
\[
 \boxed{\quad
 N=\frac{\epsilon\phi}{\eta\sqrt\lambda}E\widehat J
       +\epsilon\phi uE\left(P+\frac Q\lambda\right)
       +\frac{\epsilon x}{\eta c}E\xi .
 \quad}                                                       \tag{CJ5}
\]
The only state inverse here is the pure-testing $Z^{-1}$.  No inverse
of $u$, $v$, $W$, or a mixed expression occurs.  No radial row, source
pivot, or completed square is duplicated.

\paragraph{Bounded allocations and infinitesimal coefficients.}
Define two dependent roots
\[
 \mathcal E=\sqrt{s/\eta}\,E,
 \qquad \mathcal P=P+Q/\lambda.
\]
Both have bounded allocations on the existing source roots
$(N_s,N_T,N_Q)$.  Indeed the exact DG9/LD10 inverse identities give
coefficients $O(\sqrt\eta)$, $O(\sqrt\eta)$ and
$O(\sqrt{\eta/s})$ for $P$.  Hence the coefficients for
$\sqrt{\eta s}\,P$ are bounded (in fact infinitesimal), whereas
\[
 \mathcal E=\sqrt{\eta s}\,P+\omega N_Q,
 \qquad
 Q/\lambda=\frac{\omega}{\Lambda}\sqrt{\eta/s}\,N_Q.
                                                               \tag{CJ6}
\]
Here $\eta/s\le1$ is a pure-source inequality.  These are bounded
allocations of dependent roots, not new positive pivots.
For completeness, the other two terms of the same normal have
the particularly simple allocations
\[
 \epsilon uvQ\xi-N
   =-\frac{\epsilon}{s}u\phi\mathcal E^2
     -\frac{\epsilon uv}{\sqrt B}PN_x.               \tag{CJ6a}
\]
Their coefficients are infinitesimal, since $u,v$ are local
infinitesimals and $P$ has a bounded source-root allocation.

In these roots CJ5 is exactly
\[
\begin{split}
 N={}&\frac{\epsilon}{s}\frac{\sqrt s}{\sqrt\Lambda}
                  \phi\mathcal E\widehat J\\
 &+\epsilon\sqrt{\eta/s}\,u\phi\mathcal E\mathcal P
   +\frac{\epsilon}{s}
       \frac{x\sqrt s}{\sqrt a\sqrt B}\mathcal E N_x .
\end{split}                                                    \tag{CJ7}
\]
All three coefficients are infinitesimal bounded tensor elements.
For the first two use $\epsilon/s$ bounded, $s$ infinitesimal,
$\eta/s\le1$, and bounded
$\phi=X/Z+hu$.  The last coefficient is also infinitesimal:
\[
 \frac{x\sqrt s}{\sqrt a}=
 \begin{cases}
 O(x\sqrt s),&a\text{ a bounded positive unit},\\
 \sqrt a\,L\sqrt s,&x=aL,\ a\text{ infinitesimal}.
 \end{cases}                                                   \tag{CJ8}
\]
Also $1/\sqrt B$ is bounded, even at literal $\beta=0$ and when
$B$ is unbounded.  The apparent inverse $\sqrt a$ is canceled in
the large-$\zeta$ case before taking bounds.  Every coefficient is
a finite separated expression; $\phi$ is expanded as $X/Z+hu$.
No source rate is compared with any testing rate.

\paragraph{Full local reassembly.}
Add the already present row $\widehat J^2$ to the strict frame
CJ2 and allocate CJ7 to it and the source roots.  This produces an
infinitesimal Gram perturbation of the previous strict-residue
matrix plus a unit $\widehat J$ diagonal entry.  The separated
bounded-Gram lemma therefore applies to the full remainder.
Retain each other positive row from DG3/LD9 exactly once:
\[
 \widehat R^2,\quad c_UUQ^2,\quad e_4UF_t^2,\quad BV\xi^2,
 \quad(2x/c)uv\tau^2,\quad(G_0x/c)D_0^2WZ.
\]
The last two are literal tensor squares on the unchanged sign chart
$u,v\ge0$, $WZ>0$.  The complete angular and radial packets remain
inside CJ2; all full-$D$, full-$Y$ and higher normals are those of
the exact original-parent identity, not truncated expansions.

Consequently the entire local original parent is tensor SOS when
$Z\ne0$ and $X/Z$ is bounded, with no condition on $s/\eta$,
$u^2W/Z$, or $Y/Z$.  The proof is valid on literal local angle axes
because it uses no angle inverse.  At original contact $E,P,Q,\xi$
and $\widehat J$ vanish, so this reassembly preserves equality and
all original contact normals exactly.

\paragraph{The local state join and the remaining global bookkeeping.}
The reviewed pure-testing half-core exit covers $WZ\le0$, including
$Z=0$.  When $Z\ne0$ and $X/Z$ is unbounded, DS3 in Section~\ref{component:p8cfada778490} is a direct
original-parent tensor-SOS exit.  The above argument covers every
remaining real state on the local angle chart.  Applying these real
certificates to real and imaginary state parts gives the full complex
parent.  Thus the shared local C14.3a/b obligation is closed, with
all source gaps and $\beta=0$ retained.  The SB remote theorem in Section~\ref{component:p6d2f659c0fde} and the unit-$\eta$ theorem in Section~\ref{component:p60eca5442075} supply the other testing and source charts. The complete source partition is printed in Section~\ref{component:pfd0dd7bb10cf}.




\section{\texorpdfstring{C14.4 at large $\zeta$: every unmatched Bob residue}{C14.4 at large zeta: every unmatched Bob residue}}
\label{component:p1a6cc124f5db}

\label{p1a6cc124f5db:sec:r18-c144-bob-residue}

\paragraph{Precise admission.}
Retain the entire original C14.4 source: $x,L>0$ infinitesimal,
$\eta$ unbounded, $\zeta=L/x$ noninfinitesimal, all original physical
intersections, every positive opening, and either permitted literal
$D=0$ or $\beta=0$. This section additionally assumes $\zeta$
unbounded. Put
\[
 \sigma=1/\eta,\qquad S=\sigma+L,\qquad
 \nu=\sigma/S,\qquad \ell=L/S=1-\nu.
\]
For a real projective Bob test choose its normalized homogeneous
coordinates $(c,s)$, so $c^2+s^2=1$; on the finite chart these are
$c=(1+v^2)^{-1/2}$, $s=vc$. The complete original parent is tensor SOS
at every independent test and state satisfying
\[
                         \overline{s^2}\ne\overline\ell.       \tag{EB1}
\]
Bars mean real residues of bounded elements in their respective fields.
This compares real residues only, not independent infinitesimal rates.
The adjacent-$C,A$ theorem already covers noninfinitesimal Alice
tests. Hence only $u$ infinitesimal is proved below. No all-tests
closure of C14.4 is asserted: the matched residue and bounded-unit
$\zeta$ source remain.

\paragraph{Source limits, including both sphere endpoints.}
Use ES1--ES10 without changing their source entry. Since
$x/S\le x/L=1/\zeta$ is infinitesimal, ES7 gives
$e=xh/\sigma$ infinitesimal, and hence $E\sim1$.
Moreover $M=1+xr-\theta/\zeta\sim1$. ES8 then gives $m$ infinitesimal.
The exact equation ES6, divided by $S$, is
\[
 e(B/S)=(x/S)H_0+\nu\mathcal H b_0^2.
\]
Both summands are nonnegative; $B/S$ and $\mathcal H$ tend to one.
Consequently $\nu b_0^2$ is infinitesimal. Therefore ES5 yields
\[
 \frac DS-\frac{\nu^2}{2}\in\operatorname{Inf},\qquad
 \frac\alpha S-\frac{\ell^2}{2}\in\operatorname{Inf},
 \qquad m\in\operatorname{Inf}_{>0}.                         \tag{EB2}
\]
For the second formula, divide the exact ES9 expression by $S$:
\[
 \frac\alpha S=
 \frac{(1+\sigma r)\{\ell R e+\ell^2ME\}}
 {2Q(P/S)}.
\]
All displayed denominator factors tend to one. These are additive
residue assertions, valid also when $\nu$ or $\ell$ is infinitesimal.
In particular no division by either sphere coordinate is used.

The opening-cleared ES8 bounds give
\[
 c_*=SC(1-r)\sim\tfrac12,\qquad
 a_*=SAr\sim\tfrac12.
\]
Define the pure-source harmonic quantities
\[
 H=S(C+A),\qquad \rho=C/(C+A),\qquad
 a_h=SCA/(C+A)=\frac{c_*a_*}{c_*r+a_*(1-r)}.
                                                               \tag{EB3}
\]
Then $1/\sqrt H$ is bounded, $0<\rho<1$, and $a_h\sim1/2$,
uniformly through both $r$ openings. No upper bound on $C$ or $A$
has been assumed.

\paragraph{An exact all-state frame for every Bob chart.}
Let $D_u=1+u^2$ and $A_u=1+u^2s^2$, both positive testing units.
Use the separated roots
\[
 p=\frac{uW}{\sqrt S},\quad q=\sqrt S Z,\quad
 y=\frac Y{\sqrt S},\quad N_y=\sqrt H(y-\rho p),
 \quad N_m=\sqrt m\,sW,
\]
\[
 N_b=\sqrt{A_u}\left(X+\frac{uc^2}{A_u}Z\right),
 \qquad N_r=\frac{\sqrt{D_u}\,usc}{\sqrt{A_u}}Z.
                                                               \tag{EB4}
\]
Only pure-source and pure-testing radicals or inverses occur.
In particular no inverse of $c$, $s$, $u$, or a state coordinate is
needed. The inverse harmonic allocation
$y=\rho p+N_y/\sqrt H$ has bounded coefficients.

Write $b=X+uZ$. The exact original signed core $K/D_v$ in finite
Bob coordinates equals the following homogeneous expression:
\[
\begin{split}
 K_B={}&N_b^2+N_r^2-2D_usc\,pq
       +Sc^2(p-y)^2+SD_us^2y^2\\
 &+\frac{2\sqrt S\,usc}{\sqrt{A_u}}(p-y)N_b
  +\frac{2\sqrt S\,us^2\sqrt{D_u}}{\sqrt{A_u}}(p-y)N_r.
\end{split}                                                    \tag{EB5}
\]
Here $D_us$ denotes $D_u\,s$, and similarly in the other products.
One derivation starts with CE2 divided by $D_v$, sets
$R=usZ$, and completes its entire $(b,R)$ block:
\[
 A_ub^2-2D_us\,bR+D_uR^2=N_b^2+N_r^2.
\]
The cross term from
$[\sqrt S\,c(p-y)+usb]^2$ is allocated by
$cb=cN_b/\sqrt{A_u}+\sqrt{D_u}sN_r/\sqrt{A_u}$.
This proves EB5 without deleting a first state normal.

The original positive event roots have the exact bounded expressions
\[
 \begin{split}
 M_B={}&N_m-\frac{\sqrt m\,c}{\sqrt{A_u}}N_b
       -\frac{\sqrt{mD_u}\,s}{\sqrt{A_u}}N_r
       +\sqrt{mS}\,us\,y,\\
 D_B={}&\sqrt{DS}\,c(p-y)
       +\frac{\sqrt D\,us}{\sqrt{A_u}}N_b
       -\sqrt{D/S}\,\frac{D_us}{A_u}q,\\
 A_B={}&\sqrt{\alpha D_uS}\,s y
       -\sqrt{\alpha D_u/S}\,c q.
 \end{split}                                                   \tag{EB6}
\]
They are, respectively, the normalized original $m,D,\alpha$ event
amplitudes. The formulas are valid at $D=0$ as well. In the $D$ row,
the identity
$usb-D_usZ=usN_b/\sqrt{A_u}-D_usZ/A_u$
avoids an apparent reciprocal-$c$ pole.

Consequently the full original quarter-parent, divided by $D_v$ on
the finite chart, is exactly
\[
 \mathcal P_B=N_y^2+a_hp^2+ASu^2y^2
              +M_B^2+D_B^2+A_B^2+\beta b^2+\tfrac12K_B.
                                                               \tag{EB7}
\]
The two explicit positive rows $ASu^2y^2$ and $\beta b^2$ can be
retained unchanged. All other rows in EB7 use EB4 once. The formula
is a full four-state identity and is not a restriction to a state ray.

\paragraph{Strictness precisely off the matched residue.}
Every coefficient of EB5 except the printed $pq$ term is either
its fixed positive pivot or infinitesimal after EB3. In EB6, the
$m$ row is $N_m$ plus infinitesimal coefficients; the $D$ and
$\alpha$ rows have only their $q$ coefficients at residue.
Use EB2, $A_u,D_u\sim1$, and EB3. The resulting residue Gram on
$(N_y,N_m,N_b,N_r,p,q)$ is the direct sum of
$\operatorname{diag}(1,1,1/2,1/2)$ and
\[
 \frac12
 \begin{pmatrix}
 1&-\bar s\bar c\\
 -\bar s\bar c&\bar\nu^2\bar s^2+\bar\ell^2\bar c^2
 \end{pmatrix}.                                                \tag{EB8}
\]
The determinant of this last block is
\[
             \frac14(\bar s^2-\bar\ell)^2.                    \tag{EB9}
\]
Thus EB1 makes the entire bounded finite Gram strictly positive.
The separated bounded-Gram lemma gives tensor SOS, and restoring the
two retained positive rows proves the complete parent. The roots are
allowed to be dependent or even identically zero at an endpoint; the
claim concerns the coefficient Gram, so this does not affect the
argument.

The same homogeneous formulas allow $c=0$ and $s=0$ literally.
They include every Bob projective endpoint, every sign, and every
local Alice axis. Nonlocal Alice charts use the inherited adjacent
$C,A$ theorem. The real and imaginary state parts are treated
separately. Every normalization is pure source or pure testing.

\paragraph{The exact remaining obstacle.}
The leading scalar has the useful two-square identity
\[
 \tfrac12p^2-sc\,pq
 +\tfrac12(\nu^2s^2+\ell^2c^2)q^2
 =\tfrac12(p-scq)^2+\tfrac12(s^2-\ell)^2q^2.                  \tag{EB10}
\]
It does not complete the matched-residue chart: there the second
pivot collapses, and the actual source corrections must be allocated
through exact contact conductors. Arbitrary small positive $D$ or
$\beta$ openings do not supply a missing fixed pivot. Unit $\zeta$
also requires new source limits because $e$ can have nonzero residue.
These two obligations are retained, not declared empty.

The focused checker verifies EB5--EB7 as complete homogeneous
quadratic identities and EB9--EB10 algebraically. Source estimates
EB2--EB3 are the analytical consequences of the already proved
ES identities given above. This is a partial-domain admission only.


\section{\texorpdfstring{C14.2: bounded Bob tests at every local Alice angle}{C14.2: bounded Bob tests at every local Alice angle}}
\label{component:pda4a21937f64}

\label{pda4a21937f64:sec:r18-c142-bounded-bob}

\begin{proposition}
Retain the unchanged full C14.2 source and every original intersection
as in the preceding noninfinitesimal-Alice proposition.  For every
infinitesimal or zero Alice test $u$ and every bounded Bob test $v$,
the entire original parent is finite tensor SOS for all real or
complex states.  Every source opening and both $D=0$, $\beta=0$
are included.  Together with that proposition this leaves exactly
$u$ infinitesimal and $v$ projectively unbounded, equivalently
$u,w=1/v$ infinitesimal, as the remaining C14.2 testing corner.
\end{proposition}

\paragraph{The same source weights and a pure-source strict matrix.}
Keep the opening-safe positive source weights from LA1,
\[
 C_0=\frac1{2L(1-r)},\quad
 A_0=\frac{1+L(1-r)}{2Lr},\quad
 a_0=\frac L{2(1-Lr)}.
\]
Let $c=(1+v^2)^{-1/2}$, $s=vc$, so $c^2+s^2=1$.
Since $v$ is bounded, $c$ is a positive testing unit.  Write
$s_0=\operatorname{res}s\in(-1,1)$ and $c_0^2=1-s_0^2>0$.
These are real numbers, not source/testing comparisons.  Consider
the pure-source quadratic form in three placeholder variables
\[
 M_0(P,Y,T)=
  (C_0+c_0^2/2)(P-Y)^2+(A_0+s_0^2/2)Y^2
          +a_0(s_0Y-T)^2-s_0PT.                    \tag{LB1}
\]
Its Gram has first leading minor $C_1=C_0+c_0^2/2>0$,
second leading minor $C_1[A_0+s_0^2(a_0+1/2)]>0$, and
\[
 \det M_0=
 \frac{c_0^2\{1+L(1-r)\}^2}
                {8Lr(1-r)(1-Lr)}>0.               \tag{LB2}
\]
This is an exact source identity, not just a limiting determinant.
Moreover $M_0^{-1}$ has bounded entries through every opening.
Indeed every cofactor is a sum of bounded multiples of
$C_0A_0,C_0a_0,A_0a_0,C_0,A_0,a_0,1$; terms involving a square
of a single divergent weight cancel.  Multiplication by
$8L^2r(1-r)(1-Lr)$ makes every cofactor polynomial in bounded
$L,r,s_0$.  Division by LB2 therefore leaves only the source-unit
denominators $L$, $c_0^2$, and $1+L(1-r)$.
Explicitly, with $e=1-r$ and $B_L=1+Le$,
\[
 M_0^{-1}=\frac1{c_0^2B_L}
 \begin{pmatrix}
 2L&2Lr&2s_0\\
 2Lr&2Lr(1-s_0^2e)&2s_0r(1+Le c_0^2)\\
 2s_0&2s_0r(1+Le c_0^2)&
       2(1-Lr c_0^2)(1+Le c_0^2)/L
 \end{pmatrix}.                                      \tag{LB2a}
\]

Choose any exact positive pure-source congruence
$M_0=N_1^2+N_2^2+N_3^2$.  Its inverse allocation of $P,Y,T$
on these three roots is bounded: each inverse-row norm squared is
the corresponding diagonal of $M_0^{-1}$.  This uses only a
pure-source Gram, with the fixed real number $s_0$ as parameter;
no mixed matrix is inverted or square-rooted.

\paragraph{Freeze the only divergent Bob-dependent event exactly.}
Use the legal raw roots
\[
 P=uW,\quad X,\quad Y,\quad
 T=cZ-(s-s_0)Y,
 \qquad d_s=s-s_0\in\operatorname{Inf}_{\rm test}.
                                                               \tag{LB3}
\]
In particular the entire alpha event satisfies
$sY-cZ=s_0Y-T$ exactly.  Thus no product of an unbounded alpha
weight with $d_s$ is ever estimated.  No inverse of $c$, $u$, or
a state coordinate is needed in these root definitions.

Let $B_0=c_0^2(P-Y)^2+s_0^2Y^2-2s_0PT$.
The complete original half-core after Bob normalization is
$K_B/2$, where
\[
 K_B=\{c(P-Y)+usX-sZ\}^2
       +(1+u^2)(sY-cZ)^2+(X+uZ)^2-(1+u^2)Z^2.
\]
It has the following exact, division-free allocation on LB3:
\[
\begin{split}
 K_B={}&X^2+B_0
       +d_sP\{-(s+s_0)P+2s_0Y-2T\}\\
       &+2ucX(sP-s_0Y+T)
       +u^2s^2X^2+u^2(s_0Y-T)^2.
\end{split} \tag{LB4}
\]
This keeps the full signed core and all $u$-dependent hard normals.

\paragraph{Exact actual-parent reassembly.}
Use exactly LA4's pure-source factor
$f=1-(|C/C_0-1|+|A/A_0-1|+|\alpha/a_0-1|)>0$.
Retain the actual original event excesses
$C-fC_0$, $A-fA_0$, $\alpha-fa_0$, and the whole $m,D,\beta$
packets.  On the finite Bob chart the original quarter-parent divided
by $1+v^2$ is then the sum of those positive packets and exactly
\[
 S=f\{C_0(P-Y)^2+A_0(1+u^2)Y^2
                         +a_0(1+u^2)(s_0Y-T)^2\}+K_B/2.
\]
Substitution of LB1 and LB4 gives
\[
\begin{split}
 S={}&fM_0+\tfrac12X^2
       +fA_0u^2Y^2+fa_0u^2(s_0Y-T)^2\\
    &+\tfrac12u^2s^2X^2+\tfrac12u^2(s_0Y-T)^2\\
    &+\tfrac12(1-f)B_0
       +\tfrac12d_sP\{-(s+s_0)P+2s_0Y-2T\}\\
    &+ucX(sP-s_0Y+T).
\end{split} \tag{LB5}
\]
The first two lines consist of a strict source-normalized Gram on
$(N_1,N_2,N_3,X)$ and exact retained positive squares.  All terms
in the last two lines have infinitesimal bounded coefficients on
that same frame: $1-f$ is source infinitesimal, $d_s,u$ are testing
infinitesimals, $c,s,s_0$ bounded, and LB2 supplies a bounded inverse
allocation for all three raw roots.  Since $f\sim1$, the resulting
complete Gram has residue $\operatorname{diag}(1,1,1,1/2)$.
The separated bounded-Gram lemma and restoration of every retained
packet prove the full original parent.  The normalization factor
$1+v^2$ is pure testing and positive.

Both coefficient-boundary packets were retained rather than capped
or inverted.  The argument works literally at $u=0$ or $v=0$,
allows signed testing angles, and extends to complex states by
real/imaginary splitting.  It gives no claim for $c$ infinitesimal:
there $c_0=0$, the determinant in LB2 vanishes, and the reciprocal
two-small-angle corner remains a separate obligation.

The new focused checks establish LB2, opening-regular cofactors,
the complete original-core identity LB4 and the entire selected-parent
identity LB5.  They do not rerun prior source identities or assert
the remaining C14.2 corner closed.


\section{\texorpdfstring{The exact C14.2 remainder and its centered C extraction}{The exact C14.2 remainder and its centered C extraction}}
\label{component:p156da4167158}


\paragraph{Quantified remainder after the two new test admissions.}
The propositions in Section~\ref{component:pb0d9d28d0849} and
Section~\ref{component:pda4a21937f64} admit all C14.2 tests except $u,w=1/v$
infinitesimal.  Literal $u=0$ or $w=0$ has a direct original-core
SOS, so after the exact sign gauges a remaining test has
\[
 u,w>0\text{ infinitesimal},\qquad
 L>0\text{ a bounded source unit},\quad
 x,\sigma=1/\eta>0\text{ infinitesimal}.             \tag{LC1}
\]
All original strict gaps, coefficient signs, total-scale and prior
exclusions remain; $D=0$ and $\beta=0$ are not removed.
The pure-testing core-sign exit allows restriction to $WZ>0$,
with every state magnitude otherwise unrestricted.  No assertion
that this residual intersection is nonempty is made.

Use the exact BE2--BE8 source sphere and the entire generic JR5/JB1
parent in the form SP6 from
Section~\ref{component:pfeca652aef94}.
Those identities are algebraic for $z=L^{-1}$ a bounded unit as
well as for large $L$.  On C14.2, the original $r\to1$ opening
may make $C$ unbounded; the large-$L$ boundedness assertion for
SP5 therefore cannot be imported.  The following exact extraction
removes that particular pole without changing the source.

\paragraph{The entire C-dependent block.}
Retain all SP/JR roots and definitions, notably
\[
 q^2=z,\quad b=Lr,\quad D_w=1+w^2,\quad
 J_C=N_y+bR,\quad\Gamma=\sigma Z-uwW.
\]
With all other actual source coefficients held fixed, the complete
SP6 matrix is affine in the displayed coefficient $C$.  Inspection
of the full $H_y,K_r,R_y,R_\gamma$ formulas gives its entire
$C$-dependent part as
\[
 4CzD_wJ_C^2
       +8C(r-1)J_C(z\sigma R-qw\Gamma).            \tag{LC2}
\]
This includes both complete normal couplings and the $R\Gamma$
coupling.  Neither the scalar radial diagonal nor a first normal has
been silently removed.  No other displayed SP6 coefficient depends
on the formal variable $C$ once the other actual source coefficients
are held fixed.

Choose the bounded positive pure-source cap
\[
 C_*=\frac C{1+C},\quad d_C=C-C_*\ge0,
 \quad t_C=d_C(1-r).
\]
The original opening-cleared estimate $C(1-r)\asymp z$ proves
$t_C$ bounded, uniformly through the $r\to1$ opening.  No source
case split is needed.  The whole parent core therefore satisfies
the exact identity
\[
 \boxed{\quad
 G=G\big|_{C=C_*}+4zD_wd_CJ_C^2
              -8t_CJ_C(z\sigma R-qw\Gamma).
 \quad}                                                       \tag{LC3}
\]
Formal substitution in the first term does not assert that a capped
source is physical.  The equality is a full-parent matrix identity,
and every other original retained square remains unchanged.
The extracted row is centered at the original contact, since
$N_y=R=0$ there.  It retains its whole event dependence at every
state and angle; this is not a Hessian-only extraction.

For reference, the existing exact SP12 conductor gives
\[
 z\sigma R-qw\Gamma=qu(w^2-\sigma z)W.             \tag{LC4}
\]
Both LC3 forms have only separated source/testing coefficients and
no new state, angle, or mixed inverse.

\paragraph{Use in the full completion.}
LC3 removes the $C$-opening pole: $C_*,t_C$ are bounded and the correction carries $\sigma$ or $w$. This extraction alone does not give a positive Gram. The $\Gamma$ direction can lose its pivot at $D=0$; its normals are jointly completed with the source sphere, radial row and full $h/\sigma$ packet by the CF, HG, HP and CT certificates. Their exhaustive C14.2 join, on LC1 with all original intersections, is printed in Section~\ref{component:pfd0dd7bb10cf}.


\section{\texorpdfstring{C14.2: the complete noninfinitesimal-Alice complement}{C14.2: the complete noninfinitesimal-Alice complement}}
\label{component:pb0d9d28d0849}

\label{pb0d9d28d0849:sec:r18-c142-remote}

\begin{proposition}
Keep the unchanged physical C14.2 source: $d=1$, $x>0$
infinitesimal, $L=x\zeta$ a positive bounded unit, $\eta$ unbounded,
all original strict gaps, nonnegative event weights, total-scale and
earlier-exclusion intersections.  Then the complete original parent
is finite tensor SOS for every noninfinitesimal Alice angle $u$,
including $u=\infty$, every independent Bob angle, and all real or
complex states.  Literal $D=0$ and $\beta=0$ are retained.  Thus
only infinitesimal or zero $u$, with arbitrary Bob type, remains in
C14.2.  No entire-source admission is asserted.
\end{proposition}

\paragraph{Opening-safe source input.}
Use LH1--LH3 in the preserved R13 file
Section~\ref{component:p57d751e20564} and BE1--BE9
in Section~\ref{component:p10ee3d96a3f1}.
Put $e_r=1-r$, $Q=1-Lr>0$ and define positive pure-source weights
\[
 C_0=\frac1{2Le_r},\qquad
 A_0=\frac{1+Le_r}{2Lr},\qquad
 a_0=\frac L{2Q}.
                                                               \tag{LA1}
\]
Their actual-weight ratios satisfy
$C/C_0,A/A_0,\alpha/a_0\sim1$ uniformly through $r$, $e_r$,
and $Q$ openings.  Each leading and actual weight has a positive
real lower bound, since $L$ is a positive bounded unit and
$0<r,e_r,Q<1$.  No upper bound for the individual weights is assumed.
The BE source sphere and both its endpoints remain unchanged; no
coefficient $D$, $\beta$, $d_0$ or $b_0$ is divided by.

\paragraph{The reciprocal Bob strip with a unit Alice angle.}
Let $u$ be a bounded nonzero unit and $w=1/v$ infinitesimal or zero.
Use four legal raw roots
\[
 P=uW,\quad X,\quad Y,\quad T=wZ,
 \qquad U=u^2,\quad D_u=1+U,\quad D_w=1+w^2.
\]
The exact original signed half-core $K_w/2$ is CE6 of the preserved
R13 file Section~\ref{component:p89fd3edb3edd}.
Written in these roots it has the exact split
\[
 K_w=K_0+w^2\{(P-Y)^2+X^2\}
                      +2uwX(T+P-Y),                         \tag{LA2}
\]
\[
 K_0=D_uX^2+D_uY^2+UT^2-2UTY-2PT.
\]
This is the complete core, including its sign; no state graph is used.
The selected positive events are
$CD_w(P-Y)^2$, $AD_uD_wY^2$, and
$\alpha D_u(Y-T)^2$ in the quarter-parent.  The whole original
$m,D,\beta$ packets are retained separately.

The LH4 contact identity implies the following exact identity for
the leading selected form at $w=0$:
\[
\begin{split}
 Q_0={}&C_0(P-Y)^2+A_0D_uY^2+a_0D_u(Y-T)^2+K_0/2\\
 ={}&\frac{D_u}{2}X^2
       +C_0\{P-Y-Le_rT\}^2
       +(A_0+a_0+1/2)\{Y-LrT\}^2\\
 &+U\{A_0Y^2+(a_0+1/2)(Y-T)^2\}.
\end{split}                                                   \tag{LA3}
\]
Every displayed coefficient is nonnegative pure source or a positive
pure-testing factor.  The source coefficients in LA3 can be clipped
at positive real lower bounds, retaining their exact positive square
excesses.  The resulting bounded Gram on $(P,X,Y,T)$ has strictly
positive residue: $X$ is controlled directly, the two $U$ terms
control $Y,T$ because $u$ is a unit, and the first contact square
then controls $P$.  The coefficients $Le_r,Lr$ in these roots are
bounded pure source, including all original openings.

\paragraph{Exact common-factor accounting for the actual weights.}
To transfer LA3 to the actual parent, define solely in the source field
\[
 \delta=|C/C_0-1|+|A/A_0-1|+|\alpha/a_0-1|,
 \qquad f=1-\delta.
\]
Then $\delta$ is nonnegative infinitesimal, $f>0$ is a unit,
and each of $C-fC_0$, $A-fA_0$, $\alpha-fa_0$ is nonnegative.
Retain their original event packets as exact squares.  The complete
remaining selected polynomial is, identically,
\[
\begin{split}
 S={}&f\{C_0D_w(P-Y)^2+A_0D_uD_wY^2
                          +a_0D_u(Y-T)^2\}+K_w/2\\
  ={}&fQ_0+fw^2\{C_0(P-Y)^2+A_0D_uY^2\}\\
    &+(1-f)K_0/2+(K_w-K_0)/2.
\end{split} \tag{LA4}
\]
The second line's $w^2$ packets are retained positive squares with
their complete actual source scales.  The third line, including the
signed $(1-f)K_0/2$ term, is a bounded infinitesimal Gram on the
four raw roots, by LA2 and $1-f=\delta\in\operatorname{Inf}$.
Clip the completed source-square coefficients of $fQ_0$ at positive
real lower bounds as above.  Its retained bounded Gram remains
strict under this entire infinitesimal perturbation.  The separated
bounded-Gram lemma gives tensor SOS; restoring all square excesses
proves the whole original parent.  This is exact polynomial
accounting in the tensor ring, not an inequality that orders mixed
source and testing quantities.

\paragraph{The other Bob types and the unbounded Alice endpoint.}
If $u$ is a bounded nonzero unit and $v$ is a bounded nonzero unit,
use the inherited noncommuting-residue theorem on the unchanged
physical parent.  If $u$ is a bounded nonzero unit and $v$ is
infinitesimal or zero, clip the original $C,A,\alpha$ at positive
real lower bounds $c_*,a_*,a'_*$ and retain every excess and the
whole original $m,D,\beta$ packets.  The complete selected residue is
\[
 (c_*+1/2)(uW-Y)^2+a_*D_uY^2
          +a'_*D_uZ^2+\tfrac12(X+uZ)^2\succ0.     \tag{LA5}
\]
It is strict on all four original states when $u$ is a unit; every
finite-angle correction has bounded infinitesimal coefficients.
This covers all bounded Bob types without invoking LA4 outside its
stated reciprocal strip.

For $u$ unbounded, put $p=1/u$ and use the original homogeneous
parent.  The exact CE4--CE5 raw-frame proof, with any positive real
caps for $C,A$, covers every bounded $v$.  Its strict selected
residue is
\[
 4c_*(1+v^2)W^2+4a_*(1+v^2)Y^2
                      +2(W+R)^2+2(B-vY)^2,
 \quad B=pX+Z,\quad R=vX.                         \tag{LA6}
\]
For unbounded $v$, the exact twice-homogeneous CE7--CE8 proof has
strict residue
\[
 4c_*W^2+4a_*Y^2+2X^2+2(T-Y)^2,
 \qquad T=wZ.                                    \tag{LA7}
\]
These two homogeneous proofs need only positive real lower bounds
for $C,A$ at these endpoints.  They do not use the stronger
``$C,A$ unbounded'' hypothesis needed for the finite-unit-angle part
of the adjacent-event theorem.  All unselected original packets and
positive excesses are retained exactly.  Literal $p=0$ and $w=0$
are the polynomial values of these original homogeneous parents.

The finite unit-Alice and homogeneous unbounded-Alice cases exhaust
all noninfinitesimal $u$.  Exact sign gauges and real/imaginary splitting
give both signs and complex states.  The possible additional $D=0$
null direction at $u=w=0$ is untouched, since the proposition does
not include the infinitesimal-Alice chart.  The sole remaining C14.2
testing obligation is now $u$ infinitesimal or zero, with arbitrary
Bob types and the unchanged full source intersection.

The focused checker verifies only the new LA2--LA4 algebra, including
the entire signed core and the exact common-factor remainder.  No
historical checker is rerun and no branch closure is inferred from
this testing-complement result.


\section{\texorpdfstring{C14.1b: every state when the first testing angle is a unit}{C14.1b: every state when the first testing angle is a unit}}
\label{component:pcc84202479d2}

\label{pcc84202479d2:sec:r18-sphere-unit-u}

Retain the entire original C14.1b source intersection and the exact
SP6 parent from revision 17.  Assume $u$ is a bounded nonzero unit
and $w$ is infinitesimal.  The argument below proves the whole parent
for every testing state, including the previously unadmitted
$W/Z$-unbounded states.  It does not cover infinitesimal $u$.
All source gaps, $D=0$, both source sphere endpoints and the literal
coefficient boundaries remain unchanged.

\paragraph{No state chart is needed.}
All roots below are homogeneous polynomials in the original testing
state, with pure-source coefficients.  In particular the proof
uses no state-coordinate inverse, no bound on $W/Z$, and no sign
restriction.  The only testing inverse is that of the unit $u$.
Write $U=u^2$, $s=x/\sigma$, and preserve the exact SP/BS coefficients
\[
 L_h=\lambda\sigma/h,\quad \rho=\lambda x/\sigma,
 \quad d=4D/\sigma^2,\quad K_A=4\kappa h/\sigma-\lambda.
\]
The source identities give $s\in\mathrm{Inf}$,
$\sigma/\sqrt h\in\mathrm{Inf}$, $L_h\sim2$, $\nu\sim1$, and
\[
 d+\rho\sim2,\quad K_A-(2-4h)\in\mathrm{Inf},\quad
 \chi-(4-4h)\in\mathrm{Inf},\quad
 (c_G+zK_r)-2(1+U)\in\mathrm{Inf}.
 \tag{SU2}
\]
The last infinitesimal is a finite separated expression for bounded
$u$ and infinitesimal $w$.  For the only new source assertion in
SU2, the sphere identities give exactly
\[
 d=\frac{2\mathcal H}{PJ}d_0^2,\qquad
 \rho=\frac{2xH_0}{\sigma B}
       +\frac{2\mathcal H}{B}b_0^2+\lambda_0s.
 \tag{SU3}
\]
Thus $d+\rho\sim2$ follows from $d_0^2+b_0^2=1$, with no inverse
of either sphere coordinate.  All other bounds are the reviewed
SP9--SP11 and BS1 source identities.

\paragraph{Shift the troublesome angular product into the normal block.}
Set $A_0=qW$, $A=A_\theta=(U-x)A_0$, and retain $R$.
There is an exact conductor:
\[
 A(A+uR)=u(U-x)wWZ-x(U-x)A_0^2.
 \tag{SU4}
\]
Consequently the angular/radial part of BS2 is
\[
\begin{split}
 K_AA^2+\chi uAR+(c_G+zK_r)R^2
 ={}&2(A+uR)^2+2R^2-4hu(U-x)wWZ\\
 &+4hx(U-x)A_0^2+\delta_A A^2+
   \delta_\chi uAR+\delta_RR^2,
\end{split} \tag{SU5}
\]
where the three $\delta$ coefficients are the infinitesimals in
SU2.  This step keeps the signed product exactly; it is not an
inequality based on the sign of $U-x$.

\paragraph{A complete strict normal block.}
Introduce the legal separated roots
\[
 H=\sqrt h\,uZ,\qquad S=uwW-hu^2Z,\qquad N=N_x,
 \qquad Z_g=uwW=S+u\sqrt h\,H.
\]
Their exact allocations are
\[
\begin{split}
 \Gamma&=-S+\left(\frac{\sigma}{u\sqrt h}-u\sqrt h\right)H,\\
 P_h:=\sqrt h P_\beta&=(1-sh)H-\frac{s\sqrt h}{u}S,\\
 X_e&=N-\sqrt h P_h.
\end{split} \tag{SU6}
\]
Every coefficient is bounded; the factors $u^{-1}$ are pure testing
units.  The full signed packet in BS2 also satisfies exactly
\[
 \mathscr P=L_hP_h^2+
      \rho(1-x/U)Z_g^2-\rho\sigma(U-x)A_0^2.
 \tag{SU7}
\]
The apparently signed factor $1-x/U$ is simply a bounded Gram
coefficient of residue one.  No mixed inverse is used.

Combine the $wWZ$ term in SU5 with every $N,\Gamma,P_h,Z_g$ term:
\[
\begin{split}
 \mathscr N={}&[2(1+w^2)+4mw^2]N^2
     +2(2D+1)U X_e^2+d\Gamma^2-4u\nu N\Gamma\\
 &+L_hP_h^2+\rho(1-x/U)Z_g^2
       -4\sqrt h\,\frac{U-x}{u}\,Z_gH .
\end{split} \tag{SU8}
\]
This is an exact equality for the entire indicated normal block.
Its coefficients in $(N,S,H)$ are bounded separated expressions.
Using SU2 and SU6, its residue is exactly
\[
 2(1+u_0^2)N^2+4u_0NS+2S^2+2H^2.
 \tag{SU9}
\]
The $N,S$ determinant is $4$, so this residue is strictly positive.
The $H$ cross terms cancel: the $X_e$ square cancels the leading
$N\Gamma$ coupling to $H$, and the angular shift in SU5 cancels
the remaining $S H$ and $hH^2$ terms.  This is where the sphere
sum $d_0^2+b_0^2=1$ is essential.

The remaining angular/radial block is
\[
\begin{split}
 \mathscr A={}&2\{(U-x)A_0+uR\}^2+2R^2\\
 &+(4hx-\rho\sigma)(U-x)A_0^2
   +\delta_A(U-x)^2A_0^2
   +\delta_\chi u(U-x)A_0R+\delta_RR^2.
\end{split} \tag{SU10}
\]
Its residue on $(A_0,R)$ has determinant $4u_0^4>0$.

\paragraph{Every other original-parent term.}
The unchanged hard $N_y$ diagonal has residue $2(1+u_0^2)$.
All SP5 terms other than its two $\Gamma$ terms already have
infinitesimal coefficients on $(N,N_y,A_0,R)$, by SP's complete
opening-uniform coefficient bounds and $A=(U-x)A_0$.
The two $\Gamma$ terms have a factor $qw$ and become infinitesimal
crosses after the bounded allocation SU6.  Thus the complete Gram
on $(N,N_y,S,H,A_0,R)$ has strictly positive residue, the direct
sum of SU9, the $N_y$ entry and SU10.

The remaining exact nonnegative rows are retained once: both
centered cap excesses, $4\beta D_wN^2$, and
$4bpzU\Gamma^2$.  The $d\Gamma^2$ part and
$2(2D+1)(uX_e)^2$ have already been used in SU8 and are not
retained a second time.  All scalar, hard and first-normal terms
of the original SP6 parent are therefore included.

The bounded separated-Gram criterion proves tensor SOS for all real
states and both testing signs, including literal $Z=0$ and $w=0$.
Applying the real identity separately to real and imaginary parts
gives all complex states.  Neither source sphere endpoint is inverted.
The precise remaining C14.1b testing domain has $u$ infinitesimal
(with the earlier state and sign admissions still retained).

The focused checker verifies SU4, SU6--SU8 and the whole reassembled
SP6 parent, then the two strict residue determinants.  It does not
replay the unchanged source formulas or claim a whole-source closure.


\section{\texorpdfstring{C14.2: the complete unbounded-$H$, unbounded-$W/Z$ chart}{C14.2: the complete unbounded-H, unbounded-W/Z chart}}
\label{component:p831dd6874dcf}
% Edited scoped mathematical packet, 22 October 2026.

\begin{proposition}
Keep the exact physical C14.2 source, $d=1$, all original strict
gaps and residual intersections, and $D,\beta\ge0$, including their
literal zero faces.  Suppose $H=h/\sigma$ is unbounded.  The entire
parent is tensor SOS on the remaining double-small-angle chart with
$W/Z$ unbounded.  This includes every state ratio on this chart:
the extreme ratios are admitted by CS1--CS4, and every
positive bounded unit ratio is admitted below.  No mixed inverse or
source/testing rate comparison occurs.
\end{proposition}

\paragraph{Exact source and state data.}
Use BE2--BE9 and the entire SP6 identity, restoring the original
centered $a,\ell$ excesses.  Thus $q^2=z=L^{-1}$ is a bounded
positive source unit, $b=Lr$, $Q=1-b$, $p=ab$, $l=\ell Q$,
$a=Az$, $\ell=\alpha z$.  All original hard and radial normals
remain present.  Use the actual SP3 coefficient
\[
 \lambda=2H+\lambda_0,
 \qquad K_A=4\kappa H-\lambda=\alpha_\kappa-\lambda_0.
                                                               \tag{CT1}
\]
SP9--SP10 give $\lambda_0,\alpha_\kappa\in\mathcal O_c$.
Consequently $\lambda>0$ is unbounded, $\lambda\sigma\sim2h$,
$\lambda x\in\operatorname{Inf}_c$, and
$\rho=\lambda x/\sigma$ is bounded: use BE7 and
$x/\sigma\in\operatorname{Inf}_c$ on $H$ unbounded.
There is no assumption that $\rho$ is a unit.

After CS1--CS4 and the sign gauges, take $W,Z,u,w>0$ and define solely
in the testing field
\[
 f=Z/W\in\operatorname{Inf}_{t,>0},\quad
 c=\frac{u}{wf}=\frac{uW}{wZ}>0\text{ a bounded unit},\quad
 F=fc^{1/3},\quad V=wc^{1/3}.
\]
Both $F,V$ are infinitesimal.  Set
\[
 k_c=(zc)^{1/3},\quad
 F_0=\sqrt{x/\sigma}\,z^{1/6},\quad
 V_0=\sqrt\sigma\,z^{1/6},\quad
 g_0=\sqrt{\lambda x}\,z^{1/6}.
                                                               \tag{CT2}
\]
Here $F_0,V_0$ are infinitesimal pure-source elements, and $k_c$ is
a bounded separated unit.  Every fractional power of a displayed
separated product means the product of the two pure-field powers.

\paragraph{The scalar packet and its strict circuit rows.}
The complete SP6 angular packet is exactly
\[
 K_A A_\theta^2+S+\chi uA_\theta R,
 \qquad S=\lambda\sigma P_\beta^2+
                    \rho(u^2-x)(w^2-\sigma z)W^2.              \tag{CT3}
\]
This follows from SP7 and does not discard a scalar or normal term.
Put
\[
 A_1=\sqrt{\lambda\sigma}\,VF^2W,\quad
 A_2=\sqrt\rho\,V^2FW,\quad
 A_3=\sqrt\lambda\,xqW,\quad
 G=g_0VFW,\quad d_i=A_i-G.
\]
All four quantities are nonnegative separated monomials and
$A_1A_2A_3=G^3$.  Let $r_c=k_c^{-1}$.  The exact identity is
\[
\begin{split}
 S+\lambda xR^2={}&\sum_{i=1}^3A_i^2-3G^2\\
 &+\lambda xz\{uW(r_c-1)r_c(r_c+1)\}^2
   +2\lambda xz\{uW(r_c-1)r_c\}^2.             \end{split} \tag{CT4}
\]
To see the finite SOS and its strict $d_i$ rows, let
$p_i=A_i^{1/3}$, with separated cube roots.  Exactly
\[
 \sum_iA_i^2-3G^2
 =\sum_i d_i^2+
   \sum_{j=1}^3\sum_{1\le i<k\le3}
       \{\sqrt{Gp_j}(p_i-p_k)\}^2.             \tag{CT5}
\]
The square roots also separate.  Thus CT4 supplies three diagonal
pivots equal to one on $d_1,d_2,d_3$, and nine further squares.
The two additional CT4 squares are retained whole.  Its sole cost
$-\lambda xR^2$ has an infinitesimal pure-source coefficient.

\paragraph{The radial row without a mixed inverse.}
Define the legal root
\[
 J_c=(k_c-1)VFW.
\]
The exact radial relation is
\[
 R=-\frac{k_c^2+k_c+1}{q c^{2/3}}J_c.           \tag{CT6}
\]
Only the displayed bounded separated coefficient is used; the
polynomial $k_c^2+k_c+1$ is never inverted.
Restoring the full $a,\ell$ weights, their complete occurrence in
the radial normal block is
\[
 4D_u\{aD_w(N_y+bR)^2+\ell(N_y-QR)^2\}.
\]
The full constant radial base, including the other source terms, is
\[
 M=[4a+2z(2C+1)](N_y+bR)^2+4\ell(N_y-QR)^2.    \tag{CT7}
\]
Retain the positive $u^2,w^2$ excesses of the preceding $a,\ell$
packet and the exact $4Cz w^2(N_y+bR)^2$ packet.  Clip the two
pure-source coefficients in CT7 at positive real lower bounds,
retaining their positive square excesses.  After CT6 the clipped
Gram on $(N_y,J_c)$ is bounded and strictly positive at residue:
the two rows differ by $R$, and the coefficient of $J_c$ in CT6
has nonzero real residue.  This remains true at $b=0$ or $Q=0$
in residue.  For an explicit audit, put $t=u^2$, $\xi=w^2$,
$g=h-Q$, $a_\beta=\beta\theta^2$ and
\[
 k_1=2(1+2m)g^2+8Dhb+8pb\sigma+4a_\beta.
\]
After the displayed positive radial packets have been extracted,
the exact three remaining coefficients of $N_y^2,R^2,N_yR$ are
\[
\begin{split}
 E_{yy}={}&2z\{(1+2m)t+(1+2D)\xi\},\\
 E_{RR}={}&z\{8bD+8bp\sigma
       +4b\sigma(2C+2D+1)(r-1)+k_1t
       +2b^2(2D+1)\xi\},\\
 E_{yR}={}&4z\{t[2Dh+2p\sigma+(1+2m)g]
       +b(2D+1)\xi+\sigma(2C+1)(r-1)\\
       &\hspace{19mm}+2D[1+\sigma(r-1)]+2p\sigma\}.
\end{split}                                                     \tag{CT7a}
\]
These are bounded infinitesimal coefficients.  The only apparent
$C$ pole is its full LC2 correction,
whose coefficient is $C(1-r)$ times $\sigma$ or $w$ and is bounded
by BE9.  No entrywise cap of the original radial matrix is made.

\paragraph{The two sphere-weighted Gamma conductors.}
The monomial relations give exactly
\[
 (F-F_0)G=F_0d_1,\qquad (V-V_0)G=V_0d_2.
                                                               \tag{CT8}
\]
Write $b_0,d_0$ for the BE6 sphere coordinates, not the source
parameter $b$.  The pure-source quantity
\[
 B_\Gamma=b_0V_0/g_0=b_0\sqrt{\sigma/(\lambda x)}
                                                               \tag{CT9}
\]
is bounded.  Indeed BE7 gives
$xhB=\sigma xH_0+\sigma^2\mathcal H b_0^2$, and
$\lambda\sigma\sim2h$; this proves the bound directly and never
divides by $b_0$.  With $j=c^{1/3}$ the exact conductors are
\[
\begin{split}
 w b_0\Gamma={}&-\frac{B_\Gamma(V_0+V)}{k_cj}d_2
                       -\frac{b_0V^2}{k_cj}J_c,\\
 u b_0\Gamma={}&-\frac{B_\Gamma jF(V_0+V)}{k_c}d_2
                       -\frac{b_0jV^2F}{k_c}J_c.              \end{split} \tag{CT10}
\]
Every coefficient in these two arrays is infinitesimal and bounded.
Let $\Gamma_d=d_0\Gamma$.  The original diagonal is exactly
\[
 \frac{4D}{\sigma^2}\Gamma^2
       =\frac{2\mathcal H}{PJ}\Gamma_d^2,       \tag{CT11}
\]
with a positive-unit pure-source coefficient.  In CT11 the $J$ in
the denominator is the original source $J$ of BE2.  Finally
\[
 u\Gamma=d_0u\Gamma_d+b_0(u b_0\Gamma),\quad
 w\Gamma=d_0w\Gamma_d+b_0(w b_0\Gamma).         \tag{CT12}
\]
Thus every first normal carrying $u\Gamma$ or $w\Gamma$ has an
infinitesimal bounded allocation on $(\Gamma_d,d_2,J_c)$.
This includes $d_0=0$ and $b_0=0$ algebraically, with no false
strict pivot on an unweighted $\Gamma$.

\paragraph{The angular row at either $h$ type.}
If $h$ is infinitesimal, SP10 gives
$\alpha_\kappa\in\operatorname{Inf}_c$ and SP9 gives
$\lambda_0\sim-2H_0$; hence $K_A\sim2H_0>0$.  Retain
$K_AA_\theta^2$ as a strict additional row.
If $h$ is a positive bounded unit, $\lambda\sigma$ is a unit and
the following exact conductor replaces that row:
\[
 A_\theta=
 \frac{qcw}{\sqrt{\lambda\sigma}}d_1
 +\frac{k_c^2}{\sqrt\lambda}d_2
 -\frac1{\sqrt\lambda}d_3
 +(k_c+1)\sqrt x\,z^{1/6}J_c.                 \tag{CT13}
\]
All four coefficients are infinitesimal.  Since $K_A$ is bounded,
its entire possibly signed square is then an infinitesimal bounded
perturbation of the circuit/radial frame.  CT13 uses CT8 and
$(k_c-1)G=g_0J_c$, not a division by a contact difference.

\paragraph{Full-parent reassembly.}
The strict bounded principal frame consists of $N_x$, the two raw
radial roots $(N_y,J_c)$ with the clipped CT7 Gram, $\Gamma_d$,
and the three circuit roots $d_i$.  Add $A_\theta$ when $h$ is
infinitesimal.  The other nine CT5 squares and both final CT4
squares are retained.  The following original positive packets
are also retained whole: $4\beta D_wN_x^2$,
$2(2D+1)(uX_e)^2$, the nonconstant positive $N_x^2$ excess,
$4bpzu^2\Gamma^2$, and all radial source-weight excesses above.

There are no unallocated SP6 terms.  Its remaining Gamma normals
are exactly $-4u\nu N_x\Gamma$ and
$qwR_\gamma(N_y+bR)\Gamma$; their coefficients $\nu,R_\gamma$
are bounded, and CT10--CT12 allocate them infinitesimally.
The $N_xN_y,N_xR,N_xA_\theta,N_yA_\theta$ entries have factors
$uw,uw,wB_x,u$, respectively, with bounded remaining coefficients;
$B_x=8\kappa e$ is bounded.  The last angular-radial term is the
complete $\chi uA_\theta R$, with $\chi$ bounded.  CT6 and,
when needed, CT13 give bounded allocations for all these roots.
The remaining radial correction is bounded infinitesimal by CT7
and LC2, and the circuit cost is $-\lambda xR^2$ with
$\lambda x$ infinitesimal.  Consequently the complete remaining
Gram is a bounded infinitesimal perturbation of the stated strictly
positive residue Gram.  The separated bounded-Gram lemma proves
finite tensor SOS.  Restore every retained square to obtain the
entire original parent.

All identities are exact at the sphere endpoints; no coefficient
has been inferred positive by ordering a mixed element.  Source
openings occur only in retained positive weights or the bounded
opening-cleared combinations.  The proof is pointwise in the
testing experiment, so separated cube roots and state ratios are
legal.  Real/imaginary splitting supplies complex states, with the
appropriate extreme-ratio or positive-unit chart chosen for each part.
The bounded-$W/Z$ part of the unbounded-$H$ source and the
bounded-$H$ source are separate obligations and are not claimed here.


\section{\texorpdfstring{C14.2: pure-testing state exits at the two small angles}{C14.2: pure-testing state exits at the two small angles}}
\label{component:p8655d641d4a7}
% Edited scoped mathematical packet, 22 October 2026.

Keep the unchanged C14.2 source and the exact quarter-parent BE10,
with $d=1$, $L$ a positive bounded source unit, $x,\sigma>0$
infinitesimal, and every original opening.  In particular $D,\beta\ge0$
may vanish.  The R18 LA and LB admissions leave $u,w>0$
infinitesimal.  Literal axes have the original direct core SOS.
The original pure-testing core-sign exit leaves $WZ>0$; consequently
$P=uW$ and $T=wZ$ are nonzero and have the same sign.

\paragraph{The two extreme state ratios.}
Choose positive real numbers $c_*,a_*,a'_* $ below the actual
source weights $C,A,\alpha$, respectively; this is possible by LA1.
Retain all positive original excesses and the entire $m,D,\beta$
packets.  The selected polynomial is exactly
\[
 H_*=c_*D_w(P-Y)^2+a_*D_uD_wY^2
       +a'_*D_u(Y-T)^2+K_w/2.                     \tag{CS1}
\]
Here LA2 gives the complete signed core, without approximation:
\[
 K_w=D_uX^2+D_uY^2+u^2T^2-2u^2TY-2PT
       +w^2\{(P-Y)^2+X^2\}+2uwX(T+P-Y).          \tag{CS2}
\]
If $P/T$ is unbounded, write $T=\epsilon P$ with the
\emph{pure-testing} infinitesimal $\epsilon=T/P$.  The bounded
Gram of CS1 on $(X,Y,P)$ has residue
\[
 \tfrac12X^2+c_*(P-Y)^2+(a_*+a'_*+\tfrac12)Y^2,
                                                               \tag{CS3}
\]
which is strictly positive.  If $P/T$ is infinitesimal, write
$P=\epsilon T$.  On $(X,Y,T)$ the residue is
\[
 \tfrac12X^2+(c_*+a_*+\tfrac12)Y^2+a'_*(Y-T)^2,  \tag{CS4}
\]
also strictly positive.  All coefficients of CS1 are pure testing
and real, so the usual finite positive-Gram lift is directly legal.
Restoring the retained original squares proves the complete parent
on both charts, for arbitrary remaining state coordinates.


The remaining positive testing ratio $uW/(wZ)$ is a bounded
unit.  On the unbounded-$H$, unbounded-$W/Z$ source/test chart,
the full CT certificate applies to every such unit ratio; its proof
does not require matching its residue to $L$.  The bounded-$H$
and bounded-$W/Z$ alternatives are the retained HP/HG and CF
certificates.  These uses are detailed in the complete C1.4 join.


\section{\texorpdfstring{Two complementary C14.2 certificates}{Two complementary C14.2 certificates}}
\label{component:p98e371f5b846}

\label{p98e371f5b846:sec:r19-c142-complements}

Throughout this section the source is the unchanged physical C14.2
source, with $z=L^{-1}$ a positive bounded unit, $x,\sigma>0$
infinitesimal, and all original openings and intersections retained.
The remaining tests have $u,w$ infinitesimal. We use the full SP6
identity, its retained squares, and the exact opening-uniform radial
decomposition CT7--CT7a. In particular $N_y,R$ have a strict
pure-source radial form through every opening. Every inverse below is
pure source or pure testing.

\subsubsection*{Unbounded $H=h/\sigma$ and bounded $W/Z$}

Suppose $H$ is unbounded and $Z\ne0$ with the pure-testing quotient
$f=W/Z$ bounded. Write
\[
 s=x/\sigma,\quad L_h=\lambda/H,\quad
 K_A=4\kappa H-\lambda,\quad \rho=\lambda x/\sigma.
\]
BE7 and SP3 imply $s$ infinitesimal, $\lambda>0$ unbounded,
$L_h\sim2$, and $\rho$ bounded positive. The complete last line of
SP6, together with its angular diagonal, is exactly
\[
 K_A A^2+\chi uAR+\mathscr P,\qquad
 \mathscr P=\lambda\sigma P_\beta^2
       +\rho(u^2-x)(w^2-\sigma z)W^2.                 \tag{CF1}
\]
Here $A=q(u^2-x)W$, $q^2=z$. The source polynomial divisibility
SP10, now with $z$ a unit, gives the exact residue formulas
\[
 (K_A)_0=2\{H_0-h_0(H_0+1-r)_0\},\qquad
 \chi_0=4\{H_0-h_0(1+z-2r)_0\}.                    \tag{CF2}
\]
The notation on the right means that every bounded source quantity
is evaluated in its real residue. In particular $K_A,\chi$ are
bounded, and $h$ infinitesimal implies $(K_A)_0=2(H_0)_0>0$.

Put
\[
 U_h=\sqrt h\,uZ,\quad Q_\lambda=\sqrt\rho\,uwW,
 \quad M_\lambda=\sqrt\lambda\,xqW.
\]
The following is the complete scalar identity, including both
negative terms:
\[
\begin{split}
 \mathscr P={}&L_h\{(1-szf^2)U_h-\sqrt h\,sqfR\}^2
       +Q_\lambda^2+M_\lambda^2-L_hszf^2U_h^2\\
 &\hspace{7mm}-L_hs^2\{\sqrt h\,qfR+zf^2U_h\}^2.
\end{split}                                                   \tag{CF3}
\]
All corrections to the leading $L_hU_h^2$ on the $(U_h,R)$
frame have bounded infinitesimal coefficients. This uses $s$ infinitesimal
and $q,z,f$ bounded; it does not require $z$ infinitesimal.

Every Gamma normal is also retained. Set
$B_f=(\sigma-xzf^2)/\sqrt h$. The source bounds
$\sigma/\sqrt h$ and $x/\sqrt h=s\sigma/\sqrt h$ infinitesimal
show that $B_f$ is bounded infinitesimal. Exactly,
\[
\begin{split}
 u\Gamma&=B_fU_h-uqfA-u^2qfR,\\
 qw\Gamma&=qzfB_fU_h-uz^2f^2A
                 +z(\sigma-ufw-u^2zf^2)R.
\end{split}                                                   \tag{CF4}
\]
All coefficients in CF4 are infinitesimal because $u,w$ are
infinitesimal. These identities allocate both hard first normals,
the radial Gamma coupling, and the complete centered-$C$ correction.
The nonnegative diagonal
$(4D/\sigma^2+4bpzu^2)\Gamma^2$ remains as exact squares.

If $h$ is infinitesimal, retain $Q_\lambda^2+M_\lambda^2$
separately. On $(N_x,\text{two radial roots},U_h,A)$ the complete
bounded Gram has strictly positive residue: the first pivot is $2$,
the radial form is strict, the $U_h$ pivot is $2$, and the $A$
pivot is $2(H_0)_0$. Every other cross coefficient is infinitesimal.
If $h$ is a positive unit, use instead the exact identity
\[
 A=\frac{qfu}{\sqrt h}U_h-\frac{1}{\sqrt\lambda}M_\lambda.
                                                               \tag{CF5}
\]
Both coefficients are infinitesimal. Keep $M_\lambda^2$ in the Gram
and retain $Q_\lambda^2$ separately. The bounded signed angular
packet in CF1 then contributes only infinitesimal Gram entries.
The new $M_\lambda$ pivot is $1$.

For completeness, the remaining full SP5 terms are $N_xN_y$ with
coefficient $O(uw)$, $N_xR$ with coefficient $O(uw)$,
$N_xA$ with coefficient $qwB_x$ where $B_x=8\kappa e$ is
infinitesimal, and $N_yA$ with coefficient $uA_y$, where $A_y$ is
bounded. The entire $N_yR$ coefficient is included in CT7--CT7a.
Thus all coefficients just asserted small have explicit source or
testing small factors after the opening poles have been extracted.
Restore $2(2D+1)(uX_e)^2$, the original beta packet, and every
radial excess. The separated bounded-Gram lemma proves the entire
parent. The literal $Z=0$ chart has the inherited core-sign exit.
This proves every bounded-$W/Z$ state for the unbounded-$H$ source.

\subsubsection*{Bounded $M_\beta=\beta\theta/\sigma^2$}

Suppose $M_\beta$ is bounded. BE3--BE7 imply, uniformly through
all openings,
\[
 H\sim H_0,\qquad \theta\sim1,\qquad
 D/\sigma^2\sim H_0/2,\qquad \mu=\kappa H\sim H_0/2.
                                                               \tag{CF6}
\]
Use the full generic JR5/JB1 identity before the SP6 pole extraction,
with $V=uw(u^2-x)W$. Its coefficients are all bounded after the
same exact radial extraction. The only cross coefficient which can
have nonzero residue outside the radial block is $S\Gamma V$.
Write
\[
 S_*=\frac{2}{\sigma^2}
 \{h^2(2D+1)x+4Dh-2bp\sigma^2z
                         +2\beta\theta^2+h^2\}.
                                                               \tag{CF7}
\]
This is a bounded pure-source number, and
$S-S_*=2h^2(2D+1)u^2/\sigma^2$ is infinitesimal.
Add the exact zero JR6 multiplied by the positive bounded
pure-source number $T_*=1+S_*^2$:
\[
 T_*\{V^2-\sigma u^2A^2
       -u\sigma(u^2-x)AR+(u^2-x)V\Gamma\}=0.        \tag{CF8}
\]
All new coefficients except the $V^2$ pivot are infinitesimal.
The $(\Gamma,V)$ residue Gram is
\[
 \begin{pmatrix}
 2(H_0)_0&(S_*)_0/2\\
 (S_*)_0/2&1+(S_*)_0^2
 \end{pmatrix},
\quad
 \det=2(H_0)_0+
          \{2(H_0)_0-1/4\}(S_*)_0^2>0.              \tag{CF9}
\]
Here $H_0=1+z(1-b)\ge1$. The $A$ residue pivot is $2(H_0)_0$,
the $N_x$ pivot is $2$, and the exact radial form is strict.
All remaining generic JR/JB crosses carry $u$, $w$ or a source
infinitesimal: in particular $uA_x$, $Z_x$ and $u\sigma B$ are
infinitesimal, by CF6 and SP4. The retained beta and cap squares
are unchanged. Hence the complete original parent is tensor SOS
on this source domain, for every remaining state and test.

Finally, if $H$ is bounded and the exact SP3 coefficient $\lambda\le0$,
its defining identity implies
\[
 4\beta\theta^2+2h^2+8Dh\le4bp\sigma^2z.
\]
Since $\theta\sim1$, this implies $M_\beta$ bounded, so CF6--CF9
cover that whole case, including $\lambda=0$. No inverse of
$\lambda$, $D$ or $\beta$ is introduced.

All these identities apply with the same source to each real state
part. Together with the sign exits, they restore arbitrary complex
states. The proofs do not compare a source rate to a testing rate.


\section{\texorpdfstring{C14.2: bounded $h/\sigma$ and the unbounded $W/Z$ state chart}{C14.2: bounded h/sigma and the unbounded W/Z state chart}}
\label{component:pcb3428d66a71}

\label{pcb3428d66a71:sec:r19-c142-bounded-H}

Retain the full original C14.2 source, with $z=L^{-1}$ a positive
bounded unit, $x,\sigma$ infinitesimal and every original opening,
coefficient boundary and residual intersection unchanged.  Suppose
$H=h/\sigma$ is bounded.  Then $H$ is a positive unit by BE7 and
$h$ is infinitesimal.  The following full-parent certificate applies
to $u,w$ infinitesimal and to every real state with $W\ne0$ and
$t=Z/W$ infinitesimal, including $t=0$ and both signs.  It makes
no claim for the complementary state chart or unbounded $H$.

Use the exact generic SP6 parent, with $A=A_\theta$ and
\[
 d=4D/\sigma^2,\qquad \rho=\lambda x/\sigma.
\]
All source coefficients in the following calculations are actual
coefficients; formal caps are only exact square extractions.

\paragraph{The sphere restores a strict Gamma direction.}
The identity $uP_\beta=(u^2-x)Z+(x/\sigma)\Gamma$ gives exactly
\[
 \lambda u\Gamma P_\beta
       =\rho\Gamma^2+\frac{\lambda t}{q}A\Gamma.
 \tag{HG1}
\]
There is no state inverse beyond the pure-testing coordinate $W^{-1}$.
The complete angular/Gamma packet of SP6 is consequently
\[
 (d+\rho)\Gamma^2+4\kappa H A^2
       +\frac{\lambda t}{q}A\Gamma
       +(\lambda+\chi)uAR.
 \tag{HG2}
\]
No signed scalar product was discarded in this transformation.

To control the formerly missing Gamma pivot uniformly, BE4 implies
the exact source identity
\[
 \rho=\frac D{\sigma^2}(8hx-4\theta J)
       +\frac{2xh^2}{\sigma^2}
       +\frac{2\theta\mathcal H}{P}-4bp xz.
 \tag{HG3}
\]
Also, directly from BE7,
\[
 \frac{xh^2}{\sigma^2}
    =\frac{eH_0}{B}+\frac{h\mathcal H}{B}b_0^2,
 \qquad e=xh/\sigma\in\mathrm{Inf}.
 \tag{HG4}
\]
Since $\operatorname{res}(D/\sigma^2)=
\tfrac12\operatorname{res}(H_0d_0^2)$,
$\mathcal H/P\sim H_0$, $J\sim1$ and
$d_0^2+b_0^2=1$, these identities give, even without bounded $H$,
\[
 \operatorname{res}\rho=2\operatorname{res}(H_0b_0^2),\qquad
 \operatorname{res}(d+\rho)=2\operatorname{res}H_0>0.
 \tag{HG5}
\]
In particular no ratio to the possibly vanishing $b_0$ is asserted,
and neither $D$ nor $b_0$ is inverted.  In the bounded-$H$ chart,
$\lambda,\chi$ are bounded,
$4\kappa H$ is a positive unit, and $q$ is a unit.  Both cross
coefficients in HG2 are therefore infinitesimal.

\paragraph{The full radial and hard block through the C opening.}
Restore the centered $A,\alpha$ cap excesses into the $N_y$ diagonal.
Put $a=Az$, $\ell=\alpha z$, $Q=1-b$,
$J_C=N_y+bR$ and $J_\ell=N_y-QR$.
Use the exact LC3 cap $C_*=C/(1+C)$, retaining
$4zD_w(C-C_*)J_C^2$ and its complete correction
\[
 -8(C-C_*)(1-r)J_C(z\sigma R-qw\Gamma).
 \tag{HG6}
\]
The correction has infinitesimal coefficients on
$(N_y,R,\Gamma)$, because $(C-C_*)(1-r)$ is bounded.
The exact leading pure-source radial form is
\[
 M=\{4a+2z(2C_*+1)\}J_C^2+4\ell J_\ell^2.
 \tag{HG7}
\]
Both coefficients have positive real lower bounds; their potentially
unbounded excesses are retained as squares.  The inverse allocations
$R=J_C-J_\ell$ and $N_y=QJ_C+bJ_\ell$ are bounded.  Therefore
$M$ supplies a strict bounded Gram on $(N_y,R)$ through every
$r$, $1-r$, and $Q$ opening.

For clarity, HG7 is obtained from the entire formal radial/hard
coefficient collection by putting $u=w=\sigma=D=m=\beta\theta=h=0$
while retaining the actual pure-source $a,\ell,p=ab,l=\ell Q,z,b,C_*$.
The omitted $u^2$ and $w^2$ multiples of $a,\ell$ are retained
exactly as the nonnegative packets from
\[
 4(1+u^2)\{a(1+w^2)J_C^2+\ell J_\ell^2\}.
\]
All remaining coefficient differences are bounded infinitesimals:
$p,l,z,b,C_*$ are bounded, $m,D,\beta\theta,h,\sigma$ are
infinitesimal, and $u,w$ are testing infinitesimals.  This assertion
is made only after the divergent $C$ correction HG6 has been
accounted for, and after retaining the displayed $a,\ell$ packets.

The $N_x$ diagonal has residue $2$.  Its Gamma normal is
$-4u\nu N_x\Gamma$ with bounded $\nu$; both other Gamma normals
carry $qw$.  All remaining SP5 hard and first-normal coefficients
are infinitesimal after the radial source form HG7 is restored.
The separated Gram on $(N_x,J_C,J_\ell,\Gamma,A)$ consequently has
strict positive residue.  HG2 supplies the last two positive pivots.
The original beta packet, the $2(2D+1)(uX_e)^2$ square, the
$4bpzu^2\Gamma^2$ square, and every source cap excess remain once.
Thus the entire original SP6 parent is included, not only its scalar
restriction.  The bounded separated-Gram lemma proves tensor SOS.

All formulas are literal at zero testing angles and $Z=0$; the
case $W=0$ has the inherited original-core sign exit.  Application
to real and imaginary state parts gives the corresponding complex
state admission when each part lies in this chart or an existing
admission.  This is a source/state subchart result, not a complete
C14.2 closure.


\section{\texorpdfstring{C14.2: bounded $h/\sigma$ with positive $\lambda$, all states}{C14.2: bounded h/sigma with positive lambda, all states}}
\label{component:p021ed4a9e0db}

\label{p021ed4a9e0db:sec:r19-c142-positive-lambda}

Retain the original C14.2 source and all its intersections, and the
complete small-angle corner $u,w\in\mathrm{Inf}$.  Suppose
$H=h/\sigma$ is bounded and $\lambda>0$.  This section proves the
entire original parent for every testing state.  It retains literal
$D=0$, every source-sphere opening, and every allowed coefficient
boundary.  The source case $\lambda\le0$ is a separate obligation.

Use the exact SP6/BS2 parent, with the complete C-opening extraction
LC3 and the actual radial frame HG7.  The source estimates proved
in HG3--HG5 and the SP definitions give
\[
 \rho=\lambda x/\sigma\ge0,\quad \rho\in\mathcal O_c,
 \quad d=4D/\sigma^2,\quad
 d+\rho\sim2H_0,
 \quad K_A=4\kappa H-\lambda\sim2H_0>0.
 \tag{HP1}
\]
Here $\lambda,\chi,\nu$ are bounded and $\lambda x=\rho\sigma$
is infinitesimal.  The source $q=\sqrt z$ is a unit.  The complete
parent outside already positive rows consists of its strict hard/radial
frame, $d\Gamma^2+K_A A_\theta^2+\chi uA_\theta R$, the complete
Gamma normals, and the signed packet
\[
 \mathscr P=\lambda\sigma P_\beta^2+
      \rho(u^2-x)(w^2-\sigma z)W^2.
 \tag{HP2}
\]
No source or testing rate is compared across the two fields.

\paragraph{A ternary circuit that supplies an actual difference pivot.}
The original-core sign exit and exact testing gauges reduce the
remaining real states to $u,w,W,Z>0$.  Literal axes and $WZ\le0$
already have direct original-core SOS.  Let
\[
 a=\sqrt{\lambda\sigma}\,uZ,\qquad
 b=\sqrt\rho\,uwW,\qquad
 c=\sqrt{\lambda z}\,xW,\qquad g=(abc)^{1/3},
 \qquad r^3=\frac{wZ}{zuW}.
\]
Every radical and quotient is a product of explicitly positive
pure-source and pure-testing factors.  Direct expansion gives
\[
 \boxed{\quad
 \mathscr P+\lambda xR^2
 =a^2+b^2+c^2-3g^2
  +\lambda xzu^2W^2 r^2(r-1)^2(r^2+2r+3).
 \quad} \tag{HP3}
\]
The last term is a finite legal SOS since $r^2+2r+3=(r+1)^2+2$.
There is a stronger exact decomposition of the first term:
\[
 a^2+b^2+c^2-3g^2
 =\tfrac12\{(a-b)^2+(b-c)^2+(c-a)^2\}
       +(ab+bc+ca-3g^2).
 \tag{HP4}
\]
The final parenthesis is itself a finite legal SOS.  Indeed, put
$\alpha=(ab)^{1/3}$, $\beta=(bc)^{1/3}$,
$\gamma=(ca)^{1/3}$.  It is exactly
\[
 \tfrac12(\alpha+\beta+\gamma)
    \{(\alpha-\beta)^2+(\beta-\gamma)^2+(\gamma-\alpha)^2\}.
 \tag{HP5}
\]
Expand the positive prefactor into its three monomial terms before
taking square roots; there is no radical of a mixed sum.  Thus HP3
supplies a genuine unit-weight pivot
$D_1=(a-b)/\sqrt2$, with every other displayed row retained.
The small subtraction $-\lambda xR^2$ is an infinitesimal Gram
perturbation of the strict radial frame.

\paragraph{The source-sphere split.}
If $d_0$ is a positive unit, $d\Gamma^2$ supplies a strict Gamma
pivot.  Every complete Gamma normal in SP6/LC3 carries $u$, $w$
or $\sigma$ and has a bounded remaining coefficient.  Hence all
these normals are infinitesimal entries on the strict hard/radial,
Gamma and angular frame.  Retain all of HP3 as separate SOS.
Together with $K_A\sim2H_0$ this closes this source case for every
state, without any state quotient.

It remains that $d_0$ is infinitesimal, including $D=0$.
Then $\rho$ is a positive unit by HP1.  If $W/Z$ is unbounded,
the full-parent HG certificate applies.  Thus consider bounded
$f=W/Z$.  The next argument supplies every missing Gamma normal
without imposing a positive pivot on $\Gamma$.

\paragraph{A quadratic domination conductor, without a mixed inverse.}
Put $k=\sqrt x$, and define
\[
 T=q u(u-k)W,\qquad \delta=w/q>0,
 \qquad T_1=\sqrt\delta\,T.
\]
All these roots are separated and $\delta$ is infinitesimal.
The exact circuit and angular identities are
\[
 u\Gamma=\frac{\sqrt2 k}{\sqrt\rho}D_1-\delta T,
 \tag{HP6}
\]
\[
 A_\theta^2-T^2
   =q^2(u-k)^2W^2(k^2+2uk)
   =(qk(u-k)W)^2+2(q\sqrt u\sqrt k(u-k)W)^2.
 \tag{HP7}
\]
Reserve the infinitesimal part $\delta A_\theta^2$ of the strict
angular diagonal.  Exactly,
\[
 K_A A_\theta^2=(K_A-\delta)A_\theta^2
         +T_1^2+\delta(A_\theta^2-T^2).
 \tag{HP8}
\]
The final term is a legal SOS by HP7; the first coefficient has
unchanged positive residue.  Thus $T_1$ is a genuine additional
unit-weight pivot, paid for exactly once from the angular row.
HP6 becomes the bounded infinitesimal allocation
\[
 u\Gamma=\frac{\sqrt2 k}{\sqrt\rho}D_1
                   -\sqrt\delta\,T_1.
 \tag{HP9}
\]
This avoids division by $u+k$.  A pointwise bound on
$u/(u+k)$ alone would not be a legal tensor allocation.

The other full Gamma normal has the exact state conductor
\[
 qw\Gamma=qzf\,u\Gamma+z(\sigma-uwf)R.
 \tag{HP10}
\]
Its first coefficient is bounded and the second is infinitesimal.
Together HP9--HP10 allocate both $u\Gamma$ and $qw\Gamma$ with
infinitesimal coefficients on $(D_1,T_1,R)$.  The only state
inverse used is the pure-testing $Z^{-1}$ in the bounded quotient
$f$; no source/testing expression is inverted.

\paragraph{Full-parent reassembly.}
The entire Gamma coupling in SP6 is
\[
 -4\nu N_x(u\Gamma)+R_\gamma(N_y+bR)(qw\Gamma).
\]
LC3 adds only its recorded bounded-$t_C$ correction to the same
$qw\Gamma$ normal and to $\sigma R$.  Substitution of HP9--HP10
makes every one of these couplings infinitesimal on the strict frame
\[
 (N_x,J_C,J_\ell,A_\theta,D_1,T_1).
\]
HG7 supplies the radial/hard positive source rows; all divergent
$a,\ell$ angle-dependent excesses and the centered C excess are
retained as in HG.  The other complete SP5 couplings and
$\chi uA_\theta R$ are bounded infinitesimals.  The subtraction
$-\lambda xR^2$ is infinitesimal by HP1.  Consequently the whole
frame has strict positive residue.

Retain exactly once every remaining positive row: $d\Gamma^2$,
$4bpzu^2\Gamma^2$, $2(2D+1)(uX_e)^2$, the original beta packet,
the unused HP3--HP5 circuit rows, the HP7 remainder from HP8, and
all source-cap excesses.  In particular the $(a-b)^2/2$ row and
the angular fraction used for $T_1$ are not counted twice.
The separated bounded-Gram lemma now proves the complete parent.

The two sphere cases exhaust the source, and the bounded/unbounded
$W/Z$ alternatives exhaust the states after the direct sign/zero
exits.  Real/imaginary splitting gives all complex states.  Thus all
remaining C14.2 tests are admitted on bounded $H$ and $\lambda>0$.
The result is conditional only on the explicitly stated source split,
not on an unproved state-rate comparison.


\section{\texorpdfstring{C14.4: extend the unmatched-Bob certificate to the entire small-$e$ source}{C14.4: extend the unmatched-Bob certificate to the entire small-e source}}
\label{component:p8c30238a5f7d}

\label{p8c30238a5f7d:sec:recovery-c144-small-e}

This is a partial-domain theorem. It retains the original physical C14.4
source, all strict gaps, nonnegative event weights, the original total-scale
and prior-complement intersections, and independent source and testing fields.
Its source has $x,L>0$ infinitesimal, $\zeta=L/x$ noninfinitesimal and
$\eta$ unbounded. Write $\sigma=1/\eta$, $S=\sigma+L$ and
\[
 e=\frac{xh}{\sigma},\quad E=1-e,\quad
 \kappa=\frac{x}{S},\quad \nu=\frac{\sigma}{S},\quad
 \ell=\frac{L}{S},\quad \nu+\ell=1.
 \tag{SE1}
\]
The letter $\kappa$ in this section denotes $x/S$, not the different
coefficient denoted by $\kappa$ in the R19 reciprocal-$L$ chart.
All source inequalities are taken in the coefficient field. Real residues
are compared only after taking residues separately in the two fields.

\paragraph{The exact source input.}
Use ES1--ES10 from the retained R13 source proof. In particular
$1/18\le E<1$, $0<h,r<1$, $\theta=1-h$, and
\[
\begin{gathered}
 J=1+xh-\sigma(1-r),\qquad M=1+xr-\theta/\zeta>0,\\
 Q=1-Lr,\quad H_0=1+L(1-r),\quad B=P+x(J-\sigma),\\
 P=\sigma(1+L+2xh+rxL+rx^2h)+(1+x)L+x^2h\theta,\\
 \mathcal H=(1+xr)H_0+(1-r)xh(2+xr)
       -e\{2+x(1+h+r\theta)\}+(1+x)e^2.
\end{gathered} \tag{SE2}
\]
The opening-safe sphere identities are
\[
\begin{gathered}
 2\theta P\beta=hB-\sigma H_0,\qquad
 d_0^2+b_0^2=1,\quad d_0,b_0\ge0,\\
 D=\frac{\sigma^2\mathcal H}{2PJ}d_0^2,\qquad
 x\theta\beta=\frac{\sigma^2\mathcal H}{2P}b_0^2,\\
 e\frac BS=\kappa H_0+\nu\mathcal H b_0^2,\qquad
 e^2\frac BS=\kappa(eH_0+h\mathcal H b_0^2).
\end{gathered}\tag{SE3}
\]
Here $P/S\sim1$, $J,Q,H_0\sim1$, $B/S-(1+\kappa)$ is
infinitesimal, and $\mathcal H-E^2$ is infinitesimal. In particular
$B/S$ and $\mathcal H$ are positive bounded units. These estimates
are the already admitted ES bounds; they do not use bounded $\eta$.

\paragraph{Source dichotomy.}
On the entire C14.4 source,
\[
 e\text{ is infinitesimal}\quad\Longleftrightarrow\quad
 \kappa=x/S\text{ is infinitesimal}.\tag{SE4}
\]
Indeed the first nonnegative sum in SE3 gives $\kappa=O(e)$,
and the second gives $e^2=O(\kappa)$. Both implications are
ordered-field statements entirely inside the source field.

Suppose $e$ is infinitesimal. If $\zeta$ is unbounded, then
$\theta/\zeta$ is infinitesimal and $M\sim1$ directly. If $\zeta$
is a positive bounded unit, then the exact relation
\[
                  \ell=\zeta\kappa\tag{SE5}
\]
gives $\ell$ infinitesimal and $\nu\sim1$. Consequently
$B/\sigma\sim1$. Since $\beta\ge0$, the first identity of SE3
implies $hB\ge\sigma H_0$, whence $h\sim1$ and $\theta$ is
infinitesimal. Thus $M\sim1$ in this case as well. We have proved,
uniformly on the whole small-$e$ source,
\[
                    E\sim1,\quad M\sim1,\quad m\in\operatorname{Inf}_{>0}.
                    \tag{SE6}
\]
For the last statement use the exact original coefficient
\[
 m=\frac{(1+x)(1+\sigma r)eS}{2M\zeta P};\tag{SE7}
\]
$1/\zeta$ is bounded and every other coefficient apart from $e$ is
bounded, including the reciprocal of $M$ just proved to be a unit.
No assumption on the ratio of a source infinitesimal to a testing
infinitesimal occurs.

The two nonnegative terms in the first sum in SE3 imply
$\nu b_0^2$ infinitesimal. Hence the exact sphere formula and the
exact formula for $\alpha$ yield
\[
 \frac DS-\frac{\nu^2}{2}\in\operatorname{Inf},\qquad
 \frac\alpha S-\frac{\ell^2}{2}\in\operatorname{Inf},\qquad
 SC(1-r)\sim\tfrac12,\qquad SAr\sim\tfrac12.\tag{SE8}
\]
For clarity the second formula is obtained without dividing by $\ell$:
\[
 \frac\alpha S=
 \frac{(1+\sigma r)\{\ell(1+xr)e+\ell^2ME\}}
      {2Q(P/S)}.\tag{SE9}
\]
These are additive bounded residue estimates, including $D=0$ whenever
that boundary has points in the stated domain. They do not assert a
relative estimate at a vanishing pivot.

\paragraph{Theorem: every unmatched Bob residue on the small-$e$ source.}
Let $e$ be infinitesimal. In the testing field choose normalized homogeneous
Bob coordinates $(c,s)$ with $c^2+s^2=1$. For every independent projective
test and every real or complex state satisfying
\[
                         \overline{s^2}\ne\overline\ell,\tag{SE10}
\]
the complete original C14.4 parent is a finite sum of squares in
$K_c\otimes_{\mathbb R}K_t$.

\paragraph{Proof with the complete parent retained.}
Noninfinitesimal Alice tests use the unchanged adjacent-$C,A$ theorem:
SE8 and $S$ infinitesimal imply that both $C$ and $A$ are unbounded.
For infinitesimal or zero $u$, apply the exact EB3--EB7 homogeneous
frame from the retained R18 proof. Those formulas are algebraic identities;
their only source-limit inputs to the bounded-Gram proof are precisely
$m$ infinitesimal and SE8, now established on the larger source domain.
To make the imported implication explicit, put
\[
 H=S(C+A),\quad \rho=\frac C{C+A},\quad a_h=\frac{SCA}{C+A},
 \qquad p=\frac{uW}{\sqrt S},\ q=\sqrt S Z,\ y=\frac Y{\sqrt S}.
\]
Then $1/\sqrt H$ is bounded, $0<\rho<1$, $a_h\sim1/2$, and
the harmonic normal $N_y=\sqrt H(y-\rho p)$ can be substituted
by $y=\rho p+N_y/\sqrt H$ with bounded source coefficients.
The remaining normals are the EB4 $N_m,N_b,N_r$, with only pure-field
radicals and the positive testing denominators $1+u^2$ and $1+u^2s^2$.
The exact full quarter-parent is EB7,
\[
 N_y^2+a_hp^2+ASu^2y^2+M_B^2+D_B^2+A_B^2+\beta(X+uZ)^2
                         +\tfrac12 K_B.\tag{SE11}
\]
All EB5 core terms and all EB6 event amplitudes are retained. In particular
no state graph is imposed and neither first-normal coupling is removed.
The two positive rows $ASu^2y^2$ and $\beta(X+uZ)^2$ are left whole.
On the other six roots $(N_y,N_m,N_b,N_r,p,q)$ the bounded Gram has
residue equal to $\operatorname{diag}(1,1,1/2,1/2)$ plus the scalar block
\[
 \frac12\begin{pmatrix}
 1&-\bar s\bar c\\
 -\bar s\bar c&\bar\nu^{\,2}\bar s^{\,2}
                            +\bar\ell^{\,2}\bar c^{\,2}
 \end{pmatrix}.
\tag{SE12}
\]
Its determinant is $(\overline{s^2}-\bar\ell)^2/4$, so SE10
makes this finite coefficient Gram strictly positive definite. The
separated bounded-Gram lemma supplies tensor SOS; restoring the two
whole rows gives the complete parent. The exact homogeneous formulas
allow $c=0$, $s=0$ and $u=0$. Signs and real/imaginary state parts are
handled exactly as in EB, and the original positive source and testing
normalizations are pure-field factors. This proves the stated domain.

\paragraph{New part, and exact retained complement.}
Unbounded $\zeta$ already forces $\kappa\le1/\zeta$ infinitesimal,
so the new theorem contains the entire R18 unmatched-residue admission.
Its additional domain is $\zeta$ a positive bounded unit, $e$
infinitesimal, and Bob sine square with nonzero residue; this includes
Bob infinity. In this added domain $\ell$ has zero residue, so the
remaining Bob chart is exactly finite infinitesimal $v$ (including zero,
which may be handled by independent literal-axis exits).

After the unmatched-residue admissions, the complementary C14.4 domain is contained in the union of (i) $e$ infinitesimal and matched Bob residue, and (ii) $\zeta$ a positive bounded unit and $e$ a positive bounded unit. All original intersections remain in both. In case (ii), SE7 shows that $m$ has a positive real lower bound, while its upper bound may diverge through $M=0$. This is a source estimate. The complete certificates for both complementary domains, with their exact source partition, are joined in Section~\ref{component:pfd0dd7bb10cf}.


\section{\texorpdfstring{State-free normal conductors at the large-$L$ corner}{State-free normal conductors at the large-L corner}}
\label{component:p3be5a8005801}

\label{p3be5a8005801:sec:c141-state-free}

\paragraph{Claims and retained input.}
Keep the original physical stationary C1 source, its positive $d$,
strict gaps, nonnegative event coefficients, original infinitesimal
total scale, and the complement of the previously admitted domains.
Normalize $d=1$. Source quantities belong to a real closed extension
$K_c$ and testing angles and states to an independent extension $K_t$.
The conclusion is a finite sum of squares in
$K_c\otimes_{\mathbb R}K_t$, not pointwise positivity after identifying
the fields. In this section $x,z=1/L,\sigma=1/\eta$ are all positive
infinitesimals. All original source openings are retained, including
the allowed literal faces $D=0$ and $\beta=0$.

The first theorem closes the entire original C14.1a residual. The
second closes a proper part of C14.1b. The precise remaining C14.1b
complement is stated at the end. No global theorem is claimed.

\begin{theorem}[C14.1a]
If $H=h/\sigma$ is bounded, the complete original parent is tensor
SOS for every independent projective test and every real or complex
state. In particular the previously unresolved C14.1a domain,
$M_\beta=\beta\theta/\sigma^2$ unbounded, is admitted completely.
\end{theorem}

\begin{theorem}[A further C14.1b subdomain]
If $H$ is unbounded and the real residue of $h$ is strictly less
than $1/2$, the same full-parent conclusion holds. This includes
every infinitesimal $h$ and every positive unit $h$ with residue
less than $1/2$, but does not include residue $1/2$ or greater.
\end{theorem}

The proof uses the exact complete P4--P8/SP6 parent in the inherited
Section~\ref{component:p472f434581f3}, the BE1--BE9 source identities, and the
already admitted remote and literal-angle exits. The new argument
settles the remaining $u,w=1/v$ infinitesimal corner directly; it does
not reclassify any source outside the original C14.1 domains.

\subsubsection*{Source bounds with $z$ infinitesimal}

Write $q=\sqrt z$, $k=\sqrt x$, $b=Lr$, $Q=1-b$,
$\theta=1-h$, $e=xh/\sigma$, and retain the exact source letters
$J,M,P,B,\mathcal H,p,l,\kappa$ from P1--P2. Since $r=bz$ and
$0<b<1$, $r$ is infinitesimal and $1-r$ is a unit. BE1 applies
because it requires only that $L$ have a positive real lower bound;
it gives $e=O(x+\sigma)$ and $E=1-e\sim1$. Consequently
\[
 P,J,M,B,\mathcal H\sim1,\qquad
 p,l\sim\tfrac12,\qquad C=O(z),\qquad\kappa\sim\tfrac12.
 \tag{SF1}
\]
The possible $b^{-1},Q^{-1}$ poles occur only in the two original
positive source weights $a=Az=p/b$ and $\ell=\alpha z=l/Q$.
Both have positive real lower bounds. They need not be bounded above.
No cap of an individual unbounded matrix entry is used below.

Let $d_0,b_0$ be the exact source sphere from P2, and put
\[
 d=4D/\sigma^2,\qquad
 \lambda=2H+\lambda_0,\qquad
 \rho=\lambda x/\sigma,\qquad
 K_A=4\kappa H-\lambda.
 \tag{SF2}
\]
The identities SP9--SP10 show that $\nu,\lambda_0,\chi,K_A$ are
bounded even when $H$ is unbounded. This is a polynomial
divisibility argument, not multiplication of an uncontrolled error
by $H$. More precisely their source residues give
\[
 (\nu)_0=1,\qquad (\lambda_0)_0=-2(1-h_0),\qquad
 (K_A)_0=2-4h_0.
 \tag{SF3}
\]
For example the SP10 polynomial $F_\sigma$ has residue $-1$ at
$x=z=\sigma=0$, whereas $eF_x$ has zero residue. Thus
$\alpha_\kappa=2H(2\kappa-1)$ has residue $-2h_0$, which proves
the last equality without assuming a rate for $2\kappa-1$.

The source-sphere calculations HG3--HG5 are algebraic for bounded
$z$ and do not require it to be a unit. They give
\[
 \rho\text{ is bounded},\qquad
 (\rho)_0=2(b_0^2)_0,\qquad (d+\rho)_0=2.
 \tag{SF4}
\]
One way to check the delicate assertion is to use both exact identities
\[
 \rho=\frac D{\sigma^2}(8hx-4\theta J)
       +\frac{2xh^2}{\sigma^2}
       +\frac{2\theta\mathcal H}{P}-4bpxz,
 \qquad
 \frac{xh^2}{\sigma^2}=\frac{eH_0}{B}
                         +\frac{h\mathcal H}{B}b_0^2.
 \tag{SF5}
\]
These include $D=0$ and $\beta=0$. Throughout the positive-$\lambda$
case $\lambda x=\rho\sigma$ is infinitesimal.

When $H$ is bounded, $h$ is infinitesimal, $K_A\sim2$, and
$\lambda$ is bounded. If $\lambda\le0$, its exact defining formula
implies
\[
 8Dh+4\beta\theta^2+2h^2\le4bp\sigma^2z.
\]
Since $\theta\sim1$, this forces bounded $M_\beta$; the inherited
JB bounded-$M_\beta$ all-tests theorem applies. On the actual
C14.1a remainder $M_\beta$ is unbounded, so $\lambda>0$.

If $d_0$ is infinitesimal, including literal $D=0$, SF4 makes
$\rho$ a positive bounded unit. The positive sphere identity
\[
 H=H_0/B+(\sigma/x)(\mathcal H/B)b_0^2
 \tag{SF6}
\]
then gives the source bounds needed for the two theorems:
\[
 \begin{array}{ll}
 H\text{ bounded}:&\sigma/x\text{ bounded};\\
 H\text{ unbounded}:&x/\sigma\text{ infinitesimal},\quad
                    xh/\sigma^2\text{ a positive unit},\quad
                    \sigma/k\asymp\sqrt h\text{ bounded}.
 \end{array}
 \tag{SF7}
\]
For the second line use $\lambda/H\sim2$ and $\rho\sim2$.
None of these is a comparison between a source rate and a testing rate.

\subsubsection*{The full source-opening-safe radial block}

The roots are the original ones:
\[
 N_x=X+uhZ,\quad N_y=(Y-bwZ)/q,\quad
 \Gamma=\sigma Z-uwW,\quad R=wZ/q-uqW,\quad
 A_\theta=q(u^2-x)W.
\]
Put $J_C=N_y+bR$ and $J_\ell=N_y-QR$. The exact CT7--CT7a
radial extraction, with every $a,\ell$ excess restored, has leading
pure-source form
\[
 [4a+2z(2C+1)]J_C^2+4\ell J_\ell^2.
 \tag{SF8}
\]
Retain the nonconstant positive $u^2,w^2$ packets and all positive
source-weight excesses. Clipping the two coefficients in SF8 at
positive real lower bounds gives a strict bounded Gram on $(N_y,R)$,
since
\[
 R=J_C-J_\ell,\qquad N_y=QJ_C+bJ_\ell,
 \qquad b+Q=1.
\]
This is uniform through $b\to0$ and $Q\to0$. The remaining three
radial coefficients are precisely CT7a. They are infinitesimal
because $z$ is infinitesimal, $C=O(z)$, and all their other source
letters are bounded. In particular the $z$-unit argument's possible
$C$ pole does not occur here.

The $N_x^2$ diagonal has residue $2$. The complete remaining
non-Gamma cross terms of SP6 are
\[
 \begin{split}
 &-4uqw(2m+2D+1)N_xN_y+uqwR_xN_xR\\
 &\hspace{4mm}+qwB_xN_xA_\theta+uA_yN_yA_\theta
                         +\chi uA_\theta R.
 \end{split}
 \tag{SF9}
\]
Every coefficient is a bounded infinitesimal on $u,w$ infinitesimal.
Here $m,D,C,\beta\theta,p,l$ are bounded, $B_x=8\kappa e$ is
infinitesimal, and $A_y$ is bounded. The entire $N_yR$ term was
already included in SF8 and CT7a. No first normal was suppressed.

\subsubsection*{A scalar circuit and two normal conductors with no state quotient}

The original-core sign exit and exact testing sign gauges leave
$u,w,W,Z>0$. Literal $u=0,w=0,W=0,Z=0$ and the opposite-sign
cases have their inherited direct core SOS. Suppose $\lambda>0$.
Define separated nonnegative monomials
\[
 a_c=\sqrt{\lambda\sigma}\,uZ,\qquad
 b_c=\sqrt\rho\,uwW,\qquad
 c_c=\sqrt\lambda\,xqW,\qquad
 g_c=(a_cb_cc_c)^{1/3},\qquad
 s_c^3=\frac{wZ}{zuW}.
 \tag{SF10}
\]
Each radical is taken separately in its source and testing factors.
The exact full angular packet in SP6 is
\[
 K_AA_\theta^2+\chi uA_\theta R+\mathscr P,
 \qquad
 \mathscr P=\lambda\sigma P_\beta^2
       +\rho(u^2-x)(w^2-\sigma z)W^2.
\]
The same algebraic HP3 identity is valid for infinitesimal $z$:
\[
 \mathscr P+\lambda xR^2
 =a_c^2+b_c^2+c_c^2-3g_c^2
   +\lambda xzu^2W^2s_c^2(s_c-1)^2(s_c^2+2s_c+3).
 \tag{SF11}
\]
No bound on $s_c$ or any state ratio is asserted or required. The last
term is a legal SOS after writing $s_c^2+2s_c+3=(s_c+1)^2+2$.
Furthermore
\[
 a_c^2+b_c^2+c_c^2-3g_c^2
 =\tfrac12\{(a_c-b_c)^2+(b_c-c_c)^2+(c_c-a_c)^2\}
      +(a_cb_c+b_cc_c+c_ca_c-3g_c^2).
 \tag{SF12}
\]
The last parenthesis is the elementary cubic AM--GM SOS applied to
$(a_cb_c)^{1/3},(b_cc_c)^{1/3},(c_ca_c)^{1/3}$. Expand its
positive monomial prefactor before taking square roots. Thus
\[
 D_1=(a_c-b_c)/\sqrt2,\qquad
 D_2=(b_c-c_c)/\sqrt2
\]
are two genuine unit-weight pivots; all other SF11--SF12 squares
are retained separately.

The new root is
\[
 T=kq(u-k)W.
\]
Its legal quadratic domination is exact:
\[
 A_\theta^2-T^2=q^2(u-k)^2W^2(u^2+2uk).
 \tag{SF13}
\]
Both summands on the right are squares, using the separated root
$\sqrt u\sqrt k$. If $K_A$ has positive residue, allocate its
diagonal exactly once as
\[
 K_AA_\theta^2=\tfrac12K_AA_\theta^2+\tfrac12K_AT^2
                        +\tfrac12K_A(A_\theta^2-T^2).
 \tag{SF14}
\]
The last term is retained SOS and the first two supply strict
positive source pivots for $A_\theta,T$.

Here are the decisive two identities, valid for arbitrary states
without a state division:
\[
 \boxed{\begin{aligned}
 u\Gamma={}&\frac{\sqrt2 k}{\sqrt\rho}D_1
       -\frac{\sqrt2(u-k)}{\sqrt\rho}D_2-\sqrt\sigma T,\\
 qw\Gamma={}&\sigma zR
       -\frac{\sqrt2(qw+z\sqrt\sigma)}{\sqrt\rho}D_2
       +\left(\frac{\sigma z}{k}
                  +\frac{\sqrt\sigma qw}{k}\right)T.
 \end{aligned}}
 \tag{SF15}
\]
All divisions are pure source. In particular neither $q$ nor a
testing coordinate is inverted in a conductor coefficient. These
replace the $\lambda(Z/W)/q$ and $w/q$ coefficients that obstruct a
literal transport of the R19 $z$-unit proof.

\subsubsection*{Completion when $H$ is bounded}

Here $K_A\sim2$ and $\lambda$ is bounded. If $d_0$ has positive
residue, the original $d\Gamma^2$ diagonal gives a strict Gamma
pivot. Its two complete normals are
\[
 -4\nu N_x(u\Gamma)
       +R_\gamma(N_y+bR)(qw\Gamma).
 \tag{SF16}
\]
Their coefficients are infinitesimal on $(N_x,N_y,R,\Gamma)$.
Keep SF11 entirely as SOS; its sole cost $-\lambda xR^2$ is an
infinitesimal bounded radial correction. The frame
$(N_x,N_y,R,\Gamma,A_\theta)$ has strict positive residue.

If $d_0$ is infinitesimal, $\rho$ is a unit by SF4 and
$\sigma/x$ is bounded by SF7. All coefficients in both lines of
SF15 are bounded infinitesimals. Indeed $\sqrt\sigma/k$ is
bounded, $\sigma/k=k(\sigma/x)$ is infinitesimal, and $q,w,z$
are bounded infinitesimals. Use SF14 and the two unit circuit
pivots. After substituting SF15 into SF16, the complete frame
\[
 (N_x,N_y,R,A_\theta,T,D_1,D_2)
 \tag{SF17}
\]
has strictly positive residue. Its radial block is SF8, its angular
and $T$ pivots come from SF14, and the circuit pivots from SF12.
Every correction is a bounded infinitesimal. The diagonal
$d\Gamma^2$ is retained separately, including at $d=0$.

In both cases retain exactly once $4\beta(1+w^2)N_x^2$,
$2(2D+1)(uX_e)^2$, $4bpzu^2\Gamma^2$, the unused circuit rows,
and all positive radial and coefficient-cap excesses. Equations
SP6, SF8--SF16 list every original scalar, normal and hard coupling.
The separated bounded-Gram lemma proves the entire parent.

The $d_0$ split exhausts the source sphere. The bounded-$M_\beta$
and bounded nonzero-unit Alice-angle exits are the inherited JB
theorems; their CE remote join includes all other projective tests.
Together with the literal-axis and sign exits, this closes the
complete original C14.1a residual and the full bounded-$H$ source
domain. Apply the real argument independently to real and imaginary
state parts and restore only positive pure-source and pure-testing
normalizations.

\subsubsection*{A source factor restores the second normal when $H$ is unbounded}

Assume now $H$ unbounded and $h_0<1/2$. SF3 gives a positive-unit
$K_A$, so the same angular allocation is available. Also
$\lambda>0$, $\lambda/H\sim2$, $\rho$ bounded and
$\lambda x=\rho\sigma$ infinitesimal. The positive-$d_0$ case is
closed exactly as before: $\lambda$ occurs only inside the retained
scalar circuit, not as an unbounded coefficient of the remaining Gram.

For infinitesimal $d_0$, $\rho$ is a unit, but $\sqrt\sigma/k$
can be unbounded. One must use the actual second normal in SF16,
not claim boundedness of its factor $qw\Gamma$ by itself. Exactly
\[
 R_\gamma=R_0-8pu^2,\qquad
 R_0=-4\{2(C+D)(r-1)+2p+r-1\}.
 \tag{SF18}
\]
The required source bound is $R_0/\sigma$ bounded. Here is a
reproducible polynomial proof that retains every opening. In the
P1 formulas temporarily leave $e$ independent, and put
$N=JP R_0$. The cancellation of $1-r$ in $C(r-1)$ shows that $N$
is a polynomial in $x,\sigma,z,b,h,e$. Moreover
$N|_{x=\sigma=e=0}=0$. Define the polynomial quotients
\[
 F_x=\frac{N-N|_{x=0}}x,\quad
 F_\sigma=\frac{N|_{x=0}-N|_{x=\sigma=0}}\sigma,\quad
 F_e=\frac{N|_{x=\sigma=0}}e.
\]
All their source arguments are bounded and the denominators $JP$
are units. Restoring $e=xh/\sigma$ gives the exact expression
\[
 \frac{R_0}{\sigma}
 =\frac{(x/\sigma)F_x+F_\sigma+(xh/\sigma^2)F_e}{JP}.
 \tag{SF19}
\]
The two ratios on the right are bounded by SF7; this proves the
claim. The displayed finite differences define all three quotients
explicitly as polynomials, so no external coefficient list is required.

Use the first conductor in SF15 unchanged. For the entire other
normal use the exact identity
\[
 qwR_\gamma\Gamma=R_0(qw\Gamma)-8p qwu(u\Gamma).
 \tag{SF20}
\]
Substitute both SF15 expressions into SF20. Every resulting
coefficient is infinitesimal. The only apparently dangerous ones are
\[
 R_0\frac{\sigma z}{k},\qquad
 R_0\frac{\sqrt\sigma qw}{k}.
\]
Write $R_0=\sigma C_0$ with bounded pure-source $C_0$. SF7 makes
$\sigma/k$ bounded, so these are respectively
$C_0\sigma z(\sigma/k)$ and
$C_0\sqrt\sigma qw(\sigma/k)$, both infinitesimal. The other
term in SF20 has the explicit infinitesimal factor $qwu$.

Thus SF17 again has strict positive residue, now using SF20 for the
complete second Gamma normal. All retained squares and every other
SP6 term are accounted for as in the bounded-$H$ proof. The inherited
R17 bounded-state and R18 unit-Alice admissions, together with the
unchanged remote joins, complete all independent projective tests on
this stated source subdomain. The argument never inverts a sphere
coordinate and remains literal at $D=0$.

\subsubsection*{The complementary source and its completion}

C14.1a has no remaining mathematical domain after this proof is
combined with the stated testing exits. The complementary C14.1b
source is the original source intersected with
\[
 H\text{ unbounded},\qquad \operatorname{res}(h)\ge\tfrac12.
 \tag{SF21}
\]
The earlier bounded-$W/Z$, sign, literal-axis and nonzero-unit-$u$
admissions still remove their stated parts. Thus its real residual
still has $u,w>0$ infinitesimal, $WZ>0$ and $W/Z$ unbounded,
as well as every original source intersection.

On SF21 the angular coefficient $K_A$ has residue $2-4h_0$, which is zero or negative. The SF14 allocation therefore supplies no strict angular or $T$ pivot there. The restored angular budget in Section~\ref{component:p4e5ba027f187}, together with the US, BC, GA and bounded-state exits, completes this source complement. Their exact join is printed in Section~\ref{component:pfd0dd7bb10cf}.




\section{\texorpdfstring{C14.1b: the bounded testing-ratio complement at every unit $h$}{C14.1b: the bounded testing-ratio complement at every unit h}}
\label{component:p5cd5fa89abcc}

\label{p5cd5fa89abcc:sec:c141b-bounded-testing-ratio}

This is a further partial-domain admission, independent of the positive
angular-pivot restriction in the preceding state-free proof. Keep the
entire original physical C14.1b source, all strict gaps, original scale
and prior-complement intersections, and all allowed coefficient openings.
Thus $x,z=1/L,\sigma=1/\eta$ are positive infinitesimals, $H=h/\sigma$
is unbounded, and all source quantities belong to $K_c$, independently
of tests and states in $K_t$. We assume here that $h$ is a positive
bounded unit. The claim is about the complete original P4/SP6 parent.

After the inherited literal-axis, state-sign and testing-sign exits,
consider $u,w,W,Z>0$, with $u,w$ infinitesimal, and define exclusively
in the testing field
\[
                 c=\frac{uW}{wZ}>0.
                 \tag{BC1}
\]
The new assertion is that every real state and test with $c$ bounded
is admitted. In particular this includes infinitesimal $c$, with no
positive real lower bound on $c$. In the actual retained C14.1b
complement one can additionally keep $W/Z$ unbounded; nothing below
needs that condition. A complex state is covered whenever each of its
real and imaginary parts belongs to this domain or to an already
admitted domain. No condition is imposed on a nonexistent complex
order or complex state ratio.

\paragraph{Exact bounded source input.}
Use SF1--SF7 and the complete radial extraction SF8. They were proved
with $z$ infinitesimal and do not require a positive lower bound on $z$.
In particular
\[
 \lambda>0,\quad \lambda\sigma\sim2h,\quad
 \rho=\lambda x/\sigma\text{ bounded},\quad
 \lambda x\text{ infinitesimal},\quad K_A\text{ bounded}.
 \tag{BC2}
\]
No positivity of $K_A$ is assumed. The exact sphere split gives
a positive original Gamma pivot when $d_0$ has positive residue.
When $d_0$ is infinitesimal, $\rho$ is a positive unit and
$\sigma/k$ is bounded, where $k=\sqrt x$. The actual second
Gamma normal has the source factorization SF18--SF20, with
$R_0/\sigma$ bounded. These facts include literal $D=0$.

\paragraph{The complete scalar circuit and a new radial normalization.}
Write $q=\sqrt z$ and use exactly the three nonnegative separated
monomials from SF10:
\[
 a_1=\sqrt{\lambda\sigma}\,uZ,\qquad
 a_2=\sqrt\rho\,uwW,\qquad
 a_3=\sqrt\lambda\,xqW,
 \qquad G=(a_1a_2a_3)^{1/3},\quad d_i=a_i-G.
 \tag{BC3}
\]
Every fractional power is the product of its pure-source and
pure-testing powers. The exact CT4 scalar identity applies without
changing the source or ordering a mixed ratio: the scalar packet
$\mathscr P$ of SF11 plus $\lambda xR^2$ is
$\sum_i a_i^2-3G^2$ plus the two printed nonnegative CT4 squares.
For $p_i=a_i^{1/3}$, its exact strict circuit decomposition is
\[
 \sum_i a_i^2-3G^2
 =\sum_i d_i^2+
   \sum_{j=1}^3\sum_{i<k}
       \{\sqrt{Gp_j}(p_i-p_k)\}^2.
 \tag{BC4}
\]
The nine supplementary squares and the two CT4 squares remain whole.
Thus the circuit supplies three unit diagonal pivots on $d_1,d_2,d_3$;
its only cost is the infinitesimal source correction $-\lambda xR^2$.

Put
\[
 k_c=z^{1/3}c^{1/3},\qquad
 \overline J=(k_c-1)\frac{wZ}{q}.
 \tag{BC5}
\]
The bar in $\overline J$ is part of the root's name, not a residue.
It is a finite separated expression. Its exact relation to the
original radial root is
\[
 R=\frac{wZ}{q}-uqW
       =-(1+k_c+k_c^2)\overline J.
 \tag{BC6}
\]
No mixed expression on the right is inverted. Because $c$ is bounded,
$k_c$ and $k_c^2$ are bounded separated coefficients with zero real
residue. Therefore the coefficient in BC6 has residue $-1$.
After BC6 the clipped SF8 radial form is a bounded strict Gram on
$(N_y,\overline J)$, including the $b$ and $Q$ opening endpoints
in residue. All exact positive radial excesses are retained.

This is precisely the rescaling
$\overline J=J_c/(q c^{2/3})$ of CT6, but definition BC5 avoids
an inverse of $c$ even in the displayed root. It is valid when
$c$ is infinitesimal.

\paragraph{The signed angular packet has a small exact conductor.}
The following whole-state identity is the rescaled CT13 conductor:
\[
\begin{split}
 A_\theta={}&
 \frac{qcw}{\sqrt{\lambda\sigma}}d_1
 +\frac{k_c^2}{\sqrt\lambda}d_2
 -\frac1{\sqrt\lambda}d_3\\
 &+(k_c+1)\sqrt x\,z^{2/3}c^{2/3}\overline J
 =:\mathcal V_A.
\end{split}\tag{BC7}
\]
All four coefficients are finite sums of bounded separated factors
with zero residue. The first uses bounded $c$, infinitesimal $q,w$
and the positive source unit $\lambda\sigma$. The second and third
use $1/\sqrt\lambda$ infinitesimal. For the last, expand it as
$\sqrt x\,z^{2/3}c^{2/3}+\sqrt x\,zc$.
No bound involving an ordered product of independent rates occurs.

One can therefore replace the possibly signed angular coefficient
by a strict reserve, using an identity rather than an inequality.
For the SF13 root $T=kq(u-k)W$, exactly
\[
\begin{split}
 K_AA_\theta^2={}&\tfrac12 A_\theta^2+\tfrac12 T^2
       +\tfrac12(A_\theta^2-T^2)
       +(K_A-1)\mathcal V_A^2,\\
 A_\theta^2-T^2={}&q^2(u-k)^2W^2(u^2+2uk).
\end{split}\tag{BC8}
\]
The last expression in the second line is an exact separated SOS.
The first two rows in the first line give strict pivots $1/2$ on
$A_\theta,T$. The signed final term has bounded infinitesimal Gram
entries on $(d_1,d_2,d_3,\overline J)$ by BC7 and bounded $K_A$.
In particular this replacement remains valid at $\operatorname{res}h=1/2$
and when $K_A$ has negative residue. It is paid exactly once.

\paragraph{Every normal of the original parent.}
Keep the original $N_x$ pivot and the exact SF8 radial extraction.
If $d_0$ has positive residue, retain the strict original Gamma
pivot and use its two full normals SF16 directly. Their coefficients
are bounded infinitesimals, including after BC6. If $d_0$ is
infinitesimal, use SF15 and the complete second-normal factorization
SF20, just as in the state-free proof. Their circuit roots are now
\[
 D_1=(d_1-d_2)/\sqrt2,\qquad
 D_2=(d_2-d_3)/\sqrt2.
 \tag{BC9}
\]
Thus no new division by a sphere coordinate is needed. SF7 and SF19
make every resulting coefficient bounded infinitesimal; $A_\theta,T$
have the strict reserves just supplied by BC8. The positive original
$4D\Gamma^2/\sigma^2$ row is retained whole in this latter case.

All other SP6 terms are precisely the five complete crosses in SF9,
the infinitesimal radial remainder CT7a, and the cost
$-\lambda xR^2$. Substitute BC6 where $R$ occurs. They remain
bounded infinitesimal perturbations of the strict frame
\[
 (N_x,N_y,\overline J,A_\theta,T,d_1,d_2,d_3),
 \tag{BC10}
\]
with $\Gamma$ adjoined in the positive-$d_0$ case. The source-unit
radial block, fixed $N_x$ pivot, BC8 angular pivots, and BC4 circuit
pivots give a positive-definite real residue. Retain exactly once
$4\beta D_wN_x^2$, $2(2D+1)(uX_e)^2$, $4bpzu^2\Gamma^2$,
all nonconstant positive radial packets and source-weight excesses,
the nine BC4 squares, the two CT4 squares, and the positive BC8
domination remainder. No full state normal has been discarded.
The separated bounded-Gram lemma now gives the complete parent
as tensor SOS on the asserted domain.

The proof concerns every independent test/state in BC1 with $c$
bounded. Literal axes and opposite-sign configurations keep their
inherited admissions. Applying the real proof or an earlier exit
separately to the real and imaginary parts gives the complex-state
join stated at the start; it does not enlarge the real domain by
assuming an order on complex amplitudes.

\paragraph{Exact improvement and retained obstruction.}
Combining this result with the previously accepted source subdomain
$\operatorname{res}h<1/2$, the real C14.1b remainder can be narrowed
to its original retained source/test intersections with
\[
 \operatorname{res}h\ge1/2,\qquad
 u,w>0\text{ infinitesimal},\quad WZ>0,\quad W/Z\text{ unbounded},
 \quad \frac{uW}{wZ}\text{ unbounded}.
 \tag{BC11}
\]
This complement still contains the difficult source-contact scale:
on the actual source contact $uW/(wZ)=L=1/z$, which is unbounded.
For unbounded pure-testing $c$, BC7's coefficients involving $c$
cannot be declared bounded separated coefficients, and the present
argument does not cover that case. No full C14.1b closure or
nonemptiness assertion is made.

The narrow checker verifies the exact radial relation, the complete
angular conductor, the circuit decomposition and conversion, and
the once-paid angular reserve. The source bounds, every remaining
normal and the independent-field Gram argument are the analytical
proof above, with precisely the SF and CT dependencies identified.


\section{\texorpdfstring{C14.1b: the remaining source at bounded Alice-to-reciprocal-Bob ratio}{C14.1b: the remaining source at bounded Alice-to-reciprocal-Bob ratio}}
\label{component:p50b4df0bb619}

\label{p50b4df0bb619:sec:c141b-gamma-angular}

\paragraph{Exact additional scope.}
Keep the full original C14.1b source, physicality, strict gaps,
nonnegative coefficients, total-scale and prior-complement
intersections. Use the independent fields and the complete P4/SP6
parent of the inherited Section~\ref{component:p472f434581f3}. In this
section $h$ is a positive bounded unit and $a_0=k/\sigma$ is
infinitesimal. The unit-$a_0$ source was already admitted by
US1--US15; no conclusion about that proof is changed here.

\begin{theorem}
On the stated source, the entire original parent is tensor SOS
for every projective testing configuration whose remaining
double-small-angle chart has $u,w>0$ infinitesimal and $u/w$
bounded, and for every real or complex state. Include the exact
sign gauges, literal endpoints and other testing exits already
admitted in SF1--SF21. There is no restriction on a state ratio
or on $c=uW/(wZ)$ in this theorem.
\end{theorem}

This is an additional testing-domain admission. It does not close
the remaining chart $u/w$ unbounded. Unlike BC1--BC11, its testing
predicate is independent of the state, so the same predicate holds
for the real and imaginary state components.

\paragraph{Available source pivots.}
The source bounds SF1--SF5 and US1--US3 give
\[
 L_0=\lambda\sigma\sim2h,\qquad
 \lambda^{-1}=\sigma/L_0\text{ infinitesimal},\qquad
 \rho=L_0a_0^2\text{ infinitesimal},\qquad
 d=4D/\sigma^2\sim2.
 \tag{GA1}
\]
Indeed US2 first bounds $a_0$, and its exact residue identity makes
$b_0$ infinitesimal when $a_0$ is; the sphere then gives $d_0\sim1$.
Thus the original Gamma diagonal supplies a strict positive-unit
pivot. Also $\lambda x=\rho\sigma$ is infinitesimal. The source
coefficients $K_A,\chi,\nu,R_\gamma,R_x,A_y$ are bounded. No
lower cap on $\beta$, $b$ or $Q$ is imposed. A literal $D=0$ face
is incompatible with GA1, rather than being removed from a domain
where it could occur.

\paragraph{A paid angular domination from the original Gamma square.}
The unchanged sign and state-axis exits leave $u,w,W,Z>0$; their
ratios may otherwise be arbitrary. Set $t_s=\sqrt\sigma$ (a source
letter, unrelated to the earlier notation $t=u^2$), and define
\[
 E_0=quW-t_s Z,\quad V=uE_0,\quad
 \alpha_t=q\frac uw,\quad
 U=\frac{B_1-B_3}{\sqrt\lambda}
   =\frac{d_1-d_3}{\sqrt\lambda}.
 \tag{GA2}
\]
Here $B_i,G,d_i=B_i-G$ are exactly the four monomials and circuit
differences in US7--US10. Every root in GA2 is separated. The only
new testing quotient is $u/w$, which is bounded by the theorem's
predicate. Hence $\alpha_t$ is a bounded infinitesimal coefficient.
Exactly
\[
 \Gamma+t_sR=-\frac{w+t_sq}{q}E_0,
 \qquad A_\theta=U+V.
 \tag{GA3}
\]
The first identity is used as a polynomial identity; its mixed
factor $w+t_sq$ is never inverted. In fact it yields the explicit
separated-SOS difference
\[
\begin{split}
 \Delta&:=\alpha_t^2(\Gamma+t_sR)^2-V^2\\
 &=u^2E_0^2\left\{\frac{2t_sq}{w}
                            +\left(\frac{t_sq}{w}\right)^2\right\}\\
 &=\left(\sqrt2\sqrt{t_s}\sqrt q\,
                  \frac{u}{\sqrt w}E_0\right)^2
   +\left(\frac{t_squ}{w}E_0\right)^2.
\end{split}\tag{GA4}
\]
All square roots and inverses in the last line belong to their
individual fields. The roots themselves need not be bounded.

Choose a positive real $C_*$ so that the bounded pure-source
coefficient $K_A+C_*$ has positive real residue. The following
once-paid identity allocates the complete signed angular diagonal:
\[
\boxed{\begin{aligned}
 K_AA_\theta^2={}&(K_A+C_*)A_\theta^2
  -\frac{2C_*}{\lambda}(d_1-d_3)^2
  -2C_*\alpha_t^2(\Gamma+t_sR)^2\\
 &+C_*(U-V)^2+2C_*\Delta.
\end{aligned}}\tag{GA5}
\]
It follows simply from $A_\theta=U+V$ and GA4. The last two terms
are retained SOS. The two negative terms have bounded
infinitesimal coefficients on the circuit and $(\Gamma,R)$ frames.
This constructs a strict angular pivot without requiring a positive
residue of $K_A$, and without the collapsed $\lambda k$ pivots of
US11.

\paragraph{Complete scalar, every state normal, and all excesses.}
Use the whole four-monomial allocation once:
\[
       \mathscr S=\sum_{i=1}^4d_i^2+\mathcal E_4
                                           -\lambda xR^2.
       \tag{GA6}
\]
Here $\mathcal E_4$ is the explicit separated SOS in US10. Do not
also allocate US6 or US11. The single radial cost in GA6 is an
infinitesimal pure-source coefficient.

Use SF8/CT7--CT7a with the full source-opening-safe radial squares
and every positive excess restored. The resulting strict bounded
radial frame is $(N_y,R)$, uniformly through the $b,Q$ openings.
Together with the original $N_x$ pivot, $d\Gamma^2$ from GA1,
GA5 and GA6, the principal frame is
\[
       (N_x,N_y,R,\Gamma,A_\theta,d_1,d_2,d_3,d_4).
       \tag{GA7}
\]
Its residue is strictly positive definite: the $N_x$ residue is
$2$, the Gamma residue is $2$, the angular coefficient is
$K_A+C_*$ with positive residue, and the four circuit coefficients
are $1$. All costs in GA5--GA6 and all CT7a remainders are bounded
infinitesimal matrix entries on this frame.

The two complete Gamma normals remain exactly
\[
      -4\nu N_x(u\Gamma)
       +R_\gamma(N_y+bR)(qw\Gamma).
       \tag{GA8}
\]
Their coefficients are infinitesimal directly because $u,q,w$ are
infinitesimal and the remaining factors are bounded. No division
by $\sqrt\rho$ is used. The five remaining non-Gamma crosses are
precisely SF9, with bounded infinitesimal coefficients on GA7.
The whole $N_yR$ coefficient has already been included in the
exact radial extraction. Thus every original SP6 term, including
both first state normals, is accounted for.

Keep separately and exactly once $4\beta D_wN_x^2$,
$2(2D+1)(uX_e)^2$, the nonconstant positive $N_x$ terms,
$4bpzu^2\Gamma^2$, all radial and source-weight excesses,
$\mathcal E_4$, and the last two SOS terms of GA5. The separated
bounded-Gram lemma applied to GA7 proves the complete parent.

The exact inherited literal, sign and remote exits supply the
complementary charts specified in the theorem. For complex states,
apply the same real theorem separately to each real state part;
the hypothesis $u/w$ bounded depends only on the test and is
unchanged. Restore only positive pure-source and pure-testing
normalizations. No source/testing rate is identified or compared.

\paragraph{The complementary testing chart.}
Combining this theorem with SF, BC and US leaves C14.1b only on
the original retained complement intersected with
\[
\begin{gathered}
 \operatorname{res}h\ge\tfrac12,\quad k/\sigma\text{ infinitesimal},
 \quad b_0\text{ infinitesimal},\quad d_0\sim1,\\
 u,w>0\text{ infinitesimal},\quad WZ>0,\quad W/Z\text{ unbounded},\\
 c=uW/(wZ)\text{ unbounded},\qquad u/w\text{ unbounded}.
\end{gathered}\tag{GA9}
\]
On GA9 the coefficient $\alpha_t=q(u/w)$ in GA2 has an unbounded testing factor. It cannot be declared small by comparing the independent source $q$ to the testing reciprocal $w/u$. At contact $u_c/w_c=a_0\sqrt\sigma/q$. The restored angular budget in Section~\ref{component:p4e5ba027f187} covers this complement without requiring the GA2 coefficient to be bounded. Its exact join with the present bounded-ratio theorem and the other C14.1b exits is printed in Section~\ref{component:pfd0dd7bb10cf}.


\section{\texorpdfstring{C14.1b: the entire unit beta-sphere source, including the contact ratio}{C14.1b: the entire unit beta-sphere source, including the contact ratio}}
\label{component:p0931b3af19ec}

\label{p0931b3af19ec:sec:c141b-unit-beta-sphere}

\paragraph{The precise whole-source subdomain.}
Keep the entire original physical C14.1b source, its strict gaps,
nonnegative event weights, original total-scale and prior-complement
intersections. In particular $x,z=1/L,\sigma=1/\eta$ are positive
infinitesimals. Write $k=\sqrt x$, $q=\sqrt z$ and retain all
P1--P8/SP6 notation. Source and testing fields are independent.

\begin{theorem}
Suppose $h$ and $k/\sigma$ are positive bounded units. Then the
complete original parent is finite tensor SOS at every independent
projective testing pair and every real or complex state. All permitted
literal coefficient boundaries and all positive opening rates are
retained. There is no bound on $uW/(wZ)$ or any other testing/state
ratio in this theorem.
\end{theorem}

This includes the source-contact scale $uW/(wZ)=1/z$, which was not
covered by the bounded-testing-ratio theorem. The theorem is a whole
source-subdomain admission, not closure of all C14.1b.

\paragraph{Source entry and its sphere interpretation.}
The bounds SF1--SF5 remain valid. Since $h$ is a unit and
$\sigma$ infinitesimal, $H$ is unbounded and
\[
 \begin{gathered}
 \lambda>0,\quad L_0:=\lambda\sigma\sim2h,\quad
 a_0:=k/\sigma\text{ a positive unit},\\
 \rho=\lambda x/\sigma=L_0a_0^2\text{ a positive unit},
 \quad \lambda k=L_0a_0\text{ a positive unit}.
 \end{gathered}
 \tag{US1}
\]
Also $\lambda x=\rho\sigma$ is infinitesimal. The subscript in
$a_0$ is part of this section's source letter, not a residue.
The bounded source quantities $K_A,\chi,\nu,R_x,A_y,R_\gamma$
have the exact meanings in SP6; $K_A$ may have either sign.

On the entire unit-$h$ source, the earlier exact source bound
$e=xh/\sigma=O(x+\sigma)$ first gives $x/\sigma=e/h$
infinitesimal. SF5 gives the exact identity
\[
 h a_0^2=\frac{x}{\sigma}\frac{H_0}{B}
                    +\frac{\mathcal H}{B}b_0^2.
 \tag{US2}
\]
Since $h$ is a positive unit and the right-hand factors in US2 are
bounded, this identity bounds $a_0^2$ without assuming it bounded
in advance. It also shows
\[
        \operatorname{res}(h a_0^2)=\operatorname{res}(b_0^2).
 \tag{US3}
\]
Thus this theorem covers precisely the unit-$h$ source with the
beta-sphere coordinate $b_0$ a positive unit. Literal $D=0$ has
$b_0=1$ and is included. There is no division by $D,\beta$ or
either sphere coordinate. A source face incompatible with these
exact hypotheses is not assumed to have points.

\paragraph{Two exact allocations of the same complete scalar packet.}
The inherited remote and literal/sign exits leave $u,w,W,Z>0$ with
$u,w$ infinitesimal. No size restriction on $W/Z$ is needed below.
The entire scalar contact packet in SP6 is
\[
 K_AA_\theta^2+\chi uA_\theta R+\mathscr S,
 \qquad
 \mathscr S=\lambda\sigma P_\beta^2
       +\rho(u^2-x)(w^2-\sigma z)W^2.
 \tag{US4}
\]
Define legal separated roots
\[
 J_A=q(u-k)W,\quad
 F_A=u(Z-a_0wW),\quad
 E_u=\sqrt u\,J_A,\quad E_k=\sqrt k\,J_A.
 \tag{US5}
\]
The following identity is exact for the complete scalar, with no
restriction to a state graph:
\[
 \boxed{\mathscr S=L_0F_A^2
       +2\lambda k E_u^2+\lambda k E_k^2
       +2\lambda k\sqrt u\,E_uR.}
 \tag{US6}
\]
One direct derivation expands $P_\beta=uZ-\sigma a_0^2wW$ and
uses $wZ/q=R+uqW$. The coefficient of $q^2W^2$ that remains is
\[
 a_0(2u^3-3ku^2+k^3)=a_0(2u+k)(u-k)^2,
\]
which gives US6. Its diagonal coefficients in US1 are positive
units, and its sole cross coefficient contains the testing
infinitesimal $\sqrt u$.

For a complementary allocation set
\[
 \begin{gathered}
 B_1=\sqrt{\lambda\sigma}\,uZ,\quad
 B_2=\sqrt\rho\,uwW,\quad
 B_3=\sqrt\lambda\,xqW,\quad
 B_4=\sqrt{\lambda x}\,\frac{wZ}{q},\\
 G=(B_1B_2B_3B_4)^{1/4}
       =\sqrt{\lambda x}\sqrt{uwWZ},\qquad d_i=B_i-G.
 \end{gathered}
 \tag{US7}
\]
All radicals separate; in particular $\sqrt{uwWZ}$ is entirely
testing. The pure-source inverse $1/q$ in a root is allowed. Its
root need not be bounded, and it is not used in a coefficient.
Exactly
\[
 \boxed{\mathscr S+\lambda xR^2
               =\sum_{i=1}^4 B_i^2-4G^2.}
 \tag{US8}
\]
The four-term circuit has strict difference pivots:
\[
 \sum_iB_i^2-4G^2=\sum_i d_i^2+\mathcal E_4,
 \tag{US9}
\]
where the following is explicitly tensor SOS:
\[
 \mathcal E_4=2G\left\{
 (\sqrt{B_1}-\sqrt{B_2})^2+
 (\sqrt{B_3}-\sqrt{B_4})^2+
 2\bigl((B_1B_2)^{1/4}-(B_3B_4)^{1/4}\bigr)^2
 \right\}.
 \tag{US10}
\]
The prefactor $G$ is a nonnegative separated monomial, so it too
has a separated square root. The circuit does not assert an order
on a mixed coefficient.

Use each allocation with weight one-half. The exact paid budget is
\[
\begin{split}
 \mathscr S={}&\tfrac12L_0F_A^2+
       \lambda kE_u^2+\tfrac12\lambda kE_k^2
       +\lambda k\sqrt u\,E_uR\\
     &+\tfrac12\sum_i d_i^2+\tfrac12\mathcal E_4
                             -\tfrac12\lambda xR^2.
\end{split}\tag{US11}
\]
There is no repeated use of a scalar square. The only radial cost
is the infinitesimal pure-source coefficient $-\lambda x/2$.

\paragraph{The complete angular and Gamma conductors.}
The angular root has the exact identity
\[
             A_\theta=\sqrt u\,E_u+\sqrt k\,E_k.
             \tag{US12}
\]
Its coefficients are bounded infinitesimals. Therefore the entire
possibly signed $K_AA_\theta^2$, and the complete
$\chi uA_\theta R$ term, give bounded infinitesimal corrections
to the strict frame supplied by US11 and the radial block.
This works also when $\operatorname{res}h\ge1/2$ and $K_A$ has
zero or negative residue.

Both first Gamma normals have small exact conductors on the same
frame:
\[
\boxed{\begin{aligned}
 u\Gamma={}&\frac{k}{\sqrt\rho}(d_1-d_2)
       -\frac{u-k}{\sqrt\rho}(d_2-d_3)
       -\sqrt{\sigma k}\,E_k,\\
 qw\Gamma={}&\frac{qw+z\sqrt\sigma}{\sqrt\rho}(d_4-d_2)
                         -\sqrt\sigma\,qwR.
\end{aligned}}\tag{US13}
\]
The first is SF15 with $T=kJ_A=\sqrt k E_k$; the second follows
directly from the two monomials $B_4,B_2$. Every displayed coefficient
is a finite sum of bounded separated factors with zero residue.
No state quotient, testing inverse, source/test rate comparison,
or sphere-coordinate inverse occurs. In particular these formulas
do not use a lower bound on the original Gamma diagonal.

For comparison, US7 also yields the exact global identity
$T=(B_4-B_3)/\sqrt\lambda-kR$. This is another small conductor
on the present source and explains why the extra fourth monomial
removes the coefficient failure of the earlier three-term circuit.
The proof uses US13 directly.

\paragraph{Full-parent restoration and strictness.}
Use the exact opening-safe SF8/CT7--CT7a radial decomposition. Its
clipped pure-source Gram on $(N_y,R)$ is strict through the $b,Q$
openings, with all excess squares retained. Keep the original
$N_x$ pivot, whose residue is $2$. The principal frame is
\[
           (N_x,N_y,R,E_u,E_k,d_1,d_2,d_3,d_4).
           \tag{US14}
\]
US11 gives positive-unit pivots $\lambda k$ and $\lambda k/2$
on $E_u,E_k$, and pivots $1/2$ on the four circuit differences.
The $F_A$ square and $\mathcal E_4/2$ are retained separately.

The two entire Gamma normals of SP6 are
\[
       -4\nu N_x(u\Gamma)
            +R_\gamma(N_y+bR)(qw\Gamma).
\]
Their remaining coefficients are bounded, so US13 allocates them
with infinitesimal coefficients on US14. Retain the entire
nonnegative $(4D/\sigma^2+4bpzu^2)\Gamma^2$, including its literal
zero cases. The five non-Gamma crosses in SF9 are all bounded
infinitesimal after US12; the whole $N_yR$ coefficient was already
included in the radial extraction. The radial remainders CT7a are
infinitesimal on $z$ infinitesimal, as in SF8. US11 adds only its
printed infinitesimal radial cost and infinitesimal $E_uR$ cross.

Thus every full-parent term is either an explicitly retained SOS
or a bounded Gram entry on US14, whose real residue is strictly
positive definite. The separated bounded-Gram lemma gives tensor
SOS. Restore exactly once the original $4\beta D_wN_x^2$,
$2(2D+1)(uX_e)^2$, the positive nonconstant $N_x$ terms, every
radial/source-weight excess, the original Gamma diagonal, and the
two retained nonnegative terms of US11. No hard state normal or
signed scalar term is omitted.

This proves the complete double-small-angle chart for every state
left by the exact sign/axis exits, including $D=0$. The unchanged
remote, nonzero-unit-Alice and other projective exits listed in the
SF dependency index cover its complement. Hence the whole original
source subdomain stated in the theorem is covered at every test.
The real theorem applies separately to arbitrary real and imaginary
state parts, with no state-ratio predicate, proving the full complex
claim as well.

\paragraph{The exact remaining source and testing obstruction.}
Combine this theorem with the source subdomain
$\operatorname{res}h<1/2$ and the bounded-$c$ theorem BC1--BC11.
The remaining real C14.1b complement is contained in the original
retained source/test domain with
\[
 \operatorname{res}h\ge1/2,\quad k/\sigma\text{ infinitesimal},
 \quad b_0\text{ infinitesimal},\quad d_0\sim1,
 \quad \frac{uW}{wZ}\text{ unbounded},
 \tag{US15}
\]
and the previous $u,w>0$ infinitesimal, $WZ>0$, unbounded-$W/Z$
conditions. Here the new angular pivots $\lambda k=L_0k/\sigma$
lose their real residue, and the US13 coefficients involving
$1/\sqrt\rho$ are no longer bounded. The original Gamma diagonal
now has a positive-unit pivot, which is a concrete possible reserve
for a further completion, but this proof does not allocate it to
the lost angular roots. No nonemptiness or full C14.1b closure is
asserted.

The focused checker verifies US6, US8--US13, the exact half-and-half
budget and the fourth-monomial $T$ conductor. Source bounds and
all-test restoration are the analytical proof above, with the exact
inherited SF/CT inputs stated. Review remains same-system review.


\section{\texorpdfstring{A fresh angular budget and completion of the retained C14.1b residual}{A fresh angular budget and completion of the retained C14.1b residual}}
\label{component:p4e5ba027f187}

\label{p4e5ba027f187:sec:c141b-residual-rederived}

\paragraph{Scope and provenance.}
This is a new proof, not a restoration of the missing late-R21 source.
Keep the original physical stationary C1 source, its strict gaps,
nonnegative event coefficients, original infinitesimal total scale and
every earlier-admission complement in the surviving C14.1b statement.
Normalize the original positive coefficient to $d=1$ as in P1--P4.
Source parameters lie in a real closed extension $K_c$ and testing
angles and states in an independent extension $K_t$ of $\mathbb R$.
All conclusions are finite sums of squares in
$\mathcal T=K_c\otimes_{\mathbb R}K_t$. No order or general inverse
of a mixed tensor is used.

The exact target is the complete reciprocal-Bob polynomial P4/SP6
in the surviving Section~\ref{component:p472f434581f3}.
Use its notation, and set $k=\sqrt x$, $q=\sqrt z$, $z=1/L$,
$\sigma=1/\eta$, $b=Lr$, $Q=1-b$, $\theta=1-h$.
The source part of the retained residual GA9 is
\[
 x,z,\sigma>0\text{ infinitesimal},\quad
 \operatorname{res}h\ge\tfrac12,\quad k/\sigma\text{ infinitesimal},
 \quad b_0\text{ infinitesimal},\quad d_0\sim1.
 \tag{AR1}
\]
Its real testing part has $u,w>0$ infinitesimal, $WZ>0$,
$W/Z$ and $u/w$ unbounded, and $uW/(wZ)$ unbounded. These
intersections are retained. The proof below in fact needs only
$u,w,W,Z>0$ with $u,w$ infinitesimal after the inherited sign
gauges; it assumes no upper or lower bound on any state ratio.

\begin{theorem}[Fresh completion of the retained residual]
On AR1 and the exact retained testing/source intersections just
stated, the complete original parent is a finite tensor sum of
squares. Together with the surviving SF, BC, US and GA admissions
and their recorded literal, sign and projective exits, this closes
the entire original scoped C14.1b domain. It does not by itself
close C14.4 or the global C1 entry and union.
\end{theorem}

\subsubsection*{Source coefficients and the original scalar packet}

The source estimates SF1--SF5 and US2--US3 were proved for
$z$ infinitesimal. On AR1 they imply
\[
 \begin{gathered}
 L_0:=\lambda\sigma\sim2h\text{ is a positive bounded unit},
 \qquad \lambda^{-1}\in\operatorname{Inf}_c,\\
 \rho=\lambda x/\sigma=L_0(k/\sigma)^2\in\operatorname{Inf}_c,
 \qquad \lambda x\in\operatorname{Inf}_c,
 \qquad d_G:=4D/\sigma^2\sim2. 
 \end{gathered} \tag{AR2}
\]
Here $h$ is a positive unit even when $\theta=1-h$ is
infinitesimal. In particular $\lambda>0$ and $d_G$ is a strict
original Gamma pivot. The exact sphere identity P2 proves these
bounds without an inverse of $b_0$ or $d_0$.

For clarity, the other source bounds used below are
\[
 \begin{gathered}
 P,J,M,B,\mathcal H\sim1,\quad p,l,\kappa\sim\tfrac12,
 \quad C=O(z),\quad D=O(\sigma^2),\\
 e=xh/\sigma\in\operatorname{Inf}_c,\quad
 m=\kappa xze\in\operatorname{Inf}_c,\quad
 \beta\theta\in\mathcal O_c,\quad
 \nu,K_A,\chi,R_x,A_y\in\mathcal O_c,
 \quad B_x=8\kappa e\in\operatorname{Inf}_c.
 \end{gathered} \tag{AR3}
\]
These are SF1--SF5, the P1 formulas and the SP9--SP10 polynomial
cancellations. In particular $K_A=4\kappa h/\sigma-\lambda$
is bounded although its residue $2-4\operatorname{res}h$ may be
zero or negative. Also
\[
 R_\gamma=-4\{2(C+D)(r-1)+2p+r-1+2pu^2\}
 \tag{AR4}
\]
is a bounded separated coefficient. None of AR2--AR4 contains
$u/w$, a state ratio, or a comparison of a source and a testing
rate. The only possible source poles in $a=Az=p/b$ and
$\ell=\alpha z=l/Q$ remain in positive packets below. The
possibly unbounded $\beta$ itself also remains in its original
nonnegative packet; no upper bound on $\beta$ is asserted.

Keep the exact roots
\[
 N_x=X+uhZ,\quad N_y=(Y-bwZ)/q,\quad
 \Gamma=\sigma Z-uwW,\quad R=wZ/q-uqW,
 \quad A_\theta=q(u^2-x)W.
 \tag{AR5}
\]
The whole angular packet of SP6 is
\[
 K_AA_\theta^2+\chi uA_\theta R+\mathscr S,
 \qquad
 \mathscr S=\lambda\sigma\{uZ-(x/\sigma)wW\}^2
       +\rho(u^2-x)(w^2-\sigma z)W^2.
 \tag{AR6}
\]
This is the exact P8 transformation, not a restriction of the state.
Define the four positive separated monomials
\[
 B_1=\sqrt{\lambda\sigma}\,uZ,\quad
 B_2=\sqrt\rho\,uwW,\quad
 B_3=\sqrt\lambda\,xqW,\quad
 B_4=\sqrt{\lambda x}\,wZ/q,
 \qquad G_4=(B_1B_2B_3B_4)^{1/4}.
 \tag{AR7}
\]
Every fractional power means the product of the two pure-field
powers. With $d_i=B_i-G_4$, the complete circuit is
\[
 \mathscr C:=\mathscr S+\lambda xR^2
       =\sum_{i=1}^4B_i^2-4G_4^2
       =\sum_{i=1}^4d_i^2+\mathcal E_4,
 \tag{AR8}
\]
where the explicit surviving US10 remainder is
\[
 \mathcal E_4=2G_4\left\{
 (\sqrt{B_1}-\sqrt{B_2})^2+
 (\sqrt{B_3}-\sqrt{B_4})^2+
 2\big((B_1B_2)^{1/4}-(B_3B_4)^{1/4}\big)^2\right\}.
 \tag{AR9}
\]
Its nonnegative monomial prefactors have separated square roots.

\subsubsection*{An eight-square budget for the angular row}

The following elementary identity is the new step. Let $U,V,Z_*$
be positive separated monomials, let $a_*>0$ be pure source, and set
\[
 \begin{gathered}
 \Psi=1+U^2V^2+U^2Z_*^2+V^2Z_*^2-4UVZ_*,\\
 T=U^2-1,\qquad Q_*=VZ_*-U,\qquad
 R_*=Z_*-UV,\qquad G_*=U(V-Z_*).
 \end{gathered} \tag{AR10}
\]
There are two exact polynomial identities
\[
 \Psi=G_*^2+(1-VZ_*)^2+2VZ_*(U-1)^2,
 \qquad
 \Psi=T^2+Q_*^2+UTQ_*-UR_*^2-(U+1)G_*R_*.
 \tag{AR11}
\]
Consequently
\[
\boxed{\begin{aligned}
 &[2+a_*(U+1)]\Psi+U^2Q_*^2
       +\left[2U+\frac{U+1}{a_*}\right]R_*^2-T^2\\
 &\quad=(T+UQ_*)^2+2Q_*^2
       +(U+1)\left(\sqrt{a_*}G_*
                          -R_*/\sqrt{a_*}\right)^2\\
 &\qquad\quad+a_*(U+1)(1-VZ_*)^2
                     +2a_*(U+1)VZ_*(U-1)^2.
\end{aligned}} \tag{AR12}
\]
The right side consists of eight legal squares: expand each
factor $U+1$ into its two summands before taking roots. Only
radicals of positive separated monomials occur. In particular
neither a square root nor an inverse of $1+U$ is used. The
identity imposes no bound on $U,V,Z_*$ or on $a_*^{-1}$.

Apply this with
\[
 U=\frac uk,\qquad V=\frac w{\sqrt\sigma q},\qquad
 Z_* =\frac{\sqrt\sigma}{kq}\frac ZW,
 \qquad a_*=\lambda k,\qquad M_*=k^2qW.
 \tag{AR13}
\]
Each quotient is taken within a pure field before tensoring.
The only state inverse is the nonzero pure-testing $W^{-1}$.
The other inverses are positive pure-source factors. These
expressions are allowed to be unbounded inside a square root row.
Directly,
\[
 \mathscr C=\lambda M_*^2\Psi,\qquad
 A_\theta=M_*T,\qquad R=kqWQ_*,\qquad
 \Gamma=k\sqrt\sigma qWR_*.
 \tag{AR14}
\]
Multiplying AR12 by the square $M_*^2$ therefore proves
\[
\boxed{\quad
 \Delta_A:=\varepsilon_C\mathscr C+u^2R^2
                 +\varepsilon_G\Gamma^2-A_\theta^2
                         \ \in\ \Sigma\mathcal T^2,
 \quad} \tag{AR15}
\]
with the exact coefficients
\[
 \varepsilon_C=\frac2\lambda+u+k,\qquad
 \varepsilon_G=2\frac k\sigma u+
                         \frac{u+k}{\lambda\sigma}.
 \tag{AR16}
\]
Both are finite sums of positive bounded separated factors with
zero real residue. This follows from AR2, bounded $k/\sigma$,
and the independent facts $u\in\operatorname{Inf}_t$ and
$k,\lambda^{-1}\in\operatorname{Inf}_c$. It uses no bound on
$q u/w$, $uW/(wZ)$, or any normalized variable in AR13.
The identity also gives a direct expression for each of its eight
roots: multiply the eight expanded roots on the right of AR12
by $M_*$. Thus AR15 is a finite certificate, not a pointwise
inequality imported from an ordered compositum.

\subsubsection*{Paying the signed angular term exactly once}

Choose a fixed positive real $c_*$ such that
$\operatorname{res}(K_A+c_*)>0$, possible because $K_A$ is
bounded. The exact allocation is
\[
\begin{split}
 K_AA_\theta^2+\mathscr S
  ={}&(K_A+c_*)A_\theta^2
       +(1-c_*\varepsilon_C)\mathscr C
       -c_*\varepsilon_G\Gamma^2\\
    &-(\lambda x+c_*u^2)R^2+c_*\Delta_A.
\end{split} \tag{AR17}
\]
AR17 uses AR8 and AR15 once each. The coefficient of the whole
circuit has residue one; its four $d_i$ rows remain strict.
The three AR9 rows with this same coefficient are retained
separately. The elementary positive-residue lemma makes each
of these bounded coefficients a tensor SOS; multiplying its
finite square decomposition by the corresponding row preserves
the finite SOS property. There is no mixed square root of
$1-c_*\varepsilon_C$.

\subsubsection*{The two complete state normals and all original hard terms}

Use the exact CT7--CT7a radial extraction. Its constant base is
\[
 [4a+2z(2C+1)](N_y+bR)^2+4\ell(N_y-QR)^2.
 \tag{AR18}
\]
Since $p,l\sim1/2$ and $0<b,Q<1$, both coefficients exceed
one. Keep their positive pure-source excesses and retain
$(N_y+bR)^2+(N_y-QR)^2$ in the principal frame. Its bounded
two-by-two Gram has determinant $(b+Q)^2=1$ and positive
definite residue, including $b=0$ or $Q=0$ in residue. This is
a cap of whole positive squares, not a cap of unrelated matrix
entries. Retain exactly the nonconstant positive $a,\ell,C$
packets specified in CT7, before using its remaining coefficients.

The three exact remainders $E_{yy},E_{RR},E_{yR}$ in CT7a are
bounded infinitesimal: every source factor in their braces is
bounded by AR3, $C=O(z)$, $D=O(\sigma^2)$, and they have the
printed factor $z$ together with $u^2,w^2,\sigma$ or $D$.
No estimate involving a state coordinate or $u/w$ is present.

The two entire Gamma normal terms of SP6 are
\[
 -4\nu uN_x\Gamma
          +qwR_\gamma(N_y+bR)\Gamma.
 \tag{AR19}
\]
They are bounded infinitesimal entries on the original
$(N_x,N_y,R,\Gamma)$ frame, directly by AR3--AR4 and
$u,q,w$ infinitesimal. In this source sector the Gamma diagonal
is $d_G\Gamma^2$ with $d_G\sim2$, so no sphere-coordinate
inverse, Gamma conductor or weak Gamma pivot is required.
The five other complete cross terms are exactly
\[
 \begin{split}
 &-4uqw(2m+2D+1)N_xN_y+uqwR_xN_xR\\
 &\qquad+qwB_xN_xA_\theta+uA_yN_yA_\theta
                                      +\chi uA_\theta R.
 \end{split} \tag{AR20}
\]
Every displayed coefficient is infinitesimal and bounded by
AR3. The entire $N_yR$ term has already been included in
AR18 and CT7a. Thus both unrestricted state normals and every
remaining mixed term of SP6 have been allocated.

Use exactly $2N_x^2$ in the principal frame and retain exactly
$[2w^2+4mw^2]N_x^2$ as its positive excess. After AR17 the
principal frame is
\[
          (N_x,N_y,R,\Gamma,A_\theta,d_1,d_2,d_3,d_4).
 \tag{AR21}
\]
Its real residue is strictly positive definite: the $N_x$ pivot
is two, the radial block is AR18's clipped block, the Gamma
pivot is two, the angular pivot is $\operatorname{res}(K_A+c_*)>0$,
and the four circuit pivots are one. All other entries are
bounded infinitesimals, including the AR17 Gamma/radial costs.
The separated bounded-Gram lemma therefore proves finite SOS.

Finally restore, each exactly once, $4\beta(1+w^2)N_x^2$,
$2(2D+1)(uX_e)^2$, $[2w^2+4mw^2]N_x^2$,
$4bpzu^2\Gamma^2$, all CT7 positive source and testing
excesses, the AR9 remainder with its positive-residue coefficient,
and $c_*\Delta_A$. This restores precisely P4/SP6. The roots
need not be bounded; only the Gram entries have been bounded.

\subsubsection*{Exact join and literal boundaries}

The surviving SF theorem covers $\operatorname{res}h<1/2$ in
C14.1b. US covers its unit-$h$, unit-$k/\sigma$ source for
all states and tests; US2 bounds $k/\sigma$ throughout the
unit-$h$ source. BC and GA cover their stated bounded-ratio
testing domains. Their residual is exactly contained in AR1
with the testing predicates printed above, now covered by
AR17--AR21. The inherited bounded-state, nonzero-unit-Alice,
remote, literal-axis and sign exits remain unchanged. The bounded-state and unit-Alice bodies are Sections~\ref{component:pff0a29ebf284} and~\ref{component:pcc84202479d2}; the common remote and sign bodies are in Section~\ref{component:p89fd3edb3edd}. The SF, US, BC and GA bodies and their source predicates are joined explicitly in Section~\ref{component:pfd0dd7bb10cf}.
This is a finite union on the original scoped C14.1b source;
it does not erase any earlier-admission complement.

The residual itself has $D>0$ by AR2. Literal $D=0$ remains
covered by the US sphere endpoint, which divides by neither
$D$ nor a sphere coordinate. Literal $\beta=0$ is incompatible
with $H=h/\sigma$ unbounded here: P2 gives $H=H_0/B$ on
that face, a bounded quantity. Hence no nonempty C14.1b face
has been dropped. Arbitrarily small positive $b,Q,\theta,D$
or other original openings are kept wherever compatible with
the printed source sector. For complex states apply the real
routing independently to their real and imaginary parts;
each part satisfies an earlier exit or the new residual proof.
Restore only the original positive pure-source and pure-testing
normalizations.




\section{\texorpdfstring{C1.4: the finite source and testing cover}{C1.4: the finite source and testing cover}}
\label{component:pfd0dd7bb10cf}
% Edited scoped mathematical packet, 22 October 2026.

\label{pfd0dd7bb10cf:sec:fresh-c144-c145-join}

\paragraph{Common input.}
Use the original native C1 source, $d=1$, $\lambda=n=0$,
$m,C,A,\alpha>0$, $D,\beta\ge0$, $D+\beta>0$, and all
strict native gaps.  On the C1.4 row, $x=k^2>0$ is infinitesimal
and $\zeta>0$ is noninfinitesimal.  Put $L=x\zeta$.
Retain physicality, the actual finite noncommuting stationary contact,
and $\delta_{\rm tot}\in\operatorname{Inf}_{c,>0}$.
All alternatives below retain these predicates and any shared upstream
intersection.  Each named packet is printed in the certificate appendix.

\paragraph{The C14.4 source partition.}
When $L$ and $\sigma=1/\eta$ are infinitesimal, set
\[
 S=\sigma+L,\quad\kappa=x/S,\quad\nu=\sigma/S,
 \quad\ell=L/S,\quad\nu+\ell=1,\quad\zeta=\ell/\kappa.
                                                               \tag{JC1}
\]
The actual FM source, NC5 and ES/SE bounds give $E\ge1/18$,
$M>0$, and the exact identity
\[
 \kappa Q_s+\ell M+\nu E=1,\qquad
 Q_s=1-\ell Sr\sim1,\qquad
 0<\kappa Q_s<1-\nu/18.                                \tag{JC2}
\]
Thus $\kappa$ is bounded.  The fractions $\nu,\ell$ cannot both
be infinitesimal.  If $\kappa$ is a unit, then $\ell=\kappa\zeta$
is a unit.  This gives precisely the following six source cells:
\[
\begin{array}{c|c|c|c|l}
\kappa&\nu&\ell&\text{additional predicate}&\text{full certificate}\\\hline
\text{unit}&\text{unit}&\text{unit}& &\mathrm{FU+UKR}\\
\text{unit}&\text{infinitesimal}&\text{unit}& &\mathrm{NU}\\
\text{infinitesimal}&\text{unit}&\text{unit}& &\mathrm{FM+MB}\\
\text{infinitesimal}&\text{unit}&\text{infinitesimal}& &\mathrm{EL}\\
\text{infinitesimal}&\text{infinitesimal}&\text{unit}&\tau\text{ bounded}&\mathrm{BN}\\
\text{infinitesimal}&\text{infinitesimal}&\text{unit}&\tau\text{ unbounded}&\mathrm{UN}
\end{array}                                                     \tag{JC3}
\]
In its last two rows define only there
\[
                q_\nu=\nu/\ell,\quad\tau=h/q_\nu,
                       \quad n_s=\beta/q_\nu^2.          \tag{JC4}
\]
BN covers every bounded $\tau$, including bounded $n_s$ and
literal $\beta=0$; a separate SN terminal is unnecessary.
The exact SN dictionary is retained for BN and UN.  SE proves
$e$ infinitesimal exactly when $\kappa$ is infinitesimal, so
its complete offmatched-Bob raw frames EB/SE apply on the four
small-$\kappa$ cells.  CE supplies their noninfinitesimal-Alice
charts.  FU and NU supply their own finite and reciprocal raw
joins.  The local certificates retain every original state normal
and every allowed $D=0$ or $\beta=0$ face.  Consequently JC3
covers all independent testing angles and states on C14.4.

\paragraph{The large-$L$, large-$\eta$ source.}
Put $\sigma=1/\eta$ and $H=h/\sigma$.  For bounded $H$,
the complete SF theorem in Section~\ref{component:p3be5a8005801}, with BE/CE and JR/JB, treats all
states and tests.  It uses JR/JB on bounded
$M_\beta=\beta(1-h)/\sigma^2$ and its full normal conductor
on the unbounded part.  Thus literal $\beta=0$ is included.
For unbounded $H$, SF first treats $\operatorname{res}h<1/2$.
US treats the unit beta-sphere source.  BS, BC, GA, SU and CE
treat the bounded-state, bounded testing-ratio, unit-Alice, sign
and endpoint alternatives.  Their remaining common domain lies in
\[
 \operatorname{res}h\ge1/2,\quad k/\sigma\in\operatorname{Inf},
 \quad b_0\in\operatorname{Inf},\quad d_0\sim1,
 \quad u,w>0\text{ infinitesimal}.                        \tag{JC5}
\]
AR treats the entire remaining domain with its restored angular
budget and every SP first normal, with no additional bound on
$W/Z$, $u/w$ or $uW/(wZ)$.  This is the complete C14.1 join.

\paragraph{The C14.2 testing and state partition.}
Here $L$ is a positive source unit and $\eta$ is unbounded.
LA treats every noninfinitesimal Alice test; LB treats bounded Bob
tests at local Alice.  Literal axes and the $WZ\le0$ core sign
are direct.  It remains that $u,w>0$ are infinitesimal and
$WZ>0$.  The positive source ratio $H=h/\sigma$ has a real
lower bound by BE7.  If it is bounded, CF treats its
$\lambda\le0$ part, while HP treats its $\lambda>0$ part,
using HG when the testing state ratio is unbounded.  Here
$\lambda$ has exactly the definition in the HP sphere dictionary.
If $H$ is unbounded and $W/Z$ is bounded, CF applies.
If both are unbounded, split the positive testing ratio
$c=uW/(wZ)$ into infinitesimal, positive unit or unbounded.
CS1--CS4 treat the extreme ratios.  CT treats every positive unit
$c$: all coefficients in CT6, CT10 and CT13 remain bounded or
infinitesimal under that hypothesis, without using
$\operatorname{res}c=\operatorname{res}L$.
This supplies every C14.2 source/test/state tuple, including all
original coefficient-zero faces.  No mixed inverse is introduced.

\paragraph{The C14.3 source and state partition.}
Here $x,L$ are infinitesimal and $\eta$ is bounded.
The full SB remote theorem treats the complement of both testing
angles infinitesimal.  The literal axes and the $WZ\le0$ core
sign are direct.  When $\zeta$ is unbounded and $\eta$ is a
unit, EU treats all remaining tests and states.  Otherwise the
remaining C14.3 source has $\eta$ infinitesimal and either unit
$\zeta$ or unbounded $\zeta$.  Use respectively the actual
DG or LD regular source formulas and their common full damped
identity, with $s=x+\eta$ or $s=L+\eta$.
The positive source deficits, full angular block and joint radial
block do not require a bound on $x/\eta$ or $L/\eta$.
Their complete remainder has one remaining first normal
$\epsilon uvQ\xi$, with $\epsilon=O(s)$.
For $|X/Z|$ unbounded, DS3 is a direct complete-parent exit.
For bounded $X/Z$, CJ conducts the whole remaining normal using
the already retained $\widehat J$ row, with no bound on $W/Z$
or $Y/Z$.  Its coefficients are finite separated sums: expand
$X/Z+hu$, use $\eta/s\le1$, and cancel
$x/\sqrt a=\sqrt aL$ before estimating the large-$\zeta$
source.  These exact DG/LD/CJ identities include both literal
local axes and preserve both state normals.  They therefore prove
C14.3 at all relative source/testing rates.  The earlier restricted
DG/LD and bounded-state admissions are not separate terminal
requirements.

\paragraph{Exhaustion of the entered C1.4 source.}
If $L$ is unbounded, $1-Lr>0$ forces $r$ infinitesimal.
The native relation $1-r=a_{\rm nat}\eta$ with
$0<a_{\rm nat}<2$ excludes infinitesimal $\eta$.
Unit $\eta$ is the full large-L/unit-eta theorem in Section~\ref{component:pd4b92bf7de71};
unbounded $\eta$ is the C14.1 join just proved.
If $L$ is a unit, infinitesimal, unit and unbounded $\eta$
are respectively LC in Section~\ref{component:p080aa62de16d}, LU in Section~\ref{component:p686209b87229}, and C14.2.
If $L$ is infinitesimal, unbounded $\eta$ is C14.4.
For bounded $\eta$, unbounded $\zeta$ is C14.3a, and unit
$\zeta$ with infinitesimal $\eta$ is C14.3b.  The only
remaining case has both $\eta,\zeta$ positive units.  Its actual
positive gap is
\[
 j=\eta(1+xh)-1+r>0,\quad j_0=\eta+r-1,
                   \quad j-j_0=\eta xh\in\operatorname{Inf}.
                                                               \tag{JC6}
\]
If $j$ is a unit, $j_0$ is a positive unit and RO applies.
If $j$ is infinitesimal, $\eta+r\sim1$; unit $r$ gives JO
and infinitesimal $r$ gives HO.  Their full scalar dictionaries,
complete normal allocations and reciprocal testing joins are retained.
This does not compare $j$ to $x$ or any source rate to a testing rate.

These alternatives give seventeen terminal source families after
expanding JC3 into its six rows.  Every one has its complete
source, testing and state certificate.  Thus, for the unchanged
entered predicate $\mathcal C$,
\[
 \mathcal C\Longrightarrow
 \bigvee_{\text{the seventeen displayed source cells}}
      \text{complete separated full-parent SOS}.               \tag{JC7}
\]
Apply each real-state certificate independently to the real and
imaginary state parts, and restore the positive pure-source and
pure-testing homogenization factors.  This proves the C1.4 join
with the original signs, strict gaps, allowed literal zero weights
and every projective endpoint.


\section{\texorpdfstring{The entire bounded-$h/q$ part of the doubly small endpoint}{The entire bounded-h/q part of the doubly small endpoint}}
\label{component:p80b0c5b9879f}

\label{p80b0c5b9879f:sec:new-c144-bounded-h-ratio}

\paragraph{Source scope and result.}
Retain the original entered physical C14.4 source, all strict gaps and
all original prior-complement intersections. Assume
\[
 S=\sigma+L,\ \kappa=x/S,\ \nu=\sigma/S
       \in\operatorname{Inf}_{>0},\qquad
 \ell=1-\nu,\quad q=\nu/\ell,\quad
 \tau=h/q\ \hbox{is bounded}.
 \tag{BN1}
\]
The complete original parent is a finite separated tensor SOS for
every test and every real or complex state on BN1. The quantity
$\beta/q^2$ is unrestricted, including zero and unbounded values.
The proof below is a fresh circuit completion of SN, not an import of
a missing historical endpoint proof.

\paragraph{Exact source estimates.}
Use the definitions of SN2, except that $n_s=\beta/q^2$ need not be
bounded. The exact sphere SN3 remains valid. Put $b=\beta/q$ and
$z=\kappa\beta/q^2$. Then
\[
\begin{gathered}
 h=q\tau\in\operatorname{Inf},\quad\theta\sim1,\quad
 e=\kappa\tau/\ell\in\operatorname{Inf},\quad M\sim1,\\
 b=\frac{\tau(B/S)-\ell H_0}{2\theta(P/S)},\qquad
 \mu\ell^2=\frac{\tau(1+x)(1+\sigma r)}{2M(P/S)},\\
 Jd_s+\theta z=\frac{\ell^2\mathcal H}{2(P/S)}\sim\tfrac12.
\end{gathered}\tag{BN2}
\]
In particular $b,z,d_s,\mu,\tau$ are bounded, $\tau$ and $\mu$
have positive real lower bounds, and at least one of $d_s,z$ is a
positive bounded unit. The exact inequality $\beta\ge0$ supplies
the lower bound on $\tau$. Also
$c_0,a_0,a_1\sim1/2$, $D,\alpha,m,\beta$ are infinitesimal.
Although the two terms below need not have small difference before
subtraction, their displayed bounded formulas in BN2 give
\[
                  A_{00}:=\mu\ell^2-b\sim\tfrac12.
 \tag{BN3}
\]
This argument takes residues only of bounded source coefficients and
never uses an estimate divided by $q$. It allows arbitrary relative
rates of $S,\kappa,q$.

\paragraph{A retained circuit and two actual positive rows.}
Use the full reciprocal parent SN5, the two unrestricted normals SN6,
and the shifted exact scalar dictionary SN9. Let $\zeta=\ell/\kappa$,
$a=u^2,c=w^2$, $x=\kappa S$, and assume $u,w\ge0$ by the original
sign gauges. Both may be zero. For every state define
\[
\begin{array}{ll}
 A_1=\sqrt\beta\,uZ,&A_2=\sqrt\beta\,W/(\zeta\sqrt q),\\
 A_3=\sqrt\beta\sqrt{x/q}\,wZ,
 &A_4=\sqrt\beta\,uwW/(\zeta q\sqrt x),
\end{array}\tag{BN4}
\]
\[
 \Psi=\sum_{i=1}^4 A_i^2-\frac{4\beta}{\zeta q}uwWZ,
 \qquad E_1=A_1-A_2,\quad E_2=A_3-A_4.
 \tag{BN5}
\]
Every monomial and its radical is a product of a pure-source and a
pure-testing element. This remains true if $x/q$ is unbounded.
The exact inequality meant here is a separated SOS certificate:
\[
                    \Psi-\tfrac12(E_1^2+E_2^2)\in\Sigma T^2.
 \tag{BN6}
\]
Indeed, if $WZ\ge0$, then $A_1A_2,A_3A_4\ge0$ and
\[
 \Psi=E_1^2+E_2^2+
              2(\sqrt{A_1A_2}-\sqrt{A_3A_4})^2.
\]
The radicals displayed here split into pure-source and pure-testing
radicals. If $WZ<0$, use instead
\[
 \Psi-\tfrac12(E_1^2+E_2^2)
 =\tfrac12(A_1+A_2)^2+\tfrac12(A_3+A_4)^2
                  +\frac{4\beta}{\zeta q}|uwWZ|.
\]
The last positive separated monomial is a square. Thus BN6 is valid
for all states, including $WZ=0$. At the actual source contact every
$A_i$ agrees, so $\Psi$ and its first contact derivatives vanish.

\paragraph{The exact scalar after circuit subtraction.}
Write $N_0=W/\zeta-uwZ=N/\ell$, where $N$ is SN7, and keep
$G,B,B_u,R$ from SN7. Thus
\[
 R=\frac{\sqrt S}{\zeta}G-uN_0.
 \tag{BN7}
\]
Subtract the entire circuit BN5 from the exact scalar graph in SN9.
Define
\[
\begin{gathered}
 A=\zeta^2 A_N-b,\qquad A_P'=A_P-zc/\ell^2,\\
 K'=K-2\beta/(\zeta q),\quad J'=K'-\ell A_P',
 \quad L'=J'-A/\zeta,\\
 H'=F_Z-\beta a-\beta xc/q-\ell^2ScA_P'
                          +ac(A-2\zeta J').
\end{gathered}\tag{BN8}
\]
Then $F_0-\Psi=A N_0^2+A_P'G^2-2\zeta uwZL'N_0+H'Z^2$.
All coefficients of $A$ and $A_P'$ are bounded:
\[
 A=\mu\ell^2(1+ra+ec)^2
       +(\beta+\tfrac12)\tau^2c(1+c)-b,
 \qquad \overline A=\overline{A_P'}=\tfrac12.
 \tag{BN9}
\]
Here and below the residue is taken at infinitesimal $u,w$.

The circuit subtraction has especially simple exact effects on the
SN10 quotient dictionary:
\[
\begin{gathered}
 L'_a=L_a,\qquad L'_c=L_c+\kappa n_s/\ell,\qquad
 H'_{20}=H_{20},\quad H'_{02}=H_{02},\quad
 H'_{11}=H_{11}-2n_s,\\
 H'=S\delta^2(d_s+cH_{02})
               +c(a-x)\delta H'_{11}+c(a-x)^2H_{20}.
\end{gathered}\tag{BN10}
\]
In particular no coefficient $n_s$ remains on its own in the final
conductors. With $C_R$ and $L_a$ from SN11--SN13 put
\[
 C_G'=\kappa(H_{11}-2n_s)/\ell,\qquad
 C_N'=-H_{11}-2\ell L_c/\kappa.
 \tag{BN11}
\]
These are bounded. The cancellation for the second coefficient is
explicit:
\[
 C_N'=-Sa_0\ell^2r+(\beta+m+\tfrac12)\tau^2(c+1/\ell)
                 +2\ell\kappa\mu\tau-m\tau^2.
 \tag{BN12}
\]
For $C_G'$ use the bounded formula SN13 and $\kappa n_s=z$.
The exact full scalar is consequently
\[
\begin{split}
 F_0={}&\Psi+A N_0^2+A_P'G^2+d_sB^2+H_{02}(wB)^2
       +C_RR^2-2\sqrt S\,L_aGR\\
 &+C_G'G(wB)+C_N'N_0(wB_u).
\end{split}\tag{BN13}
\]
This is an identity in unrestricted state coordinates, not a
restriction to a contact graph in the full parent.

\paragraph{Free helper rows across the whole source sphere.}
First reserve the positive rows $E_1,E_2$ from BN6. If $d_s$ is a
unit, $wB$ has infinitesimal cost on $\sqrt{d_s}B$. Otherwise
$z$ is a positive bounded unit. The exact identity
\[
 E_2=-\sqrt z\,\{(w/\ell)G+
                          \sqrt S\,w(w-\sqrt q)Z\}
\]
then gives
\[
 wB=-\frac{w+\sqrt q}{\sqrt z}E_2
                   -\frac{w(w+\sqrt q)}\ell G.
 \tag{BN14}
\]
Every coefficient in BN14 is infinitesimal, so any fixed real pivot
on $wB$ is free at infinitesimal Gram cost. Only the chosen pure-source
unit $\sqrt z$ is inverted. Literal $\beta=0$ uses the first branch.

After this pivot has been purchased, the normalized SN8 identity is
\[
 \delta^2N_0^2+\frac{2\kappa\delta}{\ell}G(wB)
          +2\kappa(wB)^2-(wB_u)^2
                         =(\delta W/\zeta)^2.
 \tag{BN15}
\]
It buys any fixed pivot on $wB_u$ at infinitesimal cost. The order of
these two purchases is deliberate; no weak $B$ diagonal is needed
in the second step.

\paragraph{Both full normals and their only additional row.}
Use the complete $H_X,H_Y,H_{XY},Y_P,Y_a,Y_c$ from SN18 and its
bounded $\widehat A_X,\widehat X_a,X_c$ from SN17. These formulas
only require $\tau,\mu,\beta$ bounded, not $n_s$ bounded.
The entire first-normal part is
\[
\begin{split}
 L_X={}&\ell w\widehat A_XN_0
       +\frac{uw\sqrt S}{\ell}\widehat X_aG
       +d_s\sqrt S\,quB-bB_u+wX_c(wB_u),\\
 L_Y={}&\sqrt{r(1-r)}\{Y_PG+\sqrt S\,Y_aR+Y_c(wB)\}.
\end{split}\tag{BN16}
\]
The apparent extra $B_u$ coupling is paid by the actual circuit row
$E_1$, using the exact, zero-face-safe identity
\[
 bB_u=-(w+\sqrt q)\sqrt b\,E_1-b(w+\sqrt q)N_0.
 \tag{BN17}
\]
All coefficients on the right have zero residue because $b$ is
bounded. No inverse of $b$ or $\beta$ is used. In the $B$ term in
$L_X$, write $d_sB=\sqrt{d_s}(\sqrt{d_s}B)$; this is likewise
face-safe when $D=0$. All direct normal coefficients are
infinitesimal, and $H_X,H_Y\sim1/2$, $H_{XY}$ is infinitesimal.
There is no missing $X$ or $Y$ normal and no imposed state ratio.

\paragraph{Allocation and all-tests conclusion.}
Keep a fixed positive fraction of each circuit row in BN6.
Use the strict base rows $(N_0,G,X_e,N_y,\sqrt{d_s}B,E_1,E_2)$,
omitting any zero row. Their main four diagonal residues are $1/2$;
the last three have fixed positive pivots wherever their rows are
nonzero. Terms containing $R$ are infinitesimal matrix perturbations
by BN7. Purchase sufficiently large fixed pivots on $wB,wB_u$ by
BN14--BN15, and use them for the bounded cross terms in BN13 and
BN16. Formula BN17 pays the remaining first-normal row at
infinitesimal cost. Every purchase is taken from an explicitly
retained positive square or from the exact SOS remainder BN6.
The resulting bounded Gram has strictly positive definite real
residue. The surviving strict bounded-Gram lemma therefore gives
finite tensor SOS for SN5 with both unrestricted normals restored.

This proves all infinitesimal $u,w$, including the literal axes.
For noninfinitesimal Alice use the unchanged adjacent-$C,A$ theorem
CE1--CE10; $C,A$ are unbounded on BN1. For the other Bob charts use
the surviving small-$e$ unmatched theorem SE10--SE12, since
$\overline\ell=1$. The sign gauges and real/imaginary state split
are the original exact ones. The local argument already handles
either sign of $WZ$, so no state normalization or state-chart
exclusion is required. Both literal $D=0$ and $\beta=0$ are covered
whenever they intersect BN1.

\paragraph{The other source cells.}
The unbounded-$h/q$ source and the $\ell$-infinitesimal endpoint are proved in Sections~\ref{component:peaa56c1a3c7a} and~\ref{component:pb96c9bfcc12d}. Their source partition with the present bounded-$h/q$ theorem is printed in Section~\ref{component:pfd0dd7bb10cf}.


\section{\texorpdfstring{C14.4 at the entire $\ell$-infinitesimal source endpoint}{C14.4 at the entire ell-infinitesimal source endpoint}}
\label{component:pb96c9bfcc12d}
% Fresh rederivation. Independent analytical review is recorded separately.

\label{pb96c9bfcc12d:sec:c144-ell-endpoint-rederived}

\paragraph{Theorem and exact scope.}
Retain the original native C14.4 physical source, all strict gaps,
nonnegative event weights, total-scale conditions and prior-complement
intersections.  Let
\[
 S=\sigma+L,\quad x=\kappa S,\quad
 \nu=\sigma/S,\quad\ell=L/S,\quad\nu+\ell=1,
 \quad\zeta=L/x=\ell/\kappa.
\]
Suppose $S,\ell>0$ are infinitesimal and $\zeta$ is
noninfinitesimal.  Then the complete original parent is a finite
sum of squares in the separated ring
$K_c\otimes_{\mathbb R}K_t$ for every independently chosen
projective test and every real or complex state.  The proof includes
$\beta=0$ and unrestricted openings of $r,1-r,\theta$.
The intersection with $D=0$ is empty by the source estimate below.
No earlier missing R21 endpoint proof is used.

\paragraph{Source estimates at this endpoint.}
Use the exact actual-source formulas FM2--FM3 and SE1--SE3, with
$0<r,h<1$, $\theta=1-h$, $e=\kappa h/\nu$, $E=1-e\ge1/18$,
$U_s=(1+x)(1+\nu Sr)$, and
\[
\begin{gathered}
 J=1+xh-\nu S(1-r),\quad
 M=1+xr-\kappa\theta/\ell>0,\quad Q_s=1-\ell Sr,\\
 a_0=SAr,\qquad c_0=SC(1-r),\qquad P=SP_0.
\end{gathered}
\]
Put $g=\kappa/\ell=1/\zeta$, which is bounded.
Then $\nu\sim1$ and $\kappa\in\operatorname{Inf}$.
Since
\[
 0\le\beta\theta
  =\frac{U_s(\ell+\kappa h)-\theta P_0}{2P_0},
\]
we have $\theta=O(\ell+\kappa)=O(\ell)$ and
$\beta\theta=O(\ell)$, hence $h\sim1$,
$\beta\theta^2=O(\ell^2)$ and $M\sim1$.
The exact source formulas give
\[
\begin{gathered}
 J,Q_s,P_0,E\sim1,\quad c_0,a_0\sim\tfrac12,
 \quad D/S\sim\tfrac12,\quad m\in\operatorname{Inf}_{>0},\\
 \frac{qm}{\kappa^2}
   =\frac{U_s h}{2M\nu^2P_0}\sim\tfrac12,
 \qquad q=\ell/\nu,\\
 \frac{\alpha}{S\ell^2}
   =\frac{(1+\nu Sr)\{(1+xr)e/\ell+ME\}}{2Q_sP_0},
 \qquad
 \frac{\alpha}{S\ell^2}-\frac{1+g}{2}\in\operatorname{Inf}.
\end{gathered}                                                   \tag{EL0a}
\]
The last difference is uniform for bounded $g$, because
$e/\ell=gh/\nu$.  All statements are inside the source field;
none compares a testing rate to a source rate.  In particular
$D/S$ is a positive unit, while $\beta$ may be zero, bounded,
or unbounded.

\paragraph{The full parent and local testing chart.}
Noninfinitesimal Alice tests and unmatched Bob residues will be joined
at the end.  The remaining chart has $u,v\ge0$ infinitesimal,
including literal zero axes.  Write $a=u^2$, $c=v^2$,
$b=\sqrt q$, $j=\nu b$, and $d=D+1/2$.
The complete original quarter-parent is
\[
\begin{aligned}
 F/4={}&m(vW-X+uvY-uZ)^2+C(1+c)(uW-Y)^2\\
 &+A(1+a)(1+c)Y^2+d(uW-Y+uvX-vZ)^2\\
 &+(\alpha+\tfrac12)(1+a)(vY-Z)^2
 +(\beta+\tfrac12)(1+c)(X+uZ)^2\\
 &-\tfrac12(1+a)(1+c)Z^2.
\end{aligned}                                                     \tag{EL0b}
\]
All state coordinates are arbitrary elements of the testing field.
No state coordinate is divided by or bounded below.
Put $R_s=1+xr$, $p=\kappa\theta/(\ell R_s)$, $\mu=m(1-p)$, and
\[
 V=v(W+uY),\qquad N=X+uZ-pV,\qquad
 Y=ruW+\sqrt{Sr(1-r)}N_y.                         \tag{EL1}
\]
These are pure-source rescalings of polynomial state normals.
The original $m$ event is $(1-p)V-N$, the beta event is $N+pV$,
and the core event is
$uW-Y+puvV-(1+a)vZ+uvN$; thus no term has been dropped.
On $N=N_y=0$ define
\[
\begin{aligned}
 L_0&=\mu(1-p)+(\beta+\tfrac12)p^2,
 \qquad L_1=(\beta+\tfrac12)p^2,\\
 B_c&=1-r+pc(1+ra),\\
 \lambda(a,c)&=\frac q{\kappa^2}(1+ra)^2(L_0+L_1c),\\
 A_P&=(1+c)[c_0(1-r)+a_0r(1+a)]
       +S[dB_c^2+(\alpha+\tfrac12)(1+a)r^2c],\\
 K&=(1+a)[dB_c+(\alpha+\tfrac12)r],\\
 F_Z&=(1+a)[\alpha+Dc+dac].
\end{aligned}                                                    \tag{EL2}
\]
Here $c_0=SC(1-r)$ and $a_0=SAr$.  The complete graph scalar is
\[
 F_0=c(1+ra)^2(L_0+L_1c)W^2+A_PP^2-2vKPQ+F_ZZ^2,
 \quad P=uW/\sqrt S,\quad Q=\sqrt S Z.             \tag{EL3}
\]
Let $T=\kappa W,U=uZ$ and use the exact four-term circuit
\[
 \Psi=(\nu/\ell)cT^2+\nu^2cU^2+\kappa\nu\ell Q^2
       +\kappa P^2-4\kappa\nu vPQ.                \tag{EL4}
\]
Its four positive monomial square roots are
$A_1=(\kappa/b)vW$, $A_2=\nu vuZ$,
$A_3=\nu b\sqrt xZ$, and $A_4=\sqrt\kappa P$.
Their product is a square.  The usual two paired AM--GM identities
make $\Psi$ a separated SOS, including $WZ<0$ by sign excess.
The multiplier $\lambda$ itself is a sum of separated monomial
squares with positive pure-source coefficients; no mixed radical
of $\lambda$ is taken.

Set $\Gamma=P-\nu vQ$, $\Lambda=vP-\ell Q$, and
\[
 \widehat A=A_P-\kappa\lambda,\quad
 \widehat K=K-2\kappa\nu\lambda,\quad
 \widehat F=F_Z-\lambda(\nu^2ac+\kappa\nu\ell S).
\]
Then the following equality is purely polynomial and exact:
\[
 F_0-\lambda\Psi
   =\widehat A\Gamma^2-2vJ\Gamma Q+HZ^2,
 \quad J=\widehat K-\nu\widehat A,
 \quad H=\widehat F+\nu^2Sc\widehat A-2\nu Sc\widehat K.
                                                               \tag{EL5}
\]
The exact coefficient bounds are
$\lambda_0=qL_0/\kappa^2\sim1/2$ and
$\lambda_1=qL_1/\kappa^2=O(\ell)$.
Indeed the exact opening-safe formulas are
\[
 \lambda_1=\frac{\theta(2\beta\theta+\theta)}
                    {2\ell\nu R_s^2},\qquad
 \lambda_0=\frac{U_shM}{2\nu^2P_0R_s^2}+\lambda_1.
\]
Here $\theta/\ell$ and $(\beta\theta)/\ell$ are bounded, so
the formula for $\lambda_1/\ell$ has only bounded source factors.
No comparison between $S$ and $\ell$ is made.  The mixed term
arising from $-\nu^2\lambda_1c^2a$ is allocated in EL18 below
using an actual row of the same circuit.

\subsubsection*{Exact helper rows and the complete two-normal conductor}
The following exact budgets and complete normal conductors use a
positive real fraction of $\Psi$ and strict pivots on
$(\Gamma,\Lambda)$.  The scalar calculation below supplies precisely
these pivots; it does not use the state normals to obtain them.
Put $\delta=v-b$, $G_0=P-jQ$, $B_o=\delta Q$, and
\[
\begin{gathered}
 E_{12}=A_1-A_2=(v/b)(T-jU),\quad V_h=v\delta U,\\
 R_v=v(a-x)Z,\quad B=\sqrt S(c-q)Z,\quad
 R_0=\ell\sqrt S(a-x)Z.
\end{gathered}
                                                               \tag{EL6}
\]
The exact angular relations are
\[
 (v+b)G_0=\Lambda+b\Gamma,\qquad
 \nu(v+b)B_o=\Lambda-v\Gamma=\nu B.             \tag{EL7}
\]
In particular $B$ has a bounded linear expression in the two strict
angular rows.  Multiply MB3 by $v^2/j^2$ and use
$\eta=\kappa(v+b)/j\in\mathrm{Inf}$.  The exact SOS inequality is
\[
 V_h^2\preceq_{\rm SOS}
  \frac{\delta^2}{\nu^2}E_{12}^2
   +\eta(\Lambda+b\Gamma)^2
   +\frac{\eta+2\kappa}{\nu^2}(\Lambda-v\Gamma)^2.
                                                               \tag{EL8}
\]
For a direct verification, set $R_1=vG_0,R_2=vB_o$.
The mixed-term replacement uses the legal SOS
\[
 \eta(R_1^2+R_2^2)-2\kappa\delta R_1R_2/j
 =\frac{\kappa v}{j}(R_1-R_2)^2
  +\frac{\kappa b}{j}(R_1+R_2)^2.
\]
Replacing $v^2$ by $(v+b)^2$ leaves
$(2vb+b^2)G_0^2$ and $(2vb+b^2)B_o^2$, with positive
separated multipliers.  The remaining MB3 square is
$(v\delta T/j)^2$.  Every cost in EL8 is infinitesimal, since
$\kappa/j=(\kappa/\ell)b$ is infinitesimal.

Two further exact relations are
\[
 R_v=\frac{\kappa v\sqrt S}{\nu b}\Gamma
       -\frac u\nu E_{12}+\frac{xv(v-b)}bZ,          \tag{EL9}
\]
\[
 \frac{x^2}{qS}B^2-
       \left(\frac{xv(v-b)}bZ\right)^2
   =\frac{x^2}{q}(2vb+b^2)(v-b)^2Z^2.              \tag{EL10}
\]
Thus $R_v$ is free at infinitesimal Gram cost because
$x^2/(qS)=\kappa S(\kappa/q)$ is infinitesimal.  Finally
\[
 2\nu^2ScR_v^2+2\nu^2(a-x)^2B^2-R_0^2
       =[\nu\sqrt S(2c-q)(a-x)Z]^2               \tag{EL11}
\]
buys $R_0$ with infinitesimal costs.  None of EL6--11 bounds or
divides a state coordinate or a mixed source/testing sum.

For the exact complete $N$ half derivative set
\[
\begin{aligned}
 B_0(c)&=[-\mu+(\beta+\tfrac12)p(1+c)]/\kappa,\\
 F_a(a,c)&=r\kappa B_0(c)+d(1-r)+dpc(1+ra),\\
 D_t(a,c)&=d(1+a)-\nu S F_a(a,c).
\end{aligned}                                                     \tag{EL12}
\]
All their coefficients are bounded, including at $\beta=0$ or
unbounded $\beta$: $\mu/\kappa$ and
$(\beta+1/2)p/\kappa=(\beta\theta+\theta/2)/(\ell R_s)$
are bounded.  Write the monic polynomial differences
$D_t-D_t(x,q)=(a-x)D_a+(c-q)D_c$ and
$B_0(c)-B_0(q)=(c-q)B_{01}$.  Source contact gives
$D_t(x,q)=\nu B_0(q)$.  The entire half derivative is
\[
\begin{aligned}
 L_N={}&bB_0(c)E_{12}+uv\sqrt S F_a\Gamma-uvD_aR_v\\
 &+[-D_t(x,q)+\nu bB_{01}(v+b)-v(v+b)D_c]V_h.
\end{aligned}                                                     \tag{EL13}
\]
Hence EL8--10 allocate every $N$ cross term without a state bound.
The normal $\widehat N=\sqrt{1+\beta}N$ has diagonal at least
$(\beta+1/2)/(1+\beta)\ge1/2$; dividing the cross coefficients
by $\sqrt{1+\beta}$ only improves the displayed bounds.

The complete other half derivative is
\[
 L_y=\sqrt{r(1-r)}[Y_P\Gamma+vY_KQ],             \tag{EL14}
\]
where
\[
\begin{aligned}
 Y_P={}&(1+c)[-c_0+a_0(1+a)]\\
 &+S\{c(1+ra)(L_0+L_1c)+d(-1+pac)B_c
                   +(\alpha+\tfrac12)(1+a)rc\},\\
 Y_K={}&\nu Y_P+(1+a)(D-\alpha-dpac).
\end{aligned}
\]
The polynomial coefficients are bounded, $Y_P$ is infinitesimal,
and actual source contact gives $Y_K(x,q)=0$.  Monic division gives
$Y_K=(a-x)Y_a+(c-q)Y_c$ with bounded coefficients, so
\[
 L_y=\sqrt{r(1-r)}[Y_P\Gamma+\sqrt S Y_aR_v+vY_cB].
                                                               \tag{EL15}
\]
The exact quadratic normal coefficients are
\[
\begin{aligned}
 H_N={}&m+dac+(\beta+\tfrac12)(1+c),\\
 H_y={}&rc_0(1+c)+(1-r)a_0(1+a)(1+c)\\
 &+Sr(1-r)[d(-1+pac)^2+(\alpha+\tfrac12)(1+a)c
                         +ac(L_0+L_1c)],\\
 H_{Ny}={}&\sqrt{Sr(1-r)}uv[\kappa B_0(c)-d+dpac].
\end{aligned}                                                     \tag{EL16}
\]
After the source-only normalization of $N$, the first two diagonal
residues are respectively in $[1/2,1]$ and $1/2$, and the cross
coefficient is infinitesimal.  Equations EL13--16 include all terms
of the full original parent; they require no $W/Z$ or testing/source
rate restriction.  The scalar allocation and exact strict-Gram assembly are given next.


\subsubsection*{Complete scalar allocation, including the mixed radial term}
Both $\widehat A$ and the coefficients of $J,H$ in EL5 are
bounded polynomials in $a,c$: their source coefficients are built
from bounded $p,d,\alpha,c_0,a_0,\lambda_0,\lambda_1,\kappa,S,\nu$.
In fact $p=O(\kappa)$,
$\mu/\kappa=O(\kappa/\ell)$, and
$(\beta+1/2)p^2=O(\kappa^2)$.
Consequently $\widehat A\sim1/2$ on the local chart.
Actual source contact, checked in the original parent, gives
\[
 J(x,q)=0,\qquad H(x,q)=H_a(x,q)=H_c(x,q)=0.        \tag{EL17}
\]
These are exact identities, not residue equalities.  They also follow
by evaluating the actual contact state in EL1: all four circuit
monomials agree, $\Gamma=0$, and both graph normals vanish.
The vanishing first derivatives of the original contact parent
then give the three scalar jets.  The focused checker verifies all
of them directly from FM2--FM3.

The additional actual circuit row
\[
 E_{23}=A_2-A_3=\nu(uv-\sqrt{xq})Z
\]
resolves the otherwise unbounded mixed coefficient.  Direct expansion
using $x=\kappa S$ gives the exact identity
\[
 c(a-x)(c-q)Z^2
 =\frac{c-q}{\nu^2}E_{23}^2
  +\frac{2\sqrt{\kappa q}}{\nu}E_{23}B-\kappa B^2.
                                                               \tag{EL18}
\]
Every coefficient on the right is infinitesimal.  To see the equality
without division by a testing coordinate, use
\[
 (ac-xq)Z^2=E_{23}^2/\nu^2+2\sqrt{xq}\,ZE_{23}/\nu,
 \qquad c(a-x)=ac-xq-x(c-q).
\]
All square roots in this calculation are pure-source roots.

For completeness, the single circuit supplies positive fixed pivots
on both $E_{12}$ and $E_{23}$.  If $WZ\ge0$, use the two AM--GM
pairings having respectively these differences and average them;
then $\Psi-\tfrac14(E_{12}^2+E_{23}^2)$ is separated SOS.
If $W,Z$ are both negative, first change the common state sign;
the parent and all squared differences are unchanged.
If $WZ<0$, the negative product term of $\Psi$ is nonnegative,
and
\[
 \sum_{i=1}^4 A_i^2-\tfrac14[(A_1-A_2)^2+(A_2-A_3)^2]
\]
is a real positive-definite quadratic form in the $A_i$.
For example its leading $3\times3$ block has diagonal
$(3/4,1/2,3/4)$ and successive leading minors
$3/4,5/16,3/16$; the fourth diagonal is $1$.
This supplies a finite real Gram factorization.  The remaining
$-4\kappa\nu uvWZ$ is a separated square on this pure-testing sign
chart.  The zero-state boundaries use the same identities.
Since $\lambda_0\sim1/2$ and $\lambda_1\ge0$,
$\lambda-1/4$ is a separated SOS multiplier:
\[
 \lambda-\tfrac14=(\lambda_0-\tfrac14)
       +\lambda_0(2ra+r^2a^2)+\lambda_1c(1+ra)^2.
\]
Thus $\lambda\Psi$ leaves strict real diagonal $1/16$ on each of
$E_{12},E_{23}$ and an explicitly retained separated SOS remainder.
No mixed square root of $\lambda$ has been used.

Perform monic polynomial division in $a$ and then in $c$.
By EL17 there are bounded coefficient polynomials such that
\[
\begin{aligned}
 J&=(a-x)J_a+(c-q)J_c,\\
 H&=(a-x)^2H_{20}(a,c)
       +(a-x)(c-q)H_{11}(c)+(c-q)^2H_{02}(c).
\end{aligned}                                                     \tag{EL19}
\]
The division is by monic polynomials; it introduces no inverse of
$a-x$, $c-q$, or any mixed expression.  The following four particular
coefficient identities are exact:
\[
\begin{aligned}
 H_{20}(a,0)&=-\kappa\nu\ell S\lambda_0r^2,\\
 H_{11}(0)&=-\alpha/q+2\kappa\nu^2S\lambda_0r(1+rx),\\
 H_{02}(0)/S
   &=\nu^2\left[(1+x)\frac\alpha{S\ell^2}
                 -g\nu\lambda_0(1+rx)^2\right]\sim\tfrac12,\\
 H_{02}(c)&=H_{02}(0)+\nu^2S^2dp^2(1+rx)^2c.
\end{aligned}                                                     \tag{EL20}
\]
The first three also follow immediately by evaluating
$H(a,0)=\alpha(1+a)-\kappa\nu\ell S\lambda_0(1+ra)^2$
and using EL17; the last follows from the $c^3$ coefficient.
In particular $H_{11}(0)/(S\ell)$ is bounded and
$H_{02}(c)/S\sim1/2$.  The $g$ terms cancel in the real residue
of the third line by EL0a.  This cancellation is why the independent
ratio $S/\ell$ never enters the estimate.

Write $H_{20}=H_{20}(a,0)+c\widetilde H_{20}$ and
$H_{11}=H_{11}(0)+c\widetilde H_{11}$, with bounded coefficients.
Substitution into EL5 gives the exact scalar assembly
\[
\begin{aligned}
 F_0={}&\lambda\Psi+\widehat A\Gamma^2
                 +\frac{H_{02}(c)}S B^2\\
 &-2\sqrt S J_a\Gamma R_v-2vJ_c\Gamma B
       -g\nu\lambda_0r^2R_0^2+\widetilde H_{20}R_v^2\\
 &+\frac{H_{11}(0)}{S\ell}R_0B\\
 &+\widetilde H_{11}
       \left[\frac{c-q}{\nu^2}E_{23}^2
             +\frac{2\sqrt{\kappa q}}\nu E_{23}B-\kappa B^2\right].
\end{aligned}                                                     \tag{EL21}
\]
This includes every scalar term, including the one involving
$\lambda_1c^2a$.  It supplies strict real pivots on
$(E_{12},E_{23},\Gamma,B)$; the two radial rows have arbitrary
fixed positive pivots available at infinitesimal cost by EL9--11.
There is no discarded source-dependent signed term.

\subsubsection*{Strict bounded Gram and all-test completion}
Use the main rows
\[
 E_{12},\ E_{23},\ \Gamma,\ B,\ \widehat N=\sqrt{1+\beta}N,
 \ N_y
\]
and the three auxiliary rows $V_h,R_v,R_0$.
EL21 and EL13--16 express the entire EL0b as a bounded quadratic
form in these nine rows, together with the retained circuit SOS
remainder.  The real diagonal of the main block is bounded below
by $(1/16,1/16,1/2,1/2,1/2,1/2)$.
Its off-diagonal entries are infinitesimal except for couplings to
the auxiliary rows.  These are bounded, as are the possibly signed
auxiliary diagonal entries.

Choose a sufficiently large fixed real positive diagonal on each
auxiliary row, and purchase it using EL8--11.  This is an exact
SOS operation: if $Q-R^2$ is a displayed SOS budget, replace the
quadratic form by its difference with a real multiple of $Q-R^2$
and retain that multiple as a separate SOS.  EL8 and EL9--10 cost
only infinitesimal entries on the main rows; EL11 is purchased
last and costs an infinitesimal entry on $R_v$ as well as on $B$.
The real auxiliary diagonals can be chosen large enough that the
real Schur complement is positive definite.  This argument uses
only a finite real matrix, never a mixed-field matrix inverse.
The resulting bounded coefficient Gram has strictly positive real
residue.  The inherited strict separated bounded-Gram lemma in
Section~\ref{component:p472f434581f3} makes the complete form a finite
SOS in $K_c\otimes_{\mathbb R}K_t$.  Restoring the budget and circuit
remainders proves the original full parent.

There is no restriction on $W,Z$, their ratio, or either normal;
$\beta$ was never discarded or divided by.  The actual retained
normal diagonal handles an arbitrarily large $\beta$ before any
cross-term estimate.  The only inverse involving an arbitrarily
small opening is the pure-source scaling used to define $N_y$;
its inverse never appears in a Gram coefficient.  All helper roots
are products of source and testing factors or finite sums thereof.

The local proof includes $u=0$ and $v=0$, since it never divides
by either.  Negative $u$ is reduced by
$(u,Y,Z)\mapsto(-u,-Y,-Z)$ and negative $v$ by
$(v,X,Z)\mapsto(-v,-X,-Z)$; EL0b is invariant under each map.
Real and imaginary state parts are handled separately.
By EL0a both $C,A$ are unbounded, so the surviving adjacent-$C,A$
proof covers noninfinitesimal Alice tests.  Since $\kappa$ and $e$
are infinitesimal, the surviving SE10--SE12 theorem covers every
homogeneous Bob sine-square residue different from $\bar\ell=0$,
including Bob infinity.  The only remaining Bob chart is exactly
finite infinitesimal $v$, proved above.  Therefore all projective
tests and all states on the original $\ell$-infinitesimal source
intersection are covered.  All original source predicates are
retained; no numerical sampling or historical completion marker
is an input to this conclusion.


\section{\texorpdfstring{Fresh C14.4 completion when $\kappa$ is small and both Bob source fractions are units}{Fresh C14.4 completion when kappa is small and both Bob source fractions are units}}
\label{component:p2edb55cd1f69}

\label{p2edb55cd1f69:sec:new-c144-finite-matched}

\paragraph{Scope.}
This theorem treats the source sector stated in FM1: $S,\kappa$ are positive infinitesimals and $\nu,\ell$ are positive bounded units. The source endpoints $\nu\in\operatorname{Inf}$ and $\ell\in\operatorname{Inf}$, and the unit-$\kappa$ source, are separate cells in the complete partition of Section~\ref{component:pfd0dd7bb10cf}. All original physical C14.4 source conditions, strict gaps, nonnegative event coefficients, total-scale conditions and prior complement intersections are retained.

\paragraph{Theorem.}
Let the original native C14.4 source have
\[
 S=\sigma+L>0,\qquad x=\kappa S,\qquad
 \nu=\sigma/S,\quad\ell=L/S,\quad\nu+\ell=1,
 \qquad S,\kappa\in\operatorname{Inf}_{>0},
 \tag{FM1}
\]
with $\nu,\ell$ positive bounded units. Then the complete original
parent is a finite SOS in $K_c\otimes_{\mathbb R}K_t$ for every
independent projective test and every real or complex state.
The argument below supplies the missing finite matched-Bob chart;
the other tests use the surviving explicitly scoped SE theorem.

\paragraph{Exact opening-safe source dictionary.}
Write $\theta=1-h$, $0<r,h<1$, $e=\kappa h/\nu$, $E=1-e$,
and set
\[
\begin{split}
 J&=1+xh-\nu S(1-r),\quad
 M=1+xr-\kappa\theta/\ell,\quad Q_s=1-\ell Sr,\\
 P_0&=1+S\{\nu\ell+2\nu\kappa h+\kappa\ell+
                    \kappa^2h\theta\}
             +S^2\{\nu r\kappa\ell+\nu r\kappa^2h\}.
\end{split}\tag{FM2}
\]
These are the actual native source quantities, with $P=SP_0$.
The converted P1 coefficients, without applying P1's different
unit-$L$ asymptotic hypotheses, are exactly
\[
\begin{split}
 c_0=SC(1-r)&=\frac{(1+x)(1+\nu Sr)E}{2P_0},\\
 a_0=SAr&=\frac{(1+xh)E+(1-r)(\ell S+xh)}{2P_0},\\
 d_0=D/S&=\frac{(1+x)(1+\nu Sr)-JP_0}{2JSP_0},\\
 a_1=\alpha/S&=
 \frac{(1+\nu Sr)\{\ell(1+xr)e+\ell^2ME\}}{2Q_sP_0},\\
 b_1=\beta\theta&=
 \frac{(1+x)(1+\nu Sr)(\ell+\kappa h)-\theta P_0}{2P_0},\\
 \mu=m/\kappa^2&=
 \frac{(1+x)(1+\nu Sr)h}{2M\ell\nu P_0}.
\end{split}\tag{FM3}
\]
The numerator in the $d_0$ formula has a factor $S$, so its displayed
$S$ denominator cancels in the source polynomial ring. Thus all six
coefficients belong to the bounded source algebra generated by
$S,\kappa,h,r,\nu,\ell$ and inverses of $J,M,Q_s,P_0,\nu,\ell$.
None of these inverses is of an opening $r,1-r,\theta,D$ or $\beta$.
The original source gives
\[
\begin{gathered}
 J,M,Q_s,P_0,E\sim1,\qquad c_0,a_0\sim\tfrac12,\qquad
 d_0\sim\nu^2/2,\quad a_1\sim\ell^2/2,\\
 b_1-(h-\nu)/2\in\operatorname{Inf},\qquad
 \mu-h/(2\nu\ell)\in\operatorname{Inf},\qquad
 \bar h\ge\bar\nu>0.
\end{gathered}\tag{FM4}
\]
The last inequality follows from $b_1=\beta\theta\ge0$.
Hence $\mu$ is a positive bounded unit. These estimates are uniform
through $r,1-r,\theta\to0$; no comparison with a testing quantity
has been used. In this sector $D/S$ is a positive unit, so $D=0$
has empty intersection; $\beta=0$ is allowed.

\paragraph{Exact full parent and unrestricted normals.}
On the finite Bob chart write $D_u=1+u^2$, $D_v=1+v^2$.
The complete quarter-parent is
\[
\begin{split}
 F/4={}&m(vW-X+uvY-uZ)^2+C D_v(uW-Y)^2+A D_uD_vY^2\\
 &+(D+\tfrac12)(uW-Y+uvX-vZ)^2
 +(\alpha+\tfrac12)D_u(vY-Z)^2\\
 &+(\beta+\tfrac12)D_v(X+uZ)^2-\tfrac12D_uD_vZ^2.
\end{split}\tag{FM5}
\]
No state restriction is made. Put
\[
 N_x=X+huZ,\qquad N_y=\frac{Y-ruW}{\sqrt{Sr(1-r)}}.
 \tag{FM6}
\]
The source square root is nonzero even when its size is arbitrarily
small. It is used only to define a permitted tensor root; every
coefficient below is bounded without its inverse.

After the surviving unmatched-residue and adjacent-$C,A$ exits it is
enough to take $u$ infinitesimal and $v$ a bounded positive unit with
the same residue as $b=\sqrt{\ell/\nu}$. Set
\[
\begin{gathered}
 a=u^2,\quad j=\sqrt{\nu\ell},\quad T=\kappa W,\quad U=uZ,\\
 P=uW/\sqrt S,\quad Q=\sqrt S Z,\quad N=T-jU,\quad G=P-jQ,\\
 \delta=v-b,\quad B=\delta Q,\quad B_u=\delta U.
\end{gathered}
 \tag{FM7}
\]
Every state coordinate remains arbitrary, including zero and
unbounded coordinates. The rows satisfy the exact MB identities
\[
\begin{gathered}
 R_Z=(a-x)Z=(\kappa\sqrt S/j)G-(u/j)N,\\
 (a-x)\delta Z^2=(\kappa/j)GB-(1/j)NB_u,\\
 \delta^2N^2+2j\kappa\delta GB+2j^2\kappa B^2-j^2B_u^2
                  =(\delta T)^2.
\end{gathered}\tag{FM8}
\]
These are polynomial identities, not ordered-ring inequalities.

\paragraph{The full graph scalar.}
Let $F_0$ denote FM5 with $N_x=N_y=0$; it is used only as the
constant coefficient of the full normal expansion. Define bounded
source/testing polynomials
\[
\begin{split}
 A_T={}&\mu v^2(1+ra)^2,\\
 A_P={}&(1+v^2)\{c_0(1-r)+a_0r(1+a)\}\\
       &+S\{(D+\tfrac12)(1-r)^2+
                   (\alpha+\tfrac12)(1+a)r^2v^2\},\\
 A_Q={}&d_0v^2+a_1,\\
 A_U={}&m\theta^2+\alpha+b_1\theta+\theta^2/2\\
       &+v^2\{2Dh+h-\tfrac12+b_1\theta+\theta^2/2\}
              +(D+\tfrac12)h^2v^2a,\\
 K={}&m\theta(1+ra)+(D+\tfrac12)(1-r)(1+ha)
                         +(\alpha+\tfrac12)(1+a)r,\\
 H={}&vK-\kappa jA_T,\quad R_U=A_U-j^2A_T,\\
 g={}&jA_P-H,\quad R_Q=A_Q+j^2A_P-2jH,\\
 J_Z={}&SR_Q+aR_U.
\end{split}\tag{FM9}
\]
Direct expansion of every term of FM5 gives
\[
 F_0=A_TN^2+A_PG^2+2gGQ+J_ZZ^2.
 \tag{FM10}
\]
For example the only relation used to obtain FM10 from the
five monomials $T^2,U^2,P^2,Q^2,PQ$ is $TU=\kappa PQ$.

The actual source stationarity identities, also verified directly
from FM2--FM3, are
\[
 g(x,b)=0,\qquad
 J_Z(x,b)=\partial_aJ_Z(x,b)=\partial_vJ_Z(x,b)=0.
 \tag{FM11}
\]
These are coefficient-field identities at the actual source
contact; they impose no equation on independent tests.

\paragraph{Finite polynomial division.}
All following quotients mean monic polynomial division in formal
$a$ or $v$, performed before evaluation in the tensor ring. No
mixed inverse is introduced. There are bounded coefficient
polynomials $g_a,g_v,H_{20},H_{11},H_{02}$ defined exactly by
\[
\begin{split}
 g(a,v)&=(a-x)g_a(a,v)+(v-b)g_v(v),\\
 J_Z(a,v)&=(a-x)^2H_{20}(a,v)+(a-x)(v-b)H_{11}(v)
                                      +S(v-b)^2H_{02}(v).
\end{split}\tag{FM12}
\]
For clarity, first divide $g(a,v)-g(x,v)$ by $a-x$ and then
$g(x,v)$ by $v-b$. For the second line, first subtract the
constant and linear Taylor terms in $a-x$ and divide the remainder
by $(a-x)^2$. Divide $\partial_aJ_Z(x,v)$ by $v-b$ for $H_{11}$.
Finally divide the bounded polynomial
$R_Q(x,v)+\kappa R_U(x,v)$ by $(v-b)^2$ for $H_{02}$.
FM11 makes every division exact. In particular $H_{02}$ is
bounded without dividing a bounded error by $S$.

Substituting FM8 into FM10 now yields the exact scalar identity
\[
\begin{split}
 F_0={}&A_TN^2+A_PG^2+2g_vGB+H_{02}B^2
          +H_{11}\{(\kappa/j)GB-(1/j)NB_u\}\\
       &+H_{20}R_Z^2+2\sqrt S\,g_aGR_Z.
\end{split}\tag{FM13}
\]
Thus every coefficient is bounded and all remaining radial
contributions are infinitesimal matrix entries on $(N,G)$.

Here and below a displayed residue formula uses the real residues
of its bounded source parameters. FM3 gives explicitly
\[
 \bar A_T=\frac h{2\nu^2},\quad
 \bar A_P=\frac1{2\nu},\quad
 \bar g_v=\frac{\ell-\nu}{2},\quad
 \bar H_{02}=\frac\nu2,\quad \bar H_{11}=-\theta j.
 \tag{FM14}
\]
For example at $S=\kappa=0$ and $a=0$ the three relevant
polynomials are
$g=j(1+v^2)/2-v/2$,
$R_Q=\nu(v-b)^2/2$, and
$R_U=\theta(\ell-\nu v^2)/2$.
These establish FM14 by the finite divisions just specified.

\paragraph{Both complete first-normal terms.}
Write the full expansion of FM5 as
\[
 F/4=F_0+H_xN_x^2+H_yN_y^2+2H_{xy}N_xN_y
                  +2N_xL_x+2N_yL_y+\beta D_vN_x^2.
 \tag{FM15}
\]
The exact bounded diagonal and cross coefficients are
\[
\begin{split}
 H_x={}&m+(D+\tfrac12)u^2v^2+D_v/2,\\
 H_y={}&r c_0D_v+(1-r)a_0D_uD_v\\
       &+Sr(1-r)\{mu^2v^2+D+\tfrac12
                            +(\alpha+\tfrac12)D_uv^2\},\\
 H_{xy}={}&-uv\sqrt{Sr(1-r)}(m+D+\tfrac12).
\end{split}\tag{FM16}
\]
In particular $\bar H_x=\bar H_y=1/(2\nu)$ and
$\bar H_{xy}=0$. The last term of FM15 is retained whole;
it allows $\beta$ to be arbitrarily large or zero.

Define
\[
\begin{split}
 L_U={}&m\theta-(mj/\kappa)v(1+ra)
 -(D+\tfrac12)v^2(1+ha)+(b_1+\theta/2)(1+v^2)\\
       &+(D+\tfrac12)(1-r)vjS,\\
 Y_P={}&(1+v^2)\{-c_0+a_0(1+a)\}
       +S\{mv^2(1+ra)-(D+\tfrac12)(1-r)
                        +(\alpha+\tfrac12)(1+a)rv^2\},\\
 Y_Q={}&-m\theta a+(D+\tfrac12)(1+ha)
                                  -(\alpha+\tfrac12)(1+a),\qquad
 Y_C=jY_P+vY_Q.
\end{split}\tag{FM17}
\]
The expression $mj/\kappa=\kappa\mu j$ is bounded and
infinitesimal. The complete original normal terms are exactly
\[
\begin{split}
 L_x={}&-\kappa\mu v(1+ra)N
       +(D+\tfrac12)(1-r)vu\sqrt S\,G+L_UU,\\
 L_y={}&\sqrt{r(1-r)}\{Y_PG+Y_CQ\}.
\end{split}\tag{FM18}
\]
The actual source gives $L_U(x,b)=Y_C(x,b)=0$. Perform the same
finite monic divisions to write
\[
 L_U=(a-x)l_a+(v-b)l_v,\qquad
 Y_C=(a-x)y_a+(v-b)y_v.
 \tag{FM19}
\]
Then $L_UU=u l_aR_Z+l_vB_u$ and
$Y_CQ=\sqrt S\,y_aR_Z+y_vB$. All their coefficients are
bounded. Moreover $\bar l_v=-j$, while $Y_P$ and $y_v$ have
zero residue. This follows directly from the source-limit
identities
$L_U=(\ell-\nu v^2)/2$ at $a=0$, and
$Y_P=Y_Q=0$ at $a=0$.
Consequently $L_x+jB_u$ is an infinitesimal bounded combination
of $(N,G,B_u)$ and $L_y$ is an infinitesimal bounded combination
of $(N,G,B)$. This accounts for both unrestricted first-normal
terms, including all $\beta\theta$ terms and the independent
$r,1-r,\theta$ openings.

\paragraph{A strict residue for the full six-row Gram.}
Use FM8 with a fixed real $c=4$ to introduce a $cB_u^2$ pivot
at infinitesimal Gram cost. More explicitly subtract
\[
 (c/j^2)\delta^2N^2+(2c\kappa/j)\delta GB+2c\kappa B^2
 \tag{FM20}
\]
from the bounded part of FM15, add $cB_u^2$ to that part, and
retain the exact positive square $(\sqrt c\,\delta T/j)^2$.
Their sum is unchanged by FM8. The residue on
$(N_x,N_y,N,G,B,B_u)$ is therefore
\[
 \frac{N_x^2+N_y^2}{2\nu}+\frac h{2\nu^2}N^2
 +\frac{G^2}{2\nu}+(\ell-\nu)GB+\frac\nu2B^2
       +\theta NB_u-2jN_xB_u+4B_u^2.
 \tag{FM21}
\]
The $(G,B)$ block has determinant $\nu\ell>0$ and positive
diagonal. The remaining block has positive diagonal pivots
$1/(2\nu),1/(2\nu),h/(2\nu^2)$, and its final scalar Schur
pivot is
\[
 4-2\nu j^2-\frac{\theta^2\nu^2}{2h}>0.
 \tag{FM22}
\]
Indeed $j^2\le1/4$, $0<\nu\le1$, $h\ge\nu$ and
$0\le\theta\le1$ in residue make the two subtracted terms
at most $1/2$ each. All entries before taking residues are finite
sums of bounded separated factors. The surviving strict
bounded-Gram lemma therefore supplies a finite tensor SOS.
Restoring the single square from FM20 and the whole positive
$\beta D_vN_x^2$ row gives the entire FM5. No mixed matrix inverse
or mixed square root has been used.

\paragraph{Testing, sign, and boundary joins.}
The argument includes $u=0$ and every real state without a state
normalization or a sign condition. Negative $u$ is changed to
$-u$ together with $(Y,Z)\mapsto(-Y,-Z)$; negative $v$ is
changed to $-v$ together with $(X,Z)\mapsto(-X,-Z)$.
FM5 is exactly invariant under these transformations. On FM1,
$e=\kappa h/\nu$ is infinitesimal and $C,A$ are unbounded.
The surviving SE10--SE12 theorem covers every unmatched
homogeneous Bob residue, including Bob zero and infinity,
and invokes the unchanged adjacent-$C,A$ theorem for
noninfinitesimal Alice tests. The only remaining Bob residue
has finite nonzero $v$ and $\bar v^2=\bar\ell/\bar\nu$,
which is the matched chart just proved. Thus the union covers
all projective tests on FM1. Real and imaginary state parts
are treated separately for complex states. All divisions and
normalizations used in these joins are pure source or pure
testing operations; all original source predicates remain.

\paragraph{Dependencies and limits.}
The analytical inputs are the surviving native source formulas
P1 converted as FM2--FM3, the full original parent, the strict
separated bounded-Gram lemma in Section~\ref{component:p472f434581f3},
and Section~\ref{component:p8c30238a5f7d}
with its preserved EB and adjacent-$C,A$ dependencies.
The exact budget is independently printed in
Section~\ref{component:pe55ec4890781}. The source endpoint
and unit-$\kappa$ sectors require separate complete arguments;
this theorem makes no claim about them.


\section{\texorpdfstring{A source-free angular budget for the finite matched-Bob corner}{A source-free angular budget for the finite matched-Bob corner}}
\label{component:pe55ec4890781}
% Edited scoped mathematical packet, 22 October 2026.

\label{pe55ec4890781:sec:new-c144-matched-budget}

This exact auxiliary identity supplies the angular helper budget in FM.
Work in the separated tensor ring $K_c\otimes_{\mathbb R}K_t$.
Let $S,\kappa,\nu,\ell$ be positive source quantities, with
$\nu+\ell=1$, and put
\[
 x=\kappa S,\qquad j=\sqrt{\nu\ell},\qquad
 a_B=\sqrt{\ell/\nu}.
 \tag{MB1}
\]
For independent testing variables $u,v,W,Z$ define
\[
\begin{gathered}
 T=\kappa W,\quad U=uZ,\quad P=uW/\sqrt S,\quad Q=\sqrt S Z,\\
 N=T-jU,\quad G=P-jQ,\quad \delta=v-a_B,\\
 B=\delta Q,\quad B_u=\delta U.
\end{gathered}
 \tag{MB2}
\]
The monomial relation $TU=\kappa PQ$ is exact. It gives
\[
 \delta^2N^2+2j\kappa\delta GB+2j^2\kappa B^2-j^2B_u^2
       =(\delta T)^2.
 \tag{MB3}
\]
In particular, if $\nu,\ell$ are positive bounded units,
$\kappa$ is infinitesimal and $\delta$ has zero real residue,
then any fixed real positive diagonal on $B_u$ can be purchased
at an infinitesimal perturbation of the Gram matrix on $(N,G,B)$.
Indeed, for any fixed $c>0$ the exact SOS remainder is
\[
 \frac{c}{j^2}\delta^2N^2+
       \frac{2c\kappa}{j}\delta GB+2c\kappa B^2-cB_u^2
       =\frac{c}{j^2}(\delta T)^2.
 \tag{MB4}
\]
Every displayed coefficient is a finite sum of bounded separated
factors, with zero real residue except the constant $c$ in the
desired $B_u$ pivot. The only inverses are of positive source units.
The factor $\delta$ is a difference of a pure testing and a pure
source quantity; no inverse or radical of that difference is taken.
The right side has the displayed separated SOS root
$\sqrt c\,\delta T/j$.

Two exact radial conductors and one cross conductor are
\[
 \begin{split}
 (u^2-x)W/\sqrt S&=uG-\sqrt S N,\\
 (u^2-x)Z&=(\kappa\sqrt S/j)G-(u/j)N,\\
 (u^2-x)\delta Z^2&=(\kappa/j)GB-(1/j)NB_u.
 \end{split}
 \tag{MB5}
\]
These identities require no bound on the state coordinates or on
any ratio of source and testing infinitesimals. In the local matched
chart with $u$ infinitesimal their first two right sides have only
infinitesimal coefficients on $(N,G)$.


\section{\texorpdfstring{The full unit-$\kappa$, infinitesimal-$\nu$ source endpoint}{The full unit-kappa, infinitesimal-nu source endpoint}}
\label{component:p578719565e9d}

\label{p578719565e9d:sec:new-c144-nu-endpoint}

\paragraph{Scope.}
Retain every original physical C14.4 source predicate, strict gap,
nonnegative event coefficient, and original prior-complement intersection.
This theorem assumes
\[
 S=\sigma+L\in\operatorname{Inf}_{>0},\quad
 \kappa=x/S\text{ a positive bounded unit},\quad
 \nu=\sigma/S\in\operatorname{Inf}_{>0},\quad\ell=L/S=1-\nu.
 \tag{NU1}
\]
It proves a finite separated tensor SOS for the complete original
parent at every projective test and every real or complex state.
The sector with both $\kappa,\nu$ infinitesimal is not included.
No relative source/testing rate, state-ratio bound, or mixed inverse
is used. The original literal $D=0$ and $\beta=0$ faces are included.

\paragraph{Opening-safe source parameters and the two sphere masses.}
Use the exact FM2--FM3 native dictionary, and write
\[
\begin{gathered}
 q=\nu/\ell,\quad \zeta=\ell/\kappa=L/x,\quad
 e=\kappa h/\nu,\quad E=1-e,\\
 c_0=SC(1-r),\quad a_0=SAr,\quad d_0=D/S,\quad a_1=\alpha/S.
\end{gathered}
 \tag{NU2}
\]
In NU1, $q$ is infinitesimal and $\zeta$ is a positive bounded
unit. The exact source sphere SE3 gives
$e(B/S)=\kappa H_0+\nu\mathcal H b_0^2$.
Since $B/S\sim1+\kappa$, $H_0\sim1$, and
$\mathcal H\sim E^2$, its real residues satisfy
\[
 e=\frac\kappa{1+\kappa},\quad E=\frac1{1+\kappa},\quad
 h\in\operatorname{Inf},\quad\theta\sim1,\quad
 c_0=a_0=a_1=E/2.
 \tag{NU3}
\]
The two equalities for $e,E$ in this display are residue
equalities; the definition $e=\kappa h/\nu$ remains exact.
In particular $e,E,c_0,a_0,a_1$ are positive bounded units.

Define actual nonnegative source coefficients
\[
 d_s=\frac{D}{Sq^2},\qquad n_s=\frac\beta{q^2}.
 \tag{NU4}
\]
These are bounded even when either is zero. More explicitly the
SE3 sphere identities give
\[
 d_s=\frac{\ell^2\mathcal H}{2P_0J}d_0^{\,\mathrm{sph}\,2},
 \qquad
 n_s=\frac{\ell^2\mathcal H}{2\kappa P_0\theta}
                                   b_0^{\,\mathrm{sph}\,2},
 \qquad
 Jd_s+\kappa\theta n_s=\frac{\ell^2\mathcal H}{2P_0}.
 \tag{NU5}
\]
Here $(d_0^{\mathrm{sph}})^2+(b_0^{\mathrm{sph}})^2=1$;
the superscript distinguishes sphere coordinates from $D/S$.
Every denominator in NU5 is a positive unit. At least one of
$d_s,n_s$ is a positive unit. No sphere coordinate is inverted.

The sole possibly unbounded event coefficient needed locally is
$m$. Remove its pole by its actual contact-zero event, not by
dividing an error by $M$. With $R_s=1+xr$ set
\[
 p=\frac{\kappa\theta}{\ell R_s}\in(0,1),\qquad
 \mu_c=m(1-p)=
 \frac{\kappa e(1+x)(1+\nu Sr)}{2\ell P_0R_s}.
 \tag{NU6}
\]
Thus $\mu_c$ is a positive bounded unit and $m\ge\mu_c$.
Choose a fixed positive real $\epsilon<m$. On the reciprocal
Bob chart define $V=W+uY$ and $B_c=w(X+uZ)$. Retain the whole
positive packet $(m-\epsilon)(B_c-pV)^2$. The remaining $m$
quadratic is exactly
\[
 m_0B_c^2-2m_1B_cV+m_2V^2,\qquad
 m_0=\epsilon,\quad m_1=\mu_c+\epsilon p,\quad
 m_2=\mu_c(1+p)+\epsilon p^2.
 \tag{NU7}
\]
All three coefficients are bounded through $M\to0$ with arbitrary
rate. Their residues are
\[
 \mu_c=\frac{\kappa^2}{2(1+\kappa)},\quad p=\kappa,
 \qquad m_2=\kappa^2(\epsilon+\tfrac12)>0.
 \tag{NU8}
\]
The source $M>0$ also implies $0<\bar\kappa\le1$.

\paragraph{Local reciprocal chart and exact angular budgets.}
Assume $u,w$ are infinitesimal, including literal zero. Use
the homogeneous reciprocal parent $F_w/4$ from P4, with its
$m$ term replaced by NU7 and all other terms intact. Put
\[
\begin{gathered}
 X_e=X+ewW,\quad N_y=(Y-ruW)/\sqrt{Sr(1-r)},\quad
 a=u^2,\quad c=w^2,\quad \delta=c-q,\\
 P=uW/\sqrt S,\quad Q=\sqrt S\,wZ,\quad
 N=W-\zeta uwZ,\quad G=P-\ell Q,\\
 B=\sqrt S\,\delta Z,\quad B_u=u\delta Z,\quad
 R=w(a-x)Z=\frac{\sqrt S}{\zeta}G-\frac u\zeta N.
\end{gathered}\tag{NU9}
\]
Every row is a finite separated expression. No state coordinate
is divided by another. Since $W(uwZ)=PQ$, two exact identities are
\[
\begin{split}
 \delta^2N^2+2\zeta\delta wGB+2\zeta\ell w^2B^2
                  -\zeta^2(wB_u)^2&=(\delta W)^2,\\
 \delta^2G^2+2\ell\delta wNB_u+2\ell\zeta w^2B_u^2
                  -\ell^2(wB)^2&=(\delta P)^2.
\end{split}\tag{NU10}
\]
Consequently, from either positive sphere pivot in NU5, any fixed
positive diagonal budgets on both $wB$ and $wB_u$ can be bought
at infinitesimal bounded-Gram cost. To see this without a hidden
division, if $d_s$ is a unit, replace $B$ by
$\sqrt{d_s}B/\sqrt{d_s}$ in the first identity, divide by
$\zeta^2$, and use $wB=w(\sqrt{d_s}B)/\sqrt{d_s}$ for the
other budget. All cost coefficients are infinitesimal. If $n_s$
is a unit, use the second identity divided by $\ell^2$ and
$wB_u=w(\sqrt{n_s}B_u)/\sqrt{n_s}$. The factor
$\zeta w^2$ is infinitesimal because $\zeta$ is a unit in
NU1. This is the precise restriction preventing a silent extension
to infinitesimal $\kappa$.

\paragraph{The entire scalar graph before normal restoration.}
Write $D=Sd_0$, $\alpha=Sa_1$, and define bounded coefficient
polynomials in the independent testing squares $a,c$:
\[
\begin{split}
 A_P={}&(1+c)\{c_0(1-r)+a_0r(1+a)\}\\
 &+S\{(D+\tfrac12)(1-r)^2c
                  +(\alpha+\tfrac12)(1+a)r^2\},\\
 A_N={}&m_2(1+ra)^2+2m_1ec(1+ra)+m_0e^2c^2
                  +(\beta+\tfrac12)e^2c(1+c)\\
 &-(D+\tfrac12)e\{2(1-r)-e\}ac,\\
 K={}&m_1(1+ra)+m_0ec+(D+\tfrac12)(1-r-e)\\
 &+(\alpha+\tfrac12)(1+a)r+(\beta+\tfrac12)e(1+c),\\
 F_Z={}&D+\alpha c+\beta a+(m_0+\alpha+\beta+\tfrac12)ac,\\
 J_1={}&K-\ell A_P,\qquad L_1=J_1-\zeta A_N,\\
 H={}&F_Z-\ell^2ScA_P+ac(\zeta^2A_N-2\zeta J_1).
\end{split}\tag{NU11}
\]
For $X=-ewW$, $Y=ruW$, the exact graph scalar is
\[
 F_0=A_NN^2+A_PG^2-2uwZL_1N+HZ^2.
 \tag{NU12}
\]
This graph is used to compute coefficients; the unrestricted
normals are restored explicitly below. The actual source contact
has $a=x,c=q,W=\zeta uwZ$, so $N=G=X_e=N_y=0$.
The retained row $B_c-pV$ also vanishes there. Direct source
substitution in NU11 yields
\[
 L_1(x,q)=0,\qquad H(x,q)=H_a(x,q)=H_c(x,q)=0,
 \qquad H(a,0)=D+\beta a.
 \tag{NU13}
\]
These are exact rational source identities; no independent
testing angle is identified with a source angle in the certificate.

All divisions in the next display are finite monic polynomial
divisions at the bounded pure-source contacts $x,q$. They define
bounded coefficient polynomials $L_a,L_c,H_{20},H_{11},H_{02}$:
\[
\begin{split}
 L_1(a,c)&=(a-x)L_a(a,c)+(c-q)L_c(c),\\
 H(a,c)&=H(x,c)+(a-x)H_a(x,c)+c(a-x)^2H_{20}(a,c),\\
 H_a(x,c)&=(c-q)\{-\beta/q+cH_{11}(c)\},\\
 H(x,c)&=S(c-q)^2\{d_s+\kappa n_s+cH_{02}(c)\}.
\end{split}\tag{NU14}
\]
Here the factor $c$ in the quadratic $a$ remainder follows from
$H(a,0)=D+\beta a$. The two constant terms follow by putting
$c=0$: $H_a(x,0)=\beta$ and $H(x,0)=D+\beta x$.
In particular the displayed apparent divisions by $q$ are
bounded: $\beta/q=n_sq$. Finally $H(x,c)/S$ is a bounded
coefficient polynomial directly from NU11, because
$x=\kappa S$, $D/S=d_0$, and $\alpha/S=a_1$ are bounded.
Thus dividing it by the monic $(c-q)^2$ introduces no unbounded
coefficient. This verifies all asserted coefficient bounds without
using relative orders of $S,q,u,w$.

The useful mixed identity is
\[
 c(a-x)(c-q)Z^2=\frac1\zeta\{G(wB)-N(wB_u)\}.
 \tag{NU15}
\]
Substituting NU14 into NU12 gives the complete exact scalar
allocation
\[
\begin{split}
 F_0={}&A_NN^2+A_PG^2+d_sB^2+n_sB_u^2
       +H_{02}(wB)^2+H_{20}R^2\\
 &+\frac{H_{11}-n_s}{\zeta}\{G(wB)-N(wB_u)\}
       -2uL_aNR-2L_cN(wB_u).
\end{split}\tag{NU16}
\]
The residues of $A_N,A_P$ are respectively
$\kappa^2(\epsilon+1/2)$ and $E/2$, both positive.
The row $R$ has infinitesimal coefficients on $(N,G)$ by NU9.
Every remaining signed term is a bounded quadratic term on
$N,G,wB,wB_u$. The angular rows are retained as the complete
nonnegative squares $(\sqrt{d_s}B)^2$ and
$(\sqrt{n_s}B_u)^2$; neither is prematurely discarded.

\paragraph{Both complete state normals.}
Put $s_y=\sqrt{Sr(1-r)}$ and $r_y=\sqrt{r(1-r)}$.
Define
\[
\begin{split}
 A_x={}&-m_0ec-m_1(1+ra)+(D+\tfrac12)a(1-r-e)
                          -e(\beta+\tfrac12)(1+c),\\
 T_x={}&\zeta cA_x+\beta-D+c(m_0+\beta+\tfrac12),\\
 Y_P={}&(1+c)\{-c_0+a_0(1+a)\}\\
 &+S\{m_2(1+ra)+m_1ec-(D+\tfrac12)c(1-r-e)
                             +(\alpha+\tfrac12)(1+a)r\},\\
 Y_Q={}&-m_1a+(D+\tfrac12)-(\alpha+\tfrac12)(1+a),\qquad
 Y_C=\ell Y_P+Y_Q.
\end{split}\tag{NU17}
\]
The entire capped quarter-parent, before any allocation, is
\[
 F_0+H_xX_e^2+H_yN_y^2+2H_{xy}X_eN_y+2L_xX_e+2L_yN_y,
 \tag{NU18}
\]
where
\[
\begin{split}
 H_x={}&m_0c+(D+\tfrac12)a+(\beta+\tfrac12)(1+c),\\
 H_y={}&rc_0(1+c)+(1-r)a_0(1+a)(1+c)\\
 &+Sr(1-r)\{m_2a+(D+\tfrac12)c+(\alpha+\tfrac12)(1+a)\},\\
 H_{xy}={}&-uws_y(m_1+D+\tfrac12),\\
 L_x={}&wA_xN+uZT_x,\qquad
 L_y=r_y\{Y_PG+Y_CQ\}.
\end{split}\tag{NU19}
\]
These formulas retain every quadratic state coefficient and both
first normals. Their real hard residues are
$H_x=1/2$, $H_y=E/2$, $H_{xy}=0$.

Exact source substitution gives $T_x(x,q)=Y_C(x,q)=0$.
Moreover $T_x(a,0)=\beta-D$, and its $a$ dependence has a
factor $c$. Thus bounded monic divisions give
\[
\begin{split}
 T_x(a,c)&=c(a-x)X_a+(c-q)\{(D-\beta)/q+cX_c(c)\},\\
 Y_C(a,c)&=(a-x)Y_a(a,c)+(c-q)Y_c(c).
\end{split}\tag{NU20}
\]
No coefficient here needs an inverse of a sphere coordinate.
Using NU4 and NU9, the full normal conductors become
\[
\begin{split}
 L_x={}&wA_xN+uwX_aR
       +d_s\sqrt S\,q u B-n_sqB_u+wX_c(wB_u),\\
 L_y={}&r_y\{Y_PG+\sqrt S\,Y_aR+Y_c(wB)\}.
\end{split}\tag{NU21}
\]
The two potentially weak terms are face-safe:
$d_s\sqrt S\,quB=(\sqrt{d_s}\sqrt S\,qu)(\sqrt{d_s}B)$
and $n_sqB_u=(\sqrt{n_s}q)(\sqrt{n_s}B_u)$.
Their coefficients are infinitesimal, even when either sphere
mass is zero. Also $Y_P$ has zero residue by NU3 and NU17.
All other terms in NU21 have infinitesimal coefficients on
$N,G,R$ or bounded coefficients on the two angular helper rows.

\paragraph{Strict complete Gram and restoration.}
Begin with the six rows
\[
 (N,G,X_e,N_y,\sqrt{d_s}B,\sqrt{n_s}B_u).
 \tag{NU22}
\]
Their real diagonal residue in NU16--NU21 is
\[
 \operatorname{diag}\bigl(
 \kappa^2(\epsilon+\tfrac12),E/2,1/2,E/2,1,1\bigr)>0.
 \tag{NU23}
\]
Every cross coefficient among these rows is infinitesimal after
the displayed $R$ substitution. Literal zero rows are harmless:
the coefficient matrix may be strict even when the chosen rows
are linearly dependent or vanish.

Use NU10 to add sufficiently large fixed real pivots for the
two further rows $wB,wB_u$, at infinitesimal cost to NU22.
The signed scalar terms in NU16 and the remaining normal terms
in NU21 have bounded coefficients on these helper rows.
A sufficiently large fixed pair of diagonals therefore gives
a strictly positive real eight-row Gram, by an ordinary real
Schur complement. All entries are finite sums of bounded separated
factors. The separated bounded-Gram lemma supplies a finite SOS.
Restore the exact positive square remainders in NU10 and the
whole excess $(m-\epsilon)(B_c-pV)^2$ in NU7. This gives the
entire original parent, with both unrestricted state normals.

\paragraph{All projective testing joins on NU1.}
Noninfinitesimal Alice tests use CE1--CE10 because $C,A$ are
unbounded. For infinitesimal $u$, the local proof above covers
every infinitesimal reciprocal $w$, including zero. It remains
to check finite Bob coordinates $v$.

For $v$ infinitesimal use the exact FU22--FU23 frame; its source
inputs here are only $m$ strong, $c_0,a_0$ with common positive
residue $E/2$, and $a_1$ positive, all proved in NU3--NU6.
For a finite positive unit $v$ use the complete FU6--FU12
bounded parent and real Schur calculation, capping the original
$m$ event at a finite real only if $m$ is unbounded. In the
present source residue $\beta=d_0=0$, $a_1=E/2$, and, if
$\kappa<1$,
\[
 m=\frac{\kappa^2}{2(1-\kappa^2)},\qquad E=\frac1{1+\kappa}.
 \tag{NU24}
\]
The determinant polynomial FU14 becomes exactly
\[
 \mathcal P(y,c)=p_{21}c(y-\kappa)^2+p_{10}y>0,
 \qquad p_{21}=E(m+\tfrac12),\quad
 p_{10}=(m+\tfrac12)E^2,\quad y,c>0.
 \tag{NU25}
\]
It gives the strict four-term circuit inequality FU18. If
$\bar\kappa=1$, then $m$ is unbounded. The normalized limit is
\[
 \mathcal P(y,c)/m\longrightarrow
              \tfrac12c(y-1)^2+\tfrac14y>0.
 \tag{NU26}
\]
The Schur coefficients have finite positive limits, so a
sufficiently large finite cap preserves strictness, exactly as
in FU19--FU20. The whole positive original event excess is
restored. These finite-Bob arguments use no lower bound on
$\nu$, and do not import the source scope of the FU theorem.

The exact sign transformations FU31 cover negative angles;
all displayed frames include literal axes and projective
endpoints. Complex states split into real and imaginary parts.
No chart normalizes a state coordinate or bounds a state ratio.
This completes the all-tests theorem on precisely NU1.


\section{\texorpdfstring{The doubly infinitesimal endpoint with bounded $\beta/q^2$}{The doubly infinitesimal endpoint with bounded beta/q2}}
\label{component:p1de5b12d1927}
% Edited scoped mathematical packet, 22 October 2026.

\label{p1de5b12d1927:sec:new-c144-small-nu-bounded-beta}

Keep the original C14.4 source, its full native contact, and all
original coefficient signs.  The common doubly small endpoint has
\[
 S=\sigma+L\in\operatorname{Inf}_{>0},\quad\kappa=x/S\in\operatorname{Inf}_{>0},
 \quad\nu=\sigma/S\in\operatorname{Inf}_{>0},\quad
 \ell=L/S=1-\nu,\quad q=\nu/\ell,\quad n_s=\beta/q^2.
                                                               \tag{SN1}
\]
Here $n_s$ may be bounded or unbounded.  The formulas below are
exact identities; the complete BN and UN proofs supply their distinct
normalized bounds on the bounded and unbounded $\tau$ domains.
Set
\[
 c_0=SC(1-r),\quad a_0=SAr,\quad d_s=D/(Sq^2),\quad
 a_1=\alpha/S,\quad\mu=m/\kappa^2,\quad
 e=\kappa h/\nu,\quad\tau=h/q.                         \tag{SN2}
\]
The actual source-sphere identities are
\[
 d_s=\frac{\ell^2\mathcal H d_0^2}{2(P/S)J},\qquad
 n_s=\frac{\ell^2\mathcal H b_0^2}{2\kappa(P/S)\theta},\qquad
 Jd_s+\kappa\theta n_s=\frac{\ell^2\mathcal H}{2(P/S)}.   \tag{SN3}
\]
The source symbols in SN3 have the definitions in FM2--FM3 and
SE1--SE3.  In particular $e=\kappa\tau/\ell$ exactly.
\paragraph{The complete reciprocal parent.}
For infinitesimal Alice coordinate $u$ and reciprocal Bob coordinate
$w$, including either literal zero, use $D_u=1+u^2,D_w=1+w^2$ and
\[
\begin{split}
 F_w/4={}&m(W-wX+uY-uwZ)^2+CD_w(uW-Y)^2+AD_uD_wY^2\\
 &+(D+\tfrac12)(uwW+uX-wY-Z)^2
 +(\alpha+\tfrac12)D_u(Y-wZ)^2\\
 &+(\beta+\tfrac12)D_w(X+uZ)^2-\tfrac12D_uD_wZ^2.
\end{split}\tag{SN5}
\]
Keep both unrestricted state normals
\[
 X_e=X+ewW,\qquad N_y=(Y-ruW)/\sqrt{Sr(1-r)}.
 \tag{SN6}
\]
Every inverse and radical in these expressions is pure-source or
pure-testing. Write $a=u^2,c=w^2,x=\kappa S,\delta=c-q$ and set
\[
\begin{gathered}
 N=\kappa W-\ell uwZ,\quad G=uW/\sqrt S-\ell\sqrt S\,wZ,\\
 B=\sqrt S\,\delta Z,\quad B_u=u\delta Z,\quad
 R=w(a-x)Z=\frac{\kappa\sqrt S\,G-uN}{\ell}.
\end{gathered}\tag{SN7}
\]
The last equality is the exact homogeneous radial conductor.

\paragraph{Two free helper pivots.}
The exact polynomial identity
\[
 \delta^2N^2+2\kappa\ell\delta wGB
       +2\kappa\ell^2w^2B^2-\ell^2(wB_u)^2
                  =(\kappa\delta W)^2
 \tag{SN8}
\]

is an exact SOS identity.  Whenever a current strict frame controls
$N,G,B$, it buys a prescribed real pivot on $wB_u$ at
infinitesimal Gram cost.  BN and UN give the required ordered helper
constructions even when the original $d_sB^2$ pivot is absent.

\paragraph{Exact scalar dictionary and the $1/\kappa$ cancellations.}
Put $d=D+1/2$ and define the following polynomials in $a,c$:
\[
\begin{split}
 A_P={}&(1+c)\{c_0(1-r)+a_0r(1+a)\}
       +S\{d(1-r-e)^2c+(\alpha+\tfrac12)(1+a)r^2\},\\
 A_N={}&m(1+ra+ec)^2+(\beta+\tfrac12)e^2c(1+c),\\
 K={}&m(1+ra+ec)+d(1-r-e)
       +(\alpha+\tfrac12)(1+a)r+(\beta+\tfrac12)e(1+c),\\
 F_Z={}&D+\alpha c+\beta a+(m+\alpha+\beta+\tfrac12)ac,\\
 J_1={}&K-\ell A_P,\quad L_1=J_1-(\ell/\kappa)A_N,\\
 H={}&F_Z-\ell^2ScA_P
                 +ac\{(\ell/\kappa)^2A_N-2(\ell/\kappa)J_1\}.
\end{split}\tag{SN9}
\]
These are the exact graph coefficients at $X=-ewW,Y=ruW$;
the negative $dac,e\{2(1-r)-e\}$ contribution has been placed
in $A_P$, not left in $A_N$. The displayed combination is used intact before taking any
normalized bounds.

The actual source satisfies
$L_1(x,q)=H(x,q)=H_a(x,q)=H_c(x,q)=0$ and
$H(a,0)=D+\beta a$. Finite monic division therefore gives
\[
\begin{gathered}
 L_1=(a-x)L_a+\delta L_c,\\
 H=S\delta^2(d_s+\kappa n_s+cH_{02})
        +(a-x)\delta(-qn_s+cH_{11})+c(a-x)^2H_{20}.
\end{gathered}\tag{SN10}
\]
The quotients mean literal polynomial differences at $a=x,c=q$,
not division by a mixed tensor. For example $L_a$ is
$(L_1(a,c)-L_1(x,c))/(a-x)$, and $L_c$ is the quotient of
$L_1(x,c)-L_1(x,q)$ by $c-q$.

Individually $H_{20},H_{11},L_c/\kappa$ need not be bounded.
The complete scalar uses the following combined coefficients:
\[
 C_R=H_{20}+2\ell L_a/\kappa,\quad
 C_G=\kappa(H_{11}-n_s)/\ell,\quad
 C_N=(n_s-H_{11})/\ell-2L_c/\kappa.
 \tag{SN11}
\]
Substitution of $e=\kappa\tau/\ell$ gives the exact cancellations
\[
\begin{split}
 C_R={}&-\ell\mu r(a\ell r+2c\kappa\tau+2\ell),\\
 C_N={}&n_s/\ell-Sa_0\ell r
       +(\beta+m+\tfrac12)\tau^2(c/\ell+1/\ell^2)
                   +2\kappa\mu\tau-m\tau^2/\ell,\\
 H_{02}={}&\kappa\tau^2(\beta+m+\tfrac12).
\end{split}\tag{SN12}
\]
For further pole cancellation, write
$A_P=P_0(a)+cP_1(a)$. Direct differentiation in SN9 gives
\[
\begin{split}
 \kappa H_{11}={}&2\ell^2P_1(x)+\ell^2x a_0r
 +2\ell^2\kappa\mu e(1+2rx)
 +\frac{\ell^2(\beta+1/2)e^2}{\kappa}
 -2\ell(m+\beta+\tfrac12)e\\
 &+(c+q)\frac{\ell^2(m+\beta+1/2)e^2}{\kappa}.
\end{split}\tag{SN13}
\]
These are identities of the complete coefficient combinations.
Their respective normalized estimates are proved in BN and UN.

The entire scalar graph is now
\[
\begin{split}
 F_0={}&(A_N/\kappa^2)N^2+A_PG^2+d_sB^2+n_sB_u^2
       +H_{02}(wB)^2+C_RR^2-2\sqrt S\,L_aGR\\
 &+C_GG(wB)+C_NN(wB_u).
\end{split}\tag{SN14}
\]
For verification the mixed identity is
$c(a-x)\delta Z^2=\{\kappa G(wB)-N(wB_u)\}/\ell$.
Combining this with $uN=\kappa\sqrt S\,G-\ell R$ proves
SN14 directly from SN9--SN10.

\paragraph{Both complete state normals.}
Define
\[
\begin{split}
 A_X={}&-m(1+ra+ec)+da(1-r-e)-e(\beta+\tfrac12)(1+c),\\
 T_X={}&(\ell/\kappa)cA_X+\beta-D+c(m+\beta+\tfrac12),\\
 Y_P={}&(1+c)\{-c_0+a_0(1+a)\}
 +S\{m(1+ra+ec)-dc(1-r-e)
                    +(\alpha+\tfrac12)(1+a)r\},\\
 Y_C={}&\ell Y_P-ma+d-(\alpha+\tfrac12)(1+a).
\end{split}\tag{SN15}
\]
Their actual contact values satisfy $T_X(x,q)=Y_C(x,q)=0$.
Use finite polynomial quotients
\[
 T_X=c(a-x)X_a+\delta\{(D-\beta)/q+cX_c\},\qquad
 Y_C=(a-x)Y_a+\delta Y_c.
 \tag{SN16}
\]
Although $X_a,A_X/\kappa$ separately may be unbounded, their
complete conductor has
\[
\begin{gathered}
 \widehat A_X=A_X/\kappa-aX_a/\ell
 =-\kappa\mu-\frac\tau\ell
                  \{(\beta+\tfrac12)(1+c)+mc\},\\
 \widehat X_a=\kappa X_a=\ell\{d(1-r-e)-mr\},\qquad
 X_c=-\tau(m+\beta+\tfrac12).
\end{gathered}\tag{SN17}
\]
The exact quarter-parent is $F_w/4=F_0+H_XX_e^2+H_YN_y^2+
2H_{XY}X_eN_y+2L_XX_e+2L_YN_y$, with
\[
\begin{split}
 H_X={}&mc+da+(\beta+\tfrac12)(1+c),\\
 H_Y={}&rc_0(1+c)+(1-r)a_0(1+a)(1+c)
       +Sr(1-r)\{ma+dc+(\alpha+\tfrac12)(1+a)\},\\
 H_{XY}={}&-uw\sqrt{Sr(1-r)}(m+d),\\
 L_X={}&w\widehat A_XN+
        \frac{uw\sqrt S}{\ell}\widehat X_aG
        +d_s\sqrt S\,quB-n_sqB_u+wX_c(wB_u),\\
 L_Y={}&\sqrt{r(1-r)}\{Y_PG+\sqrt S\,Y_aR+Y_c(wB)\}.
\end{split}\tag{SN18}
\]

Every formula retains the original $r$ and $1-r$ openings and both
complete state normals.  In particular, the term $n_sqB_u$ may be
written $q\sqrt{n_s}(\sqrt{n_s}B_u)$ when that row is used;
literal $n_s=0$ needs no division.  BN and UN allocate these actual
rows, with their own weights, instead of treating separate uncanceled
coefficients as bounded.

\paragraph{Explicit finite divisions.}
All quotient symbols above can be defined directly in the polynomial
ring in $a,c$ over the source field.  In addition to $L_a,L_c$
specified after SN10, use
\[
\begin{aligned}
 H_{20}&=\frac{H(a,c)-H(x,c)-(a-x)H_a(x,c)}{c(a-x)^2},\\
 H_{11}&=\frac{H_a(x,c)/(c-q)+qn_s}{c},\\
 H_{02}&=\frac{H(x,c)/(S(c-q)^2)-d_s-\kappa n_s}{c},\\
 X_a&=\frac{T_X(a,c)-T_X(x,c)}{c(a-x)},\qquad
 X_c=\frac{T_X(x,c)/(c-q)-(D-\beta)/q}{c},\\
 Y_a&=\frac{Y_C(a,c)-Y_C(x,c)}{a-x},\qquad
 Y_c=\frac{Y_C(x,c)-Y_C(x,q)}{c-q}.
\end{aligned}
\]
These displays specify polynomial division, not evaluation of an
inverse of $a-x$, $c-q$ or $c$.  Divisibility follows from
$H(x,q)=H_a(x,q)=H_c(x,q)=0$, $H(a,0)=D+\beta a$,
$T_X(x,q)=Y_C(x,q)=0$, and the explicitly affine dependence of
$T_X$ on $a$.  Thus all exact coefficients are defined without an
external program.  Pure-source divisions such as $1/S$ or $1/q$
are retained until their cancellations in the printed combinations.


\section{\texorpdfstring{The unbounded-$h/q$ part of the doubly small endpoint}{The unbounded-h/q part of the doubly small endpoint}}
\label{component:peaa56c1a3c7a}

\label{peaa56c1a3c7a:sec:new-c144-unbounded-h-ratio}

\paragraph{Exact scope.}
Retain the full original entered C14.4 source, all its strict gaps and
prior-complement intersections. Suppose $S,\kappa,\nu$ are positive
infinitesimals, $\ell=1-\nu$, $q=\nu/\ell$, and
\[
 \tau=h/q\ \hbox{is unbounded},\qquad
 \theta=1-h>0,\quad e=\kappa h/\nu\in\operatorname{Inf},\quad
 \zeta=\ell/\kappa.
 \tag{UN1}
\]
This section proves finite separated tensor SOS for the complete
parent at every projective test and every state on UN1. The opening
$\theta\in\operatorname{Inf}$ is included. There is no state-ratio
restriction or source/testing relative-rate assumption.

\paragraph{Normalized source estimates and a canceled radial mass.}
Keep $c_0=SC(1-r)$, $a_0=SAr$, $d_s=D/(Sq^2)$,
$a_1=\alpha/S$, $P_0=P/S$, and the exact SE2--SE3 source sphere.
Define
\[
 d=D+\tfrac12,\quad
 n=\beta\theta^2+2Dh+\tfrac12h^2>0,\quad
 z_n=\kappa n/q^2,\quad A_{00}=m\zeta^2-n/q.
 \tag{UN2}
\]
Here $n$ is a scalar coefficient, not an omitted event weight.
It is bounded even when $\beta$ is unbounded. The exact source gives
\[
\begin{gathered}
 m\zeta^2=\frac{\tau(1+x)(1+\sigma r)}{2MP_0},\qquad
 m\zeta\in\operatorname{Inf},\quad D,\alpha,m\in\operatorname{Inf},\quad
 c_0,a_0,a_1\sim\tfrac12,\\
 e\tau=\frac{\ell\{eH_0+h\mathcal H b_0^2\}}{B/S}\ \hbox{is bounded},\\
 z_n=\frac{\ell^2\mathcal H\theta b_0^2}{2P_0}
 +\frac{\ell^2(eH_0+h\mathcal H b_0^2)}{2(B/S)}
 +2\kappa S d_s h.
\end{gathered}\tag{UN3}
\]
Consequently $d_s,z_n$ are bounded and
$d_s+z_n\sim1/2$; at least one is a positive bounded unit. This
includes literal $D=0$. Literal $\beta=0$ has no points on UN1,
because $\beta\ge0$ has the exact boundary equation
$\tau(B/S)=\ell H_0$ when $\beta=0$.

For completeness, the following finite source division proves the
bound on $A_{00}$ without dividing an error by $\theta$ or $q$.
Put $U=(1+x)(1+\sigma r)$ and
$M_*=\ell(1+xr)-\kappa\theta=\ell M$, and form the bounded
source polynomial
\[
 T=\ell JU+JM_*P_0-JM_*\kappa\theta(J-\sigma)-2M_*U.
\]
In this expression $P_0,J,U,M_*$ are their exact native polynomial
formulas in $\kappa,\nu,S,h,r$, with $\ell=1-\nu$. Thus
\[
\begin{gathered}
 T=\kappa T_\kappa+\ell\nu S T_\nu,\qquad
 T_\nu=r-1-\nu-\nu S(1-r)(r+\ell),\\
 T_\kappa=\{T-T|_{\kappa=0}\}/\kappa
             \quad\hbox{is a bounded source polynomial},\\
 A_{00}=\frac{\theta\ell H_0}{2P_0}
            +\frac{eT_\kappa+h\ell S T_\nu}{2JMP_0}.
\end{gathered}\tag{UN4}
\]
The quotient is a finite monic polynomial division in a pure source
variable. All denominator factors are positive bounded units.
Hence $A_{00}$ is bounded and $A_{00}-\theta/2$ is infinitesimal.
Moreover $n/q=m\zeta^2-A_{00}\sim\tau/2$ is positive unbounded.
If $\theta$ is infinitesimal, then $n\sim1/2$ and
$\kappa/q=e\ell/h$ is infinitesimal. These facts supply the two
source subcases used below.

\paragraph{The actual circuit and three reserved positive rows.}
Use the complete original reciprocal quarter-parent SN5. For
infinitesimal $u,w\ge0$, including zero, retain the unrestricted
normals
\[
 N_x=X+huZ,\qquad N_y=(Y-ruW)/\sqrt{Sr(1-r)}.
 \tag{UN5}
\]
Put $a=u^2,c=w^2,x=\kappa S,\delta=c-q$ and
\[
\begin{gathered}
 N_0=W/\zeta-uwZ,\quad G=uW/\sqrt S-\ell\sqrt S\,wZ,\\
 B=\sqrt S\delta Z,\quad B_u=u\delta Z,\quad
 R=w(a-x)Z=\sqrt S\,G/\zeta-uN_0,\quad
 G_h=h(a-xc/q)Z.
\end{gathered}\tag{UN6}
\]
Define $A_1,A_2,A_3,A_4$ by BN4 with $\beta$ replaced by the
positive $n$, and set
\[
 \Psi_n=\sum_iA_i^2-4nuwWZ/(\zeta q),\qquad
 E_1=A_1-A_2,\quad E_2=A_3-A_4,\quad E_3=A_1-A_3.
 \tag{UN7}
\]
The explicit separated AM--GM square decompositions in BN6 give
\[
             \Psi_n-\tfrac14(E_1^2+E_2^2+E_3^2)\in\Sigma T^2.
 \tag{UN8}
\]
To verify the simultaneous reserve, average BN6 with the pairing
$\Psi_n-(A_1-A_3)^2-(A_2-A_4)^2\in\Sigma T^2$.
For the latter pairing the products $A_1A_3,A_2A_4$ are nonnegative;
if $WZ<0$, the additional signed cross term is a positive separated
monomial. Thus UN8 holds for both signs of $WZ$, as well as its
zero face. No radical of a general mixed sum is used.

\paragraph{Exact scalar subtraction and radial cancellation.}
The complete scalar graph at $N_x=N_y=0$ is
\[
 F_0=m(1+ra)^2W^2+A_PP^2+F_ZZ^2-2KuwWZ,
 \quad P=uW/\sqrt S,
 \tag{UN9}
\]
where
\[
\begin{split}
 A_P={}&(1+c)\{c_0(1-r)+a_0r(1+a)\}
       +S\{d(1-r)^2c+(\alpha+\tfrac12)(1+a)r^2\},\\
 K={}&m\theta(1+ra)+d(1-r)(1+ha)
                         +(\alpha+\tfrac12)(1+a)r,\\
 F_Z={}&D+\alpha c+na+
          \{m\theta^2+\alpha+(\beta+\tfrac12)\theta^2\}ac
                                                 +dh^2a^2.
\end{split}\tag{UN10}
\]
Subtract $\Psi_n$, put
\[
\begin{gathered}
 A=m\zeta^2(1+ra)^2-n/q,\quad A_P'=A_P-z_nc/\ell^2,\\
 K'=K-2n/(\zeta q),\quad J'=K'-\ell A_P',\quad
 L'=J'-A/\zeta,\\
 H'=F_Z-na-nxc/q-\ell^2ScA_P'+ac(A-2\zeta J').
\end{gathered}\tag{UN11}
\]
The exact residual graph is $A N_0^2+A_P'G^2-
2\zeta uwZL'N_0+H'Z^2$. At the actual source contact
$L'(x,q)=H'(x,q)=H'_a(x,q)=H'_c(x,q)=0$.
Also $H'(a,0)=D+dh^2a^2$. Finite monic contact division gives
\[
\begin{gathered}
 L'=(a-x)L_a+\delta L_c,\\
 H'=(D+dh^2x^2)\delta^2/q^2
       +(a-x)\delta(-2dh^2x/q+cH_{11})
       +(a-x)^2(dh^2+cH_{20}).
\end{gathered}\tag{UN12}
\]
The otherwise possible third-$c$ quotient is exactly zero.
Define $P_1(a)$ to be the coefficient of $c$ in $A_P$. The complete
combined coefficients have short bounded formulas:
\[
\begin{split}
 L_a={}&m\theta r+dh(1-r)+(\alpha+\tfrac12)r
       -\ell\{(1+c)a_0r+S(\alpha+\tfrac12)r^2\}\\
       &-m\zeta r\{2+r(a+x)\},\\
 C_R={}&H_{20}+2\zeta L_a=-m\zeta^2r(2+ra),\\
 C_G={}&H_{11}/\zeta=2\ell P_1(x)+\ell x a_0r-2z_n/\ell,\\
 C_N={}&-H_{11}-2\zeta L_c=-Sa_0\ell^2r.
\end{split}\tag{UN13}
\]
Only $C_R$ may be unbounded separately. Its entire contribution
must be combined before estimating: $A+aC_R=A_{00}$ exactly.
Using $R=\sqrt S\,G/\zeta-uN_0$ therefore gives
\[
\begin{split}
 F_0={}&\Psi_n+A_{00}N_0^2+\widetilde A_PG^2
             +2u\sqrt S\,(L_a-C_R/\zeta)N_0G\\
 &+d_sB^2+dG_h^2+C_GG(wB)+C_NN_0(wB_u),\\
 \widetilde A_P={}&A_P'+SC_R/\zeta^2-2SL_a/\zeta\sim\tfrac12.
\end{split}\tag{UN14}
\]
Every coefficient in UN14 is bounded after this combination:
$C_R/\zeta=-m\zeta r(2+ra)$ is infinitesimal and
$C_R/\zeta^2=-mr(2+ra)$ is bounded. The $N_0G$ coefficient
is infinitesimal. No estimate of an unbounded isolated radial
coefficient is used.

\paragraph{Complete $X$ conductor and its exact positive completion.}
The full $N_x$ hard coefficient is
$H_X=mc+da+(\beta+1/2)(1+c)$. Its complete half-linear
coefficient, after using the actual contact identity, is
\[
\begin{split}
 L_X={}&-m\zeta wN_0+d\tau a(wN_0+B_u)
        +uw\sqrt S\{d(1-r-e)-mr\}G+t_0B_u,\\
 t_0={}&\frac{D+h/2-\beta\theta}{q}
       =D/q+\frac{\ell H_0-\kappa\tau(J-\sigma)}{2P_0}
          \sim\tfrac12.
\end{split}\tag{UN15}
\]
The exact source identity in the second line cancels the potential
$1/q$ pole. Here $D/q=S q\,d_s$ is infinitesimal. The single
unbounded $\tau$ term is removed with the actual square $dG_h^2$,
not estimated. Namely
\[
 G_h=-\tau u(wN_0+B_u)+ew\sqrt S\,G,\quad
 N_x+\tau(wN_0+B_u)=X+ewW,
\]
\[
\begin{split}
 &dG_h^2+daN_x^2+2d\tau a(wN_0+B_u)N_x\\
 &\quad=da(X+ewW)^2
       -2de\tau uw\sqrt S\,(wN_0+B_u)G+de^2Sw^2G^2.
\end{split}\tag{UN16}
\]
The remainder has bounded infinitesimal coefficients on $N_0,G,wB_u$
because $e\tau$ is bounded. Retain the entire positive square on the
right and also the whole packet $\beta(1+c)N_x^2$. The remaining
$N_x$ diagonal is $mc+(1+c)/2\sim1/2$. The remaining $L_X$ is
UN15 with its $d\tau a$ term removed; all its direct coefficients
are infinitesimal except the bounded coefficient $t_0$ on $B_u$.
The construction thus includes arbitrarily large $\beta$.

\paragraph{The unrestricted $Y$ normal and cross coefficient.}
Put
\[
\begin{split}
 Y_P={}&(1+c)\{-c_0+a_0(1+a)\}
 +S\{m(1+ra)-d(1-r)c+(\alpha+\tfrac12)(1+a)r\},\\
 Y_C={}&\ell Y_P-m\theta a+d(1+ha)-(\alpha+\tfrac12)(1+a).
\end{split}\tag{UN17}
\]
The actual source has $Y_C(x,q)=0$, so finite monic division gives
$Y_C=(a-x)Y_a+\delta Y_c$ with bounded polynomial coefficients.
The complete remaining entries are
\[
\begin{split}
 L_Y={}&\sqrt{r(1-r)}\{Y_PG+\sqrt S\,Y_aR+Y_c(wB)\},\\
 H_Y={}&rc_0(1+c)+(1-r)a_0(1+a)(1+c)\\
 &+Sr(1-r)\{ma+dc+(\alpha+\tfrac12)(1+a)\}\sim\tfrac12,\\
 H_{XY}={}&-uw\sqrt{Sr(1-r)}(m+d)\in\operatorname{Inf}.
\end{split}\tag{UN18}
\]
In particular $Y_P$ is infinitesimal. Every coefficient is bounded
through $r=0$ and $r=1$ openings; no inverse of $r$ or $1-r$ remains
after the displayed pure-source normal scaling. UN14--UN18 account
for all terms of SN5, including both independent state normals.

\paragraph{Free pivots, including the $\theta$-infinitesimal opening.}
Reserve the three actual circuit rows in UN8. In either sphere
chart, $wB$ is free at infinitesimal cost: use $\sqrt{d_s}B$ when
$d_s$ is a unit, and otherwise use the exact identity
\[
 wB=-\frac{w+\sqrt q}{\sqrt{z_n}}E_2
                 -\frac{w(w+\sqrt q)}\ell G.
 \tag{UN19}
\]
If $\theta$ is a unit, $A_{00}$ already supplies a strict $N_0$
pivot. Since $q/n$ is infinitesimal, the exact identity
\[
 B_u=-(w+\sqrt q)\sqrt{q/n}\,E_1-(w+\sqrt q)N_0
 \tag{UN20}
\]
buys any fixed $B_u$ pivot at infinitesimal cost. Then $wB_u$ is
free by multiplication by the pure-testing infinitesimal $w$.

If $\theta$ is infinitesimal, first buy the missing $N_0$ pivot
from the actual circuit, using
\[
 N_0=\sqrt{q/n}(E_2-E_1)
       -\frac{w-\sqrt q}{\sqrt n}E_3
       +\frac{\sqrt{\kappa/q}}\ell\,wG.
 \tag{UN21}
\]
All coefficients are infinitesimal: $n\sim1/2$ and $\kappa/q$
is infinitesimal by UN4. Identity UN21 holds without requiring
$z_n$ to be a unit; the apparent source factor in UN19 has
canceled exactly before estimating. After this purchase, UN20
buys $B_u$, and multiplication by $w$ buys $wB_u$. Thus the
order of purchases involves no circular use of a missing pivot.
The bounded, possibly signed $A_{00}$ in this source subcase
has zero residue and is absorbed after the strict $N_0$ pivot is
bought. No inverse of $\theta$ is introduced.

\paragraph{Final allocation and unchanged testing joins.}
Use the strict rows $G,N_x,N_y$, the applicable $N_0$ pivot, and
fixed retained fractions of $E_1,E_2,E_3,\sqrt{d_s}B$ wherever
nonzero. Purchase sufficiently large real pivots on $B_u,wB,wB_u$
by UN19--UN21. Every purchase has infinitesimal Gram cost on
the preceding strict rows. All signed terms in UN14--UN18 are
then either infinitesimal matrix perturbations or bounded
couplings to those freely purchased helper pivots. A fixed
sufficiently large choice makes the complete residue matrix
strictly positive definite. The separated strict bounded-Gram
lemma proves finite SOS. Restore the positive packets retained
in UN8 and UN16, including $\beta(1+c)N_x^2$; nothing from the
original parent is omitted.

This proves all infinitesimal $u,w$, their literal zero faces and
arbitrary states. Noninfinitesimal Alice tests use CE1--CE10 because
$C,A$ are unbounded. The other Bob tests use the surviving SE10--SE12
unmatched certificate, since $e$ is infinitesimal and
$\overline\ell=1$. The original sign gauges and real/imaginary state
split are exact; the circuit argument treats both signs of $WZ$.
These joins give the claimed all-tests theorem on UN1.

The theorem just proved has precisely the source domain UN1. Its join with the bounded-$h/q$ theorem and the other C14.4 source cells is printed in Section~\ref{component:pfd0dd7bb10cf}.


\section{\texorpdfstring{All testing charts for the unit-$\kappa,\nu,\ell$ source}{All testing charts for the unit-kappa,nu,ell source}}
\label{component:p3f02c51074be}

\label{p3f02c51074be:sec:new-c144-unit-kappa-all-tests}

\paragraph{Exact scope and conclusion.}
Retain the original physical C14.4 source, its strict source gaps,
and its original prior-complement intersections. Assume
\[
 S=\sigma+L\in\operatorname{Inf}_{>0},\qquad
 \kappa=x/S,\quad\nu=\sigma/S,\quad\ell=L/S
 \quad\hbox{are positive bounded units},\qquad\nu+\ell=1.
 \tag{FU1}
\]
Together with the matched theorem UKR1--UKR21, the argument below
proves that the complete original parent is a finite separated tensor
SOS for every projective test and every state on FU1. It includes
arbitrary openings of $M,\theta,r,1-r$, literal $D=0$ and $\beta=0$,
and independent source and testing infinitesimals. It makes no claim
at the $\nu$- or $\ell$-infinitesimal source endpoints.

All normalizations are pure-source or pure-testing operations.
All real Schur complements below prove strictness of a real residue
matrix; no such inverse is used as a tensor-ring coefficient. The
strict bounded-Gram lemma is the one in the surviving
Section~\ref{component:p472f434581f3}.

\paragraph{The full parent and opening-safe source bounds.}
Write $D_u=1+u^2,D_v=1+v^2$, and divide the parent by $4$:
\[
\begin{split}
 F/4={}&m(vW-X+uvY-uZ)^2+CD_v(uW-Y)^2+AD_uD_vY^2\\
 &+(D+\tfrac12)(uW-Y+uvX-vZ)^2
 +(\alpha+\tfrac12)D_u(vY-Z)^2\\
 &+(\beta+\tfrac12)D_v(X+uZ)^2-\tfrac12D_uD_vZ^2.
\end{split}\tag{FU2}
\]
Keep the exact FM2--FM3 source formulas, with
$c_0=SC(1-r)$, $a_0=SAr$, $d_0=D/S$, and $a_1=\alpha/S$.
These four coefficients and $\beta\theta$ are bounded, with no
inverse of the possibly small $M$ or $\theta$ in their normalized
formulas. The $D/S$ numerator has its factor $S$ canceled exactly.

The following letters refer only to real residues:
\[
 A_* =\nu-\kappa h=\nu E,\quad R=1-A_*,\quad
 \nu=A_*+\kappa h,\quad\ell=1-A_*-\kappa h.
 \tag{FU3}
\]
The exact source yields
\[
\begin{gathered}
 E\ge1/18,\qquad h-A_*\ge0,\qquad
 A_*^2-\kappa(h-A_*)\ge0,\qquad R-\kappa\ge0,\\
 c_0=a_0=E/2,\qquad
 d_0=\{A_*^2-\kappa(h-A_*)\}/2,\\
 a_1=\frac{\ell\{\kappa h+(R-\kappa)A_*\}}{2\nu},
 \qquad 2\beta\theta=h-A_*.
\end{gathered}\tag{FU4}
\]
In particular $h\ge\nu/(1+\kappa)$, so $h,E,a_1,c_0,a_0$
have positive real lower bounds. The actual source has $m$ bounded
below by a positive real, while $C,A$ are unbounded. Here $m$
can be unbounded, and so can $\beta$. They cannot both be unbounded:
$m$ unbounded forces $R=\kappa$, whence
$\ell=\kappa\theta$ makes $\theta$ a unit; $\beta$ unbounded
forces $\theta=0$, whence $R-\kappa=\ell$ makes $M$ a unit.
For bounded $m,\beta$ their real values are
\[
 m=\frac{\kappa^2h}{2\nu(R-\kappa)},\qquad
 \beta=\frac{h-A_*}{2(1-h)}.
 \tag{FU5}
\]
These formulas are used only when the indicated real denominators
are positive. Every estimate below for the actual parent uses
bounded normalized source coefficients before taking residues.

\paragraph{A complete finite-Bob frame, including the state normals.}
Assume $u$ is infinitesimal and $v$ is finite. We may cap $m$ and
$\beta$ by finite pure-source values $m_c>0,\beta_c\ge0$ and
retain the complete excesses in their original event squares.
When a weight is bounded, keep it unchanged. When it is unbounded,
the sufficiently large real cap is chosen below. Define
\[
 B_0=X+uZ,\quad N_y=\frac{Y-ruW}{\sqrt{Sr(1-r)}},\qquad
 T=\kappa W,\ U=uZ,\ P=\frac{uW}{\sqrt S},\ Q=\sqrt S Z.
 \tag{FU6}
\]
All six rows are legal separated expressions, and $TU=\kappa PQ$
holds exactly. No restriction is placed on their values or on state
ratios. Set $a=u^2$, $s_y=\sqrt{Sr(1-r)}$, and
\[
\begin{split}
 M_0&=(v/\kappa)(1+ra)T-B_0,\\
 A_P&=D_v\{c_0(1-r)+a_0r(1+a)\}\\
 &\hspace{5mm}+S\{(D+\tfrac12)(1-r)^2
       +(\alpha+\tfrac12)(1+a)r^2v^2\}.
\end{split}\tag{FU7}
\]
The entire capped parent is $F_0+2L_yN_y+H_yN_y^2$, where
\[
\begin{split}
 F_0={}&m_cM_0^2+A_PP^2+(1+a)\{d_0(1+a)v^2+a_1\}Q^2
                   +\tfrac12(1+a)v^2U^2\\
 &-2(D+\tfrac12)(1+a)v^2B_0U
 -(1+a)v\{1+2D(1-r)+2\alpha r\}PQ\\
 &+\{(\beta_c+\tfrac12)D_v+(D+\tfrac12)av^2\}B_0^2
 +2(D+\tfrac12)(1-r)uv\sqrt S\,B_0P,
\end{split}\tag{FU8}
\]
\[
\begin{split}
 L_y={}&m_cuvs_yM_0
 +\sqrt{r(1-r)}D_v\{-c_0+a_0(1+a)\}P\\
 &+S\sqrt{r(1-r)}\{-(D+\tfrac12)(1-r)
          +(\alpha+\tfrac12)(1+a)rv^2\}P\\
 &-(D+\tfrac12)uvs_yB_0
       +\sqrt{r(1-r)}v(1+a)(D-\alpha)Q,\\
 H_y={}&rc_0D_v+(1-r)a_0(1+a)D_v\\
 &+Sr(1-r)\{m_cu^2v^2+D+\tfrac12
                      +(\alpha+\tfrac12)(1+a)v^2\}.
\end{split}\tag{FU9}
\]
These are exact identities in the two unrestricted original state
coordinates $X,Y$; $B_0$ is not set to zero. Their coefficients are
bounded across all source openings. In $L_y$ every coefficient has
zero residue, since $c_0,a_0$ have the same residue, and the other
terms contain $u$ or $S$. The $N_y$ pivot has residue $ED_v/2>0$.
Thus it suffices to make the remaining five-row real residue strict.

\paragraph{Finite nonzero unmatched Bob residue.}
Take $v>0$ with a positive finite residue different from
$b=\sqrt{\ell/\nu}$. The complete real residue of FU8--FU9 is
\[
\begin{split}
 &m_c\{(v/\kappa)T-B_0\}^2+qD_vB_0^2-v^2B_0U
       +\tfrac12v^2U^2-vPQ\\
 &\hspace{8mm}+\tfrac E2D_vP^2+(d_0v^2+a_1)Q^2
             +\tfrac E2D_vN_y^2,\qquad q=\beta_c+\tfrac12.
\end{split}\tag{FU10}
\]
For the moment use the uncapped finite values $m,\beta$ from FU5.
Eliminating $B_0$ in this real matrix, and using $TU=\kappa PQ$
only to name the adjusted cross coefficient, gives
\[
 a_TT^2+a_UU^2+a_PP^2+a_QQ^2-2fPQ,
 \tag{FU11}
\]
where, for $c=v^2$ and $H=m+q(1+c)$,
\[
\begin{gathered}
 a_T=\frac{mq(1+c)c}{\kappa^2H},\quad
 a_U=\frac{c(2H-c)}{4H},\quad a_P=E(1+c)/2,\quad
 a_Q=d_0c+a_1,\\
 f=\frac{v}{2}(1+mc/H).
\end{gathered}\tag{FU12}
\]
The replacement of the real Schur cross term $TU$ by $\kappa PQ$
is implemented in the actual bounded Gram by a real multiple of
the exact zero $TU-\kappa PQ$. The elimination of $B_0$ itself
is only a real strictness test. Every diagonal is positive,
including when $D$ or $\beta$ is zero.
The reduced determinant in a formal positive real variable $y$ is
\[
 \Delta(y,c)=(a_T\kappa^2+a_Py)(a_Uy+a_Q)-f^2y.
 \tag{FU13}
\]
It obeys $4H\Delta/(1+c)=\mathcal P(y,c)$ with
\[
 \mathcal P=p_{22}y^2c^2+p_{21}y^2c+p_{12}yc^2
              +p_{11}yc+p_{10}y+p_{02}c^2+p_{01}c,
 \tag{FU14}
\]
\[
\begin{array}{ll}
 p_{22}=E\beta,&p_{21}=E(m+q),\\
 p_{12}=2m\beta+2qEd_0,&
 p_{11}=-m-q+2E\{(m+q)d_0+qa_1\},\\
 p_{10}=2(m+q)Ea_1,&p_{02}=4mqd_0,\quad p_{01}=4mqa_1.
\end{array}\tag{FU15}
\]
All coefficients except possibly $p_{11}$ are nonnegative, and
$p_{21},p_{10},p_{01}$ are strictly positive. Direct substitution of
FU3--FU5 gives the exact three contact identities
\[
 \mathcal P(\kappa,\ell/\nu)=
 \mathcal P_y(\kappa,\ell/\nu)=
 \mathcal P_c(\kappa,\ell/\nu)=0.
 \tag{FU16}
\]
These finite rational identities can also be checked without
radicals. For positive real $y,c$,
\[
 \frac{\mathcal P(y,c)}{yc}
 =p_{22}yc+p_{21}y+p_{12}c+p_{11}
                     +p_{10}/c+p_{02}c/y+p_{01}/y.
 \tag{FU17}
\]
As a real function of $(\log y,\log c)$, this is strictly convex:
each nonconstant summand has nonnegative Hessian, and
$p_{21}y+p_{10}/c$ alone has positive definite Hessian. FU16 says
its value and gradient vanish at the contact. It is consequently
positive at every other positive point. This real argument is only
a way to determine a finite residue matrix, not a use of logarithms
or mixed inverses in the tensor ring.

Since $c\ne\ell/\nu$, FU13 is positive for all $y>0$.
Minimizing its quotient by $y$ gives
\[
 f<\kappa\sqrt{a_Ta_U}+\sqrt{a_Pa_Q}.
 \tag{FU18}
\]
Indeed the minimum of
$a_Pa_Uy+a_T\kappa^2a_Q/y$ is
$2\kappa\sqrt{a_Ta_Ua_Pa_Q}$, attained at a positive $y$;
equality in FU18 would give a zero of FU13. Choose a real
$\lambda$ with $|\lambda|<\sqrt{a_Ta_U}$ and
$|f-\kappa\lambda|<\sqrt{a_Pa_Q}$. Adding the exact zero
$-2\lambda(TU-\kappa PQ)$ makes both $2\times2$ blocks in
FU11 strictly positive. Restoring the positive real $B_0$ pivot
and the $N_y$ pivot proves strictness of the complete six-row
residue. All actual Gram entries were bounded in FU8--FU9, so
the strict bounded-Gram lemma gives the tensor SOS.

\paragraph{Unbounded $m$ or $\beta$ in the unmatched chart.}
The preceding strictness persists under a sufficiently large
finite cap in the one possible unbounded weight. This assertion
has an opening-safe direct verification. If $m\to\infty$,
divide FU14--FU15 by $m$ before taking the real limit. Then
\[
 (p_{22},p_{21},p_{12},p_{11},p_{10},p_{02},p_{01})/m
 \longrightarrow
 (0,E,2\beta,-1+2Ed_0,2Ea_1,4qd_0,4qa_1).
 \tag{FU19}
\]
At $R=\kappa$, the three contact identities FU16 remain exact
for this limit polynomial. The strict convexity proof FU17 still
applies because its $y$ and $1/c$ coefficients are positive.
The Schur coefficients tend to
\[
 a_T\to q(1+c)c/\kappa^2,\qquad
 a_U\to c/2,\qquad f\to v(1+c)/2,
 \tag{FU20}
\]
with $a_P,a_Q$ unchanged. Therefore FU18 is strict at the limit,
and continuity makes it strict for a sufficiently large finite
real $m_c$.

If $\beta\to\infty$, divide instead by $q$ before taking the
limit at $h=1$. The corresponding coefficient vector is
\[
 (E,E,2m+2Ed_0,-1+2E(d_0+a_1),
                         2Ea_1,4md_0,4ma_1),
 \tag{FU21}
\]
again with the three exact contact identities and the two strictly
positive convexity coefficients. Now
$a_T\to mc/\kappa^2$, $a_U\to c/2$, $f\to v/2$.
A sufficiently large finite $\beta_c$ therefore works.
All the other normalized actual coefficients have their real
residues with bounded unit denominators; no $O(S)$ term was
divided by $M$ or $\theta$. Restore the complete excesses
$(m-m_c)(vW-X+uvY-uZ)^2$ and
$(\beta-\beta_c)D_v(X+uZ)^2$. These caps are used only on the
unmatched chart and need not vanish at the actual contact.

\paragraph{Infinitesimal or zero Bob coordinate.}
If $v$ is infinitesimal, use FU8--FU9 with the six rows
\[
 (B_0,T_v,U_v,P,Q,N_y),\qquad T_v=vT,\quad U_v=vU.
 \tag{FU22}
\]
Cap $m$ at any sufficiently small fixed positive real $m_c\le m$,
and discard no excess. One may set $\beta_c=0$ and retain its
whole event. The exact coefficients remain bounded: the $B_0U$
coefficient is now a multiple of $vB_0U_v$, and $M_0$ becomes
$(1+ra)T_v/\kappa-B_0$. Their real residue is
\[
 m_c(T_v/\kappa-B_0)^2+\tfrac12B_0^2+\tfrac12U_v^2
                      +\tfrac E2(P^2+N_y^2)+a_1Q^2>0.
 \tag{FU23}
\]
This also covers literal $v=0$. The rows need not be independent
coordinates on the state space.

\paragraph{Infinitesimal reciprocal Bob coordinate.}
Put $w=1/v$ and use the homogeneous parent $w^2F$ at $w=0$ as
well as for nonzero infinitesimal $w$. The complete signed core
in the quarter-parent is $K_w/2$, with
\[
\begin{split}
 K_w={}&u^2H^2-2u^2HY-2uHW+2uwHX+u^2w^2W^2\\
 &+2u^2wWX-2uw^2WY+(1+u^2+w^2)(X^2+Y^2)-2uwXY,
 \qquad H=wZ.
\end{split}\tag{FU24}
\]
This is the exact CE6 full core, so all state normals are present.
FU4 gives the real alternative
\[
 d_0+\kappa\beta\theta=(\nu E)^2/2>0.
 \tag{FU25}
\]
Thus either $D/S$ has a positive real lower bound, or $\beta$
has one. Both cases allow independent rates for $u,w,S$.

If $\beta$ is strong, use the four rows $(W,X,Y,U)$, $U=uZ$,
in FU24 by replacing $uH$ with $wU$. Every coefficient is bounded
and the core residue is $X^2+Y^2$. The exact $m$ event is
$(W-wX+uY-wU)^2$, and the $\beta$ event is
$D_w(X+U)^2$. Positive finite caps give the strict residue
\[
 m_cW^2+\tfrac12(X^2+Y^2)+\beta_c(X+U)^2>0.
 \tag{FU26}
\]
All other event packets and both excesses remain whole.

If $D/S$ is strong, choose positive real $c,d,m_c$ with
$c\le c_0,a_0$, $d\le D/S$, and $m_c\le m$. Select exactly
\[
 C_c=c/[S(1-r)],\qquad A_c=c/(Sr),\qquad D_c=dS.
 \tag{FU27}
\]
Retain all excesses, including the whole $A_cu^2D_wY^2$.
Writing $Y=ruW+\sqrt{Sr(1-r)}N_y$, the other two selected
$C,A$ terms have the exact harmonic identity
\[
 C_cD_w(uW-Y)^2+A_cD_wY^2=cD_w(P^2+N_y^2).
 \tag{FU28}
\]
Use six rows $(W,V,X,P,N_y,Q)$ with $V=uwZ$ and the exact
monomial identity $VW=wPQ$. Replacing $uH$ by $V$ in FU24,
the potentially non-small cross term $-2VW$ is exactly $-2wPQ$.
Every other off-diagonal core coefficient is infinitesimal after
$Y=\sqrt S\{rP+\sqrt{r(1-r)}N_y\}$; its residue is $V^2+X^2$.
The selected $m$ and $D$ events are the complete squares
\[
 m_c(W-wX+uY-V)^2,\qquad
 d\{uw\sqrt S W+u\sqrt S X-w\sqrt S Y-Q\}^2.
 \tag{FU29}
\]
The full selected residue is strictly positive:
\[
 m_c(W-V)^2+\tfrac12(V^2+X^2)+c(P^2+N_y^2)+dQ^2>0.
 \tag{FU30}
\]
The strict bounded-Gram lemma applies to the exact complete
forms, and all unselected positive packets are restored. Both
reciprocal proofs include literal $w=0$ and $u=0$.

\paragraph{Exhaustion of tests and exact joins.}
If $u$ is noninfinitesimal, $C,A$ are unbounded by FU4. The
surviving two-adjacent-events theorem CE1--CE10 in
Section~\ref{component:p89fd3edb3edd}
therefore applies to the original full parent and all states.
Its hypotheses require only nonnegative event weights and these
two unbounded coefficients; no small-$\kappa$ source statement
is imported.

For infinitesimal $u$, every positive projective $v$ is either
infinitesimal, a finite unit, or has infinitesimal reciprocal.
FU22--FU30 cover the two end cases. A finite unit has either the
matched residue, covered by UKR1--UKR21, or a different residue,
covered by FU10--FU21. The transformations
\[
 (u,Y,Z)\mapsto(-u,-Y,-Z),\qquad
 (v,X,Z)\mapsto(-v,-X,-Z)
 \tag{FU31}
\]
preserve FU2 exactly and handle both signs. Literal axes and
projective endpoints were included in the displayed raw frames.
Complex states split into real and imaginary parts. No proof in
this exhaustion bounds, normalizes, or relates the testing-state
ratios $W/Z$, $X/Z$, or $Y/Z$.

Consequently every testing chart is covered on exactly FU1,
with every original source predicate retained. The remaining
source endpoints are separate from this conclusion.


\section{\texorpdfstring{A fresh matched-Bob completion for unit $\kappa,\nu,\ell$}{A fresh matched-Bob completion for unit kappa,nu,ell}}
\label{component:pd26b89cdb3f0}

\label{pd26b89cdb3f0:sec:new-c144-unit-kappa-matched}

\paragraph{Scope.}
This theorem retains the entire original physical C14.4 source,
including its strict gaps and prior complements. Assume
$S\in\operatorname{Inf}_{>0}$ and $\kappa=x/S,\nu=\sigma/S,
\ell=L/S$ are positive bounded units, with $\nu+\ell=1$.
The testing chart is $u\in\operatorname{Inf}$ and finite $v$
whose absolute value has the same residue as
$b=\sqrt{\ell/\nu}$. All real or complex states are allowed.
The unmatched Bob tests on this source are proved in Section~\ref{component:p3f02c51074be}; the source endpoints are separate cells in Section~\ref{component:pfd0dd7bb10cf}. The present proof is confined to its stated matched-Bob chart and uses no mixed inverse.

The dictionary FM2--FM3 remains an exact source identity on this
larger domain, but its small-$\kappa$ estimates FM4 are replaced
below. Set $j=\sqrt{\nu\ell}$ and retain the rows
\[
 N_x=X+huZ,\quad N_y=(Y-ruW)/\sqrt{Sr(1-r)},\qquad
 T=\kappa W,\ U=uZ,\ P=uW/\sqrt S,\ Q=\sqrt S Z,
\]
\[
 N=T-jU,\quad G=P-jQ,\quad \delta=v-b,\quad
 B=\delta Q,\quad B_u=\delta U,\quad R_Z=(u^2-x)Z.
 \tag{UKR1}
\]
The radial and cross identities FM8 hold unchanged. Their
coefficients on $R_Z$ remain infinitesimal because $S,u$ are
infinitesimal and $\kappa,1/j$ are bounded.

\paragraph{A sharper exact angular budget.}
For any fixed real $0<\rho<1$, the monomial identity $TU=\kappa PQ$
gives
\[
 \frac{\delta^2}{j^2(1-\rho^2)}N^2+
 \frac{2\kappa\delta}{j(1+\rho)}GB+
 \frac{2\kappa}{1+\rho}B^2-B_u^2
   =\frac{(\delta T-\rho jB_u)^2}{j^2(1-\rho^2)}.
 \tag{UKR2}
\]
Its first two coefficients have zero residue. Thus a fixed
$cB_u^2$ pivot costs $q\kappa cB^2$ in residue, where
$q=2/(1+\rho)>1$ can be chosen arbitrarily close to one.
Every inverse in UKR2 is a pure source unit or a positive real.
There is no inverse of $\delta$.

\paragraph{The actual source through both coefficient openings.}
The source has $1/18\le E=1-\kappa h/\nu<1$ and
$\kappa$ is bounded because $1/\zeta$ is bounded. The exact
$\beta\theta$ identity gives a positive real lower bound on $h$.
Consequently $m$ has a positive real lower bound, although it may
be unbounded. Let
\[
 R_s=1+xr,\quad p=\frac{\kappa\theta}{\ell R_s}\in(0,1),\qquad
 \mu_c=m(1-p)=
 \frac{\kappa^2h(1+x)(1+\nu Sr)}{2\nu\ell P_0R_s}.
 \tag{UKR3}
\]
The quantity $\mu_c$ is a positive bounded unit through $M\to0$.
All of $SC(1-r),SAr,D/S,\alpha/S,\beta\theta$ are bounded
without inverses of $r,1-r,\theta,D,\beta$; this follows directly
from FM3 after the explicit factor $S$ in the $D/S$ numerator
is canceled. Their limits are evaluated below.

Only in the real residue calculation use the abbreviations
\[
 A=\nu-\kappa h=\nu E,\quad R=1-A,\quad d=h-A.
 \tag{UKR4}
\]
Here $A,R,h,\nu,\ell,\kappa$ are positive units, and the
original nonnegative weights and positive $M$ imply
\[
 d\ge0,\quad A^2-\kappa d\ge0,\quad R-\kappa\ge0,
 \qquad \nu=A+\kappa h,\quad\ell=1-A-\kappa h.
 \tag{UKR5}
\]
Indeed in residue $2\beta\theta=d$, $2D/S=A^2-\kappa d$,
and $\ell M=R-\kappa$. These are additive residue identities;
no relative error estimate is asserted at a vanishing opening.

\paragraph{Retain unbounded weights in actual contact-zero rows.}
Define
\[
 V=v(W+uY),\quad B_0=X+uZ,\quad N_m=B_0-pV.
 \tag{UKR6}
\]
At the actual source contact $u=\sqrt x$, $v=b$,
$W/Z=j\sqrt{S/\kappa}$, $X=-h\sqrt x Z$, $Y=r\sqrt x W$,
both $N_m$ and $N_x$ vanish. This is an identity at the source
contact, not an equation imposed on independent tests.

If $m$ is bounded, set $m_c=m$. If it is unbounded, choose a
positive real $m_c$ sufficiently large, as justified below.
Similarly keep $\beta_c=\beta$ when $\beta$ is bounded, and
choose a sufficiently large positive real $\beta_c$ when it is
unbounded. An unbounded positive field element exceeds either
fixed real cap. Retain the entire positive excesses
\[
 (m-m_c)N_m^2+(\beta-\beta_c)D_vN_x^2.
 \tag{UKR7}
\]
The remaining parent is exact, with its $m$ part replaced by
\[
 m_0 B_0^2-2m_1B_0V+m_2V^2,\qquad
 m_0=m_c,\quad m_1=\mu_c+m_cp,\quad
 m_2=\mu_c(1+p)+m_cp^2.
 \tag{UKR8}
\]
This follows from the identity
$m(V-B_0)^2-(m-m_c)(B_0-pV)^2$
and $\mu_c=m(1-p)$. Every coefficient of UKR8 is bounded.
The residual $\beta$ terms are exactly
\[
 \beta_cD_vN_x^2+2(\beta\theta)D_vN_xuZ
                         +(\beta\theta)\theta D_vu^2Z^2.
 \tag{UKR9}
\]
All coefficients of UKR9 are bounded. The complete cross terms
are retained, and both original excesses in UKR7 are restored
once at the end.

\paragraph{The complete bounded scalar and normal frame.}
Apply the FM9--FM19 polynomial construction to this capped
parent with the following explicit replacements: in $A_T$ replace
$m/\kappa^2$ by $m_2/\kappa^2$; in $A_U$ replace its
$m\theta^2$ term by $m_0\theta^2$; in $K$ replace its
$m\theta(1+ra)$ term by $m_1\theta(1+ra)$. In the complete
first-normal formulas replace
\[
\begin{array}{c|c}
 \text{entry}&\text{replacement}\\ \hline
 H_x\text{ first term}&m_0\\
 H_y\text{ term }mu^2v^2&m_2u^2v^2\\
 H_{xy}\text{ term }m&m_1\\
 L_x\text{ coefficient }-m/\kappa&-m_1/\kappa\\
 L_U\text{ terms }m\theta-(mj/\kappa)v(1+ra)&
                  m_0\theta-(m_1j/\kappa)v(1+ra)\\
 Y_P\text{ term }mv^2(1+ra)&m_2v^2(1+ra)\\
 Y_Q\text{ term }-m\theta a&-m_1\theta a.
\end{array}\tag{UKR10}
\]
Include $\beta_cD_v$ in $H_x$ and omit the separately retained
$\beta D_vN_x^2$ term of FM15. All other entries are unchanged.
These replacements follow by expanding UKR8--UKR9 in the two
unrestricted state normals. They are not a scalar-only operation.

The graph scalar and both complete first normals still vanish to
the orders required by FM11 and FM19: the subtracted rows in UKR7
vanish at the actual source contact. Their squares therefore
change neither the value nor any first derivative there.
The same finite monic divisions supply the exact bounded FM13
scalar and FM18--FM19 full normal conductors. The $Y$ normal has
zero-residue cross coefficients, since $SC(1-r)$ and $SAr$
both have residue $E/2$ and the remaining $Y$ coefficients
contain $S$ or $u^2$. Its diagonal residue is $E/(2\nu)>0$.

The only nontrivial real residue uses $(N_x,N,G,B,B_u)$.
Writing
\[
 H_* = h-\tfrac12+\tfrac\theta2(\ell+\kappa h),\qquad
 r_* = H_*-j^2m_2/\kappa^2,
 \tag{UKR11}
\]
its exact entries are
\[
\begin{gathered}
 H_c=\begin{pmatrix}
 m_0+(\beta_c+1/2)/\nu&-(m_1/\kappa)b\\
 -(m_1/\kappa)b&(m_2/\kappa^2)b^2
 \end{pmatrix},\qquad A_G=E/(2\nu),\\
 z=\begin{pmatrix}-j(m_1/\kappa+E)\\-r_*/\nu\end{pmatrix},\qquad
 g_*=E\ell-m_1\theta-\tfrac12+2m_2\ell/\kappa
                                        +\kappa r_*/\nu,\\
 H_B=\nu(m_1\theta+\tfrac12).
\end{gathered}\tag{UKR12}
\]
Thus the residue form is
$ (N_x,N)H_c(N_x,N)^t+A_GG^2+2g_*GB+H_BB^2
 +2z^t(N_x,N)^tB_u$ before the UKR2 pivot is added.
In particular the $\kappa r_*/\nu$ contribution to $g_*$ is
retained; it is not infinitesimal in the present source sector.

\paragraph{The exact real contact curvature and all source faces.}
First take an interior real source, $\theta>0$ and $R>\kappa$,
without capping, so
\[
 m=\frac{\kappa^2h}{2\nu(R-\kappa)},\qquad
 \beta=\frac{h-A}{2\theta},\qquad
 m_0=m_1=m_2=m.
 \tag{UKR13}
\]
The two hard pivots and the angular pivot obey the exact identities
\[
\begin{split}
 (N_x,N)H_c(N_x,N)^t
   &=\frac m{\kappa^2}(bN-\kappa N_x)^2
                                  +\frac R{2\nu\theta}N_x^2,\\
 J_*&=H_B-\frac{g_*^2}{A_G}-\kappa z^tH_c^{-1}z\\
 &=\frac{AR+\kappa(h-A)}{2AhR}
                  \{4A^2hR-\kappa(Rh-A)^2\}.
\end{split}\tag{UKR14}
\]
All inverses in UKR14 are used only in a real finite matrix
calculation to prove strictness. They are not inverses in the
tensor ring. The crucial bracket has the positive decomposition
\[
\begin{split}
4A^2hR-\kappa(Rh-A)^2={}&A^3R(4-A)+A^4(R-\kappa)\\
 &+A^2R(h-A)(3+A+2\kappa)\\
 &+R^2(h-A)\{A^2-\kappa(h-A)\}.
\end{split}\tag{UKR15}
\]
Its first term is a positive unit and all remaining terms are
nonnegative by UKR5. This includes the literal $D=0$ and
$\beta=0$ faces and both limiting openings.

For completeness the strictness is uniform toward those openings.
Put $\mu=m/\kappa^2$ and retain $H_*$ from UKR11 with uncapped
coefficients. Then the two entries of $H_c^{-1}z$ are
\[
 z_x=-\frac{2\nu\theta}{R}
                   \left(jE+\frac{\kappa H_*}{\nu b}\right),
 \qquad z_N=\frac\kappa b z_x+\nu-\frac{H_*}{\mu\ell}.
 \tag{UKR16}
\]
They stay bounded. Also
\[
 g_*=-\frac{2A^2+A h\kappa+A\kappa-A-h\kappa}{2\nu}
 \tag{UKR17}
\]
stays bounded, while $A_G$ and $J_*$ have positive real lower
bounds near any fixed allowed residue point. UKR14 gives a
positive real lower bound on $H_c$ because $\mu$ and
$R/(2\nu\theta)$ have positive lower bounds and $b$ is a unit.
Completing the real source quadratic form with the bounded
coefficients UKR16--UKR17 therefore gives a uniform positive
lower bound on
\[
 Q_{\rm phys}=(N_x,N)H_c(N_x,N)^t+A_GG^2+2g_*GB+H_BB^2
                         +2\sqrt\kappa\,z^t(N_x,N)^tB.
 \tag{UKR18}
\]
This is an auxiliary real matrix used to choose a diagonal budget.
It does not impose $B_u=\sqrt\kappa B$, or any relation between
source and testing angles, in the tensor certificate.

\paragraph{Why finite caps preserve a strict real residue.}
When $m,\beta$ are bounded no cap is needed and UKR18 is already
strict. If $m$ is unbounded, $R-\kappa=0$ in residue; since
$\ell$ is a unit and $\ell=R-\kappa h$, $\theta$ is then a unit,
so $\beta$ is bounded. If $\beta$ is unbounded, $\theta=0$ in
residue, and $R-\kappa=\ell>0$, so $m$ is bounded. Thus at
most one unbounded penalty is present.

For the $m$ opening, the real contact-zero linear form in UKR18 is
\[
 n_m=N_x-(pb/\kappa)N-(pj/\sqrt\kappa)B,
 \tag{UKR19}
\]
and $p\to1$. Its $N$ coefficient is a unit. UKR8 shows that
subtracting the whole $mn_m^2$ leaves a bounded coefficient
matrix with a well-defined limit at the opening. That limit is
strict on $\ker n_m$: approach the residue point by real
$\kappa<R$ with $A,h$ fixed, adjust any kernel vector continuously
to the exact kernel using the unit $N$ coefficient, and apply the
uniform lower bound from UKR18. The parameter inequalities UKR5
are preserved when $\kappa$ is decreased. A real symmetric matrix
strict on the kernel of a nonzero linear form becomes strictly
positive definite after adding a sufficiently large fixed real
multiple of the form's square; this follows by a real block Schur
complement. Choose $m_c$ that large in UKR7.

For the $\beta$ opening the same argument uses the exact penalty
$(\beta/\nu)N_x^2$ and its kernel $N_x=0$. Approach $h=1$ from
below with $A,\kappa$ fixed, which also preserves UKR5, and choose
$\beta_c$ sufficiently large. This is a calculation on real
residue matrices, not a mixed-field limiting certificate. The
actual capped coefficients have those residues because UKR3,
UKR8 and UKR9 have bounded denominators through the original
independently varying $S,M,\theta$ openings. In particular no
estimate divides an $O(S)$ error by $M$ or $\theta$.

We have proved that the actual caps can be chosen so that the
capped matrix UKR18 is strictly positive definite. It follows
that $H_c\succ0$, $A_G>0$, and
\[
 H_B-g_*^2/A_G-\kappa L>0,\qquad L=z^tH_c^{-1}z\ge0.
 \tag{UKR20}
\]

\paragraph{Complete separated SOS allocation.}
Choose $\rho<1$ sufficiently close to one and then a positive
real $c$ so that
\[
 L<c<\frac{H_B-g_*^2/A_G}{q\kappa},\qquad q=2/(1+\rho).
 \tag{UKR21}
\]
Such choices exist by UKR20. Use $c$ times UKR2 to introduce
the $cB_u^2$ pivot. The full real residue on
$(N_x,N,G,B,B_u)$ is strictly positive: eliminate $H_c$ to
leave $(c-L)B_u^2$, and eliminate $A_G$ to leave
$(H_B-q\kappa c-g_*^2/A_G)B^2$.
The $N_y$ pivot is $E/(2\nu)>0$ and its cross coefficients
have zero residue. Every other full normal and radial term has
already been allocated by the exact bounded FM constructions
and UKR10. The strict separated bounded-Gram lemma supplies a
finite tensor SOS. Restore the square on the right of UKR2 and
both complete excess packets UKR7. This is the full original
parent, not its restriction to a state graph.

The original sign transformations of FM5 handle negative $u,v$,
and the proof includes $u=0$. Complex states split into real and
imaginary parts. Arbitrary state ratios and independent source
and testing rates are allowed throughout. Unmatched unit-$\kappa$
Bob tests and the $\nu$- or $\ell$-infinitesimal source endpoints
remain outside this theorem.

